% Generated mechanically from the frozen algebraization.tex target.
% Do not edit this generated file; edit the bound locale source instead.

% OK, start here.
%

\setcounter{chapter}{51}
\chapter{代数几何与形式几何}



\phantomsection
\label{algebraization-section-phantom}


\section{引言}
\label{algebraization-section-introduction}

\noindent
本章继续研究形式代数几何，尤其研究形式对象是否为某个代数对象的完备化这一问题。
一部奠基性参考文献是 \cite{SGA2}。下面列出我们已经在 Stacks 项目中讨论过的结果：
\begin{enumerate}
\item 形式函数定理，见概形上同调，\ref{coherent-section-theorem-formal-functions} 节。
\item 凝聚形式模，见概形上同调，\ref{coherent-section-coherent-formal} 节。
\item Grothendieck 存在性定理，见概形上同调，
\ref{coherent-section-existence}、\ref{coherent-section-existence-proper} 和
\ref{coherent-section-existence-proper-support}.
\item Grothendieck 代数化定理，见概形上同调，\ref{coherent-section-algebraization} 节。
\item 更一般的 Grothendieck 存在性定理，见更多平坦性，
\ref{flat-section-existence} 和
\ref{flat-section-existence-derived}.
\end{enumerate}
下面概述本章内容。

\medskip\noindent
设 $X$ 是一个概形，$\mathcal{I} \subset \mathcal{O}_X$ 是有限型拟凝聚理想层。
本章的许多问题都涉及逆系统 $(\mathcal{F}_n)$，其各项是拟凝聚
$\mathcal{O}_X$-模，并且它满足
$\mathcal{F}_n = \mathcal{F}_{n + 1}/\mathcal{I}^n\mathcal{F}_{n + 1}$。
一个重要的特殊情形是：$X$ 是诺特环 $A$ 上的概形，并且
$\mathcal{I} = I \mathcal{O}_X$，其中 $I \subset A$ 是某个理想。在上同调，
\ref{cohomology-section-inverse-systems}、
\ref{cohomology-section-inverse-systems-bis} 和
\ref{cohomology-section-inverse-systems-tri} 节中，我们给出了一些一般结果。
在本章中，\ref{algebraization-section-ML-degree-zero} 和
\ref{algebraization-section-formal-functions-principal} 节包含概形与拟凝聚模所特有的结果。
在 \ref{algebraization-section-formal-sections-cd-one} 节中，我们证明：
$\lim H^p(X, \mathcal{F}_n)$ 上的极限拓扑是 $I$-进拓扑，如果
$\text{cd}(A, I) = 1$。本章的一个主题是说明：在主理想情形 $I = (f)$ 下证明的结果，
只假设 $\text{cd}(A, I) = 1$ 时仍然成立。

\medskip\noindent
在 \ref{algebraization-section-derived-completion} 节中，我们讨论环化位点
$(\mathcal{C}, \mathcal{O})$ 上的模关于有限型理想层 $\mathcal{I}$ 的导出完备化。
这一节自然延续更多代数，\ref{more-algebra-section-derived-completion} 节中
所述的交换代数导出完备化理论。第一个主要结果是导出完备化存在。
第二个主要结果是：对环化位点的态射 $f$，导出完备化与导出推出可交换：
$$
(Rf_*K)^\wedge = Rf_*(K^\wedge)
$$
前提是上方的理想层局部地由来自下方理想的截面生成；见引理
\ref{algebraization-lemma-pushforward-commutes-with-derived-completion}。需要强调的是，若所讨论的
理想都是全局有限生成的，则这两个主要结果都非常初等；本章中这一理论的所有应用
都属于这种情形。上面展示的等式是形式函数定理的“正确”版本，见
\ref{algebraization-section-formal-functions} 节中的讨论。

\medskip\noindent
设 $A$ 是诺特环，$I, J$ 是 $A$ 的两个理想。设 $M$ 是有限 $A$-模。
本章接下来的主题是映射
$$
R\Gamma_J(M) \longrightarrow R\Gamma_J(M)^\wedge
$$
它从 $M$ 的局部上同调映入其导出 $I$-进完备化。事实证明，若施加合适的深度条件，
这个映射在一定次数范围内的上同调上成为同构。在
\ref{algebraization-section-algebraization-sections-general} 节中，我们基本上在刚才所述的一般性下工作。
在 \ref{algebraization-section-algebraization-punctured} 节中，假设 $A$ 是局部环，且
$J = \mathfrak m$ 是极大理想。我们建议读者先读这一节，再读本章这一部分的另外两节。
最后，在 \ref{algebraization-section-bootstrap} 节中，我们从局部情形逐步推导，在一般情形中得到更强的结果。

\medskip\noindent
在本章的下一部分，我们利用局部上同调完备化的结果，给出关于凝聚模完备化之
上同调的一组并非穷尽的结果。更准确地说，设 $A$ 是诺特环，$I \subset A$
是理想，$U \subset \Spec(A)$ 是开子概形。若 $\mathcal{F}$ 是凝聚
$\mathcal{O}_U$-模，则可以考虑映射
$$
H^i(U, \mathcal{F}) \longrightarrow \lim H^i(U, \mathcal{F}/I^n\mathcal{F})
$$
并问在一定次数范围内是否得到同构。在 \ref{algebraization-section-algebraization-sections} 节中，
我们处理一些 $U$ 是局部环穿孔谱的例子。在
\ref{algebraization-section-algebraization-sections-coherent} 节中，我们讨论一般情形。
在 \ref{algebraization-section-connected} 节中，我们把所得的一些结果应用于代数几何中的连通性问题。

\medskip\noindent
本章余下各节专门讨论凝聚形式模的代数化。换言之，给定如上的凝聚模逆系统
$(\mathcal{F}_n)$，它位于 $U$ 上并满足
$\mathcal{F}_n = \mathcal{F}_{n + 1}/I^n\mathcal{F}_{n + 1}$
我们要问：是否存在一个凝聚 $\mathcal{O}_U$-模 $\mathcal{F}$，使得
$\mathcal{F}_n = \mathcal{F}/I^n\mathcal{F}$
对所有 $n$ 都成立。我们建议读者阅读 \ref{algebraization-section-algebraization-modules} 节，
其中给出了这一问题的精确表述、一个有用的一般结果
（引理 \ref{algebraization-lemma-when-done}）以及一个非平凡的应用
（引理 \ref{algebraization-lemma-algebraization-principal-variant}）。要证明实质上超出这一情形的结果，
还需要发展相当多的理论。本章所得此类结果中最强的那些，见
\ref{algebraization-section-algebraization-modules-conclusion} 节。



\section{形式截面，I}
\label{algebraization-section-ML-degree-zero}

\noindent
我们建议先阅读《上同调》，第 \ref{cohomology-section-inverse-systems} 节。

\begin{lemma}
\label{algebraization-lemma-properties-system}
设 $X$ 是概形。设 $\mathcal{I} \subset \mathcal{O}_X$
是拟凝聚理想层。设
$$
\ldots \to \mathcal{F}_3 \to \mathcal{F}_2 \to \mathcal{F}_1
$$
是拟凝聚 $\mathcal{O}_X$-模的一个逆系统，并且
$\mathcal{F}_n = \mathcal{F}_{n + 1}/\mathcal{I}^n\mathcal{F}_{n + 1}$.
置 $\mathcal{F} = \lim \mathcal{F}_n$。则
\begin{enumerate}
\item $\mathcal{F} = R\lim \mathcal{F}_n$，
\item 对任意仿射开集 $U \subset X$，有
$H^p(U, \mathcal{F}) = 0$，其中 $p > 0$，并且
\item 对每个 $p$ 都有短正合列
$0 \to R^1\lim H^{p - 1}(X, \mathcal{F}_n) \to
H^p(X, \mathcal{F}) \to \lim H^p(X, \mathcal{F}_n) \to 0$。
\end{enumerate}
此外，若 $\mathcal{I}$ 是有限型的，则
\begin{enumerate}
\item[(4)]
$\mathcal{F}_n = \mathcal{F}/\mathcal{I}^n\mathcal{F}$，并且
\item[(5)]
$\mathcal{I}^n \mathcal{F} = \lim_{m \geq n} \mathcal{I}^n\mathcal{F}_m$。
\end{enumerate}
\end{lemma}

\begin{proof}
(1)、(2) 与 (3) 是具有满射过渡映射的拟凝聚模逆系统的一般事实；见
《概形的导出范畴》，引理 \ref{perfect-lemma-Rlim-quasi-coherent}
与《上同调》，引理 \ref{cohomology-lemma-RGamma-commutes-with-Rlim}。
其次，设 $\mathcal{I}$ 是有限型的。
设 $U \subset X$ 是仿射开集。写 $U = \Spec(A)$，并设 $\mathcal{I}|_U$
对应于 $I \subset A$。注意 $I$ 是有限生成理想。
利用拟凝聚 $\mathcal{O}_U$-模与 $A$-模之间的范畴等价
（《概形》，引理 \ref{schemes-lemma-equivalence-quasi-coherent}），
可知 $M_n = \mathcal{F}_n(U)$ 是 $A$-模的一个逆系统，且
$M_n = M_{n + 1}/I^nM_{n + 1}$。因此
$$
M = \mathcal{F}(U) = \lim \mathcal{F}_n(U) = \lim M_n
$$
是一个 $I$-进完备模，并且由《代数》，引理
\ref{algebra-lemma-limit-complete}，有 $M/I^nM = M_n$。这证明了 (4)。
(5) 转化为命题
$\lim_{m \geq n} I^nM/I^mM = I^nM$。
由于 $I^mM = I^{m - n} \cdot I^nM$，这正是说
$I^mM$ 是 $I$-进完备的。这由《代数》，引理
\ref{algebra-lemma-hathat-finitely-generated}
以及 $M$ 的完备性得出。
\end{proof}





\section{形式截面，II}
\label{algebraization-section-formal-functions-principal}

\noindent
我们建议先阅读《上同调》，第
\ref{cohomology-section-inverse-systems-bis} 节与第
\ref{cohomology-section-inverse-systems-tri} 节。

\begin{lemma}
\label{algebraization-lemma-equivalent-f-good}
设 $X$ 是概形。设 $f \in \Gamma(X, \mathcal{O}_X)$。设
$$
\ldots \to \mathcal{F}_3 \to \mathcal{F}_2 \to \mathcal{F}_1
$$
是拟凝聚 $\mathcal{O}_X$-模的一个逆系统。下列条件等价：
\begin{enumerate}
\item 对所有 $n \geq 1$，映射
$f : \mathcal{F}_{n + 1} \to \mathcal{F}_{n + 1}$ 通过
$\mathcal{F}_{n + 1} \to \mathcal{F}_n$ 分解，并给出短正合列
$0 \to \mathcal{F}_n \to \mathcal{F}_{n + 1} \to \mathcal{F}_1 \to 0$；
\item 对所有 $n \geq 1$，映射
$f^n : \mathcal{F}_{n + 1} \to \mathcal{F}_{n + 1}$
通过 $\mathcal{F}_{n + 1} \to \mathcal{F}_1$ 分解，并给出短正合列
$0 \to \mathcal{F}_1 \to \mathcal{F}_{n + 1} \to \mathcal{F}_n \to 0$；
\item 存在一个 $\mathcal{O}_X$-模 $\mathcal{G}$，它是
$f$-可除的，并且 $\mathcal{F}_n = \mathcal{G}[f^n]$；
\item 存在一个 $\mathcal{O}_X$-模 $\mathcal{F}$，它无 $f$-挠，并且
$\mathcal{F}_n = \mathcal{F}/f^n\mathcal{F}$。
\end{enumerate}
\end{lemma}

\begin{proof}
(1)、(2)、(3) 的等价性以及蕴涵 (4) $\Rightarrow$ (1)，已在《上同调》，
引理 \ref{cohomology-lemma-equivalent-f-good} 中证明。
设 (1) 成立。置
$\mathcal{F} = \lim \mathcal{F}_n$。由引理 \ref{algebraization-lemma-properties-system}
的 (4)，有 $\mathcal{F}_n = \mathcal{F}/f^n\mathcal{F}$。
设 $U \subset X$ 是开集，且
$s = (s_n) \in \mathcal{F}(U) = \lim \mathcal{F}_n(U)$。
选取 $n \geq 1$。若 $fs = 0$，则由条件 (1)，$s_{n + 1}$ 属于
$\mathcal{F}_{n + 1} \to \mathcal{F}_n$ 的核。
故 $s_n = 0$。由于 $n$ 是任意的，可知 $s = 0$。
因此 $\mathcal{F}$ 无 $f$-挠。
\end{proof}

\begin{lemma}
\label{algebraization-lemma-formal-functions-principal}
\begin{reference}
这是 \cite[Lemma 1.6]{Bhatt-local} 的略加改进版本。
\end{reference}
设 $A$ 是环且 $f \in A$。设 $X$ 是 $A$ 上的概形。
设 $\mathcal{F}$ 是拟凝聚 $\mathcal{O}_X$-模。
假设 $\mathcal{F}[f^n] = \Ker(f^n : \mathcal{F} \to \mathcal{F})$
稳定。则
$$
R\Gamma(X, \lim \mathcal{F}/f^n\mathcal{F}) =
R\Gamma(X, \mathcal{F})^\wedge
$$
其中右边表示关于理想 $(f) \subset A$ 的导出完备化。因此，对
$p \in \mathbf{Z}$，得到交换图
$$
\xymatrix{
& 0 & 0 \\
0 \ar[r] &
\widehat{H^p(X, \mathcal{F})} \ar[r] \ar[u] &
\lim H^p(X, \mathcal{F}/f^n\mathcal{F}) \ar[r] \ar[u] &
T_f(H^{p + 1}(X, \mathcal{F})) \ar[r] &
0 \\
0 \ar[r] &
H^0(H^p(X, \mathcal{F})^\wedge) \ar[r] \ar[u] &
H^p(X, \lim \mathcal{F}/f^n\mathcal{F}) \ar[r] \ar[u] &
T_f(H^{p + 1}(X, \mathcal{F})) \ar[r] \ar@{=}[u] &
0 \\
&
R^1\lim H^p(X, \mathcal{F})[f^n] \ar[u] \ar[r]^\cong &
R^1\lim H^{p - 1}(X, \mathcal{F}/f^n\mathcal{F}) \ar[u] \\
& 0 \ar[u] & 0 \ar[u]
}
$$
其各行各列均正合，其中
$\widehat{H^p(X, \mathcal{F})} =
\lim H^p(X, \mathcal{F})/f^n H^p(X, \mathcal{F})$
是通常的 $f$-进完备化，而 $T_f(-)$ 表示 $f$-进 Tate 模，
如《代数进阶》，例
\ref{more-algebra-example-spectral-sequence-principal} 中所定义。
\end{lemma}

\begin{proof}
由引理 \ref{algebraization-lemma-properties-system}，有
$\lim \mathcal{F}/f^n\mathcal{F} = R\lim \mathcal{F}/f^n \mathcal{F}$.
其余结论由《上同调》，例
\ref{cohomology-example-formal-functions-principal} 得出。
\end{proof}









\section{形式截面，III}
\label{algebraization-section-formal-sections-cd-one}

\noindent
本节证明引理 \ref{algebraization-lemma-topology-I-adic}。在诺特概形与凝聚模的情形下，
当不假设理想 $I$ 是主理想，而是假设
$\text{cd}(A, I) = 1$ 时，该引理是《上同调》，引理
\ref{cohomology-lemma-topology-I-adic-f} 的类似结果。

\begin{lemma}
\label{algebraization-lemma-cd-one}
设 $I = (f_1, \ldots, f_r)$ 是诺特环 $A$ 的理想。
若 $\text{cd}(A, I) = 1$，则存在 $c \geq 1$ 与映射
$\varphi_j : I^c \to A$，使得 $\sum f_j \varphi_j : I^c \to I$
是包含映射。
\end{lemma}

\begin{proof}
由于 $\text{cd}(A, I) = 1$，补集 $U = \Spec(A) \setminus V(I)$
是仿射的（《局部上同调》，引理 \ref{local-cohomology-lemma-cd-is-one}）。
写 $U = \Spec(B)$。于是 $IB = B$，并且可以写成
$1 = \sum_{j = 1, \ldots, r} f_j b_j$，
其中 $b_j \in B$。由《概形上同调》，引理
\ref{coherent-lemma-homs-over-open}，可将 $b_j$ 表示为映射
$\varphi_j : I^c \to A$，其中 $c \geq 0$。
再由同一引理，$\sum f_j \varphi_j : I^c \to I \subset A$
是典范嵌入；这里必要时把 $c$ 换成更大的整数。
\end{proof}

\begin{lemma}
\label{algebraization-lemma-cd-one-extend}
设 $I = (f_1, \ldots, f_r)$ 是诺特环 $A$ 的理想，且
$\text{cd}(A, I) = 1$。设 $c \geq 1$ 以及 $\varphi_j : I^c \to A$，
$j = 1, \ldots, r$，如引理 \ref{algebraization-lemma-cd-one} 中所述。
则存在唯一的分次 $A$-代数映射
$$
\Phi : \bigoplus\nolimits_{n \geq 0} I^{nc} \to A[T_1, \ldots, T_r]
$$
使得 $\Phi(g) = \sum \varphi_j(g) T_j$，其中 $g \in I^c$。
此外，$\Phi$ 与映射
$A[T_1, \ldots, T_r] \to \bigoplus_{n \geq 0} I^n$,
$T_j \mapsto f_j$ 的复合是包含映射
$\bigoplus_{n \geq 0} I^{nc} \to \bigoplus_{n \geq 0} I^n$.
\end{lemma}

\begin{proof}
对每个 $j$ 与 $m \geq c$，$\varphi_j$ 在
$I^m$ 上的限制是映射 $\varphi_j : I^m \to I^{m - c}$。
给定 $j_1, \ldots, j_n \in \{1, \ldots, r\}$，我们断言复合
$$
\varphi_{j_1} \ldots \varphi_{j_n} :
I^{nc} \to I^{(n - 1)c} \to \ldots \to I^c \to A
$$
与指标 $j_1, \ldots, j_n$ 的次序无关。事实上，若
$g = g_1 \ldots g_n$ 且 $g_i \in I^c$，则可见
$$
(\varphi_{j_1} \ldots \varphi_{j_n})(g) =
\varphi_{j_1}(g_1) \ldots \varphi_{j_n}(g_n)
$$
与次序无关，因为 $A$ 中的乘法可交换。因此，可以如下定义 $\Phi$：
把 $g \in I^{nc}$ 映到
$$
\Phi(g) = \sum\nolimits_{e_1 + \ldots + e_r = n}
(\varphi_1^{e_1} \circ \ldots \circ \varphi_r^{e_r})(g)
T_1^{e_1} \ldots T_r^{e_r}
$$
直接验证可知，这是具有所需性质的分次 $A$-代数同态。
唯一性与最后一个性质都是立即可见的。这就证明了引理。
\end{proof}

\begin{lemma}
\label{algebraization-lemma-cd-one-extend-to-module}
设 $I = (f_1, \ldots, f_r)$ 是诺特环 $A$ 的理想，且
$\text{cd}(A, I) = 1$。设 $c \geq 1$ 以及 $\varphi_j : I^c \to A$，
$j = 1, \ldots, r$，如引理 \ref{algebraization-lemma-cd-one} 中所述。
设 $A \to B$ 是环映射，其中 $B$ 诺特，并设 $N$ 是有限
$B$-模。必要时增大 $c$ 并相应调整 $\varphi_j$ 之后，
存在唯一的分次 $B$-模映射
$$
\Phi_N : \bigoplus\nolimits_{n \geq 0} I^{nc}N \to N[T_1, \ldots, T_r]
$$
使得 $\Phi_N(g x) = \Phi(g) x$，其中 $g \in I^{nc}$ 且 $x \in N$，
这里 $\Phi$ 如引理
\ref{algebraization-lemma-cd-one-extend} 中所述。$\Phi_N$ 与映射
$N[T_1, \ldots, T_r] \to \bigoplus_{n \geq 0} I^nN$,
$T_j \mapsto f_j$ 的复合是包含映射
$\bigoplus_{n \geq 0} I^{nc}N \to \bigoplus_{n \geq 0} I^nN$.
\end{lemma}

\begin{proof}
由该公式以及引理 \ref{algebraization-lemma-cd-one-extend} 中 $\Phi$ 的唯一性，
唯一性是清楚的。考虑诺特 $A$-代数
$B' = B \oplus N$，其中 $N$ 是平方为零的理想。为了证明
$\Phi_N$ 存在，由引理 \ref{algebraization-lemma-cd-one}，只需证明 $\varphi_j$
可延拓为映射 $\varphi'_j : I^cB' \to B'$；这里必要时把 $c$
增大到某个 $c'$（并把 $\varphi_j$ 换成包含映射
$I^{c'} \to I^c$ 与 $\varphi_j$ 的复合）。回忆，$\varphi_j$ 对应于截面
$$
h_j \in \Gamma(\Spec(A) \setminus V(I), \mathcal{O}_{\Spec(A)})
$$
；见《概形上同调》，引理 \ref{coherent-lemma-homs-over-open}。
（事实上，在引理 \ref{algebraization-lemma-cd-one} 的证明中，我们正是这样选取
$\varphi_j$ 的。）利用同一引理，把拉回
$$
h'_j \in \Gamma(\Spec(B') \setminus V(IB'), \mathcal{O}_{\Spec(B')})
$$
（即 $h_j$ 的拉回）表示为 $B'$-线性映射
$\varphi'_j : I^{c'}B' \to B'$，其中 $c' \geq c$。
再次应用该引理可知，它将与 $\varphi_j$ 相符，只要 $c'$
充分大：事实上，只需在 $I^{c'}$
的一组有限生成元上检验二者相符。略去这个小细节。
\end{proof}

\begin{lemma}
\label{algebraization-lemma-cd-is-one-for-system}
设 $I = (f_1, \ldots, f_r)$ 是诺特环 $A$ 的理想，且
$\text{cd}(A, I) = 1$。设 $c \geq 1$ 以及 $\varphi_j : I^c \to A$，
$j = 1, \ldots, r$，如引理 \ref{algebraization-lemma-cd-one} 中所述。
设 $X$ 是 $\Spec(A)$ 上的诺特概形。设
$$
\ldots \to \mathcal{F}_3 \to \mathcal{F}_2 \to \mathcal{F}_1
$$
是凝聚 $\mathcal{O}_X$-模的一个逆系统，并且
$\mathcal{F}_n = \mathcal{F}_{n + 1}/I^n\mathcal{F}_{n + 1}$。
置 $\mathcal{F} = \lim \mathcal{F}_n$。
必要时增大 $c$ 并相应调整 $\varphi_j$ 之后，
存在唯一的分次 $\mathcal{O}_X$-模映射
$$
\Phi_\mathcal{F} :
\bigoplus\nolimits_{n \geq 0} I^{nc}\mathcal{F}
\longrightarrow
\mathcal{F}[T_1, \ldots, T_r]
$$
使得 $\Phi_\mathcal{F}(g s) = \Phi(g) s$，其中 $g \in I^{nc}$，
$s$ 是 $\mathcal{F}$ 的局部截面，而 $\Phi$ 如引理
\ref{algebraization-lemma-cd-one-extend} 中所述。$\Phi_\mathcal{F}$ 与映射
$\mathcal{F}[T_1, \ldots, T_r] \to \bigoplus_{n \geq 0} I^n\mathcal{F}$,
$T_j \mapsto f_j$
的复合是典范包含
$\bigoplus_{n \geq 0} I^{nc}\mathcal{F} \to
\bigoplus_{n \geq 0} I^n\mathcal{F}$.
\end{lemma}

\begin{proof}
由 $\mathcal{O}_X$-线性以及要求
$\Phi_\mathcal{F}(g s) = \Phi(g) s$ 对 $g \in I^{nc}$ 和
$s$ 是 $\mathcal{F}$ 的局部截面这一条件成立，唯一性立即可见。
因此可设 $X = \Spec(B)$ 是仿射的。注意 $(\mathcal{F}_n)$ 是范畴
$\textit{Coh}(X, I\mathcal{O}_X)$ 的对象；该范畴在《概形上同调》，
第 \ref{coherent-section-coherent-formal} 节中引入。
设 $B' = B^\wedge$ 是相对于 $I$ 而取的 $B$ 的完备化。
由《概形上同调》，引理 \ref{coherent-lemma-inverse-systems-affine}，
对象 $(\mathcal{F}_n)$ 对应于有限 $B'$-模 $N$；其含义是，
$\mathcal{F}_n$ 是与有限 $B$-模 $N/I^n N$ 相伴的凝聚模。
把引理 \ref{algebraization-lemma-cd-one-extend-to-module} 应用于
$I \subset A \to B'$ 与 $N$，可知必要时增大 $c$ 并相应调整
$\varphi_j$ 之后，得到唯一的映射
$$
\Phi_N : \bigoplus\nolimits_{n \geq 0} I^{nc}N \to N[T_1, \ldots, T_r]
$$
并具有相应性质。注意，在次数 $n$，得到映射
$N/I^mN \to \bigoplus_{e_1 + \ldots + e_r = n}
N/I^{m - nc}N \cdot T_1^{e_1} \ldots T_r^{e_r}$ 的逆系统，其中
$m \geq nc$。翻译回凝聚层，可知 $\Phi_N$ 对应于映射系统
$$
\Phi^n_m :
I^{nc}\mathcal{F}_m
\longrightarrow
\bigoplus\nolimits_{e_1 + \ldots + e_r = n}
\mathcal{F}_{m - nc} \cdot T_1^{e_1} \ldots T_r^{e_r}
$$
，其中 $m \geq nc$ 与 $n \geq 1$ 变化。对这些映射关于 $m$
取逆极限，得到 $\Phi_\mathcal{F} = \bigoplus_n \lim_m \Phi^n_m$。
注意 $\lim_m I^t\mathcal{F}_m = I^t \mathcal{F}$；例如在仿射开集上求值
即可看出。不过事实上并不需要这一点，因为显然有映射
$I^t\mathcal{F} \to \lim_m I^t\mathcal{F}_m$。
\end{proof}

\begin{lemma}
\label{algebraization-lemma-topology-I-adic}
设 $I$ 是诺特环 $A$ 的理想。设 $X$ 是 $\Spec(A)$ 上的诺特概形。设
$$
\ldots \to \mathcal{F}_3 \to \mathcal{F}_2 \to \mathcal{F}_1
$$
是凝聚 $\mathcal{O}_X$-模的一个逆系统，并且
$\mathcal{F}_n = \mathcal{F}_{n + 1}/I^n\mathcal{F}_{n + 1}$。
若 $\text{cd}(A, I) = 1$，则对所有 $p \in \mathbf{Z}$，
$\lim H^p(X, \mathcal{F}_n)$ 上的极限拓扑是 $I$-进拓扑。
\end{lemma}

\begin{proof}
首先，显然 $I^t \lim H^p(X, \mathcal{F}_n)$ 在
$H^p(X, \mathcal{F}_t)$ 中映为零。因此，$I$-进拓扑比极限拓扑更细。
反过来，置 $\mathcal{F} = \lim \mathcal{F}_n$，选取生成元
$f_1, \ldots, f_r$，它们生成 $I$；选取 $c \geq 1$，并选取引理
\ref{algebraization-lemma-cd-is-one-for-system} 中的 $\Phi_\mathcal{F}$。
下文将直接使用引理 \ref{algebraization-lemma-properties-system} 的结果而不再说明。
特别地，有短正合列
$$
0 \to R^1\lim H^{p - 1}(X, \mathcal{F}_n) \to H^p(X, \mathcal{F})
\to \lim H^p(X, \mathcal{F}_n) \to 0
$$
因此，可以把任意元素 $\xi$（它属于 $\lim H^p(X, \mathcal{F}_n)$）
提升为元素 $\xi' \in H^p(X, \mathcal{F})$。假设 $\xi$ 映为
$H^p(X, \mathcal{F}_{nc})$ 中的零，其中 $n$ 是某个整数；换言之，假设 $\xi$
在极限拓扑中是``小的''。有短正合列
$$
0 \to I^{nc}\mathcal{F} \to \mathcal{F} \to \mathcal{F}_{nc} \to 0
$$
因而该假设意味着可以把 $\xi'$ 提升为元素
$\xi'' \in H^p(X, I^{nc}\mathcal{F})$。应用 $\Phi_\mathcal{F}$，得到
$$
\Phi_\mathcal{F}(\xi'') = \sum\nolimits_{e_1 + \ldots + e_r = n}
\xi'_{e_1, \ldots, e_r} \cdot T_1^{e_1} \ldots T_r^{e_r}
$$
，其中某些 $\xi'_{e_1, \ldots, e_r} \in H^p(X, \mathcal{F})$。
令 $\xi_{e_1, \ldots, e_r} \in \lim H^p(X, \mathcal{F}_n)$ 为它们的像，
再利用引理 \ref{algebraization-lemma-cd-is-one-for-system} 的最后一个断言，可得
$$
\xi = \sum f_1^{e_1} \ldots f_r^{e_r} \xi_{e_1, \ldots, e_r}
$$
属于 $I^n \lim H^p(X, \mathcal{F}_n)$，如所需。
\end{proof}

\begin{example}
\label{algebraization-example-not-I-adic}
设 $k$ 是域。设 $A = k[x, y][[s, t]]/(xs - yt)$。
设 $I = (s, t)$ 且 $\mathfrak a = (x, y, s, t)$。
设 $X = \Spec(A) - V(\mathfrak a)$ 且
$\mathcal{F}_n = \mathcal{O}_X/I^n\mathcal{O}_X$。注意有理函数
$$
g = \frac{t}{x} = \frac{s}{y}
$$
在一个开邻域 $V \subset X$ 上正则，该邻域包含
$V(I\mathcal{O}_X)$。
因此，每个幂 $g^e$ 都确定一个截面
$g^e \in M = \lim H^0(X, \mathcal{F}_n)$。注意，$g^e \to 0$，
这里 $e \to \infty$，且收敛是相对于 $M$ 上的极限拓扑而言，因为 $g^e$ 在
$\mathcal{F}_e$ 中映为零。另一方面，$g^e \not \in IM$ 对任意 $e$ 成立；
读者可通过计算 $H^0(U, \mathcal{F}_n)$ 看出这一点，
此处略去计算。注意 $\text{cd}(A, I) = 2$。
因此，引理 \ref{algebraization-lemma-topology-I-adic} 的结果是最佳的。
\end{example}








\section{Mittag--Leffler 条件}
\label{algebraization-section-ML}

\noindent
对诺特局部环的极大理想取局部上同调时，我们往往自动得到
Mittag--Leffler 条件。这蕴涵：在穿孔谱上，对具有满射过渡映射的
凝聚层逆系统，其高阶上同调群也具有同一性质。

\begin{lemma}
\label{algebraization-lemma-descending-chain}
设 $(A, \mathfrak m)$ 是诺特局部环。
\begin{enumerate}
\item 设 $M$ 是有限 $A$-模。则 $A$-模
$H^i_\mathfrak m(M)$ 对任意 $i$ 都满足降链条件。
\item 设 $U = \Spec(A) \setminus \{\mathfrak m\}$ 是
$A$ 的穿孔谱。设 $\mathcal{F}$ 是凝聚 $\mathcal{O}_U$-模。
则 $A$-模 $H^i(U, \mathcal{F})$ 在 $i > 0$ 时满足降链条件。
\end{enumerate}
\end{lemma}

\begin{proof}
对 $M$ 的支撑维数作归纳来证明 (1)。当 $M = 0$ 时命题成立，
因而可以并且确实假设 $M$ 非零。

\medskip\noindent
归纳基。若 $\dim(\text{Supp}(M)) = 0$，则 $M$ 的支撑是
$\{\mathfrak m\}$，并且由局部上同调的构造可知
$H^0_\mathfrak m(M) = M$，且 $H^i_\mathfrak m(M) = 0$ 对
$i > 0$ 成立；见《对偶复形》，第
\ref{dualizing-section-local-cohomology} 节。
由于 $M$ 具有有限长度（《代数》，引理 \ref{algebra-lemma-length-finite}），
它满足降链条件。

\medskip\noindent
归纳步骤。假设 $\dim(\text{Supp}(M)) > 0$。
由归纳基，有限模 $H^0_\mathfrak m(M) \subset M$ 满足降链条件。
由《对偶复形》，引理 \ref{dualizing-lemma-divide-by-torsion}，
可以把 $M$ 换成 $M/H^0_\mathfrak m(M)$。
于是 $H^0_\mathfrak m(M) = 0$，即 $M$ 的深度 $\geq 1$；见
《对偶复形》，引理 \ref{dualizing-lemma-depth}。
选取 $x \in \mathfrak m$，使得 $x : M \to M$ 是单射。
由《代数》，引理 \ref{algebra-lemma-one-equation-module}，有
$\dim(\text{Supp}(M/xM)) = \dim(\text{Supp}(M)) - 1$，故可应用归纳假设。
选取指标 $i$ 并考虑正合列
$$
H^{i - 1}_\mathfrak m(M/xM) \to H^i_\mathfrak m(M) \xrightarrow{x}
H^i_\mathfrak m(M)
$$
，它来自短正合列 $0 \to M \xrightarrow{x} M \to M/xM \to 0$。
由此，$x$-挠子模 $H^i_\mathfrak m(M)[x]$ 是一个满足降链条件的模的商，
因而自身也满足降链条件。于是 $\mathfrak m$-挠子模
$H^i_\mathfrak m(M)[\mathfrak m]$ 满足降链条件（因而在
$A/\mathfrak m$ 上有限维）。所以由《对偶复形》，引理
\ref{dualizing-lemma-describe-categories}.
可知 $\mathfrak m$-幂挠模 $H^i_\mathfrak m(M)$ 满足降链条件。

\medskip\noindent
(2) 由 (1) 与《局部上同调》，引理
\ref{local-cohomology-lemma-finiteness-pushforwards-and-H1-local} 得出。
\end{proof}

\begin{lemma}
\label{algebraization-lemma-ML-local}
设 $(A, \mathfrak m)$ 是诺特局部环。
\begin{enumerate}
\item 设 $(M_n)$ 是有限 $A$-模的一个逆系统。则逆系统
$H^i_\mathfrak m(M_n)$ 对任意 $i$ 都满足 Mittag--Leffler 条件。
\item 设 $U = \Spec(A) \setminus \{\mathfrak m\}$ 是
$A$ 的穿孔谱。设 $\mathcal{F}_n$ 是凝聚 $\mathcal{O}_U$-模的一个逆系统。
则逆系统 $H^i(U, \mathcal{F}_n)$ 在 $i > 0$ 时
满足 Mittag--Leffler 条件。
\end{enumerate}
\end{lemma}

\begin{proof}
由引理 \ref{algebraization-lemma-descending-chain} 立即得出。
\end{proof}

\begin{lemma}
\label{algebraization-lemma-terrific}
设 $(A, \mathfrak m)$ 是诺特局部环。
设 $(M_n)$ 是有限 $A$-模的一个逆系统。
设 $M \to \lim M_n$ 是映射，其中 $M$ 是有限 $A$-模，
并且对某个 $i$，映射
$H^i_\mathfrak m(M) \to \lim H^i_\mathfrak m(M_n)$
是同构。则逆系统 $H^i_\mathfrak m(M_n)$
本质常值，其值为 $H^i_\mathfrak m(M)$。
\end{lemma}

\begin{proof}
由引理 \ref{algebraization-lemma-ML-local}，逆系统 $H^i_\mathfrak m(M_n)$
满足 Mittag--Leffler 条件。设 $E_n \subset H^i_\mathfrak m(M_n)$
是 $H^i_\mathfrak m(M_{n'})$ 在 $n' \gg n$ 时的像。
则 $(E_n)$ 是具有满射过渡映射的逆系统，并且
$H^i_\mathfrak m(M) = \lim E_n$。由引理 \ref{algebraization-lemma-descending-chain}，
$H^i_\mathfrak m(M)$ 满足降链条件，故满射
$H^i_\mathfrak m(M) \to E_n$ 中只有有限多个具有非平凡核。
因此，$E_n \to E_{n - 1}$ 对所有 $n \gg 0$ 都是同构，如所需。
\end{proof}

\begin{lemma}
\label{algebraization-lemma-local-cohomology-derived-completion}
设 $(A, \mathfrak m)$ 是诺特局部环。
设 $I \subset A$ 是理想。设 $M$ 是有限 $A$-模。则
$$
H^i(R\Gamma_\mathfrak m(M)^\wedge) = \lim H^i_\mathfrak m(M/I^nM)
$$
对所有 $i$ 成立，其中 $R\Gamma_\mathfrak m(M)^\wedge$ 表示导出
$I$-进完备化。
\end{lemma}

\begin{proof}
应用《对偶复形》，引理 \ref{dualizing-lemma-completion-local}
与引理 \ref{algebraization-lemma-ML-local}，可知各 $R^1\lim$ 项消失。
\end{proof}








\section{环化位点上的导出完备化}
\label{algebraization-section-derived-completion}

\noindent
我们强烈建议读者初次阅读时跳过本节。

\medskip\noindent
这些内容的代数版本见《代数进阶》，第
\ref{more-algebra-section-derived-completion} 节。
设 $\mathcal{O}$ 是位点 $\mathcal{C}$ 上的环层。
设 $f$ 是 $\mathcal{O}$ 的整体截面。用 $\mathcal{O}_f$ 表示与局部化预层
$U \mapsto \mathcal{O}(U)_f$ 相伴的层。

\begin{lemma}
\label{algebraization-lemma-map-twice-localize}
设 $(\mathcal{C}, \mathcal{O})$ 是环化位点。设 $f$ 是 $\mathcal{O}$ 的整体截面。
\begin{enumerate}
\item 对 $L, N \in D(\mathcal{O}_f)$，有
$R\SheafHom_\mathcal{O}(L, N) = R\SheafHom_{\mathcal{O}_f}(L, N)$.
特别地，两个 $\mathcal{O}_f$-结构在
$R\SheafHom_\mathcal{O}(L, N)$ 上相同。
\item 对 $K \in D(\mathcal{O})$ 与
$L \in D(\mathcal{O}_f)$，有
$$
R\SheafHom_\mathcal{O}(L, K) =
R\SheafHom_{\mathcal{O}_f}(L, R\SheafHom_\mathcal{O}(\mathcal{O}_f, K))
$$
特别地，
$R\SheafHom_\mathcal{O}(\mathcal{O}_f,
R\SheafHom_\mathcal{O}(\mathcal{O}_f, K)) =
R\SheafHom_\mathcal{O}(\mathcal{O}_f, K)$.
\item 若 $g$ 是 $\mathcal{O}$ 的另一个整体截面，则
$$
R\SheafHom_\mathcal{O}(\mathcal{O}_f, R\SheafHom_\mathcal{O}(\mathcal{O}_g, K))
= R\SheafHom_\mathcal{O}(\mathcal{O}_{gf}, K).
$$
\end{enumerate}
\end{lemma}

\begin{proof}
(1) 的证明。设 $\mathcal{J}^\bullet$ 是 $\mathcal{O}_f$-模 K-内射复形，
并且表示 $N$。由《位点上的上同调》，引理
\ref{sites-cohomology-lemma-K-injective-flat}，可知 $\mathcal{J}^\bullet$
也是 $\mathcal{O}$-模 K-内射复形。设 $\mathcal{F}^\bullet$ 是
$\mathcal{O}_f$-模复形，并且表示 $L$。则
$$
R\SheafHom_\mathcal{O}(L, N) =
R\SheafHom_\mathcal{O}(\mathcal{F}^\bullet, \mathcal{J}^\bullet) =
R\SheafHom_{\mathcal{O}_f}(\mathcal{F}^\bullet, \mathcal{J}^\bullet)
$$
；这是由《位点上的模》，引理
\ref{sites-modules-lemma-epimorphism-modules} 得出的，因为
$\mathcal{J}^\bullet$ 既是 $\mathcal{O}$-模 K-内射复形，
也是 $\mathcal{O}_f$-模 K-内射复形。

\medskip\noindent
(2) 的证明。设 $\mathcal{I}^\bullet$ 是 $\mathcal{O}$-模 K-内射复形，
并且表示 $K$。于是
$R\SheafHom_\mathcal{O}(\mathcal{O}_f, K)$ 由
$\SheafHom_\mathcal{O}(\mathcal{O}_f, \mathcal{I}^\bullet)$ 表示；由
《位点上的上同调》，引理 \ref{sites-cohomology-lemma-hom-K-injective} 与
\ref{sites-cohomology-lemma-K-injective-flat}，后者既是
$\mathcal{O}_f$-模 K-内射复形，也是 $\mathcal{O}$-模 K-内射复形。
设 $\mathcal{F}^\bullet$ 是 $\mathcal{O}_f$-模复形，并且表示 $L$。则
$$
R\SheafHom_\mathcal{O}(L, K) =
R\SheafHom_\mathcal{O}(\mathcal{F}^\bullet, \mathcal{I}^\bullet) =
R\SheafHom_{\mathcal{O}_f}(\mathcal{F}^\bullet,
\SheafHom_\mathcal{O}(\mathcal{O}_f, \mathcal{I}^\bullet))
$$
；这是由《位点上的模》，引理
\ref{sites-modules-lemma-adjoint-hom-restrict} 以及
$\SheafHom_\mathcal{O}(\mathcal{O}_f, \mathcal{I}^\bullet)$ 是
$\mathcal{O}_f$-模 K-内射复形得出的。

\medskip\noindent
(3) 的证明。这由以下两个事实得出：
$R\SheafHom_\mathcal{O}(\mathcal{O}_g, \mathcal{I}^\bullet)$
作为 $\mathcal{O}$-模复形是 K-内射的；并且
$\SheafHom_\mathcal{O}(\mathcal{O}_f,
\SheafHom_\mathcal{O}(\mathcal{O}_g, \mathcal{H})) = 
\SheafHom_\mathcal{O}(\mathcal{O}_{gf}, \mathcal{H})$
对所有 $\mathcal{O}$-模层 $\mathcal{H}$ 成立。
\end{proof}

\noindent
设 $K \in D(\mathcal{O})$。用 $T(K, f)$ 表示逆系统
$$
\ldots \to K \xrightarrow{f} K \xrightarrow{f} K
$$
在 $D(\mathcal{O})$ 中的一个导出极限（《导出范畴》，定义
\ref{derived-definition-derived-limit}）。

\begin{lemma}
\label{algebraization-lemma-hom-from-Af}
设 $(\mathcal{C}, \mathcal{O})$ 是环化位点。设 $f$ 是 $\mathcal{O}$
的整体截面。设 $K \in D(\mathcal{O})$。下列条件等价：
\begin{enumerate}
\item $R\SheafHom_\mathcal{O}(\mathcal{O}_f, K) = 0$,
\item $R\SheafHom_\mathcal{O}(L, K) = 0$ 对所有 $L$ 成立，
其中该对象位于 $D(\mathcal{O}_f)$ 中，
\item $T(K, f) = 0$.
\end{enumerate}
\end{lemma}

\begin{proof}
显然 (2) 蕴涵 (1)。蕴涵 (1) $\Rightarrow$ (2) 由引理
\ref{algebraization-lemma-map-twice-localize} 得出。$\mathcal{O}$-模
$\mathcal{O}_f$ 的一个自由解消为
$$
0 \to \bigoplus\nolimits_{n \in \mathbf{N}} \mathcal{O} \to
\bigoplus\nolimits_{n \in \mathbf{N}} \mathcal{O}
\to \mathcal{O}_f \to 0
$$
，其中第一个映射把局部截面 $(x_0, x_1, \ldots)$ 映到
$(x_0, x_1 - fx_0, x_2 - fx_1, \ldots)$，第二个映射把
$(x_0, x_1, \ldots)$ 映到 $x_0 + x_1/f + x_2/f^2 + \ldots$。
应用 $\SheafHom_\mathcal{O}(-, \mathcal{I}^\bullet)$，其中
$\mathcal{I}^\bullet$ 是 $\mathcal{O}$-模 K-内射复形，并且表示 $K$，
得到复形的短正合列
$$
0 \to \SheafHom_\mathcal{O}(\mathcal{O}_f, \mathcal{I}^\bullet) \to
\prod \mathcal{I}^\bullet \to \prod \mathcal{I}^\bullet \to 0
$$
，因为 $\mathcal{I}^n$ 是内射 $\mathcal{O}$-模。这些积是
$D(\mathcal{O})$ 中的积；见《内射对象》，引理
\ref{injectives-lemma-derived-products}。这意味着对象 $T(K, f)$
表示 $R\SheafHom_\mathcal{O}(\mathcal{O}_f, K)$，并且是在
$D(\mathcal{O})$ 中表示。
因此 (1) 与 (3) 等价。
\end{proof}

\begin{lemma}
\label{algebraization-lemma-ideal-of-elements-complete-wrt}
设 $(\mathcal{C}, \mathcal{O})$ 是环化位点。设 $K \in D(\mathcal{O})$。
对 $U$，令 $\mathcal{I}(U)$ 为截面 $f \in \mathcal{O}(U)$ 所成的集合，
这些截面满足 $T(K|_U, f) = 0$；这一规则给出 $\mathcal{O}$ 中的理想层。
\end{lemma}

\begin{proof}
下文直接使用引理 \ref{algebraization-lemma-hom-from-Af} 的结果而不再说明。
若 $f \in \mathcal{I}(U)$ 且 $g \in \mathcal{O}(U)$，则
$\mathcal{O}_{U, gf}$ 是 $\mathcal{O}_{U, f}$-模，因而
$R\SheafHom_\mathcal{O}(\mathcal{O}_{U, gf}, K|_U) = 0$，故
$gf \in \mathcal{I}(U)$。设 $f, g \in \mathcal{O}(U)$。则有短正合列
$$
0 \to \mathcal{O}_{U, f + g} \to
\mathcal{O}_{U, f(f + g)} \oplus \mathcal{O}_{U, g(f + g)} \to
\mathcal{O}_{U, gf(f + g)} \to 0
$$
，因为 $f, g$ 在 $\mathcal{O}(U)_{f + g}$ 中生成单位理想。
这由《代数》，引理 \ref{algebra-lemma-standard-covering}
以及最后一个箭头是满射这一简单事实得出。
因为 $R\SheafHom_\mathcal{O}( - , K|_U)$ 是三角范畴之间的正合函子，
下列三个对象的消失：
$R\SheafHom_{\mathcal{O}_U}(\mathcal{O}_{U, f(f + g)}, K|_U)$,
$R\SheafHom_{\mathcal{O}_U}(\mathcal{O}_{U, g(f + g)}, K|_U)$，以及
$R\SheafHom_{\mathcal{O}_U}(\mathcal{O}_{U, gf(f + g)}, K|_U)$,
蕴涵
$R\SheafHom_{\mathcal{O}_U}(\mathcal{O}_{U, f + g}, K|_U)$.
层条件的验证从略。
\end{proof}

\noindent
对任意环化位点，都可以作如下定义。

\begin{definition}
\label{algebraization-definition-derived-complete}
设 $(\mathcal{C}, \mathcal{O})$ 是环化位点。
设 $\mathcal{I} \subset \mathcal{O}$ 是理想层。
设 $K \in D(\mathcal{O})$。若 $K$ {\it 关于 $\mathcal{I}$ 导出完备}，
是指对每个对象 $U$（属于 $\mathcal{C}$）以及 $f \in \mathcal{I}(U)$，
对象 $T(K|_U, f)$ 在 $D(\mathcal{O}_U)$ 中都为零。
\end{definition}

\noindent
显然，由导出完备对象组成的全子范畴
$D_{comp}(\mathcal{O}) = D_{comp}(\mathcal{O}, \mathcal{I}) \subset
D(\mathcal{O})$
是饱和三角子范畴；见《导出范畴》，定义
\ref{derived-definition-triangulated-subcategory} 与
\ref{derived-definition-saturated}。该子范畴在 $D(\mathcal{O})$ 中对积与
同伦极限封闭，但一般不对可数直和封闭。

\begin{lemma}
\label{algebraization-lemma-derived-complete-internal-hom}
设 $(\mathcal{C}, \mathcal{O})$ 是环化位点。
设 $\mathcal{I} \subset \mathcal{O}$ 是理想层。
若 $K \in D(\mathcal{O})$ 且 $L \in D_{comp}(\mathcal{O})$，则
$R\SheafHom_\mathcal{O}(K, L) \in D_{comp}(\mathcal{O})$.
\end{lemma}

\begin{proof}
设 $U$ 是 $\mathcal{C}$ 的对象，并设 $f \in \mathcal{I}(U)$。回忆，
$$
\Hom_{D(\mathcal{O}_U)}(\mathcal{O}_{U, f}, R\SheafHom_\mathcal{O}(K, L)|_U)
=
\Hom_{D(\mathcal{O}_U)}(
K|_U \otimes_{\mathcal{O}_U}^\mathbf{L} \mathcal{O}_{U, f}, L|_U)
$$
；这是《位点上的上同调》，引理
\ref{sites-cohomology-lemma-internal-hom}。右边由引理
\ref{algebraization-lemma-hom-from-Af} 以及内部 Hom 与实际 Hom 之间的关系而为零；见
《位点上的上同调》，引理
\ref{sites-cohomology-lemma-section-RHom-over-U}。
对所有 $U'/U$，同样的消失成立。因此，对象
$R\SheafHom_{\mathcal{O}_U}(\mathcal{O}_{U, f},
R\SheafHom_\mathcal{O}(K, L)|_U)$ 作为 $D(\mathcal{O}_U)$ 的对象，
其第 $0$ 个上同调层消失（由上述引文）。其他上同调层亦然；也就是说，
$R\SheafHom_{\mathcal{O}_U}(\mathcal{O}_{U, f},
R\SheafHom_\mathcal{O}(K, L)|_U)$ 在 $D(\mathcal{O}_U)$ 中为零。
由引理 \ref{algebraization-lemma-hom-from-Af} 即得结论。
\end{proof}

\begin{lemma}
\label{algebraization-lemma-restriction-derived-complete}
设 $\mathcal{C}$ 是位点。设 $\mathcal{O} \to \mathcal{O}'$
是环层同态。设 $\mathcal{I} \subset \mathcal{O}$ 是理想层。
$D_{comp}(\mathcal{O}, \mathcal{I})$ 在限制函子
$D(\mathcal{O}') \to D(\mathcal{O})$ 下的原像是
$D_{comp}(\mathcal{O}', \mathcal{I}\mathcal{O}')$.
\end{lemma}

\begin{proof}
利用引理 \ref{algebraization-lemma-ideal-of-elements-complete-wrt} 可知，
$K' \in D(\mathcal{O}')$ 属于
$D_{comp}(\mathcal{O}', \mathcal{I}\mathcal{O}')$
当且仅当 $T(K'|_U, f)$ 对每个局部截面 $f \in \mathcal{I}(U)$ 都为零。
注意，$T(K'|_U, f)$ 的上同调层是在阿贝尔层范畴中计算的，
所以把 $f$ 看作 $\mathcal{O}$ 的截面，或者取 $f$ 在
$\mathcal{O}'$ 中的像，并无区别。由此及导出完备对象的定义，立即得出引理。
\end{proof}

\begin{lemma}
\label{algebraization-lemma-pushforward-derived-complete}
设 $f : (\Sh(\mathcal{D}), \mathcal{O}') \to (\Sh(\mathcal{C}), \mathcal{O})$
是环化拓扑斯态射。设 $\mathcal{I} \subset \mathcal{O}$ 与
$\mathcal{I}' \subset \mathcal{O}'$ 是理想层，并且 $f^\sharp$ 把
$f^{-1}\mathcal{I}$ 映入 $\mathcal{I}'$。则 $Rf_*$ 把
$D_{comp}(\mathcal{O}', \mathcal{I}')$ 映入
$D_{comp}(\mathcal{O}, \mathcal{I})$。
\end{lemma}

\begin{proof}
可以假设 $f$ 由环化位点态射给出，该态射对应于连续函子
$\mathcal{C} \to \mathcal{D}$（《位点上的模》，引理
\ref{sites-modules-lemma-morphism-ringed-topoi-comes-from-morphism-ringed-sites}
)。设 $U$ 是 $\mathcal{C}$ 的对象，并设 $g$ 是 $\mathcal{I}$ 在
$U$ 上的截面。必须证明
$\Hom_{D(\mathcal{O}_U)}(\mathcal{O}_{U, g}, Rf_*K|_U) = 0$
，只要 $K$ 关于 $\mathcal{I}'$ 导出完备。事实上，由《位点上的上同调》，引理
\ref{sites-cohomology-lemma-section-RHom-over-U}
，把这一点应用于 $U$ 上的所有对象及 $K$ 的所有平移，可推出
$R\SheafHom_{\mathcal{O}_U}(\mathcal{O}_{U, g}, Rf_*K|_U)$ 为零；
这又推出 $T(Rf_*K|_U, g)$ 为零（引理
\ref{algebraization-lemma-hom-from-Af}），而这正是需要证明的（定义
\ref{algebraization-definition-derived-complete}）。
设 $V$ 是 $\mathcal{D}$ 中 $U$ 的像。则
$$
\Hom_{D(\mathcal{O}_U)}(\mathcal{O}_{U, g}, Rf_*K|_U) =
\Hom_{D(\mathcal{O}'_V)}(\mathcal{O}'_{V, g'}, K|_V) = 0
$$
，其中 $g' = f^\sharp(g) \in \mathcal{I}'(V)$。第二个等号成立是因为
$K$ 导出完备；第一个等号成立是因为 $\mathcal{O}_{U, g}$ 的导出拉回是
$\mathcal{O}'_{V, g'}$，并应用《位点上的上同调》，引理
\ref{sites-cohomology-lemma-adjoint}。
\end{proof}

\noindent
下一个引理是存在导出完备化的最简单情形。

\begin{lemma}
\label{algebraization-lemma-derived-completion}
设 $(\mathcal{C}, \mathcal{O})$ 是环化位点。设 $f_1, \ldots, f_r$
是 $\mathcal{O}$ 的整体截面。设 $\mathcal{I} \subset \mathcal{O}$ 是由
$f_1, \ldots, f_r$ 生成的理想层。则包含函子
$D_{comp}(\mathcal{O}) \to D(\mathcal{O})$ 有左伴随；也就是说，
给定任意对象 $K$（它属于 $D(\mathcal{O})$），存在映射
$K \to K^\wedge$，其中 $K^\wedge$ 位于 $D_{comp}(\mathcal{O})$ 中，
使得映射
$$
\Hom_{D(\mathcal{O})}(K^\wedge, E) \longrightarrow \Hom_{D(\mathcal{O})}(K, E)
$$
在 $E$ 位于 $D_{comp}(\mathcal{O})$ 中时是双射。事实上，有
$$
K^\wedge =
R\SheafHom_\mathcal{O}
(\mathcal{O} \to \prod\nolimits_{i_0} \mathcal{O}_{f_{i_0}} \to
\prod\nolimits_{i_0 < i_1} \mathcal{O}_{f_{i_0}f_{i_1}} \to
\ldots \to \mathcal{O}_{f_1\ldots f_r}, K)
$$
，并且关于 $K$ 具有函子性。
\end{lemma}

\begin{proof}
用引理中最后一个展示公式定义 $K^\wedge$。有复形映射
$$
(\mathcal{O} \to \prod\nolimits_{i_0} \mathcal{O}_{f_{i_0}} \to
\prod\nolimits_{i_0 < i_1} \mathcal{O}_{f_{i_0}f_{i_1}} \to
\ldots \to \mathcal{O}_{f_1\ldots f_r}) \longrightarrow \mathcal{O}
$$
，它诱导映射 $K \to K^\wedge$。只需证明 $K^\wedge$ 导出完备，
并且 $K \to K^\wedge$ 在 $K$ 导出完备时是同构。

\medskip\noindent
设 $f$ 是 $\mathcal{O}$ 的整体截面。由引理
\ref{algebraization-lemma-map-twice-localize}，对象
$R\SheafHom_\mathcal{O}(\mathcal{O}_f, K^\wedge)$
等于
$$
R\SheafHom_\mathcal{O}(
(\mathcal{O}_f \to \prod\nolimits_{i_0} \mathcal{O}_{ff_{i_0}} \to
\prod\nolimits_{i_0 < i_1} \mathcal{O}_{ff_{i_0}f_{i_1}} \to
\ldots \to \mathcal{O}_{ff_1\ldots f_r}), K)
$$
若 $f = f_i$ 对某个 $i$ 成立，则 $f_1, \ldots, f_r$ 在
$\mathcal{O}_f$ 中生成单位理想，因而增广交错 {\v C}ech 复形
$$
\mathcal{O}_f \to \prod\nolimits_{i_0} \mathcal{O}_{ff_{i_0}} \to
\prod\nolimits_{i_0 < i_1} \mathcal{O}_{ff_{i_0}f_{i_1}} \to
\ldots \to \mathcal{O}_{ff_1\ldots f_r}
$$
为零（甚至同伦于零）。由此可见 $K^\wedge$ 导出完备。

\medskip\noindent
若 $K$ 导出完备，则 $R\SheafHom_\mathcal{O}(\mathcal{O}_f, K)$
对所有 $f = f_{i_0} \ldots f_{i_p}$、$p \geq 0$ 都为零。因此，
$K \to K^\wedge$ 在 $D(\mathcal{O})$ 中是同构。
\end{proof}

\noindent
下面说明导出完备化为何确实是一种完备化。

\begin{lemma}
\label{algebraization-lemma-derived-completion-koszul}
设 $(\mathcal{C}, \mathcal{O})$ 是环化位点。设 $f_1, \ldots, f_r$
是 $\mathcal{O}$ 的整体截面。设 $\mathcal{I} \subset \mathcal{O}$ 是由
$f_1, \ldots, f_r$ 生成的理想层。设 $K \in D(\mathcal{O})$。
引理 \ref{algebraization-lemma-derived-completion} 的导出完备化 $K^\wedge$ 由公式
$$
K^\wedge = R\lim K \otimes^\mathbf{L}_\mathcal{O} K_n
$$
给出，其中 $K_n = K(\mathcal{O}, f_1^n, \ldots, f_r^n)$ 是
关于 $f_1^n, \ldots, f_r^n$ 的 Koszul 复形，定义在 $\mathcal{O}$ 上。
\end{lemma}

\begin{proof}
在《代数进阶》，引理
\ref{more-algebra-lemma-extended-alternating-Cech-is-colimit-koszul}
中已经看到，增广交错 {\v C}ech 复形
$$
\mathcal{O} \to \prod\nolimits_{i_0} \mathcal{O}_{f_{i_0}} \to
\prod\nolimits_{i_0 < i_1} \mathcal{O}_{f_{i_0}f_{i_1}} \to
\ldots \to \mathcal{O}_{f_1\ldots f_r}
$$
是 Koszul 复形 $K^n = K(\mathcal{O}, f_1^n, \ldots, f_r^n)$ 的余极限，
这些复形位于次数 $0, \ldots, r$。注意，$K^n$ 是有限自由
$\mathcal{O}$-模的有限链复形，其对偶为
$\SheafHom_\mathcal{O}(K^n, \mathcal{O}) = K_n$，其中 $K_n$ 是位于次数
$-r, \ldots, 0$ 的 Koszul 上链复形（如通常约定）。由引理
\ref{algebraization-lemma-derived-completion}，函子 $E \mapsto E^\wedge$ 是从增广交错
{\v C}ech 复形取 $R\SheafHom$ 而得到的，其目标为 $E$：
$$
E^\wedge = R\SheafHom(\colim K^n, E)
$$
由《位点上的上同调》，引理
\ref{sites-cohomology-lemma-colim-and-lim-of-duals}，这等于
$R\lim (E \otimes_\mathcal{O}^\mathbf{L} K_n)$。
\end{proof}

\begin{lemma}
\label{algebraization-lemma-all-rings}
存在一种如下的构造方式：
\begin{enumerate}
\item 对每个二元组 $(A, I)$，其中 $A$ 是环而 $I \subset A$ 是有限生成理想，
构造复形 $K(A, I)$，它由 $A$-模组成；
\item 构造映射 $K(A, I) \to A$，它是 $A$-模复形之间的映射；
\item 对每个环映射 $A \to B$ 与有限生成理想 $I \subset A$，
构造复形映射 $K(A, I) \to K(B, IB)$；
\end{enumerate}
并满足：
\begin{enumerate}
\item[(a)] 对 $A \to B$ 与有限生成的 $I \subset A$，图
$$
\xymatrix{
K(A, I) \ar[r] \ar[d] & A \ar[d] \\
K(B, IB) \ar[r] & B
}
$$
交换；
\item[(b)] 对 $A \to B \to C$ 与有限生成的 $I \subset A$，映射的复合
$K(A, I) \to K(B, IB) \to K(C, IC)$ 就是映射
$K(A, I) \to K(C, IC)$；
\item[(c)] 对 $A \to B$ 与有限生成理想 $I \subset A$，诱导映射
$K(A, I) \otimes_A^\mathbf{L} B \to K(B, IB)$
在 $D(B)$ 中是同构；并且
\item[(d)] 若 $I = (f_1, \ldots, f_r) \subset A$，则有交换图
$$
\xymatrix{
(A \to \prod\nolimits_{i_0} A_{f_{i_0}} \to
\prod\nolimits_{i_0 < i_1} A_{f_{i_0}f_{i_1}} \to
\ldots \to A_{f_1\ldots f_r}) \ar[r] \ar[d] &  K(A, I) \ar[d] \\
A \ar[r]^1 & A
}
$$
，它位于 $D(A)$ 中，且横向箭头都是同构。
\end{enumerate}
\end{lemma}

\begin{proof}
设 $S$ 是形如 $A_0$ 的环所成的集合，其中
$A_0 = \mathbf{Z}[x_1, \ldots, x_n]/J$。
每个有限型 $\mathbf{Z}$-代数都同构于 $S$ 的一个元素。
设 $\mathcal{A}_0$ 为如下范畴：其对象是二元组 $(A_0, I_0)$，
其中 $A_0 \in S$ 且 $I_0 \subset A_0$ 是理想；从
$(A_0, I_0) \to (B_0, J_0)$ 的态射是环映射
$\varphi : A_0 \to B_0$，它满足 $J_0 = \varphi(I_0)B_0$。

\medskip\noindent
假设能以函子方式构造 $K(A_0, I_0) \to A_0$，这是对
$\mathcal{A}_0$ 的对象而言，并满足 (a)、(b)、(c)、(d)。此时取
$$
K(A, I) = \colim_{\varphi : (A_0, I_0) \to (A, I)} K(A_0, I_0)
$$
，其中余极限遍历环映射 $\varphi : A_0 \to A$，它们满足
$\varphi(I_0)A = I$，且 $(A_0, I_0)$ 位于 $\mathcal{A}_0$ 中。
在 $(A_0, I_0) \to (A, I)$ 与
$(A_0', I_0') \to (A, I)$ 之间的态射，由映射
$(A_0, I_0) \to (A_0', I_0')$ 给出；该映射位于 $\mathcal{A}_0$ 中，
并与到 $A$ 的映射可交换。
这些 $(A_0, I_0) \to (A, I)$ 所成的范畴是滤过的（略去细节）。
此外，$\colim_{\varphi : (A_0, I_0) \to (A, I)} A_0 = A$，
所以 $K(A, I)$ 是 $A$-模复形。
最后，给定 $\varphi : A \to B$ 与 $I \subset A$，对余极限中的每个
$(A_0, I_0) \to (A, I)$，复合
$(A_0, I_0) \to (B, IB)$ 位于 $(B, IB)$ 所对应的余极限中。
这样便得到余极限之间的映射。性质 (a)、(b)、(c)、(d) 立即由此及复形
$K(A_0, I_0)$ 的相应性质得出。

\medskip\noindent
在 $\mathcal{C}_0 = \mathcal{A}_0^{opp}$ 上赋予混沌拓扑。
在 $\mathcal{C}_0$ 上赋予环层
$\mathcal{O} : (A, I) \mapsto A$。各理想 $I$ 拼合成理想层
$\mathcal{I} \subset \mathcal{O}$。选取内射解消
$\mathcal{O} \to \mathcal{J}^\bullet$。考虑对象
$$
\mathcal{F}^\bullet = \bigcup\nolimits_n \mathcal{J}^\bullet[\mathcal{I}^n]
$$
设 $U = (A, I) \in \Ob(\mathcal{C}_0)$。
由于 $\mathcal{C}_0$ 上的拓扑是混沌的，$\mathcal{J}^\bullet(U)$
是由内射 $A$-模组成的 $A$ 的解消。因此，$\mathcal{F}^\bullet(U)$
是 $D(A)$ 的对象，表示 $R\Gamma_I(A)$ 在 $D(A)$ 中的像；见
《对偶复形》，第 \ref{dualizing-section-local-cohomology} 节。
选取 $\mathcal{O}$-模复形 $\mathcal{K}^\bullet$ 以及交换图
$$
\xymatrix{
\mathcal{O} \ar[r] & \mathcal{J}^\bullet \\
\mathcal{K}^\bullet \ar[r] \ar[u] & \mathcal{F}^\bullet \ar[u]
}
$$
，其中横向箭头是拟同构。由导出范畴 $D(\mathcal{O})$ 的构造，
这可以做到。置 $K(A, I) = \mathcal{K}^\bullet(U)$，其中
$U = (A, I)$。性质 (a)、(b) 是清楚的，性质 (c)、(d) 由
《对偶复形》，引理
\ref{dualizing-lemma-compute-local-cohomology-noetherian} 与
\ref{dualizing-lemma-local-cohomology-change-rings}.
\end{proof}

\begin{lemma}
\label{algebraization-lemma-global-extended-cech-complex}
设 $(\mathcal{C}, \mathcal{O})$ 是环化位点。设
$\mathcal{I} \subset \mathcal{O}$ 是有限型理想层。
存在映射 $K \to \mathcal{O}$，它位于 $D(\mathcal{O})$ 中，使得对每个
$U \in \Ob(\mathcal{C})$，若 $\mathcal{I}|_U$ 由
$f_1, \ldots, f_r \in \mathcal{I}(U)$ 生成，则有同构
$$
(\mathcal{O}_U \to \prod\nolimits_{i_0} \mathcal{O}_{U, f_{i_0}} \to
\prod\nolimits_{i_0 < i_1} \mathcal{O}_{U, f_{i_0}f_{i_1}} \to
\ldots \to \mathcal{O}_{U, f_1\ldots f_r}) \longrightarrow K|_U
$$
，且与到 $\mathcal{O}_U$ 的映射相容。
\end{lemma}

\begin{proof}
设 $\mathcal{C}' \subset \mathcal{C}$ 是由如下对象 $U$ 组成的全子范畴：
$\mathcal{I}|_U$ 由有限多个截面生成。则 $\mathcal{C}' \to \mathcal{C}$
是特殊余连续函子（《位点》，定义
\ref{sites-definition-special-cocontinuous-functor}）。
因此，只需在 $\mathcal{C}'$ 上工作；见《位点》，引理
\ref{sites-lemma-equivalence}。换言之，可以假设对每个对象
$U$（属于 $\mathcal{C}$），存在有限生成理想 $I \subset \mathcal{I}(U)$，使得
$\mathcal{I}|_U = \Im(I \otimes \mathcal{O}_U \to \mathcal{O}_U)$。
我们称 $I$ 生成 $\mathcal{I}|_U$。
警告：并不知道 $\mathcal{I}(U)$ 是 $\mathcal{O}(U)$ 中的有限生成理想。

\medskip\noindent
设 $U$ 是对象，并设 $I \subset \mathcal{O}(U)$ 是生成
$\mathcal{I}|_U$ 的有限生成理想。在范畴 $\mathcal{C}/U$ 上考虑预层复形
$$
K_{U, I}^\bullet : U'/U \longmapsto K(\mathcal{O}(U'), I\mathcal{O}(U'))
$$
，其中 $K(-, -)$ 如引理 \ref{algebraization-lemma-all-rings} 中所述。
我们断言其层化与 $I$ 的选择无关。事实上，若
$I' \subset \mathcal{O}(U)$ 是另一个也生成 $\mathcal{I}|_U$ 的有限生成理想，
则存在覆盖 $\{U_j \to U\}$，使得
$I\mathcal{O}(U_j) = I'\mathcal{O}(U_j)$。（提示：这是因为 $I$ 与 $I'$
均有限生成且都生成 $\mathcal{I}|_U$。）因此，
$K_{U, I}^\bullet$ 与 $K_{U, I'}^\bullet$ 在任意对象上
{\it 相同}，只要该对象位于某个 $U_j$ 之上。关于层化的断言随即得出。
用 $K_U^\bullet$ 表示这一共同值。

\medskip\noindent
对 $I$ 的选择无关还表明
$K_U^\bullet|_{\mathcal{C}/U'} = K_{U'}^\bullet$
，只要给定态射 $U' \to U$，从而给定局部化态射
$\mathcal{C}/U' \to \mathcal{C}/U$。因此，各复形 $K_U^\bullet$
可粘合成一个良定义的复形 $K^\bullet$，它由 $\mathcal{O}$-模组成。
映射 $K^\bullet \to \mathcal{O}$ 的存在性以及引理中的拟同构，
立即由引理 \ref{algebraization-lemma-all-rings} 中复形 $K(-, -)$ 的相应性质得出。
\end{proof}

\begin{proposition}
\label{algebraization-proposition-derived-completion}
设 $(\mathcal{C}, \mathcal{O})$ 是环化位点。
设 $\mathcal{I} \subset \mathcal{O}$ 是有限型理想层。
包含函子 $D_{comp}(\mathcal{O}) \to D(\mathcal{O})$ 有左伴随。
\end{proposition}

\begin{proof}
设 $K \to \mathcal{O}$ 位于 $D(\mathcal{O})$ 中，并如引理
\ref{algebraization-lemma-global-extended-cech-complex} 中所构造。设 $E \in D(\mathcal{O})$。
那么，$E^\wedge = R\SheafHom(K, E)$ 连同映射 $E \to E^\wedge$
即可满足要求。事实上，在位点 $\mathcal{C}$ 上局部地，
这恢复了引理 \ref{algebraization-lemma-derived-completion} 的伴随。
由此可知，$E^\wedge$ 总是导出完备的，并且
$E \to E^\wedge$ 在 $E$ 导出完备时是同构。
\end{proof}

\begin{remark}[与完备化的比较]
\label{algebraization-remark-compare-with-completion}
设 $(\mathcal{C}, \mathcal{O})$ 是一个环化位点。
设 $\mathcal{I} \subset \mathcal{O}$ 是一个有限型理想层。
设 $K \mapsto K^\wedge$ 为命题
\ref{algebraization-proposition-derived-completion} 的导出完备化函子。
对任意 $n \geq 1$，对象
$K \otimes_\mathcal{O}^\mathbf{L} \mathcal{O}/\mathcal{I}^n$
是导出完备的，因为它被 $\mathcal{I}$ 的局部截面的幂所零化。因此存在典范分解
$$
K \to K^\wedge \to K \otimes_\mathcal{O}^\mathbf{L} \mathcal{O}/\mathcal{I}^n
$$
，它分解了典范映射
$K \to K \otimes_\mathcal{O}^\mathbf{L} \mathcal{O}/\mathcal{I}^n$。
当 $n$ 变化时，这些映射彼此相容，于是得到比较映射
$$
K^\wedge
\longrightarrow
R\lim \left(K \otimes_\mathcal{O}^\mathbf{L} \mathcal{O}/\mathcal{I}^n\right)
$$
右端更容易被辨认为某种完备化。一般而言，这个比较映射并不是同构。
\end{remark}

\begin{remark}[局部化与导出完备化]
\label{algebraization-remark-localization-and-completion}
设 $(\mathcal{C}, \mathcal{O})$ 是一个环化位点。
设 $\mathcal{I} \subset \mathcal{O}$ 是一个有限型理想层。
设 $K \mapsto K^\wedge$ 为命题
\ref{algebraization-proposition-derived-completion} 的导出完备化函子。
由该命题证明中的构造可知，$K^\wedge|_U$
是 $K|_U$ 的导出完备化，其中 $U \in \Ob(\mathcal{C})$ 任意。
也可以如下证明这一点。由导出完备对象的定义可知，
$K^\wedge|_U$ 是导出完备的。因此得到典范映射
$a : (K|_U)^\wedge \to K^\wedge|_U$。
另一方面，若 $E$ 是 $D(\mathcal{O}_U)$ 的导出完备对象，则
$Rj_*E$ 是 $D(\mathcal{O})$ 的导出完备对象，这是引理
\ref{algebraization-lemma-pushforward-derived-complete} 的结论。这里 $j$ 是局部化态射
（《位点上的模》，\ref{sites-modules-section-localize} 节）。
因此还得到典范映射
$b : K^\wedge \to Rj_*((K|_U)^\wedge)$。
我们略去对 $b$ 的伴随映射为 $a$ 的逆映射这一（形式）事实的验证。
\end{remark}

\begin{remark}[完备张量积]
\label{algebraization-remark-completed-tensor-product}
设 $(\mathcal{C}, \mathcal{O})$ 是一个环化位点，并设
$\mathcal{I} \subset \mathcal{O}$ 是一个有限型理想层。
以 $K \mapsto K^\wedge$ 表示命题
\ref{algebraization-proposition-derived-completion} 中的伴随函子。置
$$
K \otimes^\wedge_\mathcal{O} L = (K \otimes_\mathcal{O}^\mathbf{L} L)^\wedge
$$
这个{\it 完备张量积}定义了函子
$D_{comp}(\mathcal{O}) \times D_{comp}(\mathcal{O}) \to D_{comp}(\mathcal{O})$
，使得
$$
\Hom_{D_{comp}(\mathcal{O})}(K, R\SheafHom_\mathcal{O}(L, M))
=
\Hom_{D_{comp}(\mathcal{O})}(K \otimes_\mathcal{O}^\wedge L, M)
$$
对 $K, L, M \in D_{comp}(\mathcal{O})$ 成立。注意，
$R\SheafHom_\mathcal{O}(L, M) \in D_{comp}(\mathcal{O})$；
这是引理 \ref{algebraization-lemma-derived-complete-internal-hom} 的结论。
\end{remark}

\begin{lemma}
\label{algebraization-lemma-map-identifies-koszul-and-cech-complexes}
设 $\mathcal{C}$ 是一个位点。
假设 $\varphi : \mathcal{O} \to \mathcal{O}'$ 是环层的平坦同态。
设 $f_1, \ldots, f_r$ 是 $\mathcal{O}$ 的全局截面，并满足
$\mathcal{O}/(f_1, \ldots, f_r) \cong
\mathcal{O}'/(f_1, \ldots, f_r)\mathcal{O}'$。
则增广交错 {\v C}ech 复形之间的映射
$$
\xymatrix{
\mathcal{O} \to
\prod_{i_0} \mathcal{O}_{f_{i_0}} \to
\prod_{i_0 < i_1} \mathcal{O}_{f_{i_0}f_{i_1}} \to \ldots \to
\mathcal{O}_{f_1\ldots f_r} \ar[d] \\
\mathcal{O}' \to
\prod_{i_0} \mathcal{O}'_{f_{i_0}} \to
\prod_{i_0 < i_1} \mathcal{O}'_{f_{i_0}f_{i_1}} \to \ldots \to
\mathcal{O}'_{f_1\ldots f_r}
}
$$
是拟同构。
\end{lemma}

\begin{proof}
注意，第二个复形是第一个复形与 $\mathcal{O}'$ 的张量积。
第一个增广交错 {\v C}ech 复形可以写成 Koszul 复形
$K_n = K(\mathcal{O}, f_1^n, \ldots, f_r^n)$ 的余极限；见《更多代数》，
引理 \ref{more-algebra-lemma-extended-alternating-Cech-is-colimit-koszul}。
因此，只需证明 $K_n \to K_n \otimes_\mathcal{O} \mathcal{O}'$
是拟同构。由于 $\mathcal{O} \to \mathcal{O}'$ 平坦，只需证明
$H^i \to H^i \otimes_\mathcal{O} \mathcal{O}'$ 是同构，其中
$H^i$ 是第 $i$ 个上同调层，即 $H^i = H^i(K_n)$。
这些层被 $f_1^n, \ldots, f_r^n$ 零化；见《更多代数》，
引理 \ref{more-algebra-lemma-homotopy-koszul}。
因此这些层被 $(f_1, \ldots, f_r)^m$ 零化，其中某个 $m \gg 0$。
于是 $H^i \to H^i \otimes_\mathcal{O} \mathcal{O}'$ 是同构；
这是《位点上的模》，引理
\ref{sites-modules-lemma-neighbourhood-isomorphism} 的结论。
\end{proof}

\begin{lemma}
\label{algebraization-lemma-restriction-derived-complete-equivalence}
设 $\mathcal{C}$ 是一个位点。设 $\mathcal{O} \to \mathcal{O}'$
是环层的同态，并设 $\mathcal{I} \subset \mathcal{O}$
是一个有限型理想层。
若 $\mathcal{O} \to \mathcal{O}'$ 平坦且
$\mathcal{O}/\mathcal{I} \cong \mathcal{O}'/\mathcal{I}\mathcal{O}'$,
则限制函子 $D(\mathcal{O}') \to D(\mathcal{O})$
诱导等价
$D_{comp}(\mathcal{O}', \mathcal{I}\mathcal{O}') \to
D_{comp}(\mathcal{O}, \mathcal{I})$.
\end{lemma}

\begin{proof}
由引理 \ref{algebraization-lemma-pushforward-derived-complete}，限制函子
$r : D(\mathcal{O}') \to D(\mathcal{O})$
把 $D_{comp}(\mathcal{O}', \mathcal{I}\mathcal{O}')$
映入 $D_{comp}(\mathcal{O}, \mathcal{I})$。我们将构造一个拟逆
$E \mapsto E'$。

\medskip\noindent
设 $K \to \mathcal{O}$ 是 $D(\mathcal{O})$ 中由引理
\ref{algebraization-lemma-global-extended-cech-complex} 构造的态射。置
$K' = K \otimes_\mathcal{O}^\mathbf{L} \mathcal{O}'$；这是
$D(\mathcal{O}')$ 中的对象。
则 $K' \to \mathcal{O}'$ 是 $D(\mathcal{O}')$ 中的映射；
它满足引理 \ref{algebraization-lemma-global-extended-cech-complex}
相对于 $\mathcal{I}' = \mathcal{I}\mathcal{O}'$ 的结论。
映射 $K \to r(K')$ 是拟同构，这是引理
\ref{algebraization-lemma-map-identifies-koszul-and-cech-complexes} 的结论。现在，对
$E \in D_{comp}(\mathcal{O}, \mathcal{I})$，置
$$
E' = R\SheafHom_\mathcal{O}(r(K'), E)
$$
把它视为 $D(\mathcal{O}')$ 中的对象，其中使用
$\mathcal{O}'$ 在 $K'$ 上的模结构。由于 $E$ 导出完备，我们有
$E = R\SheafHom_\mathcal{O}(K, E)$；见命题
\ref{algebraization-proposition-derived-completion} 的证明。
另一方面，由于 $K \to r(K')$ 是同构，可见
$E \to r(E')$ 是 $D(\mathcal{O})$ 中的同构。为了完成证明，
还须说明：若 $E = r(M')$，其中 $M'$ 是
$D_{comp}(\mathcal{O}', \mathcal{I}')$ 的对象，则
$E' \cong M'$。为得到一个映射，我们使用
$$
M' = R\SheafHom_{\mathcal{O}'}(\mathcal{O}', M') \to
R\SheafHom_\mathcal{O}(r(\mathcal{O}'), r(M')) \to
R\SheafHom_\mathcal{O}(r(K'), r(M')) = E'
$$
，其中第二个箭头使用映射 $K' \to \mathcal{O}'$。
要看出这是同构，只需说明对该箭头应用 $r$ 后，所得映射就是上面的同构
$E \to r(E')$。细节从略。
\end{proof}

\begin{lemma}
\label{algebraization-lemma-pushforward-derived-complete-adjoint}
设 $f : (\Sh(\mathcal{D}), \mathcal{O}') \to (\Sh(\mathcal{C}), \mathcal{O})$
是环化拓扑斯的态射。设 $\mathcal{I} \subset \mathcal{O}$
和 $\mathcal{I}' \subset \mathcal{O}'$ 是有限型理想层，并且
$f^\sharp$ 把 $f^{-1}\mathcal{I}$ 映入 $\mathcal{I}'$。
则 $Rf_*$ 把 $D_{comp}(\mathcal{O}', \mathcal{I}')$
映入 $D_{comp}(\mathcal{O}, \mathcal{I})$，并且有左伴随
$Lf_{comp}^*$；该左伴随由先应用 $Lf^*$ 再作导出完备化得到。
\end{lemma}

\begin{proof}
第一条陈述已见于引理 \ref{algebraization-lemma-pushforward-derived-complete}。
第二条陈述是有意义的，因为我们有导出完备化函子
$D(\mathcal{O}') \to D_{comp}(\mathcal{O}', \mathcal{I}')$；
这是命题 \ref{algebraization-proposition-derived-completion} 的结论。
现在设 $K \in D_{comp}(\mathcal{O}, \mathcal{I})$ 且
$M \in D_{comp}(\mathcal{O}', \mathcal{I}')$。则
$$
\Hom(K, Rf_*M) = \Hom(Lf^*K, M) = \Hom(Lf_{comp}^*K, M)
$$
，其中等式来自导出完备化的泛性质。
\end{proof}

\begin{lemma}
\label{algebraization-lemma-pushforward-commutes-with-derived-completion}
\begin{reference}
这是 \cite[Lemma 6.5.9 (2)]{BS} 的推广。在拟凝聚模与（导出）代数栈
的态射这一背景下，可与 \cite[Theorem 6.5]{HL-P} 比较。
\end{reference}
设 $f : (\Sh(\mathcal{D}), \mathcal{O}') \to (\Sh(\mathcal{C}), \mathcal{O})$
是环化拓扑斯的态射。设 $\mathcal{I} \subset \mathcal{O}$
是一个有限型理想层，并设 $\mathcal{I}' \subset \mathcal{O}'$
为由 $f^\sharp(f^{-1}\mathcal{I})$ 生成的理想。则
$Rf_*$ 与导出完备化可交换，即
$Rf_*(K^\wedge) = (Rf_*K)^\wedge$.
\end{lemma}

\begin{proof}
由命题 \ref{algebraization-proposition-derived-completion}，导出完备化函子存在。
由引理 \ref{algebraization-lemma-pushforward-derived-complete}，对象
$Rf_*(K^\wedge)$ 是导出完备的，因此由导出完备化的泛性质得到典范映射
$(Rf_*K)^\wedge \to Rf_*(K^\wedge)$。可以在 $\mathcal{C}$
上局部地检验这个映射是同构。由于导出完备化与局部化可交换
（备注 \ref{algebraization-remark-localization-and-completion}），可以假设
$\mathcal{I}$ 由全局截面 $f_1, \ldots, f_r$ 生成。
于是 $\mathcal{I}'$ 由 $g_i = f^\sharp(f_i)$ 生成。由引理
\ref{algebraization-lemma-derived-completion-koszul}，只需证明
$$
R\lim \left(
Rf_*K \otimes^\mathbf{L}_\mathcal{O} K(\mathcal{O}, f_1^n, \ldots, f_r^n)
\right)
=
Rf_*\left(
R\lim
K \otimes^\mathbf{L}_{\mathcal{O}'} K(\mathcal{O}', g_1^n, \ldots, g_r^n)
\right)
$$
由于 $Rf_*$ 与 $R\lim$ 可交换
（《位点上的上同调》，引理
\ref{sites-cohomology-lemma-Rf-commutes-with-Rlim}），只需证明
$$
Rf_*K \otimes^\mathbf{L}_\mathcal{O} K(\mathcal{O}, f_1^n, \ldots, f_r^n) =
Rf_*\left(
K \otimes^\mathbf{L}_{\mathcal{O}'} K(\mathcal{O}', g_1^n, \ldots, g_r^n)
\right)
$$
这由投影公式（《位点上的上同调》，引理
\ref{sites-cohomology-lemma-projection-formula}）以及等式
$Lf^*K(\mathcal{O}, f_1^n, \ldots, f_r^n) =
K(\mathcal{O}', g_1^n, \ldots, g_r^n)$ 得出。
\end{proof}

\begin{lemma}
\label{algebraization-lemma-formal-functions-general}
设 $A$ 是一个环，$I \subset A$ 是有限生成理想。
设 $\mathcal{C}$ 是一个位点，$\mathcal{O}$ 是 $A$-代数层。
设 $\mathcal{F}$ 是一个 $\mathcal{O}$-模层。则
$$
R\Gamma(\mathcal{C}, \mathcal{F})^\wedge =
R\Gamma(\mathcal{C}, \mathcal{F}^\wedge)
$$
在 $D(A)$ 中成立，其中 $\mathcal{F}^\wedge$ 是 $\mathcal{F}$
关于 $I\mathcal{O}$ 的导出完备化，而左端取关于 $I$ 的导出完备化。
这给出两个谱序列
$$
E_2^{i, j} = H^i(H^j(\mathcal{C}, \mathcal{F})^\wedge)
\quad\text{且}\quad
E_2^{p, q} = H^p(\mathcal{C}, H^q(\mathcal{F}^\wedge))
$$
，二者都收敛到
$H^*(R\Gamma(\mathcal{C}, \mathcal{F})^\wedge) =
H^*(\mathcal{C}, \mathcal{F}^\wedge)$
\end{lemma}

\begin{proof}
把引理 \ref{algebraization-lemma-pushforward-commutes-with-derived-completion}
应用于环化拓扑斯的态射 $(\mathcal{C}, \mathcal{O}) \to (pt, A)$，
再取上同调，即得第一条陈述。第二个谱序列是《导出范畴》，引理
\ref{derived-lemma-two-ss-complex-functor} 的第二个谱序列。
第一个谱序列则是把《更多代数》，例
\ref{more-algebra-example-derived-completion-spectral-sequence}
中的谱序列应用于 $R\Gamma(\mathcal{C}, \mathcal{F})^\wedge$ 所得。
\end{proof}

\begin{remark}
\label{algebraization-remark-local-calculation-derived-completion}
设 $(\mathcal{C}, \mathcal{O})$ 是一个环化位点。
设 $\mathcal{I} \subset \mathcal{O}$ 是一个有限型理想层。
设 $K \mapsto K^\wedge$ 为命题
\ref{algebraization-proposition-derived-completion} 的导出完备化。
设 $U \in \Ob(\mathcal{C})$ 是一个对象，并且作为理想层，
$\mathcal{I}$ 由 $f_1, \ldots, f_r \in \mathcal{I}(U)$ 生成。
置 $A = \mathcal{O}(U)$ 和 $I = (f_1, \ldots, f_r) \subset A$。
警告：未必有 $I = \mathcal{I}(U)$。于是
$$
R\Gamma(U, K^\wedge) = R\Gamma(U, K)^\wedge
$$
，其中右端是对象 $R\Gamma(U, K)$（它属于 $D(A)$）关于 $I$
的导出完备化。这是因为导出完备化与局部化可交换
（备注 \ref{algebraization-remark-localization-and-completion}），并且可应用引理
\ref{algebraization-lemma-formal-functions-general}。
\end{remark}








\section{形式函数定理}
\label{algebraization-section-formal-functions}

\noindent
我们暂时中断正文的推进，讨论概形上拟凝聚模背景中的导出完备化，
并用它给出形式函数定理的一个略有不同的证明。备注
\ref{algebraization-remark-references} 中列出了一些文献线索。

\medskip\noindent
引理 \ref{algebraization-lemma-pushforward-commutes-with-derived-completion}
是形式函数定理的一个（非常形式化的）导出版本
（《概形上同调》，定理 \ref{coherent-theorem-formal-functions}）。
更明确地说，设 $f : X \to S$ 是概形的态射，
$\mathcal{I} \subset \mathcal{O}_S$ 是有限型拟凝聚理想层，
并且 $\mathcal{F}$ 是 $X$ 上的拟凝聚层。则该引理断言
\begin{equation}
\label{algebraization-equation-formal-functions}
Rf_*(\mathcal{F}^\wedge) = (Rf_*\mathcal{F})^\wedge
\end{equation}
，其中 $\mathcal{F}^\wedge$ 是 $\mathcal{F}$ 关于
$f^{-1}\mathcal{I} \cdot \mathcal{O}_X$ 的导出完备化，右端则是
$Rf_*\mathcal{F}$ 关于 $\mathcal{I}$ 的导出完备化。要看出这确实还原了
形式函数定理，还需要做一些工作。

\begin{lemma}
\label{algebraization-lemma-sections-derived-completion-pseudo-coherent}
设 $X$ 是局部诺特概形。设 $\mathcal{I} \subset \mathcal{O}_X$
是拟凝聚理想层。设 $K$ 是 $D(\mathcal{O}_X)$ 的伪凝聚对象，
其导出完备化为 $K^\wedge$。则
$$
H^p(U, K^\wedge) = \lim H^p(U, K)/I^nH^p(U, K) =
H^p(U, K)^\wedge
$$
对任意仿射开集 $U \subset X$ 成立，其中 $I = \mathcal{I}(U)$，
而右端取关于 $I$ 的导出完备化。
\end{lemma}

\begin{proof}
写成 $U = \Spec(A)$。环 $A$ 是诺特的，因此 $I \subset A$
是有限生成理想。于是
$$
R\Gamma(U, K^\wedge) = R\Gamma(U, K)^\wedge
$$
，这是备注 \ref{algebraization-remark-local-calculation-derived-completion} 的结论。
现在，$R\Gamma(U, K)$ 是一个伪凝聚的 $A$-模复形
（《概形的导出范畴》，引理
\ref{perfect-lemma-pseudo-coherent-affine}）。
由《更多代数》，引理
\ref{more-algebra-lemma-derived-completion-pseudo-coherent}，
其第 $p$ 个上同调模（这里的复形是 $R\Gamma(U, K^\wedge)$）
等于相应模的 $I$-进完备化，也就是 $H^p(U, K)$ 的完备化。
这证明了第一个等式。
第二个（较不重要的）等式则由再次应用刚才的引理立即得到。
\end{proof}

\begin{lemma}
\label{algebraization-lemma-derived-completion-pseudo-coherent}
设 $X$ 是局部诺特概形。设 $\mathcal{I} \subset \mathcal{O}_X$
是拟凝聚理想层。设 $K$ 是 $D(\mathcal{O}_X)$ 的对象。则
\begin{enumerate}
\item 导出完备化 $K^\wedge$ 等于
$R\lim (K \otimes_{\mathcal{O}_X}^\mathbf{L} \mathcal{O}_X/\mathcal{I}^n)$.
\end{enumerate}
设 $K$ 是 $D(\mathcal{O}_X)$ 的伪凝聚对象。则
\begin{enumerate}
\item[(2)] 上同调层 $H^q(K^\wedge)$ 等于
$\lim H^q(K)/\mathcal{I}^nH^q(K)$.
\end{enumerate}
设 $\mathcal{F}$ 是凝聚 $\mathcal{O}_X$-模\footnote{例如可取
$H^q(K)$，其中 $K$ 是所讨论局部诺特概形 $X$ 上的伪凝聚对象。}。则
\begin{enumerate}
\item[(3)] 导出完备化 $\mathcal{F}^\wedge$ 等于
$\lim \mathcal{F}/\mathcal{I}^n\mathcal{F}$,
\item[(4)]
$\lim \mathcal{F}/\mathcal{I}^n \mathcal{F} =
R\lim \mathcal{F}/\mathcal{I}^n \mathcal{F}$,
\item[(5)] $H^p(U, \mathcal{F}^\wedge) = 0$ 对每个 $p \not = 0$
以及每个仿射开集 $U \subset X$ 都成立。
\end{enumerate}
\end{lemma}

\begin{proof}
(1) 的证明。存在典范映射
$$
K \longrightarrow
R\lim (K \otimes_{\mathcal{O}_X}^\mathbf{L} \mathcal{O}_X/\mathcal{I}^n),
$$
；见备注 \ref{algebraization-remark-compare-with-completion}。
导出完备化与取开子概形可交换
（备注 \ref{algebraization-remark-localization-and-completion}）。
$R\lim$ 的构造也与取开子概形可交换。因此，为检验上述映射是同构，
可以局部地工作。于是可假设 $X = U = \Spec(A)$。
写 $I = (f_1, \ldots, f_r)$，并设
$K_n = K(A, f_1^n, \ldots, f_r^n)$ 为 Koszul 复形。
由《更多代数》，引理 \ref{more-algebra-lemma-sequence-Koszul-complexes}，
前系统 $\{K_n\}$ 与 $\{A/I^n\}$ 在 $D(A)$ 中同构。
利用等价 $D(A) = D_{\QCoh}(\mathcal{O}_X)$
（《概形的导出范畴》，引理
\ref{perfect-lemma-affine-compare-bounded}），可知前系统
$\{K(\mathcal{O}_X, f_1^n, \ldots, f_r^n)\}$
与 $\{\mathcal{O}_X/\mathcal{I}^n\}$ 在
$D(\mathcal{O}_X)$ 中同构。这证明了下式中的第二个等式：
$$
K^\wedge = R\lim \left(
K \otimes_{\mathcal{O}_X}^\mathbf{L} K(\mathcal{O}_X, f_1^n, \ldots, f_r^n)
\right) =
R\lim (K \otimes_{\mathcal{O}_X}^\mathbf{L} \mathcal{O}_X/\mathcal{I}^n)
$$
第一个等式则是引理 \ref{algebraization-lemma-derived-completion-koszul}。

\medskip\noindent
假设 $K$ 伪凝聚。对仿射开集 $U \subset X$，由
$H^q(U, K^\wedge) = \lim H^q(U, K)/\mathcal{I}^n(U)H^q(U, K)$
和引理 \ref{algebraization-lemma-sections-derived-completion-pseudo-coherent} 可知该等式成立。
由于这对每个 $U$ 都成立，作为层有
$H^q(K^\wedge) = \lim H^q(K)/\mathcal{I}^nH^q(K)$。这证明了 (2)。

\medskip\noindent
(3) 是 (2) 的特殊情形。(4) 与 (5) 由《概形的导出范畴》，引理
\ref{perfect-lemma-Rlim-quasi-coherent} 得出。
\end{proof}

\begin{lemma}
\label{algebraization-lemma-formal-functions}
设 $A$ 是诺特环，$I \subset A$ 是理想。设 $X$ 是
$A$ 上的诺特概形，$\mathcal{F}$ 是凝聚 $\mathcal{O}_X$-模。
假设 $H^p(X, \mathcal{F})$ 是有限 $A$-模，并且对所有 $p$ 都如此。
则有短正合列
$$
0 \to R^1\lim H^{p - 1}(X, \mathcal{F}/I^n\mathcal{F}) \to
H^p(X, \mathcal{F})^\wedge \to \lim H^p(X, \mathcal{F}/I^n\mathcal{F}) \to 0
$$
，它们是 $A$-模的短正合列，其中 $H^p(X, \mathcal{F})^\wedge$
是通常的 $I$-进完备化。若 $f$ 是固有的，则 $R^1\lim$ 项为零。
\end{lemma}

\begin{proof}
考察引理 \ref{algebraization-lemma-formal-functions-general} 的两个谱序列。
由《更多代数》，引理
\ref{more-algebra-lemma-derived-completion-pseudo-coherent}，第一个谱序列退化。
得到 $H^p(X, \mathcal{F})^\wedge$，它位于次数 $p$。
这里用到了 $H^p(X, \mathcal{F})$ 是有限 $A$-模这一假设。
第二个谱序列退化，因为
$$
\mathcal{F}^\wedge = \lim \mathcal{F}/I^n\mathcal{F} =
R\lim \mathcal{F}/I^n\mathcal{F}
$$
由引理 \ref{algebraization-lemma-derived-completion-pseudo-coherent} 可知这是一个层。
得到 $H^p(X, \lim \mathcal{F}/I^n\mathcal{F})$，它位于次数 $p$。
由于 $R\Gamma(X, -)$ 与导出极限可交换
（《内射对象》，引理 \ref{injectives-lemma-RF-commutes-with-Rlim}），还有
$$
R\Gamma(X, \lim \mathcal{F}/I^n\mathcal{F}) =
R\Gamma(X, R\lim \mathcal{F}/I^n\mathcal{F}) =
R\lim R\Gamma(X, \mathcal{F}/I^n\mathcal{F})
$$
由《更多代数》，备注
\ref{more-algebra-remark-how-unique}，得到正合列
$$
0 \to
R^1\lim H^{p - 1}(X, \mathcal{F}/I^n\mathcal{F}) \to
H^p(X, \lim \mathcal{F}/I^n\mathcal{F}) \to
\lim H^p(X, \mathcal{F}/I^n\mathcal{F}) \to 0
$$
，它们是 $A$-模的正合列。结合以上结果，得到引理的第一条陈述。
$R^1\lim$ 项的消失由《概形上同调》，引理
\ref{coherent-lemma-ML-cohomology-powers-ideal} 得出。
\end{proof}

\begin{remark}
\label{algebraization-remark-references}
下面列出文献中讨论相关材料的一些参考文献。
对于固有映射，“导出形式函数定理”似乎可追溯到
\cite[Theorem 6.3.1]{lurie-thesis}。
\cite{dag12} 中也有讨论，尤其是第 4 章对仿射情形的讨论；另见《更多代数》，
\ref{more-algebra-section-derived-completion} 节。
\cite[Section 2.9]{G-R} 使用（ind）凝聚模讨论了固有基变换与导出完备化。
拟凝聚模复形的一个与（\ref{algebraization-equation-formal-functions}）类似的结果见
\cite[Theorem 6.5]{HL-P}。
\end{remark}







\section{局部上同调的代数化，I}
\label{algebraization-section-algebraization-sections-general}

\noindent
设 $A$ 是诺特环，$I$ 与 $J$ 是 $A$ 的两个理想。
设 $M$ 是有限 $A$-模。本节研究对象
$$
R\Gamma_J(M)^\wedge
\quad\text{属于}\quad
D(A)
$$
的上同调群，其中 ${}^\wedge$ 表示导出 $I$-进完备化。注意，在
《对偶化复形》，引理 \ref{dualizing-lemma-completion-local-H0} 中，
我们已经证明：若 $A$ 关于 $I$ 完备，则存在同构
$$
\colim H^0_Z(M) \longrightarrow H^0(R\Gamma_J(M)^\wedge)
$$
，其中（有向）余极限遍历闭子集 $Z = V(J')$，它们满足
$J' \subset J$ 和 $V(J') \cap V(I) = V(J) \cap V(I)$。
这些闭子集的并为
\begin{equation}
\label{algebraization-equation-associated-subset}
T = \{\mathfrak p \in \Spec(A) :
V(\mathfrak p) \cap V(I) \subset V(J) \cap V(I)\}
\end{equation}
这是 $\Spec(A)$ 的一个对特殊化稳定的子集。上述结果于是变为：
$$
H^0_T(M) \longrightarrow H^0(R\Gamma_J(M)^\wedge)
$$
若 $A$ 关于 $I$ 完备，则这个映射是同构；见《局部上同调》，引理
\ref{local-cohomology-lemma-adjoint-ext} 与备注
\ref{local-cohomology-remark-upshot}。
我们把这一同构推广到更高上同调群的方法基于下面的引理。

\begin{lemma}
\label{algebraization-lemma-kill-completion-general}
设 $I, J$ 是诺特环 $A$ 的理想。设 $M$ 是有限 $A$-模。
设 $\mathfrak p \subset A$ 是素理想，$s$ 与 $d$ 是整数。假设
\begin{enumerate}
\item $A$ 有对偶化复形，
\item $\mathfrak p \not \in V(J) \cap V(I)$,
\item $\text{cd}(A, I) \leq d$，且
\item 对所有素理想 $\mathfrak p' \subset \mathfrak p$，都有
$\text{depth}_{A_{\mathfrak p'}}(M_{\mathfrak p'}) +
\dim((A/\mathfrak p')_\mathfrak q) > d + s$
，其中 $\mathfrak q \in V(\mathfrak p') \cap V(J) \cap V(I)$ 任意。
\end{enumerate}
则存在 $f \in A$，$f \not \in \mathfrak p$，它零化
$H^i(R\Gamma_J(M)^\wedge)$，其中 $i \leq s$；这里 ${}^\wedge$
表示 $I$-进完备化。
\end{lemma}

\begin{proof}
我们将使用 $R\Gamma_J = R\Gamma_{V(J)}$，并对 $I + J$ 作同样处理；
见《对偶化复形》，引理
\ref{dualizing-lemma-local-cohomology-noetherian}。注意
$R\Gamma_J(M)^\wedge = R\Gamma_I(R\Gamma_J(M))^\wedge =
R\Gamma_{I + J}(M)^\wedge$；见
《对偶化复形》，引理
\ref{dualizing-lemma-complete-and-local} 与
\ref{dualizing-lemma-local-cohomology-ss}.
因此可以把 $J$ 替换为 $I + J$，并假设 $I \subset J$
以及 $\mathfrak p \not \in V(J)$。回忆
$$
R\Gamma_J(M)^\wedge = R\Hom_A(R\Gamma_I(A), R\Gamma_J(M))
$$
，这是把《更多代数》，引理
\ref{more-algebra-lemma-derived-completion} 中对导出完备化的描述
与《对偶化复形》，引理
\ref{dualizing-lemma-compute-local-cohomology-noetherian}
中对局部上同调的描述结合起来所得。假设 (3) 意味着
$R\Gamma_I(A)$ 仅在次数 $\leq d$ 上有非零上同调。
利用 $R\Gamma_I(A)$ 的典范截断，只需证明
$$
\text{Ext}^i(N, R\Gamma_J(M))
$$
被某个 $f \in A$，$f \not \in \mathfrak p$ 零化，其中
$i \leq s + d$，而所述结论对任意 $A$-模 $N$ 都成立。
再利用 $R\Gamma_J(M)$ 的典范截断，只需证明
$H^i_J(M)$ 被某个 $f \in A$，$f \not \in \mathfrak p$ 零化，
其中 $i \leq s + d$。这由《局部上同调》，引理
\ref{local-cohomology-lemma-kill-local-cohomology-at-prime} 得出。
\end{proof}

\begin{lemma}
\label{algebraization-lemma-kill-colimit-weak-general}
设 $I, J$ 是某个诺特环的理想。设 $M$ 是有限 $A$-模。
设 $s$ 与 $d$ 是整数。令 $T$ 如
（\ref{algebraization-equation-associated-subset}）中所定义，并假设
\begin{enumerate}
\item $A$ 有对偶化复形，
\item 若 $\mathfrak p \in V(I)$，则不加条件，
\item 若 $\mathfrak p \not \in V(I)$ 且 $\mathfrak p \in T$，则
$\dim((A/\mathfrak p)_\mathfrak q) \leq d$，其中存在
某个 $\mathfrak q \in V(\mathfrak p) \cap V(J) \cap V(I)$，
\item 若 $\mathfrak p \not \in V(I)$ 且 $\mathfrak p \not \in T$，则
$$
\text{depth}_{A_\mathfrak p}(M_\mathfrak p) \geq s
\quad\text{或}\quad
\text{depth}_{A_\mathfrak p}(M_\mathfrak p) +
\dim((A/\mathfrak p)_\mathfrak q) > d + s
$$
对所有 $\mathfrak q \in V(\mathfrak p) \cap V(J) \cap V(I)$ 成立。
\end{enumerate}
则存在理想 $J_0 \subset J$，满足
$V(J_0) \cap V(I) = V(J) \cap V(I)$，使得对任意 $J' \subset J_0$，若
$V(J') \cap V(I) = V(J) \cap V(I)$，则映射
$$
R\Gamma_{J'}(M) \longrightarrow R\Gamma_{J_0}(M)
$$
在次数 $\leq s$ 的上同调上诱导同构，并且这些模被 $J_0I$ 的某个幂零化。
\end{lemma}

\begin{proof}
考虑集合
$$
B = \{\mathfrak p \not \in V(I),\ \mathfrak p \in T,\text{ 且 }
\text{depth}(M_\mathfrak p) \leq s\}
$$
选择 $J_0 \subset J$，使 $V(J_0)$ 是 $B \cup V(J)$ 的闭包。

\medskip\noindent
断言 I：$V(J_0) \cap V(I) = V(J) \cap V(I)$。

\medskip\noindent
断言 I 的证明。包含关系 $\supset$ 由构造成立。
设 $\mathfrak p$ 是 $V(J_0)$ 的极小素理想。
若 $\mathfrak p \in B \cup V(J)$，则或者 $\mathfrak p \in T$，
或者 $\mathfrak p \in V(J)$；两种情形下都如所需有
$V(\mathfrak p) \cap V(I) \subset V(J) \cap V(I)$。
若 $\mathfrak p \not \in B \cup V(J)$，则
$V(\mathfrak p) \cap B$ 稠密，因而无限，可知
$\text{depth}(M_\mathfrak p) < s$；见《局部上同调》，引理
\ref{local-cohomology-lemma-depth-function}。
事实上，设
$V(\mathfrak p) \cap B = \{\mathfrak p_\lambda\}_{\lambda \in \Lambda}$.
依照 (3)，选择
$\mathfrak q_\lambda \in V(\mathfrak p_\lambda) \cap V(J) \cap V(I)$。
设 $\delta : \Spec(A) \to \mathbf{Z}$ 是与对偶化复形
$\omega_A^\bullet$（定义在 $A$ 上）相伴的维数函数。由于 $\Lambda$ 无限而 $\delta$ 有界，
存在无限子集 $\Lambda' \subset \Lambda$，使得
$\delta(\mathfrak q_\lambda)$ 在其上为常数。对
$\lambda \in \Lambda'$，有
$$
\text{depth}(M_{\mathfrak p_\lambda}) +
\delta(\mathfrak p_\lambda) - \delta(\mathfrak q_\lambda) =
\text{depth}(M_{\mathfrak p_\lambda}) +
\dim((A/\mathfrak p_\lambda)_{\mathfrak q_\lambda})
\leq d + s
$$
，这是 (3) 与 $B$ 的定义的结论。利用函数
$\text{depth} + \delta$ 的半连续性（见《概形的对偶性》，引理
\ref{duality-lemma-sitting-in-degrees}），可得
$$
\text{depth}(M_\mathfrak p) +
\dim((A/\mathfrak p)_{\mathfrak q_\lambda}) =
\text{depth}(M_\mathfrak p) + \delta(\mathfrak p) - \delta(\mathfrak q_\lambda)
\leq d + s
$$
又因为 $\mathfrak p \not \in V(I)$，由 (4) 可知
$\mathfrak p \in T$，即
$V(\mathfrak p) \cap V(I) \subset V(J) \cap V(I)$。断言 I 得证。

\medskip\noindent
断言 II：映射 $H^i_{J_0}(M) \to H^i_J(M)$ 是同构，其中
$i \leq s$，并且 $J' \subset J_0$ 满足
$V(J') \cap V(I) = V(J) \cap V(I)$。

\medskip\noindent
断言 II 的证明。选择 $\mathfrak p \in V(J')$，使其不属于 $V(J_0)$。
只需证明 $H^i_{\mathfrak pA_\mathfrak p}(M_\mathfrak p) = 0$
对 $i \leq s$ 成立；见《局部上同调》，引理
\ref{local-cohomology-lemma-isomorphism}。
注意 $\mathfrak p \in T$。因此，由于 $\mathfrak p$ 不属于 $B$，
可知 $\text{depth}(M_\mathfrak p) > s$，而这些群由《对偶化复形》，引理
\ref{dualizing-lemma-depth} 消失。

\medskip\noindent
断言 III：引理的最后一条陈述成立。

\medskip\noindent
由断言 II，对 $i \leq s$ 有
$$
H^i_T(M) = H^i_{J_0}(M) = H^i_{J'}(M)
$$
，其中 $J' \subset J_0$ 是任意满足
$V(J')  \cap V(I) = V(J) \cap V(I)$ 的理想。见《局部上同调》，引理
\ref{local-cohomology-lemma-adjoint-ext}。
现在检验《局部上同调》，命题
\ref{local-cohomology-proposition-annihilator} 对子集
$T \subset T \cup V(I)$、模 $M$ 和整数 $s$ 的假设。
须证明：给定 $\mathfrak p \subset \mathfrak q$，若
$\mathfrak p \not \in T \cup V(I)$ 且 $\mathfrak q \in T$，则
$$
\text{depth}_{A_\mathfrak p}(M_\mathfrak p) +
\dim((A/\mathfrak p)_\mathfrak q) > s
$$
若 $\text{depth}(M_\mathfrak p) \geq s$，则结论成立，因为
$(A/\mathfrak p)_\mathfrak q$ 的维数至少为 $1$。
因此可假设 $\text{depth}(M_\mathfrak p) < s$。
若 $\mathfrak q \in V(I)$，则 $\mathfrak q \in V(J) \cap V(I)$，
不等式由 (4) 成立。若 $\mathfrak q \not \in V(I)$，则由 (3)
可以选择 $\mathfrak q' \in V(\mathfrak q) \cap V(J) \cap V(I)$，使
$\dim((A/\mathfrak q)_{\mathfrak q'}) \leq d$。于是由假设 (4) 得
$$
\text{depth}_{A_\mathfrak p}(M_\mathfrak p) +
\dim((A/\mathfrak p)_{\mathfrak q'}) > s + d
$$
由于 $A$ 是链环，这蕴含所需不等式。应用《局部上同调》，命题
\ref{local-cohomology-proposition-annihilator}，得到
$J'' \subset A$，满足 $V(J'') \subset T \cup V(I)$，并且
$J''$ 零化 $H^i_T(M)$，其中 $i \leq s$。
于是可以写成 $V(J'') \cup V(J_0) \cup V(I) = V(J'I)$，
其中某个 $J' \subset J_0$ 满足
$V(J') \cap V(I) = V(J) \cap V(I)$。
把 $J_0$ 替换为 $J'$，证明完成。
\end{proof}

\begin{lemma}
\label{algebraization-lemma-kill-colimit-general}
在引理 \ref{algebraization-lemma-kill-colimit-weak-general} 中，若把空条件 (2)
替换为假设
\begin{enumerate}
\item[(2')] 若 $\mathfrak p \in V(I)$ 且
$\mathfrak p \not \in V(J) \cap V(I)$，则
$\text{depth}_{A_\mathfrak p}(M_\mathfrak p) +
\dim((A/\mathfrak p)_\mathfrak q) > s$
对所有 $\mathfrak q \in V(\mathfrak p) \cap V(J) \cap V(I)$ 成立，
\end{enumerate}
则这些条件还蕴含 $H^i_{J_0}(M)$ 是有限 $A$-模，其中 $i \leq s$。
\end{lemma}

\begin{proof}
回忆 $H^i_{J_0}(M) = H^i_T(M)$；见引理
\ref{algebraization-lemma-kill-colimit-weak-general} 的证明。因此只需检验：若
$\mathfrak p \not \in T$、$\mathfrak q \in T$ 且
$\mathfrak p \subset \mathfrak q$，则
$\text{depth}_{A_\mathfrak p}(M_\mathfrak p) +
\dim((A/\mathfrak p)_\mathfrak q) > s$；见《局部上同调》，命题
\ref{local-cohomology-proposition-finiteness}。
条件 (2') 告诉我们，这对 $\mathfrak p \in V(I)$ 成立。
由于 $H^i_T(M)$ 被 $IJ_0$ 的某个幂零化，当
$\mathfrak p \not \in V(IJ_0)$ 时该条件也成立；这是《局部上同调》，命题
\ref{local-cohomology-proposition-annihilator} 的结论。
这涵盖所有情形，证明完成。
\end{proof}

\begin{lemma}
\label{algebraization-lemma-kill-colimit-support-general}
若在引理 \ref{algebraization-lemma-kill-colimit-weak-general} 中还假设
\begin{enumerate}
\item[(6)] 若 $\mathfrak p \not \in V(I)$ 且 $\mathfrak p \in T$，则
$\text{depth}_{A_\mathfrak p}(M_\mathfrak p) > s$,
\end{enumerate}
则 $H^i_{J_0}(M) = H^i_J(M) = H^i_{J + I}(M)$ 对 $i \leq s$ 成立，
并且这些模被 $I$ 的某个幂零化。
\end{lemma}

\begin{proof}
选择 $\mathfrak p \in V(J)$ 或 $\mathfrak p \in V(J_0)$，但要求
$\mathfrak p \not \in V(J + I) = V(J_0 + I)$。
只需证明 $H^i_{\mathfrak pA_\mathfrak p}(M_\mathfrak p) = 0$
对 $i \leq s$ 成立；见《局部上同调》，引理
\ref{local-cohomology-lemma-isomorphism}。
这些群由条件 (6) 和《对偶化复形》，引理
\ref{dualizing-lemma-depth} 消失。最后一条陈述由《局部上同调》，命题
\ref{local-cohomology-proposition-annihilator} 得出。
\end{proof}

\begin{lemma}
\label{algebraization-lemma-algebraize-local-cohomology-general}
设 $I, J$ 是诺特环 $A$ 的理想。设 $M$ 是有限 $A$-模。
设 $s$ 与 $d$ 是整数。令 $T$ 如
（\ref{algebraization-equation-associated-subset}）中所定义，并假设
\begin{enumerate}
\item $A$ 是 $I$-进完备的，并且有对偶化复形，
\item 若 $\mathfrak p \in V(I)$，则不加条件，
\item $\text{cd}(A, I) \leq d$,
\item 若 $\mathfrak p \not \in V(I)$ 且 $\mathfrak p \not \in T$，则
$$
\text{depth}_{A_\mathfrak p}(M_\mathfrak p) \geq s
\quad\text{或}\quad
\text{depth}_{A_\mathfrak p}(M_\mathfrak p) +
\dim((A/\mathfrak p)_\mathfrak q) > d + s
$$
对所有 $\mathfrak q \in V(\mathfrak p) \cap V(J) \cap V(I)$ 成立，
\item 若 $\mathfrak p \not \in V(I)$、$\mathfrak p \not \in T$、
$V(\mathfrak p) \cap V(J) \cap V(I) \not = \emptyset$，并且
$\text{depth}(M_\mathfrak p) < s$，则下列条件之一成立
\footnote{我们的方法迫使我们加入这个额外条件。稍后将回到这一点
（插入未来的引用）。}：
\begin{enumerate}
\item $\dim(\text{Supp}(M_\mathfrak p)) < s + 2$\footnote{例如，若
$M$ 满足 Serre 条件 $(S_s)$，并且这一条件是在
$V(I) \cup T$ 的补集上要求的，则如此。}，或者
\item  $\delta(\mathfrak p) > d + \delta_{max} - 1$
，其中 $\delta$ 是维数函数，而 $\delta_{max}$
是 $\delta$ 在 $V(J) \cap V(I)$ 上的最大值，或者
\item $\text{depth}_{A_\mathfrak p}(M_\mathfrak p) +
\dim((A/\mathfrak p)_\mathfrak q) > d + s + \delta_{max} - \delta_{min} - 2$
对所有 $\mathfrak q \in V(\mathfrak p) \cap V(J) \cap V(I)$ 成立。
\end{enumerate}
\end{enumerate}
则存在理想 $J_0 \subset J$，满足
$V(J_0) \cap V(I) = V(J) \cap V(I)$
，使得对任意 $J' \subset J_0$，若
$V(J') \cap V(I) = V(J) \cap V(I)$，则映射
$$
R\Gamma_{J'}(M) \longrightarrow R\Gamma_J(M)^\wedge
$$
在次数 $\leq s$ 的上同调上诱导同构。
这里 ${}^\wedge$ 表示导出 $I$-进完备化。
\end{lemma}

\noindent
我们建议读者先阅读局部情形的证明
（引理 \ref{algebraization-lemma-algebraize-local-cohomology}），因为它无需处理所有这些
不等式便能解释证明的结构。

\begin{proof}
对理想 $\mathfrak a \subset A$，有
$R\Gamma_\mathfrak a = R\Gamma_{V(\mathfrak a)}$；见《对偶化复形》，引理
\ref{dualizing-lemma-local-cohomology-noetherian}。其次，注意
$$
R\Gamma_J(M)^\wedge =
R\Gamma_I(R\Gamma_J(M))^\wedge =
R\Gamma_{I + J}(M)^\wedge =
R\Gamma_{I + J'}(M)^\wedge =
R\Gamma_I(R\Gamma_{J'}(M))^\wedge =
R\Gamma_{J'}(M)^\wedge
$$
，这是《对偶化复形》，引理
\ref{dualizing-lemma-local-cohomology-ss} 与
\ref{dualizing-lemma-complete-and-local} 的结论。
这说明了如何定义引理陈述中的箭头。

\medskip\noindent
我们断言：引理 \ref{algebraization-lemma-kill-colimit-weak-general} 的假设
由当前对 $M$ 的假设蕴含。唯一需要检验的是假设 (3)。
因此设 $\mathfrak p \not \in V(I)$ 且 $\mathfrak p \in T$。
那么 $V(\mathfrak p) \cap V(I)$ 非空，因为 $I$ 包含在
$A$ 的 Jacobson 根中
（《代数》，引理 \ref{algebra-lemma-radical-completion}）。
由于 $\mathfrak p \in T$，有
$V(\mathfrak p) \cap V(I) = V(\mathfrak p) \cap V(J) \cap V(I)$。
设 $\mathfrak q \in V(\mathfrak p) \cap V(I)$ 是某个不可约分支的通用点。
有 $\text{cd}(A_\mathfrak q, I_\mathfrak q) \leq d$；这是《局部上同调》，引理
\ref{local-cohomology-lemma-cd-local} 的结论。
又有 $V(\mathfrak pA_\mathfrak q) \cap V(I_\mathfrak q) =
\{\mathfrak qA_\mathfrak q\}$，这是由 $\mathfrak q$ 的选择得到的；进而
$\dim((A/\mathfrak p)_\mathfrak q) \leq d$，见《局部上同调》，引理
\ref{local-cohomology-lemma-cd-bound-dim-local}。

\medskip\noindent
注意，引理对 $s < 0$ 成立。这并非平凡情形，因为不能先验地看出
$H^i(R\Gamma_J(M)^\wedge)$ 在 $i < 0$ 时为零。不过，这一消失已在
《对偶化复形》，引理 \ref{dualizing-lemma-completion-local} 中证明。
对于 $s \geq 0$，我们将用归纳法证明引理。

\medskip\noindent
$s = 0$ 时的引理由引理 \ref{algebraization-lemma-kill-colimit-weak-general}
的结论和《对偶化复形》，引理
\ref{dualizing-lemma-completion-local-H0} 立即得出。

\medskip\noindent
假设 $s > 0$，且引理对更小的 $s$ 已经证明。
设 $M' \subset M$ 是支集包含于 $V(I) \cup T$ 的极大子模。
则 $M'$ 是有限 $A$-模，并且其支集包含于
$V(J') \cup V(I)$，其中某个理想 $J' \subset J$ 满足
$V(J') \cap V(I) = V(J) \cap V(I)$。我们断言
$$
R\Gamma_{J'}(M') \to R\Gamma_J(M')^\wedge
$$
对 $J'$ 的任意选择都是同构。事实上，可以选择短正合列
$0 \to M_1 \oplus M_2 \to M' \to N \to 0$
，其中 $M_1$ 被 $J'$ 的某个幂零化，$M_2$ 被 $I$ 的某个幂零化，
而 $N$ 被 $I + J'$ 的某个幂零化。因此，只需对
$M_1$、$M_2$ 与 $N$ 证明该断言。对 $M_1$，有
$R\Gamma_{J'}(M_1) = M_1$；又因为 $M_1$ 是有限 $A$-模且
$I$-进完备，所以 $M_1^\wedge = M_1$。结合证明开头的说明，
这就对 $M_1$ 证明了断言。对 $M_2$，有
$H^i_J(M_2) = H^i_{I + J}(M) = H^i_{I + J'}(M) = H^i_{J'}(M_2)$
被 $I$ 的某个幂零化，因而导出完备。所以该情形中的断言也成立。
对 $N$，可以使用刚才两个论证中的任一个。考察短正合列
$0 \to M' \to M \to M/M' \to 0$
可知，只需对 $M/M'$ 证明引理。因此可假设
$\text{Ass}(M) \cap (V(I) \cup T) = \emptyset$。

\medskip\noindent
设 $\mathfrak p \in \text{Ass}(M)$ 满足
$V(\mathfrak p) \cap V(J) \cap V(I) = \emptyset$.
由于 $I$ 包含在 $A$ 的 Jacobson 根中，这蕴含
$V(\mathfrak p) \cap V(J') = \emptyset$ 对任意
$J' \subset J$ 都成立，只要 $V(J') \cap V(I) = V(J) \cap V(I)$。
置 $N = H^0_\mathfrak p(M)$，则
$R\Gamma_J(N) = R\Gamma_{J'}(N) = 0$ 对所有
$J' \subset J$ 都成立，只要 $V(J') \cap V(I) = V(J) \cap V(I)$。
特别地，$R\Gamma_J(N)^\wedge = 0$。
因此可以把 $M$ 替换为 $M/N$，因为这只在条件 (4) 或 (5)
不起作用的素理想处改变 $M$ 的结构。重复此过程，可假设
$V(\mathfrak p) \cap V(J) \cap V(I) \not = \emptyset$
对所有 $\mathfrak p \in \text{Ass}(M)$ 成立。

\medskip\noindent
假设 $\text{Ass}(M) \cap (V(I) \cup T) = \emptyset$，并且
$V(\mathfrak p) \cap V(J) \cap V(I) \not = \emptyset$
对所有 $\mathfrak p \in \text{Ass}(M)$ 成立。
设 $\mathfrak p \in \text{Ass}(M)$。我们要说明可以应用引理
\ref{algebraization-lemma-kill-completion-general}。正是在这一验证中，我们将使用补充条件 (5)。
选择 $\mathfrak p' \subset \mathfrak p$ 以及
$\mathfrak q' \subset V(\mathfrak p) \cap V(J) \cap V(I)$。
\begin{enumerate}
\item 若 $M_{\mathfrak p'} = 0$，则
$\text{depth}(M_{\mathfrak p'}) = \infty$，并且
$\text{depth}(M_{\mathfrak p'}) +
\dim((A/\mathfrak p')_{\mathfrak q'}) > d + s$.
\item 若 $\text{depth}(M_{\mathfrak p'}) < s$，则
$\text{depth}(M_{\mathfrak p'}) +
\dim((A/\mathfrak p')_{\mathfrak q'}) > d + s$，这是 (4) 的结论。
\end{enumerate}
在其余情形中，有 $M_{\mathfrak p'} \not = 0$ 且
$\text{depth}(M_{\mathfrak p'}) \geq s$。特别地，
$\mathfrak p'$ 属于 $M$ 的支集，并且可以选择
$\mathfrak p'' \subset \mathfrak p'$，使 $\mathfrak p'' \in \text{Ass}(M)$。
\begin{enumerate}
\item[(a)] 注意
$\dim((A/\mathfrak p'')_{\mathfrak p'}) \geq \text{depth}(M_{\mathfrak p'})$
，这是《代数》，引理 \ref{algebra-lemma-depth-dim-associated-primes} 的结论。
若等号成立，则
$$
\text{depth}(M_{\mathfrak p'}) + \dim((A/\mathfrak p')_{\mathfrak q'}) =
\text{depth}(M_{\mathfrak p''}) + \dim((A/\mathfrak p'')_{\mathfrak q'})
> s + d
$$
；这是把 (4) 应用于 $\mathfrak p''$ 所得，故这一情形完成。这意味着，
只有当
$\dim((A/\mathfrak p'')_{\mathfrak p'}) > \text{depth}(M_{\mathfrak p'})$.
时才有困难。这蕴含 $\dim(M_\mathfrak p) \geq s + 2$。
因此，若 (5)(a) 成立，就不会出现这种情形。
\item[(b)] 若 (5)(b) 成立，则
$$
\text{depth}(M_{\mathfrak p'}) + \dim((A/\mathfrak p')_{\mathfrak q'})
\geq s + \delta(\mathfrak p') - \delta(\mathfrak q')
\geq s + 1 + \delta(\mathfrak p) - \delta_{max}
> s + d
$$
，如所需。
\item[(c)] 若 (5)(c) 成立，则
\begin{align*}
\text{depth}(M_{\mathfrak p'}) + \dim((A/\mathfrak p')_{\mathfrak q'})
& \geq
s + \delta(\mathfrak p') - \delta(\mathfrak q') \\
& \geq
s + 1 + \delta(\mathfrak p) - \delta(\mathfrak q') \\
& =
s + 1 + \delta(\mathfrak p) - \delta(\mathfrak q) +
\delta(\mathfrak q) - \delta(\mathfrak q') \\
& >
s + 1 + (s + d + \delta_{max} - \delta_{min} - 2) +
\delta(\mathfrak q) - \delta(\mathfrak q') \\
& \geq 
2s + d - 1 \geq s + d
\end{align*}
，如所需。注意，这一论证成立是因为我们知道存在素理想
$\mathfrak q \in V(\mathfrak p) \cap V(J) \cap V(I)$。
\end{enumerate}
现在可以进行归纳步骤。

\medskip\noindent
选择理想 $J_0$，如引理 \ref{algebraization-lemma-kill-colimit-weak-general} 中所述，
并选择整数 $t > 0$，使 $(J_0I)^t$ 零化 $H^s_J(M)$。
引理 \ref{algebraization-lemma-kill-completion-general} 的假设对每个
$\mathfrak p \in \text{Ass}(M)$ 都成立（见上一段）。
因此，零化理想 $\mathfrak a \subset A$（它零化
$H^s(R\Gamma_J(M)^\wedge)$）不包含于 $\mathfrak p$，其中
$\mathfrak p \in \text{Ass}(M)$。于是可以找到
$f \in \mathfrak a(J_0I)^t$，它不属于 $M$ 的任何相伴素理想，
并且同时零化 $H^s(R\Gamma_J(M)^\wedge)$ 与 $H^s_J(M)$。
于是 $f$ 是 $M$ 上的非零因子，可以考察短正合列
$$
0 \to M \xrightarrow{f} M \to M/fM \to 0
$$
由 $f$ 的选择，得到
$$
\xymatrix{
H^{s - 1}_{J'}(M) \ar[d] \ar[r] &
H^{s - 1}_{J'}(M/fM) \ar[d] \ar[r] &
H^s_{J'}(M) \ar[d] \ar[r] & 0 \\
H^{s - 1}(R\Gamma_J(M)^\wedge) \ar[r] &
H^{s - 1}(R\Gamma_J(M/fM)^\wedge) \ar[r] &
H^s(R\Gamma_J(M)^\wedge) \ar[r] & 0
}
$$
，其中 $J' \subset J_0$ 任意，并满足
$V(J') \cap V(I) = V(J) \cap V(I)$。
因此，若选择 $J'$，使它对 $M$、$M/fM$ 和 $s - 1$ 都适用
（由归纳假设可以做到——见下一段），便可得出引理成立。

\medskip\noindent
为完成证明，须说明模 $M/fM$ 对 $s - 1$ 满足假设 (4) 与 (5)。
因此，设 $\mathfrak p$ 是 $M/fM$ 支集中的素理想，满足
$\text{depth}((M/fM)_\mathfrak p) < s - 1$，并且
$V(\mathfrak p) \cap V(J) \cap V(I)$ 非空。于是
$\dim(M_\mathfrak p) = \dim((M/fM)_\mathfrak p) + 1$，并且
$\text{depth}(M_\mathfrak p) = \text{depth}((M/fM)_\mathfrak p) + 1$。
特别地，我们知道 (4) 与 (5) 对 $\mathfrak p$ 和 $M$
在原来的值 $s$ 下成立。逐项检验即可得到所需不等式。
\end{proof}

\begin{example}
\label{algebraization-example-no-ML}
在引理 \ref{algebraization-lemma-algebraize-local-cohomology-general} 中，
我们不知道逆系统 $H^i_J(M/I^nM)$ 是否满足 Mittag--Leffler 条件。
例如，假设 $A = \mathbf{Z}_p[[x, y]]$、$I = (p)$、
$J = (p, x)$ 且 $M = A/(xy - p)$。则映射
$H^0_J(M/p^nM) \to H^0_J(M/pM)$
的像是由 $y^n$ 生成的理想；该理想位于
$M/pM = A/(p, xy)$ 中。
\end{example}





\section{局部上同调的代数化，II}
\label{algebraization-section-algebraization-punctured}

\noindent
本节重新给出 \ref{algebraization-section-algebraization-sections-general} 节的论证；
此时 $(A, \mathfrak m)$ 是局部环，并取局部上同调
$R\Gamma_\mathfrak m$，它关于 $\mathfrak m$ 定义。
与前面一样，主要工具是下面的引理。

\begin{lemma}
\label{algebraization-lemma-kill-completion}
设 $(A, \mathfrak m)$ 是诺特局部环。设 $I \subset A$ 是理想。
设 $M$ 是有限 $A$-模，$\mathfrak p \subset A$ 是素理想。
设 $s$ 与 $d$ 是整数。假设
\begin{enumerate}
\item $A$ 有对偶化复形，
\item $\text{cd}(A, I) \leq d$，并且
\item
$\text{depth}_{A_\mathfrak p}(M_\mathfrak p) + \dim(A/\mathfrak p) > d + s$.
\end{enumerate}
则存在 $f \in A \setminus \mathfrak p$，它零化
$H^i(R\Gamma_\mathfrak m(M)^\wedge)$，其中 $i \leq s$；这里 ${}^\wedge$
表示 $I$-进完备化。
\end{lemma}

\begin{proof}
根据《局部上同调》，引理
\ref{local-cohomology-lemma-sitting-in-degrees}，函数
$$
\mathfrak p' \longmapsto
\text{depth}_{A_{\mathfrak p'}}(M_{\mathfrak p'}) + \dim(A/\mathfrak p')
$$
在 $\Spec(A)$ 上下半连续。因此，该函数在
$\mathfrak p' \subset \mathfrak p$ 上的值为 $> s + d$。
所以，本引理是引理 \ref{algebraization-lemma-kill-completion-general} 的特殊情形，
只要 $\mathfrak p \not = \mathfrak m$。
若 $\mathfrak p = \mathfrak m$，则由深度与局部上同调的关系，
$H^i_\mathfrak m(M) = 0$ 对 $i \leq s + d$ 成立
（《对偶化复形》，引理 \ref{dualizing-lemma-depth}）。
因此，引理 \ref{algebraization-lemma-kill-completion-general} 证明中的论证表明，
在这个（退化）情形中 $H^i(R\Gamma_\mathfrak m(M)^\wedge) = 0$
对 $i \leq s$ 成立。
\end{proof}

\begin{lemma}
\label{algebraization-lemma-kill-colimit-weak}
设 $(A, \mathfrak m)$ 是诺特局部环。设 $I \subset A$ 是理想。
设 $M$ 是有限 $A$-模。设 $s$ 与 $d$ 是整数。假设
\begin{enumerate}
\item $A$ 有对偶化复形，
\item 若 $\mathfrak p \in V(I)$，则不加条件，
\item 若 $\mathfrak p \not \in V(I)$ 且
$V(\mathfrak p) \cap V(I) = \{\mathfrak m\}$，则
$\dim(A/\mathfrak p) \leq d$,
\item 若 $\mathfrak p \not \in V(I)$ 且
$V(\mathfrak p) \cap V(I) \not = \{\mathfrak m\}$，则
$$
\text{depth}_{A_\mathfrak p}(M_\mathfrak p) \geq s
\quad\text{或}\quad
\text{depth}_{A_\mathfrak p}(M_\mathfrak p) + \dim(A/\mathfrak p) > d + s
$$
\end{enumerate}
则存在理想 $J_0 \subset A$，满足
$V(J_0) \cap V(I) = \{\mathfrak m\}$，使得对任意 $J \subset J_0$，若
$V(J) \cap V(I) = \{\mathfrak m\}$，则映射
$$
R\Gamma_J(M) \longrightarrow R\Gamma_{J_0}(M)
$$
在次数 $\leq s$ 的上同调上诱导同构，并且这些模被 $J_0I$ 的某个幂零化。
\end{lemma}

\begin{proof}
这是引理 \ref{algebraization-lemma-kill-colimit-weak-general} 的特殊情形。
\end{proof}

\begin{lemma}
\label{algebraization-lemma-kill-colimit}
在引理 \ref{algebraization-lemma-kill-colimit-weak} 中，若把空条件 (2) 替换为假设
\begin{enumerate}
\item[(2')] 若 $\mathfrak p \in V(I)$ 且 $\mathfrak p \not = \mathfrak m$，
则 $\text{depth}_{A_\mathfrak p}(M_\mathfrak p) + \dim(A/\mathfrak p) > s$，
\end{enumerate}
则这些条件还蕴含 $H^i_{J_0}(M)$ 是有限 $A$-模，其中 $i \leq s$。
\end{lemma}

\begin{proof}
这是引理 \ref{algebraization-lemma-kill-colimit-general} 的特殊情形。
\end{proof}

\begin{lemma}
\label{algebraization-lemma-kill-colimit-support}
若在引理 \ref{algebraization-lemma-kill-colimit-weak} 中还假设
\begin{enumerate}
\item[(6)] 若 $\mathfrak p \not \in V(I)$ 且
$V(\mathfrak p) \cap V(I) = \{\mathfrak m\}$，则
$\text{depth}_{A_\mathfrak p}(M_\mathfrak p) > s$,
\end{enumerate}
则 $H^i_{J_0}(M) = H^i_J(M) = H^i_\mathfrak m(M)$ 对 $i \leq s$ 成立，
并且这些模被 $I$ 的某个幂零化。
\end{lemma}

\begin{proof}
这是引理 \ref{algebraization-lemma-kill-colimit-support-general} 的特殊情形。
\end{proof}

\begin{lemma}
\label{algebraization-lemma-algebraize-local-cohomology}
设 $(A, \mathfrak m)$ 是诺特局部环。设 $I \subset A$ 是理想。
设 $M$ 是有限 $A$-模。设 $s$ 与 $d$ 是整数。假设
\begin{enumerate}
\item $A$ 是 $I$-进完备的，并且有对偶化复形，
\item 若 $\mathfrak p \in V(I)$，则不加条件，
\item $\text{cd}(A, I) \leq d$,
\item 若 $\mathfrak p \not \in V(I)$ 且
$V(\mathfrak p) \cap V(I) \not = \{\mathfrak m\}$，则
$$
\text{depth}_{A_\mathfrak p}(M_\mathfrak p) \geq s
\quad\text{或}\quad
\text{depth}_{A_\mathfrak p}(M_\mathfrak p) + \dim(A/\mathfrak p) > d + s
$$
\end{enumerate}
则存在理想 $J_0 \subset A$，满足
$V(J_0) \cap V(I) = \{\mathfrak m\}$，使得对任意 $J \subset J_0$，若
$V(J) \cap V(I) = \{\mathfrak m\}$，则映射
$$
R\Gamma_J(M) \longrightarrow
R\Gamma_J(M)^\wedge = R\Gamma_\mathfrak m(M)^\wedge
$$
在次数 $\leq s$ 的上同调上诱导同构。
这里 ${}^\wedge$ 表示导出 $I$-进完备化。
\end{lemma}

\begin{proof}
本引理是引理 \ref{algebraization-lemma-algebraize-local-cohomology-general} 的特殊情形，
因为在 $\delta_{max} = \delta_{min} = \delta(\mathfrak m)$ 时，
条件 (4) 蕴含条件 (5)(c)。我们仍将给出这个重要特殊情形的证明，
因为它略为简单（需要检验的事项更少）。

\medskip\noindent
在当前情形中，$R\Gamma_\mathfrak a$ 与
$R\Gamma_{V(\mathfrak a)}$ 没有区别；见《对偶化复形》，引理
\ref{dualizing-lemma-local-cohomology-noetherian}。其次，注意
$$
R\Gamma_\mathfrak m(M)^\wedge =
R\Gamma_I(R\Gamma_J(M))^\wedge =
R\Gamma_J(M)^\wedge
$$
，这是《对偶化复形》，引理
\ref{dualizing-lemma-local-cohomology-ss} 与
\ref{dualizing-lemma-complete-and-local} 的结论，
也解释了引理陈述中的等号。

\medskip\noindent
注意，引理对 $s < 0$ 成立。这并非平凡情形，因为不能先验地看出
$H^s(R\Gamma_\mathfrak m(M)^\wedge)$ 在负的 $s$ 上为零。
不过，这一消失已在引理
\ref{algebraization-lemma-local-cohomology-derived-completion} 中证明。
对于 $s \geq 0$，我们将用归纳法证明引理。

\medskip\noindent
引理 \ref{algebraization-lemma-kill-colimit-weak} 的假设由《局部上同调》，引理
\ref{local-cohomology-lemma-cd-bound-dim-local} 保证。
$s = 0$ 时的引理由引理 \ref{algebraization-lemma-kill-colimit-weak} 和《对偶化复形》，
引理 \ref{dualizing-lemma-completion-local-H0} 得出。

\medskip\noindent
假设 $s > 0$，且引理对更小的 $s$ 成立。设 $M' \subset M$
是由下述元素组成的子模：其支集包含于 $V(I) \cup V(J)$，
其中某个理想 $J$ 满足 $V(J) \cap V(I) = \{\mathfrak m\}$。
则 $M'$ 是有限 $A$-模。我们断言
$$
R\Gamma_J(M') \to R\Gamma_\mathfrak m(M')^\wedge
$$
对 $J$ 的任意选择都是同构。事实上，对任意这样的模，都有短正合列
$0 \to M_1 \oplus M_2 \to M' \to N \to 0$，其中
$M_1$ 被 $J$ 的某个幂零化，$M_2$ 被 $I$ 的某个幂零化，
而 $N$ 被 $\mathfrak m$ 的某个幂零化。对 $M_1$，有
$R\Gamma_J(M_1) = M_1$；又因为 $M_1$ 是有限 $A$-模且
$I$-进完备，所以 $M_1^\wedge = M_1$。因此断言对 $M_1$ 成立。
对 $M_2$，$H^i_J(M_2)$ 被 $I$ 的某个幂零化，因而导出完备；
所以断言对 $M_2$ 也成立。对 $N$ 可用相同论证，故断言成立。
考察短正合列 $0 \to M' \to M \to M/M' \to 0$，
可知只需对 $M/M'$ 证明引理。因此可假设
$\mathfrak p \in \text{Ass}(M)$ 蕴含
$V(\mathfrak p) \cap V(I) \not = \{\mathfrak m\}$，即
$\mathfrak p$ 是 (4) 中所述的素理想。

\medskip\noindent
选择理想 $J_0$，如引理 \ref{algebraization-lemma-kill-colimit-weak} 中所述；
再选择整数 $t > 0$，使 $(J_0I)^t$ 零化 $H^s_J(M)$。
这里 $J$ 表示任意理想 $J \subset J_0$，并满足
$V(J) \cap V(I) = \{\mathfrak m\}$。
引理 \ref{algebraization-lemma-kill-completion} 的假设对每个
$\mathfrak p \in \text{Ass}(M)$ 都成立（见上一段）。因此，
零化理想 $\mathfrak a \subset A$（它零化
$H^s(R\Gamma_\mathfrak m(M)^\wedge)$）不包含于 $\mathfrak p$，
其中 $\mathfrak p \in \text{Ass}(M)$。于是可以找到
$f \in \mathfrak a(J_0I)^t$，它不属于 $M$ 的任何相伴素理想，
并且同时零化 $H^s(R\Gamma_\mathfrak m(M)^\wedge)$ 与 $H^s_J(M)$。
于是 $f$ 是 $M$ 上的非零因子，可以考察短正合列
$$
0 \to M \xrightarrow{f} M \to M/fM \to 0
$$
由 $f$ 的选择，得到
$$
\xymatrix{
H^{s - 1}_J(M) \ar[d] \ar[r] &
H^{s - 1}_J(M/fM) \ar[d] \ar[r] &
H^s_J(M) \ar[d] \ar[r] & 0 \\
H^{s - 1}(R\Gamma_\mathfrak m(M)^\wedge) \ar[r] &
H^{s - 1}(R\Gamma_\mathfrak m(M/fM)^\wedge) \ar[r] &
H^s(R\Gamma_\mathfrak m(M)^\wedge) \ar[r] & 0
}
$$
，其中 $J \subset J_0$ 任意，并满足
$V(J) \cap V(I) = \{\mathfrak m\}$。
因此，若选择 $J$，使它对 $M$、$M/fM$ 与 $s - 1$
都适用（由归纳假设可以做到），便可得出引理成立。
\end{proof}








\section{局部上同调的代数化，III}
\label{algebraization-section-bootstrap}

\noindent
本节利用第 \ref{algebraization-section-algebraization-sections-general} 节和
第 \ref{algebraization-section-algebraization-punctured} 节中的材料作进一步提升，
在如下情形中得到一个更强的结果。

\begin{situation}
\label{algebraization-situation-bootstrap}
这里 $A$ 是诺特环。给定理想 $I \subset A$、有限 $A$-模 $M$，
以及对特殊化稳定的子集 $T \subset V(I)$。另给定整数 $s$ 和 $d$。
我们假设
\begin{enumerate}
\item[(1)] $A$ 有对偶化复形，
\item[(3)] $\text{cd}(A, I) \leq d$,
\item[(4)] 对任意素理想 $\mathfrak p \subset \mathfrak r \subset \mathfrak q$，
若 $\mathfrak p \not \in V(I)$、
$\mathfrak r \in V(I) \setminus T$ 且
$\mathfrak q \in T$，则
$$
\text{depth}_{A_\mathfrak p}(M_\mathfrak p) \geq s
\quad\text{或}\quad
\text{depth}_{A_\mathfrak p}(M_\mathfrak p) +
\dim((A/\mathfrak p)_\mathfrak q) > d + s
$$
\item[(6)] 对任意 $\mathfrak q \in T$，以
$A', \mathfrak m', I', M'$ 分别表示关于 $I$ 的通常的进完备化，
它们来自 $A_\mathfrak q, \mathfrak qA_\mathfrak q, I_\mathfrak q, M_\mathfrak q$，
则
$$
\text{depth}(M'_{\mathfrak p'}) > s
$$
对所有 $\mathfrak p' \in \Spec(A') \setminus V(I')$ 中满足
$V(\mathfrak p') \cap V(I') = \{\mathfrak m'\}$ 者成立。
\end{enumerate}
\end{situation}

\noindent
下面的引理说明，在情形 \ref{algebraization-situation-bootstrap} 中，
在 (4) 中考察素理想三元组
$\mathfrak p \subset \mathfrak r \subset \mathfrak q$ 已经足够，
尽管实际的假设只涉及 $\mathfrak p$ 和 $\mathfrak q$。

\begin{lemma}
\label{algebraization-lemma-helper-bootstrap}
在情形 \ref{algebraization-situation-bootstrap} 中，设 $\mathfrak p \subset \mathfrak q$
是 $A$ 的素理想，且 $\mathfrak p \not \in V(I)$、
$\mathfrak q \in T$。若不存在满足
$\mathfrak r \in V(I) \setminus T$ 且
$\mathfrak p \subset \mathfrak r \subset \mathfrak q$ 的素理想，
则 $\text{depth}(M_\mathfrak p) > s$。
\end{lemma}

\begin{proof}
选取 $\mathfrak q' \in T$，使得
$\mathfrak p \subset \mathfrak q' \subset \mathfrak q$，
并且 $T$ 中不存在严格位于 $\mathfrak p$ 与 $\mathfrak q'$ 之间的
素理想。为证明该引理，我们可以并且确实把 $\mathfrak q$
替换为 $\mathfrak q'$。接着，令 $\mathfrak p' \subset A_\mathfrak q$
为与 $\mathfrak p$ 对应的素理想。由于不存在引理中所述的
素理想，这样便得到
$V(\mathfrak p') \cap V(IA_\mathfrak q) = \{\mathfrak q A_\mathfrak q\}$；
这里所用的正是这样的 $\mathfrak r$ 不存在这一事实。
令 $A', I', \mathfrak m', M'$ 为关于 $I$ 的进完备化，
它们分别来自
$A_\mathfrak q, I_\mathfrak q, \mathfrak qA_\mathfrak q, M_\mathfrak q$。
由于 $A_\mathfrak q \to A'$ 忠实平坦
（代数，引理 \ref{algebra-lemma-completion-faithfully-flat}），
我们可以选取 $\mathfrak p'' \subset A'$，它位于 $\mathfrak p'$ 之上，并使得
$\dim(A'_{\mathfrak p''}/\mathfrak p' A'_{\mathfrak p''}) = 0$。
于是可见
$$
\text{depth}(M'_{\mathfrak p''}) =
\text{depth}((M_\mathfrak q \otimes_{A_\mathfrak q} A')_{\mathfrak p''}) =
\text{depth}(M_\mathfrak p \otimes_{A_\mathfrak p} A'_{\mathfrak p''}) =
\text{depth}(M_\mathfrak p)
$$
这是由 $A \to A'$ 的平坦性以及对 $\mathfrak p''$ 的选取所得；参见
代数，引理 \ref{algebra-lemma-apply-grothendieck-module}。
由于 $\mathfrak p''$ 位于 $\mathfrak p'$ 之上，我们有
$V(\mathfrak p'') \cap V(I') = \{\mathfrak m'\}$。因此，
情形 \ref{algebraization-situation-bootstrap} 中的条件 (6) 蕴含
$\text{depth}(M'_{\mathfrak p''}) > s$，证明完毕。
\end{proof}

\noindent
下面这个冗长的引理说明可能施加的各组条件之间的关系。

\begin{lemma}
\label{algebraization-lemma-bootstrap-inherited}
在情形 \ref{algebraization-situation-bootstrap} 中，有
\begin{enumerate}
\item[(E)] 若 $T' \subset T$ 是更小的、对特殊化稳定的子集，则
$A, I, T', M$ 满足情形 \ref{algebraization-situation-bootstrap} 的假设；
\item[(F)] 若 $S \subset A$ 是乘法子集，则
$S^{-1}A, S^{-1}I, T', S^{-1}M$
满足情形 \ref{algebraization-situation-bootstrap} 的假设，
其中 $T' \subset V(S^{-1}I)$ 是 $T$ 的逆像；
\item[(G)] 四元组 $A', I', T', M'$
满足情形 \ref{algebraization-situation-bootstrap} 的假设，
其中 $A', I', M'$ 是通常的 $I$-进完备化，
它们分别来自 $A, I, M$，
而 $T' \subset V(I')$ 是 $T$ 的逆像。
\end{enumerate}
设 $I \subset \mathfrak a \subset A$ 是满足
$V(\mathfrak a) \subset T$ 的理想。则
\begin{enumerate}
\item[(A)] 若 $I$ 包含于 $A$ 的 Jacobson 根，则对于
$A, I, \mathfrak a, M$，引理 \ref{algebraization-lemma-kill-colimit-weak-general} 和
\ref{algebraization-lemma-kill-colimit-support-general} 的全部假设均成立；
\item[(B)] 若 $A$ 关于 $I$ 完备，则对于
$A, I, \mathfrak a, M$，引理
\ref{algebraization-lemma-algebraize-local-cohomology-general}
除可能的 (5) 外全部假设均成立；
\item[(C)] 若 $A$ 是以 $\mathfrak m = \mathfrak a$ 为极大理想的局部环，
则对于 $A, \mathfrak m, I, M$，
引理 \ref{algebraization-lemma-kill-colimit-weak} 和 \ref{algebraization-lemma-kill-colimit-support}
的全部假设均成立；
\item[(D)] 若 $A$ 是以 $\mathfrak m = \mathfrak a$ 为极大理想的局部环，
且 $I$-进完备，则对于 $A, \mathfrak m, I, M$，
引理 \ref{algebraization-lemma-algebraize-local-cohomology}
的全部假设均成立。
\end{enumerate}
\end{lemma}

\begin{proof}
证明 (E)。我们须证明情形 \ref{algebraization-situation-bootstrap} 的假设
(1)、(3)、(4)、(6) 对 $A, I, T, M$ 成立。把 $T$ 缩小为 $T'$
会减弱假设 (6) 而加强假设 (4)。不过，若有假设 (4) 中所述的
$\mathfrak p \subset \mathfrak r \subset \mathfrak q$，其中
$\mathfrak p \not \in V(I)$、$\mathfrak r \in V(I) \setminus T'$、
$\mathfrak q \in T'$，这正是 $A, I, T', M$ 的假设 (4)；那么要么可以选取
$\mathfrak r \in V(I) \setminus T$，此时可用 $A, I, T, M$ 的条件 (4)；
要么找不到这样的 $\mathfrak r$，此时由引理
\ref{algebraization-lemma-helper-bootstrap} 得到
$\text{depth}(M_\mathfrak p) > s$。
这便按要求证明了 (4) 对 $A, I, T', M$ 成立。

\medskip\noindent
证明 (F)。这是直接的，细节从略。

\medskip\noindent
证明 (G)。我们须证明情形 \ref{algebraization-situation-bootstrap} 的假设
(1)、(3)、(4)、(6) 对 $I$-进完备化
$A', I', T', M'$ 成立。请注意，
$\Spec(A') \to \Spec(A)$ 诱导同构 $V(I') \to V(I)$。

\medskip\noindent
假设 (1)：环 $A'$ 有对偶化复形；参见
对偶化复形，引理 \ref{dualizing-lemma-ubiquity-dualizing}。

\medskip\noindent
假设 (3)：由于 $I' = IA'$，这可由局部上同调，
引理 \ref{local-cohomology-lemma-cd-change-rings} 得出。

\medskip\noindent
假设 (4)：若有素理想
$\mathfrak p' \subset \mathfrak r' \subset \mathfrak q'$ 属于 $A'$，并满足
$\mathfrak p' \not \in V(I')$、
$\mathfrak r' \in V(I') \setminus T'$、
$\mathfrak q' \in T'$，则它们的像
$\mathfrak p \subset \mathfrak r \subset \mathfrak q$ 位于 $A$ 的谱中，并
满足
$\mathfrak p \not \in V(I)$、$\mathfrak r \in V(I) \setminus T$、
$\mathfrak q \in T$。
于是
$$
\text{depth}_{A_\mathfrak p}(M_\mathfrak p) \geq s
\quad\text{或}\quad
\text{depth}_{A_\mathfrak p}(M_\mathfrak p) +
\dim((A/\mathfrak p)_\mathfrak q) > d + s
$$
这是由 $A, I, T, M$ 的假设 (4) 得到的。又有
$\text{depth}(M'_{\mathfrak p'}) \geq \text{depth}(M_\mathfrak p)$，且
$\text{depth}(M'_{\mathfrak p'}) +
\dim((A'/\mathfrak p')_{\mathfrak q'}) =
\text{depth}(M_\mathfrak p) +
\dim((A/\mathfrak p)_\mathfrak q)$
由局部上同调，引理 \ref{local-cohomology-lemma-change-completion} 可得。
因此假设 (4) 对 $A', I', T', M'$ 成立。

\medskip\noindent
假设 (6)：令 $\mathfrak q' \in T'$ 位于素理想
$\mathfrak q \in T$ 之上。则 $A'_{\mathfrak q'}$
与 $A_\mathfrak q$ 有同构的 $I$-进完备化，
$M_\mathfrak q$ 与 $M'_{\mathfrak q'}$ 也同样如此。
因此，$A', I', T', M'$ 的假设 (6) 等价于
$A, I, T, M$ 的假设 (6)。

\medskip\noindent
证明 (A)。我们须对 $(A, I, \mathfrak a, M)$ 检验
引理 \ref{algebraization-lemma-kill-colimit-weak-general} 和
\ref{algebraization-lemma-kill-colimit-support-general} 的条件 (1)、(2)、(3)、(4)、(6)。
注意：这些引理陈述中的集合 $T$ 与上面的集合 $T$ 不同。

\medskip\noindent
条件 (1)：成立，因为我们在情形 \ref{algebraization-situation-bootstrap} 中假设了
$A$ 有对偶化复形。

\medskip\noindent
条件 (2)：此条件无内容。

\medskip\noindent
条件 (3)：设 $\mathfrak p \subset A$ 满足
$V(\mathfrak p) \cap V(I) \subset V(\mathfrak a)$。
由于 $I$ 包含于 $A$ 的 Jacobson 根，可见
$V(\mathfrak p) \cap V(I) \not = \emptyset$。
令 $\mathfrak q \in V(\mathfrak p) \cap V(I)$ 为一个一般点。
由于 $\text{cd}(A_\mathfrak q, I_\mathfrak q) \leq d$
（局部上同调，引理 \ref{local-cohomology-lemma-cd-local}），并且
$V(\mathfrak p A_\mathfrak q) \cap V(I_\mathfrak q) =
\{\mathfrak q A_\mathfrak q\}$，由局部上同调，
引理 \ref{local-cohomology-lemma-cd-bound-dim-local} 得到
$\dim((A/\mathfrak p)_\mathfrak q) \leq d$，从而证明 (3)。

\medskip\noindent
条件 (4)：设 $\mathfrak p \not \in V(I)$ 且
$\mathfrak q \in V(\mathfrak p) \cap V(\mathfrak a)$。
只须证明
$$
\text{depth}_{A_\mathfrak p}(M_\mathfrak p) \geq s
\quad\text{或}\quad
\text{depth}_{A_\mathfrak p}(M_\mathfrak p) +
\dim((A/\mathfrak p)_\mathfrak q) > d + s
$$
若存在素理想 $\mathfrak p \subset \mathfrak r \subset \mathfrak q$，
其中 $\mathfrak r \in V(I) \setminus T$，则这立即由情形
\ref{algebraization-situation-bootstrap} 中的假设 (4) 得出。
若不存在，则由引理 \ref{algebraization-lemma-helper-bootstrap} 有
$\text{depth}(M_\mathfrak p) > s$。

\medskip\noindent
条件 (6)：设 $\mathfrak p \not \in V(I)$ 满足
$V(\mathfrak p) \cap V(I) \subset V(\mathfrak a)$。
由于 $I$ 包含于 $A$ 的 Jacobson 根，可见
$V(\mathfrak p) \cap V(I) \not = \emptyset$。
选取 $\mathfrak q \in V(\mathfrak p) \cap V(I) \subset V(\mathfrak a)$。
显然不存在素理想
$\mathfrak p \subset \mathfrak r \subset \mathfrak q$，
使得 $\mathfrak r \in V(I) \setminus T$。
由引理 \ref{algebraization-lemma-helper-bootstrap} 有
$\text{depth}(M_\mathfrak p) > s$，这证明了 (6)。

\medskip\noindent
证明 (B)。我们须检验引理
\ref{algebraization-lemma-algebraize-local-cohomology-general} 的条件 (1)、(2)、(3)、(4)。
注意：该引理陈述中的集合 $T$ 与上面的集合 $T$ 不同。

\medskip\noindent
条件 (1)：成立，因为 $A$ 完备且有对偶化复形。

\medskip\noindent
条件 (2)：此条件无内容。

\medskip\noindent
条件 (3)：这与情形 \ref{algebraization-situation-bootstrap} 中的假设 (3) 相同。

\medskip\noindent
条件 (4)：这与引理 \ref{algebraization-lemma-kill-colimit-weak-general} 中的假设 (4)
相同，而我们已在 (A) 中证明了后者。

\medskip\noindent
证明 (C)。这是因为引理 \ref{algebraization-lemma-kill-colimit-weak} 和
\ref{algebraization-lemma-kill-colimit-support} 中的假设，与引理
\ref{algebraization-lemma-kill-colimit-weak-general} 和
\ref{algebraization-lemma-kill-colimit-support-general} 在局部情形下的假设相同，
而我们已在 (A) 中证明了这些假设成立。

\medskip\noindent
证明 (D)。这是因为引理
\ref{algebraization-lemma-algebraize-local-cohomology} 中的假设，与引理
\ref{algebraization-lemma-algebraize-local-cohomology-general}
中的假设 (1)、(2)、(3)、(4) 相同，而我们已在 (B) 中证明了这些假设成立。
\end{proof}

\begin{lemma}
\label{algebraization-lemma-algebraize-local-cohomology-bis}
在情形 \ref{algebraization-situation-bootstrap} 中，假设 $A$ 是以
$\mathfrak m$ 为极大理想的局部环，并且 $T = \{\mathfrak m\}$。则
$H^i_\mathfrak m(M) \to \lim H^i_\mathfrak m(M/I^nM)$
在 $i \leq s$ 时为同构，并且这些模都被 $I$ 的某个幂零化。
\end{lemma}

\begin{proof}
令 $A', I', \mathfrak m', M'$ 为通常的 $I$-进完备化，
分别来自 $A, I, \mathfrak m, M$。回忆我们有
$H^i_\mathfrak m(M) \otimes_A A' = H^i_{\mathfrak m'}(M')$
这是由 $A \to A'$ 的平坦性以及对偶化复形，引理
\ref{dualizing-lemma-torsion-change-rings} 得到的。
由于 $H^i_\mathfrak m(M)$ 是 $\mathfrak m$-幂挠模，故有
$H^i_\mathfrak m(M) = H^i_\mathfrak m(M) \otimes_A A'$；参见
代数进阶，引理 \ref{more-algebra-lemma-neighbourhood-equivalence}。
由此得到 $H^i_\mathfrak m(M) = H^i_{\mathfrak m'}(M')$。
完全相同的论证表明
$H^i_\mathfrak m(M/I^nM) =  H^i_{\mathfrak m'}(M'/(I')^nM')$
对所有 $n$ 和 $i$ 成立。

\medskip\noindent
由引理 \ref{algebraization-lemma-algebraize-local-cohomology}、
\ref{algebraization-lemma-kill-colimit-weak} 和
\ref{algebraization-lemma-kill-colimit-support}
以及引理 \ref{algebraization-lemma-bootstrap-inherited} 的 (C)、(D) 部分，
上述引理适用于 $A', \mathfrak m', I', M'$。
因此得到同构
$$
H^i_{\mathfrak m'}(M') \longrightarrow H^i(R\Gamma_{\mathfrak m'}(M')^\wedge)
$$
其中 $i \leq s$，而 ${}^\wedge$ 表示导出 $I'$-进完备化；
这些模都被 $I'$ 的某个幂零化。
由引理 \ref{algebraization-lemma-local-cohomology-derived-completion}
得到同构
$$
H^i_{\mathfrak m'}(M') \longrightarrow
\lim H^i_{\mathfrak m'}(M'/(I')^nM'))
$$
其中 $i \leq s$。结合已经建立的与 $A$ 上局部上同调的比较，
便得出该引理。
\end{proof}

\begin{lemma}
\label{algebraization-lemma-bootstrap-bis-bis}
设 $I \subset \mathfrak a$ 是诺特环 $A$ 的理想。
设 $M$ 是有限 $A$-模，并设 $s$ 和 $d$ 为整数。
若假设
\begin{enumerate}
\item[(a)] $A$ 有对偶化复形；
\item[(b)] $\text{cd}(A, I) \leq d$,
\item[(c)] 若 $\mathfrak p \not \in V(I)$ 且
$\mathfrak q \in V(\mathfrak p) \cap V(\mathfrak a)$，则
$\text{depth}_{A_\mathfrak p}(M_\mathfrak p) > s$，或
$\text{depth}_{A_\mathfrak p}(M_\mathfrak p) +
\dim((A/\mathfrak p)_\mathfrak q) > d + s$.
\end{enumerate}
则 $A, I, V(\mathfrak a), M, s, d$ 满足
情形 \ref{algebraization-situation-bootstrap} 的条件。
\end{lemma}

\begin{proof}
我们须证明情形 \ref{algebraization-situation-bootstrap} 的假设 (1)、(3)、(4)、(6) 成立。
显然 (a) $\Rightarrow$ (1)、
(b) $\Rightarrow$ (3)，且 (c) $\Rightarrow$ (4)。
为完成证明，下一段证明 (6) 成立。

\medskip\noindent
设 $\mathfrak q \in V(\mathfrak a)$。
用 $A', I', \mathfrak m', M'$ 表示关于 $I$ 的进完备化，
它们分别来自
$A_\mathfrak q, I_\mathfrak q, \mathfrak qA_\mathfrak q, M_\mathfrak q$。
设 $\mathfrak p' \subset A'$ 是满足
$V(\mathfrak p') \cap V(I') = \{\mathfrak m'\}$ 的非极大素理想。
由局部上同调，引理 \ref{local-cohomology-lemma-cd-bound-dim-local}，
注意到这蕴含 $\dim(A'/\mathfrak p') \leq d$。
记 $\mathfrak p \subset A$ 为 $\mathfrak p'$ 的像。我们有
$\text{depth}(M'_{\mathfrak p'}) \geq \text{depth}(M_\mathfrak p)$，且
$\text{depth}(M'_{\mathfrak p'}) +
\dim(A'/\mathfrak p') =
\text{depth}(M_\mathfrak p) +
\dim((A/\mathfrak p)_\mathfrak q)$
这是由局部上同调，引理 \ref{local-cohomology-lemma-change-completion} 得到的。
由假设 (c)，要么有
$\text{depth}(M'_{\mathfrak p'}) \geq \text{depth}(M_\mathfrak p) > s$
而结论已经得到；要么有
$\text{depth}(M'_{\mathfrak p'}) +
\dim(A'/\mathfrak p') > s + d$，结合已经证明的不等式，便蕴含
$\text{depth}(M'_{\mathfrak p'}) > s$，因为
$\dim(A'/\mathfrak p') \leq d$。两种情形都得到所需结论。
\end{proof}

\begin{lemma}
\label{algebraization-lemma-bootstrap}
在情形 \ref{algebraization-situation-bootstrap} 中，逆系统
$\{H^i_T(I^nM)\}_{n \geq 0}$ 在 $i \leq s$ 时为 pro-零。
此外，存在整数 $m_0$，使得对所有
$m \geq m_0$，都存在整数 $m'(m) \geq m$，满足：当
$k \geq m'(m)$ 时，映射
$H^{s + 1}_T(I^kM) \to H^{s + 1}_T(I^mM)$
的像单射映入 $H^{s + 1}_T(I^{m_0}M)$。
\end{lemma}

\begin{proof}
固定 $m$。设 $\mathfrak q \in T$。
由引理 \ref{algebraization-lemma-bootstrap-inherited} 和
\ref{algebraization-lemma-algebraize-local-cohomology-bis}
可知
$$
H^i_\mathfrak q(M_\mathfrak q)
\longrightarrow
\lim H^i_\mathfrak q(M_\mathfrak q/I^nM_\mathfrak q)
$$
在 $i \leq s$ 时为同构。逆系统
$\{H^i_\mathfrak q(I^nM_\mathfrak q)\}_{n \geq 0}$ 和
$\{H^i_\mathfrak q(M/I^nM)\}_{n \geq 0}$
对所有 $i$ 都满足 Mittag-Leffler 条件；参见
引理 \ref{algebraization-lemma-ML-local}。因此，考察如下长正合序列组成的逆系统：
$$
0 \to H^0_\mathfrak q(I^nM_\mathfrak q) \to
H^0_\mathfrak q(M_\mathfrak q) \to
H^0_\mathfrak q(M_\mathfrak q/I^nM_\mathfrak q) \to
H^1_\mathfrak q(I^nM_\mathfrak q) \to
H^1_\mathfrak q(M_\mathfrak q) \to \ldots
$$
可得（略去一些细节）：存在整数
$m'(m, \mathfrak q) \geq m$，使得对所有 $k \geq m'(m, \mathfrak q)$，
映射
$H^i_\mathfrak q(I^kM_\mathfrak q) \to H^i_\mathfrak q(I^mM_\mathfrak q)$
在 $i \leq s$ 时为零，并且映射
$H^{s + 1}_\mathfrak q(I^kM_\mathfrak q) \to
H^{s + 1}_\mathfrak q(I^mM_\mathfrak q)$
的像与 $k \geq m'(m, \mathfrak q)$ 的选取无关，并单射映入
$H^{s + 1}_\mathfrak q(M_\mathfrak q)$。

\medskip\noindent
假设能够证明 $m'(m, \mathfrak q)$ 的选取
与 $\mathfrak q \in T$ 无关。
那么该引理立即由局部上同调，引理 \ref{local-cohomology-lemma-zero} 和
\ref{local-cohomology-lemma-essential-image}.

\medskip\noindent
设 $\omega_A^\bullet$ 为对偶化复形。令
$\delta : \Spec(A) \to \mathbf{Z}$ 为相应的维数函数。
回忆 $\delta$ 只取有限多个值；参见
对偶化复形，引理 \ref{dualizing-lemma-universally-catenary}。
断言：对每个 $d \in \mathbf{Z}$，整数
$m'(m, \mathfrak q)$ 可以选得与满足
$\mathfrak q \in T$ 且 $\delta(\mathfrak q) = d$ 的素理想无关。
由上面的论述，该断言显然蕴含引理。

\medskip\noindent
选取 $\mathfrak q \in T$，使得 $\delta(\mathfrak q) = d$。
考察 Ext 模
$$
E(n, j) = \text{Ext}^j_A(I^nM, \omega_A^\bullet)
$$
我们将用到的一个关键性质是：这些都是有限 $A$-模。
由与对偶化复形相伴的维数函数之定义，回忆
$(\omega_A^\bullet)_\mathfrak q[-d]$ 是 $A_\mathfrak q$ 的规范化
对偶化复形；参见对偶化复形，第 \ref{dualizing-section-dimension-function} 节。
局部对偶定理（对偶化复形，引理
\ref{dualizing-lemma-special-case-local-duality}）告诉我们，
$\mathfrak qA_\mathfrak q$-进完备化后的
$E(n, -d - i)_\mathfrak q$ 与
$H^i_\mathfrak q(I^nM_\mathfrak q)$ Matlis 对偶。因此，第一段中
对 $m'(m, \mathfrak q)$ 的选取（其中 $i \leq s$）说明：
当 $k \geq m'(m, \mathfrak q)$ 且 $j \geq -d - s$ 时，映射
$$
E(m, j)_\mathfrak q \to E(k, j)_\mathfrak q
$$
为零。由于这些模是有限的，并且只对有限多个可能的 $j$ 非零
（略去一个小细节），我们可以找到开邻域
$W \subset \Spec(A)$，它包含 $\mathfrak q$，并使得
$$
E(m, j)_{\mathfrak q'} \to E(m'(m, \mathfrak q), j)_{\mathfrak q'}
$$
在 $j \geq -d - s$ 时，对所有 $\mathfrak q' \in W$ 均为零。
于是当然，映射 $E(m, j)_{\mathfrak q'} \to E(k, j)_{\mathfrak q'}$
对 $k \geq m'(m, \mathfrak q)$ 也为零。

\medskip\noindent
对于 $i = s + 1$，它对应于 $j = - d - s - 1$，
由局部对偶和第一段的结果可得
$$
K_{k, \mathfrak q} =
\Ker(E(m, -d - s - 1)_\mathfrak q \to E(k, -d - s - 1)_\mathfrak q)
$$
与 $k \geq m'(m, \mathfrak q)$ 无关，并且
$$
E(0, -d - s - 1)_\mathfrak q \to
E(m, -d - s - 1)_\mathfrak q/K_{m'(m, \mathfrak q), \mathfrak q}
$$
为满射。对 $k \geq m'(m, \mathfrak q)$，令
$$
K_k = \Ker(E(m, -d - s - 1) \to E(k, -d - s - 1))
$$
由于 $K_k$ 是有限模 $E(m, -d - s - 1)$ 的递增子模序列，
可见在稍微增大 $m'(m, \mathfrak q)$ 后，我们可以假设
$K_{m'(m, \mathfrak q)} = K_k$ 对 $k \geq m'(m, \mathfrak q)$ 成立。
必要时进一步缩小 $W$，还可假设
$$
E(0, -d - s - 1)_{\mathfrak q'} \to
E(m, -d - s - 1)_{\mathfrak q'}/K_{m'(m, \mathfrak q), \mathfrak q'}
$$
对所有 $\mathfrak q' \in W$ 都为满射（与前面一样，使用这些模的有限性，
以及该映射在 $\mathfrak q$ 处局部化后为满射）。

\medskip\noindent
任意子集，特别是
$T_d = \{\mathfrak q \in T \text{ 其中 }\delta(\mathfrak q) = d\}$,
作为赋予子空间拓扑的诺特拓扑空间 $\Spec(A)$ 的子集，
仍是诺特的，因而是拟紧的。
上面已经看到，对每个 $\mathfrak q \in T_d$，
都存在开邻域 $W$，使得 $m'(m, \mathfrak q)$
对所有 $\mathfrak q' \in T_d \cap W$ 都适用。
由此可知，可以找到整数 $m'(m, d)$，使得对所有
$\mathfrak q \in T_d$ 都有
$$
E(m, j)_\mathfrak q \to E(m'(m, d), j)_\mathfrak q
$$
在 $j \geq -d - s$ 时为零；并且令
$K_{m'(m, d)} = \Ker(E(m, -d - s - 1) \to E(m'(m, d), -d - s - 1))$
后，有
$$
K_{m'(m, d), \mathfrak q} =
\Ker(E(m, -d - s - 1)_{\mathfrak q} \to E(k, -d - s - 1)_{\mathfrak q})
$$
对所有 $k \geq m'(m, d)$ 成立，并且映射
$$
E(0, -d - s - 1)_\mathfrak q \to
E(m, -d - s - 1)_\mathfrak q/K_{m'(m, d), \mathfrak q}
$$
为满射。再次反向使用局部对偶定理，可得该断言成立。证明完毕。
\end{proof}

\begin{lemma}
\label{algebraization-lemma-final-bootstrap}
在情形 \ref{algebraization-situation-bootstrap} 中，存在整数 $m_0 \geq 0$，使得
\begin{enumerate}
\item $H^i_T(M) \to H^i_T(M/I^nM)$ 对所有 $i \leq s$
及 $n \geq m_0$ 都是单射；
\item $\{H^i_T(M/I^nM)\}_{n \geq 0}$
在 $i < s$ 时满足 Mittag-Leffler 条件；
\item $\{H^i_T(I^{m_0}M/I^nM)\}_{n \geq m_0}$
在 $i \leq s$ 时满足 Mittag-Leffler 条件；
\item $H^i_T(M) \to \lim H^i_T(M/I^nM)$
在 $i < s$ 时为同构；
\item $H^s_T(I^{m_0}M) \to \lim H^s_T(I^{m_0}M/I^nM)$
在 $i \leq s$ 时为同构；
\item $H^s_T(M) \to \lim H^s_T(M/I^nM)$ 为单射，
且其余核被 $I^{m_0}$ 零化；
\item $R^1\lim H^s_T(M/I^nM)$ 被 $I^{m_0}$ 零化。
\end{enumerate}
\end{lemma}

\begin{proof}
为证明 (1)，设 $0 \leq i \leq s$，并选取整数 $m > n > 0$，使得
$H^i_T(I^mM) \to H^i_T(I^nM)$ 为零；由引理
\ref{algebraization-lemma-bootstrap} 可以这样选取。考察各行正合的交换图
$$
\xymatrix{
H^i_T(I^nM) \ar[r] &
H^i_T(M) \ar[r] &
H^i_T(M/I^nM) \\
H^i_T(I^mM) \ar[r] \ar[u]^0 &
H^i_T(M) \ar[r] \ar[u]^1 &
H^i_T(M/I^mM) \ar[u] \\
}
$$
该图表明映射 $H^i_T(M) \to H^i_T(M/I^mM)$ 为单射。因此，
对任意充分大的 $m_0$，(1) 成立。

\medskip\noindent
考察长正合序列
$$
0 \to H^0_T(I^nM) \to H^0_T(M) \to
H^0_T(M/I^nM) \to H^1_T(I^nM) \to
H^1_T(M) \to \ldots
$$
由此及引理 \ref{algebraization-lemma-bootstrap} 得到 (2) 和 (4)。

\medskip\noindent
令 $m_0$ 和 $m'(-)$ 如引理 \ref{algebraization-lemma-bootstrap} 中所述。
对 $m \geq m_0$，考察长正合序列
$$
H^s_T(I^mM) \to H^s_T(I^{m_0}M) \to
H^s_T(I^{m_0}M/I^mM) \to H^{s + 1}_T(I^mM) \to
H^1_T(I^{m_0}M)
$$
于是，对 $k \geq m'(m)$，映射
$H^{s + 1}_T(I^kM) \to H^{s + 1}_T(I^mM)$
的像单射映入 $H^{s + 1}_T(I^{m_0}M)$。
因此，映射
$H^s_T(I^{m_0}M/I^kM) \to H^s_T(I^{m_0}M/I^mM)$
的像映为 $H^{s + 1}_T(I^mM)$ 中的零；这对所有 $k \geq m'(m)$ 成立。
由此得到 (3) 和 (5)。

\medskip\noindent
考察短正合序列
$0 \to I^{m_0}M \to M \to M/I^{m_0} M \to 0$ 和
$0 \to I^{m_0}M/I^nM \to M/I^nM \to M/I^{m_0} M \to 0$.
由此得到图表
$$
\xymatrix{
H^{s - 1}_T(M/I^{m_0}M) \ar[r] &
\lim H^s_T(I^{m_0}M/I^nM) \ar[r] &
\lim H^s_T(M/I^nM) \ar[r] &
H^s_T(M/I^{m_0}M) \\
H^{s - 1}_T(M/I^{m_0}M) \ar[r] \ar@{=}[u] &
H^s_T(I^{m_0}M) \ar[r] \ar[u]_{\cong} &
H^s_T(M) \ar[r] \ar[u] &
H^s_T(M/I^{m_0}M) \ar@{=}[u]
}
$$
其下行正合。由同调，引理
\ref{homology-lemma-apply-Mittag-Leffler}，上行在中间两处也正合。
由此得到 (6)。

\medskip\noindent
记 $B_n = H^s_T(M/I^nM)$。令 $A_n \subset B_n$ 为映射
$H^s_T(I^{m_0}M/I^nM) \to H^s_T(M/I^nM)$ 的像。
由 (3) 和同调，引理 \ref{homology-lemma-Mittag-Leffler}，
$(A_n)$ 满足 Mittag-Leffler 条件。
又有 $C_n = B_n/A_n$ 被 $I^{m_0}$ 零化。因此
$R^1\lim B_n \cong R^1\lim C_n$ 被 $I^{m_0}$ 零化，从而得到 (7)。
\end{proof}

\begin{theorem}
\label{algebraization-theorem-final-bootstrap}
在情形 \ref{algebraization-situation-bootstrap} 中，对 $i \leq s$，逆系统
$\{H^i_T(M/I^nM)\}_{n \geq 0}$ 本质上常值，其值为 $H^i_T(M)$。
特别地，$\{H^i_T(M/I^nM)\}_{n \geq 0}$ 满足 Mittag-Leffler 条件，
$H^i_T(M)$ 被 $I$ 的某个幂零化，并且
$H^i_T(M) = \lim H^i_T(M/I^nM)$.
\end{theorem}

\begin{proof}
设 $0 \leq i \leq s$。
若能证明 $\{H^i_T(M/I^nM)\}_{n \geq 0}$ 满足 Mittag-Leffler 条件并且
$H^i_T(M) = \lim H^i_T(M/I^nM)$，那么引理
\ref{algebraization-lemma-final-bootstrap} 中证明的映射
$H^i_T(M) \to H^i_T(M/I^nM)$ 在 $n \gg 0$ 时的单射性将蕴含：
$\{H^i_T(M/I^nM)\}_{n \geq 0}$ 本质上常值，其值为 $H^i_T(M)$。
略去一个小细节。

\medskip\noindent
由引理 \ref{algebraization-lemma-final-bootstrap}，我们已经知道
$\{H^i_T(M/I^nM)\}_{n \geq 0}$ 满足 Mittag-Leffler 条件，并且
$H^i_T(M) = \lim H^i_T(M/I^nM)$ 对 $i < s$ 成立。
我们还知道，对 $i = s$ 以及模 $I^{m_0}M$，这些陈述成立，
其中某个 $m_0 > 0$。为完成证明，我们将说明事实上
这些关于 $i = s$ 的断言对 $M$ 也成立。

\medskip\noindent
令 $M' = H^0_I(M)$ 且 $M'' = M/M'$，从而有短正合序列
$$
0 \to M' \to M \to M'' \to 0
$$
并且由对偶化复形，引理 \ref{dualizing-lemma-divide-by-torsion}，
$M''$ 满足 $H^0_I(M') = 0$。
由 Artin--Rees（代数，引理 \ref{algebra-lemma-Artin-Rees}），得到短正合序列
$$
0 \to M' \to M/I^n M \to M''/I^n M'' \to 0
$$
对充分大的 $n$ 成立。考察长正合序列
$$
H^s_T(M') \to
H^s_T(M/I^nM) \to
H^s_T(M''/I^nM'') \to
H^{s + 1}_T(M')
$$
现在容易看出，若逆系统 $\{H^s_T(M''/I^nM'')\}_{n \geq 0}$
满足 Mittag-Leffler 条件，则逆系统
$\{H^s_T(M/I^nM)\}_{n \geq 0}$ 也满足 Mittag-Leffler 条件。
（注意，群的逆系统 $G_n$ 是否满足 ML 条件，
只取决于该逆系统在充分大的 $n$ 处的值。）
此外，序列
$$
H^s_T(M') \to
\lim H^s_T(M/I^nM) \to
\lim H^s_T(M''/I^nM'') \to
H^{s + 1}_T(M')
$$
是正合的，因为在所需位置满足 ML 条件；参见
同调，引理 \ref{homology-lemma-apply-Mittag-Leffler}。
因此，若 $H^s_T(M'') \to \lim H^s_T(M''/I^nM'')$
为同构，则
$H^s_T(M) \to \lim H^s_T(M/I^nM)$
由五引理也为同构
（同调，引理 \ref{homology-lemma-five-lemma}）。
于是归结为下一段讨论的情形。

\medskip\noindent
假设 $H^0_I(M) = 0$。选取生成元
$f_1, \ldots, f_r$ 来生成 $I^{m_0}$，其中 $m_0$ 是引理
\ref{algebraization-lemma-final-bootstrap} 中为 $M$ 找到的整数。
考察正合序列
$$
0 \to M \xrightarrow{f_1, \ldots, f_r}
(I^{m_0}M)^{\oplus r} \to Q \to 0
$$
并以此定义 $Q$。作几点观察：第一个映射为单射，
恰恰是因为 $H^0_I(M) = 0$。该单射的余核 $Q$
是有限 $A$-模，并且对每个 $1 \leq j \leq r$ 都有
$Q_{f_j} \cong (M_{f_j})^{\oplus r - 1}$。
特别地，对素理想 $\mathfrak p \subset A$，若
$\mathfrak p \not \in V(I)$，则有
$Q_\mathfrak p \cong (M_\mathfrak p)^{\oplus r - 1}$。
类似地，给定 $\mathfrak q \in T$ 和
$\mathfrak p' \subset A' = (A_\mathfrak q)^\wedge$，
其不包含于 $V(IA')$，则有
$Q'_{\mathfrak p'} \cong (M'_{\mathfrak p'})^{\oplus r - 1}$，
其中 $Q' = (Q_\mathfrak q)^\wedge$ 且 $M' = (M_\mathfrak q)^\wedge$。
因此，情形 \ref{algebraization-situation-bootstrap} 中的条件对
$A, I, T, Q$ 成立。（注意，$Q$ 可能有非零的
$H^0_I(Q)$，但这无关紧要。）

\medskip\noindent
对任意 $n \geq 0$，令 $F^nM = M \cap I^n(I^{m_0}M)^{\oplus r}$，
从而得到短正合序列
$$
0 \to F^nM \to I^n(I^{m_0}M)^{\oplus r} \to I^nQ \to 0
$$
由 Artin--Rees（代数，引理 \ref{algebra-lemma-Artin-Rees}），
存在 $c \geq 0$，使得
$I^n M \subset F^nM \subset I^{n - c}M$ 对所有 $n \geq c$ 成立。
令 $m_0$ 为整数，并令 $m'(m)$ 为对 $m \geq m_0$ 定义的函数；
它们来自将引理 \ref{algebraization-lemma-bootstrap} 应用于 $M$ 的结果。
注意，整数 $m_0$ 与我们上面选取的整数 $m_0$ 相同
（无须核验这一点：也可以直接取两个整数的最大值）。最后，
将引理 \ref{algebraization-lemma-bootstrap} 应用于 $Q$ 可知，对每个整数 $m$，
都存在整数 $m''(m) \geq m$，使得
$H^s_T(I^kQ) \to H^s_T(I^mQ)$ 对所有 $k \geq m''(m)$ 都为零。

\medskip\noindent
固定 $m \geq m_0$。选取 $k \geq m'(m''(m + c))$。
选取 $\xi \in H^{s + 1}_T(I^kM)$，
使其在 $H^{s + 1}_T(M)$ 中的像为零。
我们要证明 $\xi$ 在 $H^{s + 1}_T(I^mM)$ 中的像为零。
这将完全如引理 \ref{algebraization-lemma-final-bootstrap} 的证明那样说明
$\{H^s_T(M/I^nM)\}_{n \geq 0}$ 满足 Mittag-Leffler 条件。
下面的图有助于直观理解该论证：
$$
\xymatrix{
&
H^{s + 1}_T(I^kM) \ar[r] \ar[d] &
H^{s + 1}_T(I^k(I^{m_0}M)^{\oplus r}) \ar[d] &
\\
H^s_T(I^{m''(m + c)}Q) \ar[r]_-\delta \ar[d] &
H^{s + 1}_T(F^{m''(m + c)}M) \ar[r] \ar[d] &
H^{s + 1}_T(I^{m''(m + c)}(I^{m_0}M)^{\oplus r}) \\
H^s_T(I^{m + c}Q) \ar[r] &
H^{s + 1}_T(F^{m + c}M) \ar[d] &
\\
&
H^{s + 1}_T(I^mM)
}
$$
$\xi$ 在 $H^{s + 1}_T(I^k(I^{m_0}M)^{\oplus r})$ 中的像
映为 $H^{s + 1}_T((I^{m_0}M)^{\oplus r})$ 中的零，
因而它也映为
$H^{s + 1}_T(I^{m''(m + c)}(I^{m_0}M)^{\oplus r})$ 中的零，
这是由 $m'(-)$ 的选取得到的。
因此，像 $\xi' \in H^{s + 1}_T(F^{m''(m + c)}M)$
映为 $H^{s + 1}_T(I^{m''(m + c)}(I^{m_0}M)^{\oplus r})$ 中的零，
从而有 $\xi' = \delta(\eta)$，其中某个
$\eta \in H^s_T(I^{m''(m + c)}Q)$。
由对 $m''(-)$ 的选取可知，$\eta$ 在
$H^s_T(I^{m + c}Q)$ 中的像为零。
这又意味着 $\xi'$ 在
$H^{s + 1}_T(F^{m + c}M)$ 中的像为零。
由于 $F^{m + c}M \subset I^mM$，结论成立。

\medskip\noindent
最后证明关于极限的陈述。考察短正合序列
$$
0 \to M/F^nM \to (I^{m_0}M)^{\oplus r}/I^n (I^{m_0}M)^{\oplus r}
\to Q/I^nQ \to 0
$$
由于这些逆系统 pro-同构，我们有
$\lim H^s_T(M/I^nM) = \lim H^s_T(M/F^nM)$。得到交换图
$$
\xymatrix{
H^{s - 1}_T(Q) \ar[r] \ar[d] &
\lim H^{s - 1}_T(Q/I^nQ) \ar[d] \\
H^s_T(M) \ar[r] \ar[d] &
\lim H^s_T(M/I^nM) \ar[d] \\
H^s_T((I^{m_0}M)^{\oplus r}) \ar[r] \ar[d] &
\lim H^s_T((I^{m_0}M)^{\oplus r}/I^n(I^{m_0}M)^{\oplus r}) \ar[d] \\
H^s_T(Q) \ar[r] &
\lim H^s_T(Q/I^nQ)
}
$$
右列正合，因为在所需位置满足 ML 条件；参见
同调，引理 \ref{homology-lemma-apply-Mittag-Leffler}。
由引理 \ref{algebraization-lemma-final-bootstrap} 的 (6)，最下面的水平箭头为单射（！）。
由引理 \ref{algebraization-lemma-final-bootstrap} 的 (5)，其上方的水平箭头为双射。
上同调次数 $\leq s - 1$ 中的箭头均为同构。
因此，由五引理（同调，引理 \ref{homology-lemma-five-lemma}），可知
$H^s_T(M) \to \lim H^s_T(M/I^nM)$ 为同构。
定理得证。
\end{proof}

\begin{lemma}
\label{algebraization-lemma-combine-two}
设 $I \subset \mathfrak a \subset A$ 是诺特环 $A$ 的理想，
并设 $M$ 是有限 $A$-模。设 $s$ 和 $d$ 为整数。假设
\begin{enumerate}
\item $A, I, V(\mathfrak a), M$ 对 $s$ 和 $d$ 满足
情形 \ref{algebraization-situation-bootstrap} 的假设；并且
\item $A, I, \mathfrak a, M$ 对 $s + 1$ 和 $d$ 满足
引理 \ref{algebraization-lemma-algebraize-local-cohomology-general}
的条件，其中 $J = \mathfrak a$。
\end{enumerate}
则存在理想
$J_0 \subset \mathfrak a$，满足 $V(J_0) \cap V(I) = V(\mathfrak a)$，
使得：对任意 $J \subset J_0$，若 $V(J) \cap V(I) = V(\mathfrak a)$，
映射
$$
H^{s + 1}_J(M) \longrightarrow \lim H^{s + 1}_\mathfrak a(M/I^nM)
$$
为同构。
\end{lemma}

\begin{proof}
具体而言，由引理 \ref{algebraization-lemma-algebraize-local-cohomology-general}，
存在 $J_0$，并且有同构
$H^{s + 1}_J(M) = H^{s + 1}(R\Gamma_\mathfrak a(M)^\wedge)$
；由对偶化复形，引理 \ref{dualizing-lemma-completion-local}，有短正合序列
$$
0 \to R^1\lim H^s_\mathfrak a(M/I^nM) \to
H^{s + 1}(R\Gamma_\mathfrak a(M)^\wedge) \to
\lim H^{s + 1}_\mathfrak a(M/I^nM) \to 0
$$
而模 $R^1\lim H^s_\mathfrak a(M/I^nM)$ 为零，因为由定理
\ref{algebraization-theorem-final-bootstrap}，
$\{H^s_\mathfrak a(M/I^nM)\}_{n \geq 0}$ 满足 Mittag-Leffler 条件。
\end{proof}








\section{形式截面的代数化，I}
\label{algebraization-section-algebraization-sections}

\noindent
本节研究局部情形下形式截面的代数化问题。
设 $(A, \mathfrak m)$ 为诺特局部环。
设 $I \subset A$ 为理想。令
$$
X = \Spec(A) \supset U = \Spec(A) \setminus \{\mathfrak m\}
$$
并用 $Y = V(I)$ 表示与 $I$ 对应的闭子概形。
设 $\mathcal{F}$ 为凝聚 $\mathcal{O}_U$-模。
本节考察极限
$$
\lim_n H^i(U, \mathcal{F}/I^n\mathcal{F})
$$
这与把 $\mathcal{F}$ 拉回到 $U$ 沿 $Y$ 的形式完备化所得之上同调密切相关；
不过，由于我们尚未引入形式概形，此处不能使用这种术语。

\begin{lemma}
\label{algebraization-lemma-compare-with-derived-completion}
设 $U$ 为诺特局部环 $A$ 的穿孔谱。
设 $\mathcal{F}$ 为凝聚 $\mathcal{O}_U$-模。
设 $I \subset A$ 为理想。则
$$
H^i(R\Gamma(U, \mathcal{F})^\wedge) =
\lim H^i(U, \mathcal{F}/I^n\mathcal{F})
$$
对所有 $i$ 成立，其中 $R\Gamma(U, \mathcal{F})^\wedge$ 表示
导出 $I$-进完备化。
\end{lemma}

\begin{proof}
由引理 \ref{algebraization-lemma-formal-functions-general} 和
\ref{algebraization-lemma-derived-completion-pseudo-coherent}，有
$$
R\Gamma(U, \mathcal{F})^\wedge =
R\Gamma(U, \mathcal{F}^\wedge) =
R\Gamma(U, R\lim \mathcal{F}/I^n\mathcal{F})
$$
因此得到短正合序列
$$
0 \to R^1\lim H^{i - 1}(U, \mathcal{F}/I^n\mathcal{F}) \to
H^i(R\Gamma(U, \mathcal{F})^\wedge) \to
\lim H^i(U, \mathcal{F}/I^n\mathcal{F}) \to 0
$$
这是由上同调，引理 \ref{cohomology-lemma-RGamma-commutes-with-Rlim} 得到的。
$R^1\lim$ 项为零，因为由引理 \ref{algebraization-lemma-ML-local}，群的逆系统
$H^i(U, \mathcal{F}/I^n\mathcal{F})$ 满足 Mittag-Leffler 条件。
\end{proof}

\begin{theorem}
\label{algebraization-theorem-algebraization-formal-sections}
\begin{reference}
本证明方法大体沿用 \cite[Theorem 1]{Faltings-algebraisation}
和 \cite[Satz 2]{Faltings-uber} 的证明方法。
所得结果几乎等同于
\cite[Theorem 1.1]{MRaynaud-paper}（补集为仿射的情形）和
\cite[Theorem 3.9]{MRaynaud-book}（补集为少数仿射开集之并的情形）。
\end{reference}
设 $(A, \mathfrak m)$ 为有对偶化复形且关于理想 $I$ 完备的诺特局部环。
令 $X = \Spec(A)$、$Y = V(I)$，且 $U = X \setminus \{\mathfrak m\}$。
设 $\mathcal{F}$ 为 $U$ 上的凝聚层。假设
\begin{enumerate}
\item $\text{cd}(A, I) \leq d$, i.e.,
$H^i(X \setminus Y, \mathcal{G}) = 0$ 对 $i \geq d$ 以及
拟凝聚 $\mathcal{G}$（在 $X$ 上）成立；
\item 对任意 $x \in X \setminus Y$，若其闭包
$\overline{\{x\}}$（在 $X$ 中）与 $U \cap Y$ 相交，则有
$$
\text{depth}_{\mathcal{O}_{X, x}}(\mathcal{F}_x) \geq s
\quad\text{或}\quad
\text{depth}_{\mathcal{O}_{X, x}}(\mathcal{F}_x)
+ \dim(\overline{\{x\}}) > d + s
$$
\end{enumerate}
则存在开集 $V_0 \subset U$，它包含 $U \cap Y$，
使得对任意开集 $V \subset V_0$，若它包含 $U \cap Y$，则映射
$$
H^i(V, \mathcal{F}) \to \lim H^i(U, \mathcal{F}/I^n\mathcal{F})
$$
在 $i < s$ 时为同构。若此外
$
\text{depth}_{\mathcal{O}_{X, x}}(\mathcal{F}_x) +
\dim(\overline{\{x\}}) > s
$
对所有 $x \in U \cap Y$ 成立，则这些上同调群是有限 $A$-模。
\end{theorem}

\begin{proof}
选取有限 $A$-模 $M$，使得 $\mathcal{F}$ 是在 $U$ 上限制所得的凝聚
$\mathcal{O}_X$-模，该模与 $M$ 相伴；参见局部上同调，
引理 \ref{local-cohomology-lemma-finiteness-pushforwards-and-H1-local}。
于是引理 \ref{algebraization-lemma-algebraize-local-cohomology}
的假设成立。
按该引理选取 $J_0$，并令 $V_0 = X \setminus V(J_0)$。
于是，开集 $V \subset V_0$ 若包含 $U \cap Y$，
便与理想以 $1$-对-$1$ 的方式对应；这里 $J \subset J_0$ 且
$V(J) \cap V(I) = \{\mathfrak m\}$。
此外，对这样的选取，有优三角
$$
R\Gamma_J(M) \to M \to R\Gamma(V, \mathcal{F}) \to
R\Gamma_J(M)[1]
$$
类似地，有优三角
$$
R\Gamma_\mathfrak m(M)^\wedge \to
M \to
R\Gamma(U, \mathcal{F})^\wedge \to
R\Gamma_\mathfrak m(M)^\wedge[1]
$$
其中涉及导出 $I$-进完备化。
由引理 \ref{algebraization-lemma-compare-with-derived-completion}，
$R\Gamma(U, \mathcal{F})^\wedge$ 的上同调群等于定理陈述中的极限。
这两个三角之间的典范映射和一些容易的论证表明，本定理可由主要引理
\ref{algebraization-lemma-algebraize-local-cohomology} 得出
（注意，此处为 $i < s$，而引理中为
$i \leq s$；这是由移位造成的）。
在附加假设下，上同调群的有限性由引理
\ref{algebraization-lemma-kill-colimit} 得出。
\end{proof}

\begin{lemma}
\label{algebraization-lemma-application-theorem}
设 $(A, \mathfrak m)$ 为有对偶化复形且关于理想 $I$ 完备的诺特局部环。
令 $X = \Spec(A)$、$Y = V(I)$，且 $U = X \setminus \{\mathfrak m\}$。
设 $\mathcal{F}$ 为 $U$ 上的凝聚层。
假设对任意相伴点 $x \in U$（相对于 $\mathcal{F}$）都有
$\dim(\overline{\{x\}}) > \text{cd}(A, I) + 1$，
其中 $\overline{\{x\}}$ 是其在 $X$ 中的闭包。则映射
$$
\colim H^0(V, \mathcal{F})
\longrightarrow
\lim H^0(U, \mathcal{F}/I^n\mathcal{F})
$$
是有限 $A$-模之间的同构，
其中余极限遍历开集 $V \subset U$，这些开集包含 $U \cap Y$。
\end{lemma}

\begin{proof}
在 $s = 1$ 时应用定理 \ref{algebraization-theorem-algebraization-formal-sections}
（同时也得到有限性）。
\end{proof}




\section{形式截面的代数化，II}
\label{algebraization-section-algebraization-sections-coherent}

\noindent
要简洁陈述第 \ref{algebraization-section-algebraization-sections-general} 节和
第 \ref{algebraization-section-bootstrap} 节中的结果对于拟仿射概形上凝聚层的上同调
及其关于理想的完备化所能给出的全部推论，有些困难。
本节给出关于 $H^0$ 的一份不穷尽的应用清单。
下一节讨论对高阶上同调的应用。

\begin{lemma}
\label{algebraization-lemma-ses-H0}
设 $I \subset \mathfrak a$ 是诺特环 $A$ 的理想。
设 $0 \to \mathcal{F}' \to \mathcal{F} \to \mathcal{F}'' \to 0$
为 $U = \Spec(A) \setminus V(\mathfrak a)$ 上凝聚模的短正合序列。
令 $\mathcal{V}$ 为开子概形 $V \subset U$ 的集合，其中各开子概形都包含
$U \cap V(I)$，并按反包含排序。考察交换图
$$
\xymatrix{
\colim_\mathcal{V} H^0(V, \mathcal{F}') \ar[d] \ar[r] &
\colim_\mathcal{V} H^0(V, \mathcal{F}) \ar[d] \ar[r] &
\colim_\mathcal{V} H^0(V, \mathcal{F}'') \ar[d] \\
\lim H^0(U, \mathcal{F}'/I^n\mathcal{F}') \ar[r] &
\lim H^0(U, \mathcal{F}'/I^n\mathcal{F}) \ar[r] &
\lim H^0(U, \mathcal{F}'/I^n\mathcal{F}'')
}
$$
若左边和右边的竖直箭头都是同构，则中间的也是同构。
若中间和左边的竖直箭头都是同构，则左边的也是同构。
\end{lemma}

\begin{proof}
图中的序列在中间正合，且第一个箭头为单射。
因此最后一个陈述由简单的追图得到。在证明的其余部分，
假设左边和右边的竖直箭头均为同构。追图表明中间的竖直箭头为单射。
只须再证明它是满射。

\medskip\noindent
可以选取有限 $A$-模 $M$ 和 $M'$，使得 $\mathcal{F}$
和 $\mathcal{F}'$ 分别是 $\widetilde{M}$ 和 $\widetilde{M}'$ 在 $U$ 上的限制；
参见局部上同调，引理
\ref{local-cohomology-lemma-finiteness-pushforwards-and-H1-local}。
把 $M'$ 替换为 $\mathfrak a^n M'$（其中某个 $n \geq 0$）后，
可以假设 $\mathcal{F}' \to \mathcal{F}$ 对应于模映射
$M' \to M$；参见概形上同调，引理 \ref{coherent-lemma-homs-over-open}。
把 $M'$ 替换为 $M' \to M$ 的像，并令 $M'' = M/M'$ 后，
可见原短正合序列对应于与短正合序列
$0 \to M' \to M \to M'' \to 0$ 相伴之凝聚模短正合序列的限制；
这里所述的是 $A$-模。

\medskip\noindent
令 $\hat s \in \lim H^0(U, \mathcal{F}/I^n\mathcal{F})$，其像为
$\hat s'' \in \lim H^0(U, \mathcal{F}''/I^n\mathcal{F}'')$。
由假设，找到 $V \in \mathcal{V}$ 以及截面
$s'' \in \mathcal{F}''(V)$，后者映为 $\hat s''$。
令 $J \subset A$ 为满足 $V(J) = \Spec(A) \setminus V$ 的理想。
由概形上同调，引理 \ref{coherent-lemma-homs-over-open}，
用 $J$ 的某个幂替换它之后，可以假设存在一个
$A$-线性映射 $\varphi : J \to M''$；它对应于 $s''$。
固定这个 $J$ 的选择；在证明的其余部分，将用 $V$ 的一个更小的
$V$ 来替换它，该开集属于 $\mathcal{V}$，即将有 $V \cap V(J) = \emptyset$。

\medskip\noindent
选取一个表现 $A^{\oplus m} \to A^{\oplus n} \to J \to 0$。
用 $g_1, \ldots, g_n \in J$ 表示 $A^{\oplus n}$ 的基向量的像，
从而 $J = (g_1, \ldots, g_n)$。设
$A^{\oplus m} \to A^{\oplus n}$ 由矩阵 $(a_{ji})$ 给出，
于是 $\sum a_{ji} g_i = 0$，其中 $j = 1, \ldots, m$。
由于 $M \to M''$ 为满射，对每个 $i$ 都可选取 $m_i \in M$，
使其映为 $\varphi(g_i) \in M''$。于是
元素 $g_i \hat s - m_i$ 属于
$\lim H^0(U, \mathcal{F}/I^n\mathcal{F})$，并落在子模
$\lim H^0(U, \mathcal{F}'/I^n\mathcal{F}')$ 中。
由假设，缩小 $V$ 后可假设存在
$s'_i \in \mathcal{F}'(V)$，$i = 1, \ldots, n$，
且 $s'_i$ 映为 $g_i \hat s - m_i$。
令 $s_i = s'_i + m_i$，这是在 $\mathcal{F}(V)$ 中。
注意，$\sum a_{ji} s_i$ 映为 $\sum a_{ji}g_i\hat s = 0$；
所用映射为
$$
\colim_\mathcal{V} \mathcal{F}(V')
\longrightarrow
\lim H^0(U, \mathcal{F}/I^n\mathcal{F})
$$
由于该映射为单射（见上文），缩小 $V$ 后可以假设
$\sum a_{ji}s_i = 0$ 在 $\mathcal{F}(V)$ 中对所有 $j = 1, \ldots, m$ 成立。
于是得到 $A$-模映射 $J \to \mathcal{F}(V)$，它把 $g_i$ 映为 $s_i$。
由 $\widetilde{J}$ 的泛性质，该 $A$-模映射对应于
$\mathcal{O}_V$-模映射 $\widetilde{J}|_V \to \mathcal{F}$。
不过，由于 $V(J) \cap V = \emptyset$，有
$\widetilde{J}|_V = \mathcal{O}_V$。因此构造出了截面
$s \in \mathcal{F}(V)$。略去验证 $s$ 经上面所示映射映为 $\hat s$ 的计算。
\end{proof}

\noindent
下面的引理将被命题 \ref{algebraization-proposition-application-H0} 取代。

\begin{lemma}
\label{algebraization-lemma-application-H0-pre}
设 $I \subset \mathfrak a$ 是诺特环 $A$ 的理想。
设 $\mathcal{F}$ 为 $U = \Spec(A) \setminus V(\mathfrak a)$ 上的凝聚模。
假设
\begin{enumerate}
\item $A$ 是 $I$-进完备的，并且有对偶化复形；
\item 若 $x \in \text{Ass}(\mathcal{F})$、$x \not \in V(I)$、
$\overline{\{x\}} \cap V(I) \not \subset V(\mathfrak a)$,
且 $z \in \overline{\{x\}} \cap V(\mathfrak a)$，则
$\dim(\mathcal{O}_{\overline{\{x\}}, z}) > \text{cd}(A, I) + 1$；
\item 下列条件之一成立：
\begin{enumerate}
\item $\mathcal{F}$ 在 $U \setminus V(I)$ 上的限制满足 $(S_1)$；
\item $V(\mathfrak a)$ 的维数至多为 $2$\footnote{这里的含义是：
在 $V(\mathfrak a)$ 上，来自 $\Spec(A)$ 上一个维数函数的最大值与
最小值之差至多为 $2$。}。
\end{enumerate}
\end{enumerate}
则得到同构
$$
\colim H^0(V, \mathcal{F})
\longrightarrow
\lim H^0(U, \mathcal{F}/I^n\mathcal{F})
$$
其中余极限遍历开集 $V \subset U$，这些开集包含 $U \cap V(I)$。
\end{lemma}

\begin{proof}
选取有限 $A$-模 $M$，使得 $\mathcal{F}$ 是在 $U$ 上的限制，
所限制的凝聚模与 $M$ 相伴；参见局部上同调，引理
\ref{local-cohomology-lemma-finiteness-pushforwards-and-H1-local}。
令 $d = \text{cd}(A, I)$。
设 $\mathfrak p$ 为 $A$ 的素理想，且不包含于 $V(I)$，
并设 $\mathfrak q \in V(\mathfrak p) \cap V(\mathfrak a)$。
那么要么 $\mathfrak p$ 不是 $M$ 的相伴素理想，因而
$\text{depth}(M_\mathfrak p) \geq 1$；
要么由 (2) 有 $\dim((A/\mathfrak p)_\mathfrak q) > d + 1$。
因此，引理 \ref{algebraization-lemma-algebraize-local-cohomology-general}
的假设对 $s = 1$ 和 $d$ 成立；这里使用了条件 (3)。
由此存在理想
$J_0 \subset \mathfrak a$，满足 $V(J_0) \cap V(I) = V(\mathfrak a)$，
使得对任意 $J \subset J_0$，若 $V(J) \cap V(I) = V(\mathfrak a)$，
则映射
$$
H^i_J(M) \longrightarrow H^i(R\Gamma_\mathfrak a(M)^\wedge)
$$
在 $i = 0, 1$ 时为同构。考察正合三角之间的态射
$$
\xymatrix{
R\Gamma_J(M)  \ar[d] \ar[r] &
M \ar[r] \ar[d] &
R\Gamma(V, \mathcal{F}) \ar[d] \ar[r] &
R\Gamma_J(M)[1]  \ar[d] \\
R\Gamma_J(M)^\wedge \ar[r] &
M \ar[r] &
R\Gamma(V, \mathcal{F})^\wedge \ar[r] &
R\Gamma_J(M)^\wedge[1] \\
R\Gamma_\mathfrak a(M)^\wedge \ar[r] \ar[u] &
M \ar[r] \ar[u] &
R\Gamma(U, \mathcal{F})^\wedge \ar[r] \ar[u] &
R\Gamma_\mathfrak a(M)^\wedge[1] \ar[u]
}
$$
其中 $V = \Spec(A) \setminus V(J)$。回忆
$R\Gamma_\mathfrak a(M)^\wedge \to R\Gamma_J(M)^\wedge$
为同构（因为 $\mathfrak a$、$\mathfrak a + I$ 和 $J + I$
截出同一个闭子概形；例如参见引理
\ref{algebraization-lemma-algebraize-local-cohomology-general} 的证明）。因此
$R\Gamma(U, \mathcal{F})^\wedge = R\Gamma(V, \mathcal{F})^\wedge$.
由此得到交换图
$$
\xymatrix{
0 \ar[r] &
H^0_J(M) \ar[r] \ar[d] &
M \ar[r] \ar[d] \ar[r] &
\Gamma(V, \mathcal{F}) \ar[d] \ar[r] &
H^1_J(M) \ar[d] \ar[r] &
0 \\
0 \ar[r] &
H^0(R\Gamma_J(M)^\wedge) \ar[r] &
M \ar[r] &
H^0(R\Gamma(V, \mathcal{F})^\wedge) \ar[r] &
H^1(R\Gamma_J(M)^\wedge) \ar[r] &
0 \\
0 \ar[r] &
H^0(R\Gamma_\mathfrak a(M)^\wedge) \ar[r] \ar[u] &
M \ar[r] \ar[u] &
H^0(R\Gamma(U, \mathcal{F})^\wedge) \ar[r] \ar[u] &
H^1(R\Gamma_\mathfrak a(M)^\wedge) \ar[r] \ar[u] &
0
}
$$
其各行正合，且下方的竖直箭头均为同构。因此得到同构
$\Gamma(V, \mathcal{F}) \to H^0(R\Gamma(U, \mathcal{F})^\wedge)$.
由引理 \ref{algebraization-lemma-formal-functions-general}
和 \ref{algebraization-lemma-derived-completion-pseudo-coherent}，有
$$
R\Gamma(U, \mathcal{F})^\wedge =
R\Gamma(U, \mathcal{F}^\wedge) =
R\Gamma(U, R\lim \mathcal{F}/I^n\mathcal{F})
$$
并且由上同调，引理 \ref{cohomology-lemma-RGamma-commutes-with-Rlim}，得到
$H^0(R\Gamma(U, \mathcal{F})^\wedge) =
\lim H^0(U, \mathcal{F}/I^n\mathcal{F})$。
\end{proof}

\noindent
现在提升前一引理，以去掉条件 (3)。

\begin{proposition}
\label{algebraization-proposition-application-H0}
设 $I \subset \mathfrak a$ 是诺特环 $A$ 的理想。
设 $\mathcal{F}$ 为 $U = \Spec(A) \setminus V(\mathfrak a)$ 上的凝聚模。
假设
\begin{enumerate}
\item $A$ 是 $I$-进完备的，并且有对偶化复形；
\item 若 $x \in \text{Ass}(\mathcal{F})$、$x \not \in V(I)$、
$\overline{\{x\}} \cap V(I) \not \subset V(\mathfrak a)$,
且 $z \in \overline{\{x\}} \cap V(\mathfrak a)$，则
$\dim(\mathcal{O}_{\overline{\{x\}}, z}) > \text{cd}(A, I) + 1$.
\end{enumerate}
则得到同构
$$
\colim H^0(V, \mathcal{F})
\longrightarrow
\lim H^0(U, \mathcal{F}/I^n\mathcal{F})
$$
其中余极限遍历开集 $V \subset U$，这些开集包含 $U \cap V(I)$。
\end{proposition}

\begin{proof}
令 $T \subset U$ 为点 $x$ 所成的集合，这些点满足
$\overline{\{x\}} \cap V(I) \subset V(\mathfrak a)$。
令 $\mathcal{F} \to \mathcal{F}'$ 为局部上同调，引理
\ref{local-cohomology-lemma-get-depth-1-along-Z} 中在 $U$ 上构造的
凝聚模满射。由于 $\mathcal{F} \to \mathcal{F}'$ 在开集
$V \subset U$ 上为同构，而该开集包含 $U \cap V(I)$，
只须证明把 $\mathcal{F}$ 替换为 $\mathcal{F}'$ 后的引理。
因此可以并且确实假设：对 $x \in U$，若
$\overline{\{x\}} \cap V(I) \subset V(\mathfrak a)$，
有 $\text{depth}(\mathcal{F}_x) \geq 1$。

\medskip\noindent
令 $\mathcal{V}$ 为开子概形 $V \subset U$ 所成的集合，
其中各开子概形包含 $U \cap V(I)$，并按反包含排序。
这是一个有向集。我们先断言
$$
\mathcal{F}(V)
\longrightarrow
\lim H^0(U, \mathcal{F}/I^n\mathcal{F})
$$
对任意 $V \in \mathcal{F}$ 都为单射（特别地，引理中的映射为单射）。
事实上，由上一段，相伴点 $x$（属于 $\mathcal{F}$）必须满足
$\overline{\{x\}} \cap U \cap Y \not = \emptyset$。
若 $y \in \overline{\{x\}} \cap U \cap Y$，则
$\mathcal{F}_x$ 是 $\mathcal{F}_y$ 的局部化，并且由 Krull 交定理
（代数，引理 \ref{algebra-lemma-intersect-powers-ideal-module-zero}），有
$\mathcal{F}_y \subset \lim \mathcal{F}_y/I^n \mathcal{F}_y$。
这证明了断言，因为核中的截面 $s \in \mathcal{F}(V)$
必有空支集，因而必为零。

\medskip\noindent
选取有限 $A$-模 $M$，使得 $\mathcal{F}$ 是 $\widetilde{M}$ 在 $U$ 上的限制；
参见局部上同调，引理
\ref{local-cohomology-lemma-finiteness-pushforwards-and-H1-local}。
可以并且确实假设 $H^0_\mathfrak a(M) = 0$。
令 $\text{Ass}(M) \setminus V(I) = \{\mathfrak p_1, \ldots, \mathfrak p_n\}$。
对 $n$ 作归纳来证明该引理。重新排序后，可以假设 $\mathfrak p_n$
是集合 $\{\mathfrak p_1, \ldots, \mathfrak p_n\}$ 关于包含关系的极小元，
即 $\mathfrak p_n$ 是 $M$ 的支集的一般点。令
$$
M' = H^0_{\mathfrak p_1 \ldots \mathfrak p_{n - 1} I}(M)
$$
并令 $M'' = M/M'$。设 $\mathcal{F}'$ 和 $\mathcal{F}''$ 为凝聚
$\mathcal{O}_U$-模，它们分别对应于 $M'$ 和 $M''$。
对偶化复形，引理 \ref{dualizing-lemma-divide-by-torsion}
蕴含 $M''$ 只有一个相伴素理想，即 $\mathfrak p_n$。
因此 $\mathcal{F}''$ 只有一个相伴点，并且可见引理
\ref{algebraization-lemma-application-H0-pre} 的条件 (3)(a) 成立；
从而映射 $\colim H^0(V, \mathcal{F}'')
\to \lim H^0(U, \mathcal{F}''/I^n\mathcal{F}'')$ 为同构。
另一方面，由于
$\mathfrak p_n \not \in V(\mathfrak p_1 \ldots \mathfrak p_{n - 1} I)$
可知 $\mathfrak p_n$ 不是 $M'$ 的相伴素理想。
因此归纳假设适用于 $M'$；注意，由于 $\mathcal{F}' \subset \mathcal{F}$，
条件 $\text{depth}(\mathcal{F}'_x) \geq 1$ 对点 $x$ 成立，
其中 $\overline{\{x\}} \cap V(I) \subset V(\mathfrak a)$；
参见代数，引理 \ref{algebra-lemma-depth-in-ses}。
因此映射
$\colim H^0(V, \mathcal{F}')
\to \lim H^0(U, \mathcal{F}'/I^n\mathcal{F}')$ 也为同构。
由引理 \ref{algebraization-lemma-ses-H0} 得出结论。
\end{proof}

\begin{lemma}
\label{algebraization-lemma-application-H0}
设 $I \subset \mathfrak a$ 是诺特环 $A$ 的理想。
设 $\mathcal{F}$ 为 $U = \Spec(A) \setminus V(\mathfrak a)$ 上的凝聚模。
假设
\begin{enumerate}
\item $A$ 是 $I$-进完备的，并且有对偶化复形；
\item 若 $x \in \text{Ass}(\mathcal{F})$、$x \not \in V(I)$、
$\overline{\{x\}} \cap V(I) \not \subset V(\mathfrak a)$，且
$z \in V(\mathfrak a) \cap \overline{\{x\}}$，则
$\dim(\mathcal{O}_{\overline{\{x\}}, z}) > \text{cd}(A, I) + 1$；
\item 对 $x \in U$，若 $\overline{\{x\}} \cap V(I) \subset V(\mathfrak a)$，
则 $\text{depth}(\mathcal{F}_x) \geq 2$。
\end{enumerate}
则得到同构
$$
H^0(U, \mathcal{F})
\longrightarrow
\lim H^0(U, \mathcal{F}/I^n\mathcal{F})
$$
\end{lemma}

\begin{proof}
令 $\hat s \in \lim H^0(U, \mathcal{F}/I^n\mathcal{F})$。
由命题 \ref{algebraization-proposition-application-H0}，可知 $\hat s$ 是某个元素
$s \in \mathcal{F}(V)$ 的像，其中 $V \subset U$ 是包含
$U \cap V(I)$ 的开集。不过，条件 (3) 表明
$\text{depth}(\mathcal{F}_x) \geq 2$ 对所有 $x \in U \setminus V$ 成立，
因而由除子，引理 \ref{divisors-lemma-depth-2-hartog} 得到
$\mathcal{F}(V) = \mathcal{F}(U)$，证明完毕。
\end{proof}

\begin{lemma}
\label{algebraization-lemma-alternative-colim-H0}
设 $A$ 为诺特环。设 $f \in \mathfrak a \subset A$
为 $A$ 的一个理想中的元素。设 $M$ 为有限 $A$-模。假设
\begin{enumerate}
\item $A$ 是 $f$-进完备的；
\item $f$ 在 $M$ 上是非零因子；
\item $H^1_\mathfrak a(M/fM)$ 是有限 $A$-模。
\end{enumerate}
则令 $U = \Spec(A) \setminus V(\mathfrak a)$，映射
$$
\colim_V \Gamma(V, \widetilde{M})
\longrightarrow
\lim \Gamma(U, \widetilde{M/f^nM})
$$
为同构，其中余极限遍历开集 $V \subset U$，这些开集包含
$U \cap V(f)$。
\end{lemma}

\begin{proof}
令 $\mathcal{F} = \widetilde{M}|_U$。
$H^1_\mathfrak a(M/fM)$ 的有限性蕴含
$H^0(U, \mathcal{F}/f\mathcal{F})$ 有限；参见局部上同调，引理
\ref{local-cohomology-lemma-finiteness-pushforwards-and-H1-local}。
由上同调，引理 \ref{cohomology-lemma-limit-finite}
（该引理适用，因为 $f$ 在 $\mathcal{F}$ 上是非零因子），可知
$N = \lim H^0(U, \mathcal{F}/f^n\mathcal{F})$
是有限 $A$-模、无 $f$-挠，并且
$N/fN \subset H^0(U, \mathcal{F}/f\mathcal{F})$。
另一方面，有映射 $M \to N$ 以及相容映射
$$
M/fM \longrightarrow H^0(U, \mathcal{F}/f\mathcal{F})
$$
对 $g \in \mathfrak a$，可见 $(M/fM)_g$ 同构地映至
$H^0(U \cap D(f), \mathcal{F}/f\mathcal{F})$，因为
$\mathcal{F}/f\mathcal{F}$ 是 $\widetilde{M/fM}$ 在 $U$ 上的限制。
由此可知 $M_g \to N_g$ 诱导同构
$$
M_g/fM_g = (M/fM)_g \to (N/fN)_g = N_g/fN_g
$$
由于 $f$ 在 $N$ 和 $M$ 上都是非零因子，可知
$M_g \to N_g$ 在 $f$-进完备化上诱导同构，
进而蕴含 $M_g \to N_g$ 在 $V(f) \cap D(g)$ 的某个开邻域上为同构。
由于 $g \in \mathfrak a$ 是任意的，可知 $M$ 和 $N$
在引理陈述中所述的某个开集 $V$ 上确定同构的凝聚模。证明完毕。
\end{proof}

\begin{proposition}
\label{algebraization-proposition-alternative-colim-H0}
设 $A$ 为诺特环。设 $f \in \mathfrak a \subset A$
为 $A$ 的一个理想中的元素。设 $\mathcal{F}$ 为
$U = \Spec(A) \setminus V(\mathfrak a)$ 上的凝聚模。假设
\begin{enumerate}
\item $A$ 是 $f$-进完备的，并且有对偶化复形；
\item 若 $x \in \text{Ass}(\mathcal{F})$、$x \not \in V(f)$、
$\overline{\{x\}} \cap V(f) \not \subset V(\mathfrak a)$,
且 $z \in \overline{\{x\}} \cap V(\mathfrak a)$，则
$\dim(\mathcal{O}_{\overline{\{x\}}, z}) > 2$.
\end{enumerate}
则映射
$$
\colim_V \Gamma(V, \mathcal{F})
\longrightarrow
\lim \Gamma(U, \mathcal{F}/f^n\mathcal{F})
$$
为同构，其中余极限遍历开集 $V \subset U$，这些开集包含
$U \cap V(f)$。
\end{proposition}

\begin{proof}[第一种证明]
回忆 $A$ 是泛链状的，并且有 Gorenstein 形式纤维；参见对偶化复形，引理
\ref{dualizing-lemma-dualizing-gorenstein-formal-fibres} 和
\ref{dualizing-lemma-universally-catenary}。因此可以考察映射
$\mathcal{F} \to \mathcal{F}'$，它由局部上同调，引理
\ref{local-cohomology-lemma-make-S2-along-Z} 针对闭子集
$V(f) \cap U$（属于 $U$）构造。注意
\begin{enumerate}
\item $\mathcal{F} \to \mathcal{F}'$ 的核和余核支撑在 $V(f) \cap U$ 上。
\item 模 $\mathcal{F}'$ 无 $f$-挠，因为其茎的深度 $\geq 1$，
这对 $V(f) \cap U$ 的所有点成立；换言之，$\mathcal{F}'$
在 $V(f) \cap U$ 中没有相伴点。
\item 若 $y \in V(f) \cap U$ 是
$\mathcal{F}'/f\mathcal{F}'$ 的相伴点，则
$\text{depth}(\mathcal{F}'_y) = 1$
并且由 $\mathcal{F}'$ 的构造，存在立即特殊化 $x \leadsto y$，
其中 $x \not \in V(f)$ 是 $\mathcal{F}$ 的相伴点。
由假设 (2) 可知，$y$ 不可能在 $\Spec(A)$ 中立即特殊化到点
$z \in V(\mathfrak a)$。
\item 由 (3) 可知 $H^0(U, \mathcal{F}'/f\mathcal{F}')$
是有限 $A$-模；参见局部上同调，引理
\ref{local-cohomology-lemma-finiteness-Rjstar}。
\end{enumerate}
这些观察将使我们能够完成证明。

\medskip\noindent
首先，断言该引理对 $\mathcal{F}'$ 成立。
具体而言，选取有限 $A$-模 $M'$，使得 $\mathcal{F}'$ 是在 $U$ 上的限制，
所限制的凝聚模与 $M'$ 相伴；参见局部上同调，引理
\ref{local-cohomology-lemma-finiteness-pushforwards-and-H1-local}。
由于 $\mathcal{F}'$ 无 $f$-挠，可以假设 $M'$ 也无 $f$-挠。
上面的观察 (4) 表明 $H^1_\mathfrak a(M')$ 是有限 $A$-模；
参见局部上同调，引理
\ref{local-cohomology-lemma-finiteness-pushforwards-and-H1-local}。
因此该断言由引理 \ref{algebraization-lemma-alternative-colim-H0} 得到。

\medskip\noindent
其次，注意该引理对任意凝聚 $\mathcal{O}_U$-模都显然成立，
只要它支撑在 $V(f) \cap U$ 上。
令 $\mathcal{K}$、$\mathcal{G}$、$\mathcal{Q}$ 分别为映射
$\mathcal{F} \to \mathcal{F}'$ 的核、像、余核。短正合序列
$0 \to \mathcal{G} \to \mathcal{F}' \to \mathcal{Q} \to 0$
和引理 \ref{algebraization-lemma-ses-H0} 表明结论对 $\mathcal{G}$ 成立。
然后对短正合序列再次作同样处理：
$0 \to \mathcal{K} \to \mathcal{F} \to \mathcal{G} \to 0$
从而完成证明。
\end{proof}

\begin{proof}[第二种证明]
该命题是命题 \ref{algebraization-proposition-application-H0} 的特殊情形。
\end{proof}

\begin{lemma}
\label{algebraization-lemma-alternative-H0}
设 $A$ 为诺特环。设 $f \in \mathfrak a \subset A$
为 $A$ 的一个理想中的元素。设 $M$ 为有限 $A$-模。假设
\begin{enumerate}
\item $A$ 是 $f$-进完备的；
\item $H^1_\mathfrak a(M)$ 和 $H^2_\mathfrak a(M)$ 被 $f$ 的某个幂零化。
\end{enumerate}
则令 $U = \Spec(A) \setminus V(\mathfrak a)$，映射
$$
\Gamma(U, \widetilde{M})
\longrightarrow
\lim \Gamma(U, \widetilde{M/f^nM})
$$
为同构。
\end{lemma}

\begin{proof}
可以把引理 \ref{algebraization-lemma-formal-functions-principal} 应用于 $U$ 以及
$\mathcal{F} = \widetilde{M}|_U$，因为 $\mathcal{F}$ 是凝聚
$\mathcal{O}_U$-模范畴中的诺特对象。
由于 $H^1(U, \mathcal{F}) = H^2_\mathfrak a(M)$
（局部上同调，引理
\ref{local-cohomology-lemma-finiteness-pushforwards-and-H1-local}）
被 $f$ 的某个幂零化，可知其 $f$-进 Tate 模为零。
因此该引理表明 $\lim H^0(U, \mathcal{F}/f^n \mathcal{F})$
等于通常的 $f$-进完备化，所完备化的是 $H^0(U, \mathcal{F})$。
考察短正合序列
$$
0 \to M/H^0_\mathfrak a(M) \to H^0(U, \mathcal{F}) \to H^1_\mathfrak a(M) \to 0
$$
它来自局部上同调，引理
\ref{local-cohomology-lemma-finiteness-pushforwards-and-H1-local}。
由于 $M/H^0_\mathfrak a(M)$ 是有限 $A$-模，它是完备的；
参见代数，引理 \ref{algebra-lemma-completion-tensor}。
由于 $H^1_\mathfrak a(M)$ 被 $f$ 的某个幂零化，
由代数，引理 \ref{algebra-lemma-completion-differ-by-torsion} 可知
$H^0(U, \mathcal{F})$ 也完备。证明完毕。
\end{proof}







\section{形式截面的代数化，III}
\label{algebraization-section-algebraization-sections-coherent-III}

\noindent
下一节给出了一个非穷尽的应用列表：将局部上同调完备化方面的材料
应用于拟仿射概形上凝聚模的高阶上同调，以及这些上同调关于某个理想的完备化。

\begin{proposition}
\label{algebraization-proposition-application-higher}
设 $I \subset \mathfrak a$ 为诺特环 $A$ 的理想。
设 $\mathcal{F}$ 是
$U = \Spec(A) \setminus V(\mathfrak a)$ 上的凝聚模。
设 $s \geq 0$。假设
\begin{enumerate}
\item $A$ 是 $I$-进完备的，并且有对偶化复形；
\item 若 $x \in U \setminus V(I)$，则
$\text{depth}(\mathcal{F}_x) > s$，或者
$$
\text{depth}(\mathcal{F}_x) +
\dim(\mathcal{O}_{\overline{\{x\}}, z}) > \text{cd}(A, I) + s + 1
$$
对所有 $z \in V(\mathfrak a) \cap \overline{\{x\}}$ 成立；
\item 下列条件之一成立：
\begin{enumerate}
\item $\mathcal{F}$ 在 $U \setminus V(I)$ 上的限制满足 $(S_{s + 1})$；或者
\item $V(\mathfrak a)$ 的维数至多为 $2$\footnote{这里的含义是：
该函数在 $V(\mathfrak a)$ 上所取最大值与最小值之差（这里该维数函数定义在
$\Spec(A)$ 上）至多为 $2$。}。
\end{enumerate}
\end{enumerate}
则映射
$$
H^i(U, \mathcal{F})
\longrightarrow
\lim H^i(U, \mathcal{F}/I^n\mathcal{F})
$$
在 $i < s$ 时为同构。此外，还有同构
$$
\colim H^s(V, \mathcal{F})
\longrightarrow
\lim H^s(U, \mathcal{F}/I^n\mathcal{F})
$$
其中余极限取遍开子概形 $V \subset U$，这些开子概形包含 $U \cap V(I)$。
\end{proposition}

\begin{proof}
可以假设 $s > 0$，因为 $s = 0$ 的情形已在
命题 \ref{algebraization-proposition-application-H0} 中处理。

\medskip\noindent
取一个有限 $A$-模 $M$，使得 $\mathcal{F}$ 是某个凝聚模
在 $U$ 上的限制，而该凝聚模与 $M$ 相伴；参见局部上同调，引理
\ref{local-cohomology-lemma-finiteness-pushforwards-and-H1-local}。
令 $d = \text{cd}(A, I)$。
设 $\mathfrak p$ 是 $A$ 的一个不包含在 $V(I)$ 中的素理想，
并设 $\mathfrak q \in V(\mathfrak p) \cap V(\mathfrak a)$。
于是或者 $\text{depth}(M_\mathfrak p) \geq s + 1 > s$，
或者由 (2) 有 $\dim((A/\mathfrak p)_\mathfrak q) > d + s + 1$。
由引理 \ref{algebraization-lemma-bootstrap-bis-bis} 可知，
对于 $A, I, V(\mathfrak a), M, s, d$，情形 \ref{algebraization-situation-bootstrap}
的假设均成立。另一方面，引理
\ref{algebraization-lemma-algebraize-local-cohomology-general}
对于 $s + 1$ 和 $d$ 的假设也成立；条件 (3) 正是在这里使用的。

\medskip\noindent
应用引理 \ref{algebraization-lemma-algebraize-local-cohomology-general}，
可得存在理想
$J_0 \subset \mathfrak a$，满足 $V(J_0) \cap V(I) = V(\mathfrak a)$，
使得对于任何 $J \subset J_0$，若 $V(J) \cap V(I) = V(\mathfrak a)$，
映射
$$
H^i_J(M) \longrightarrow H^i(R\Gamma_\mathfrak a(M)^\wedge)
$$
在 $i \leq s + 1$ 时为同构。

\medskip\noindent
当 $i \leq s$ 时，由引理 \ref{algebraization-lemma-bootstrap-inherited} 和
\ref{algebraization-lemma-kill-colimit-support-general}，映射
$H^i_\mathfrak a(M) \to H^i_J(M)$ 为同构。
利用上同调与局部上同调之间的比较
（局部上同调，引理 \ref{local-cohomology-lemma-local-cohomology}），
可推出映射 $H^i(U, \mathcal{F}) \to H^i(V,\mathcal{F})$
在 $V = \Spec(A) \setminus V(J)$ 且 $i < s$ 时为同构。

\medskip\noindent
由定理 \ref{algebraization-theorem-final-bootstrap}，等式
$H^i_\mathfrak a(M) = \lim H^i_\mathfrak a(M/I^nM)$
对 $i \leq s$ 成立。
由引理 \ref{algebraization-lemma-combine-two}，有
$H^{s + 1}_\mathfrak a(M) = \lim H^{s + 1}_\mathfrak a(M/I^nM)$。

\medskip\noindent
由以上结论和命题 \ref{algebraization-proposition-application-H0}，得到同构
$H^0(U, \mathcal{F}) = H^0(V, \mathcal{F}) =
\lim H^0(U, \mathcal{F}/I^n\mathcal{F})$。
当 $0 < i < s$ 时，利用局部上同调与上同调之间的关系，
以相同方式得到所需同构
$H^i(U, \mathcal{F}) = H^i(V, \mathcal{F}) =
\lim H^i(U, \mathcal{F}/I^n\mathcal{F})$；
这比 $i = 0$ 的情形更容易，因为当 $i > 0$ 时有
$$
H^i(U, \mathcal{F}) = H^{i + 1}_\mathfrak a(M),
\quad
H^i(V, \mathcal{F}) = H^{i + 1}_J(M),
\quad
H^i(R\Gamma(U, \mathcal{F})^\wedge) = 
H^{i + 1}(R\Gamma_\mathfrak a(M)^\wedge)
$$
最后一个断言的证明也类似。
\end{proof}

\begin{lemma}
\label{algebraization-lemma-alternative-higher}
设 $A$ 为诺特环。设 $f \in \mathfrak a \subset A$
是 $A$ 的一个理想中的元素。设 $M$ 为有限 $A$-模。
设 $s \geq 0$。假设
\begin{enumerate}
\item $A$ 是 $f$-进完备的；
\item $H^i_\mathfrak a(M)$ 被 $f$ 的某个幂零化，其中 $i \leq s + 1$。
\end{enumerate}
则令 $U = \Spec(A) \setminus V(\mathfrak a)$，映射
$$
H^i(U, \widetilde{M})
\longrightarrow
\lim H^i(U, \widetilde{M/f^nM})
$$
在 $i < s$ 时为同构。
\end{lemma}

\begin{proof}
对 $s$ 作归纳。若 $s = 0$，则断言为空。若 $s = 1$，则结论即为
引理 \ref{algebraization-lemma-alternative-H0}。假设 $s > 1$。由归纳法，
只需对 $i = s - 1 \geq 1$ 证明结论。
可以把引理 \ref{algebraization-lemma-formal-functions-principal}
应用于 $U$ 和 $\mathcal{F} = \widetilde{M}|_U$，
因为 $\mathcal{F}$ 是凝聚 $\mathcal{O}_U$-模范畴中的诺特对象。
由局部上同调，引理
\ref{local-cohomology-lemma-finiteness-pushforwards-and-H1-local}，
有 $H^j(U, \mathcal{F}) = H^{j + 1}_\mathfrak a(M)$，其中 $j$ 任意。
因此，当 $j = s = (s - 1) + 1$ 时，按假设该模被 $f$ 的某个幂零化。
于是由引理 \ref{algebraization-lemma-formal-functions-principal}，
$\lim H^{s - 1}(U, \mathcal{F}/f^n\mathcal{F})$
是关于 $f$ 的通常进完备化，所完备化的模为
$H^{s - 1}(U, \mathcal{F})$。
再次利用该模被 $f$ 的某个幂零化，可见此完备化恰好等于
$H^{s - 1}(U, \mathcal{F})$，这正是所需结论。
\end{proof}





\section{连通性的应用}
\label{algebraization-section-connected}

\noindent
本节讨论 Grothendieck 连通性定理及其变体；原始版本见
\cite[Exposee XIII, Theorem 2.1]{SGA2}。文献中有一个称为
Faltings 连通性定理的版本；我们猜测这指的是
\cite[Theorem 6]{Faltings-some}。下面陈述并证明
\cite[Theorem 1.6]{Varbaro} 中关于完备局部环的最优版本。

\begin{lemma}
\label{algebraization-lemma-punctured-still-connected}
\begin{reference}
\cite[Theorem 1.6]{Varbaro}
\end{reference}
设 $(A, \mathfrak m)$ 为诺特完备局部环。
设 $I$ 为 $A$ 的真理想。
令 $X = \Spec(A)$ 且 $Y = V(I)$。记
\begin{enumerate}
\item $d$ 为 $X$ 的不可约分支的维数最小值；
\item $c$ 为闭子集 $Z \subset X$ 的维数最小值，其中
$X \setminus Z$ 不连通。
\end{enumerate}
则对于闭子集 $Z \subset Y$，$Y \setminus Z$ 在
$\dim(Z) < \min(c, d - 1) - \text{cd}(A, I)$ 时连通。
特别地，$A/I$ 的穿孔谱在 $\text{cd}(A, I) < \min(c, d - 1)$ 时连通。
\end{lemma}

\begin{proof}
先证明最后一个断言。首先，若 $A/I$ 的穿孔谱为空，则由
局部上同调，引理 \ref{local-cohomology-lemma-cd-bound-dim-local}，
$X$ 的每个不可约分支的维数均
$\leq \text{cd}(A, I)$，从而 $\min(c, d - 1) - \text{cd}(A, I) < 0$；
这说明该引理在此情形下成立。因此可以假设
$U \cap Y$ 非空，其中 $U = X \setminus \{\mathfrak m\}$
是 $A$ 的穿孔谱。可以将 $A$ 替换为其既约化。
注意 $A$ 有对偶化复形
（对偶化复形，引理 \ref{dualizing-lemma-ubiquity-dualizing}），
并且 $A$ 关于 $I$ 完备
（代数，引理 \ref{algebra-lemma-complete-by-sub}）。
若假设 $d - 1 > \text{cd}(A, I)$，则应用
引理 \ref{algebraization-lemma-application-theorem} 可得
$$
\colim H^0(V, \mathcal{O}_V)
\longrightarrow
\lim H^0(U, \mathcal{O}_U/I^n\mathcal{O}_U)
$$
为同构，其中余极限取遍开子概形 $V \subset U$，且这些开子概形包含
$U \cap Y$。若 $U \cap Y$ 不连通，则它的第 $n$ 个无穷小邻域在
$U$ 中对所有 $n$ 都不连通，因而右端有一个非平凡幂等元
（这里使用了 $U \cap Y$ 非空这一点）。
因此可以找到不连通的 $V$。
令 $Z = X \setminus V$。由局部上同调，引理
\ref{local-cohomology-lemma-cd-bound-dim-local}，
可知 $Z$ 的每个不可约分支的维数均
$\leq \text{cd}(A, I)$。因此 $c \leq \text{cd}(A, I)$，
这确实证明了最后一个断言。

\medskip\noindent
现在可以从刚刚证明的结论推出引理的断言。设 $Z \subset Y$ 为闭子集，
$Y \setminus Z$ 不连通，且 $\dim(Z) = e$。回忆按照约定，连通空间非空。
因此，或者 (a) $Y = Z$，或者 (b)
$Y \setminus Z = W_1 \amalg W_2$，其中 $W_i$ 在 $Y \setminus Z$ 中
非空且既开又闭。在情形 (b) 中，可以选取 $w_i \in W_i$，
使其在 $U$ 中为闭点；参见态射，引理
\ref{morphisms-lemma-ubiquity-Jacobson-schemes}。
于是可以找到 $f_1, \ldots, f_e \in \mathfrak m$，使得
$V(f_1, \ldots, f_e) \cap Z = \{\mathfrak m\}$，
并且在情形 (b) 中可以假设 $w_i \in V(f_1, \ldots, f_e)$。
具体而言，利用素理想回避，可以归纳地选取 $f_i$，使得
$\dim V(f_1, \ldots, f_i) \cap Z = e - i$，
并使得在情形 (b) 中有 $w_1, w_2 \in V(f_i)$。
由此可知 $A/I + (f_1, \ldots, f_e)$ 的穿孔谱不连通
（略去一个小细节）。由于由局部上同调，引理
\ref{local-cohomology-lemma-cd-sum} 和
\ref{local-cohomology-lemma-bound-cd} 有
$\text{cd}(A, I + (f_1, \ldots, f_e)) \leq \text{cd}(A, I) + e$，故
$$
\text{cd}(A, I) + e \geq \min(c, d - 1)
$$
由证明的第一部分可得此式。这推出
$e \geq \min(c, d - 1) - \text{cd}(A, I)$，这正是所需结论。
\end{proof}

\begin{lemma}
\label{algebraization-lemma-connected}
设 $I \subset \mathfrak a$ 为诺特环 $A$ 的理想。假设
\begin{enumerate}
\item $A$ 是 $I$-进完备的，并且有对偶化复形；
\item 若 $\mathfrak p \subset A$ 是一个不包含在
$V(I)$ 中的极小素理想，并且 $\mathfrak q \in V(\mathfrak p) \cap V(\mathfrak a)$，则
$\dim((A/\mathfrak p)_\mathfrak q) > \text{cd}(A, I) + 1$；
\item 任何非空开子集 $V \subset \Spec(A)$，只要其包含
$V(I) \setminus V(\mathfrak a)$，就是连通的\footnote{例如，当 $A$ 为整环时。}。
\end{enumerate}
则 $V(I) \setminus V(\mathfrak a)$ 或为空，或为连通的。
\end{lemma}

\begin{proof}
可以将 $A$ 替换为其既约化。于是对于 $A$ 的所有相伴素理想，
(2) 中的不等式均成立。由命题 \ref{algebraization-proposition-application-H0} 可知
$$
\colim H^0(V, \mathcal{O}_V) = \lim H^0(T_n, \mathcal{O}_{T_n})
$$
其中余极限取遍 (3) 中的开子集 $V$，而 $T_n$ 表示第 $n$ 个
无穷小邻域；这是 $T = V(I) \setminus V(\mathfrak a)$ 在
$U = \Spec(A) \setminus V(\mathfrak a)$ 中的无穷小邻域。
因此 $T$ 或为空，或为连通的；否则右端将有非平凡幂等元，
而我们已假设左端没有。略去若干细节。
\end{proof}

\begin{lemma}
\label{algebraization-lemma-H0}
设 $A$ 为有对偶化复形的诺特整环，并且关于某个非零元
$f \in A$ 完备。设 $f \in \mathfrak a \subset A$，并把这里的中间项取为理想。
假设 $Z = V(\mathfrak a)$ 的每个不可约分支的
余维均 $> 2$，这里是在 $X = \Spec(A)$ 中计算；也就是说，假设 $Z$ 的每个
不可约分支的余维均 $> 1$，这里是在 $Y = V(f)$ 中计算。则 $Y \setminus Z$ 连通。
\end{lemma}

\begin{proof}
这是引理 \ref{algebraization-lemma-connected} 的特殊情形（该引理的证明依赖
命题 \ref{algebraization-proposition-application-H0}）。下面用较容易的命题
\ref{algebraization-proposition-alternative-colim-H0} 来证明它。

\medskip\noindent
令 $U = X \setminus Z$。由命题 \ref{algebraization-proposition-alternative-colim-H0}，
有同构
$$
\colim \Gamma(V, \mathcal{O}_V) \to
\lim_n \Gamma(U, \mathcal{O}_U/f^n \mathcal{O}_U)
$$
其中余极限取遍开子集 $V \subset U$，这些开子集包含 $U \cap Y$。
因此，若 $U \cap Y$ 不连通，则对于某个 $V$，
$\Gamma(V, \mathcal{O}_V)$ 中存在非平凡幂等元。这不可能，
因为 $V$ 是整概形，而 $X$ 是整环的谱。
\end{proof}






\section{完备化函子}
\label{algebraization-section-completion}

\noindent
设 $X$ 为诺特概形。设 $Y \subset X$ 为闭子概形，
其拟凝聚理想层为 $\mathcal{I} \subset \mathcal{O}_X$。
本节考虑凝聚 $\mathcal{O}_X$-模 $(\mathcal{F}_n)$ 的逆系统，其中
$\mathcal{F}_n$ 被 $I^n$ 零化，并且转移映射诱导同构
$\mathcal{F}_{n + 1}/I^n\mathcal{F}_{n + 1} \to \mathcal{F}_n$。
这些逆系统所成的范畴在概形上同调，第
\ref{coherent-section-coherent-formal} 节中记为
$$
\textit{Coh}(X, \mathcal{I})
$$
。这个范畴等价于 $X$ 沿 $Y$ 的形式完备化上的凝聚模范畴；
但是，我们还没有引入形式概形或其上的凝聚模，
所以这里不能使用这一术语。我们特别关注完备化函子
$$
\textit{Coh}(\mathcal{O}_X)
\longrightarrow
\textit{Coh}(X, \mathcal{I}),\quad
\mathcal{F} \longmapsto \mathcal{F}^\wedge
$$
；参见概形上同调，等式 (\ref{coherent-equation-completion-functor})。

\begin{lemma}
\label{algebraization-lemma-completion-fully-faithful}
设 $X$ 为诺特概形，$Y \subset X$ 为闭子概形。
设 $Y_n \subset X$ 为第 $n$ 个无穷小邻域；这是 $Y$ 在 $X$ 中的邻域。
考虑下列条件：
\begin{enumerate}
\item $X$ 是拟仿射的，并且
$\Gamma(X, \mathcal{O}_X) \to \lim \Gamma(Y_n, \mathcal{O}_{Y_n})$
为同构；
\item $X$ 有充足可逆模 $\mathcal{L}$，并且
$\Gamma(X, \mathcal{L}^{\otimes m}) \to
\lim \Gamma(Y_n, \mathcal{L}^{\otimes m}|_{Y_n})$
对所有 $m \gg 0$ 都是同构；
\item 对每个有限局部自由 $\mathcal{O}_X$-模
$\mathcal{E}$，映射
$\Gamma(X, \mathcal{E}) \to \lim \Gamma(Y_n, \mathcal{E}|_{Y_n})$
为同构；
\item 完备化函子
$\textit{Coh}(\mathcal{O}_X) \to \textit{Coh}(X, \mathcal{I})$
在有限局部自由对象所成的全子范畴上是全忠实的。
\end{enumerate}
则 (1) $\Rightarrow$ (2) $\Rightarrow$ (3) $\Rightarrow$ (4)，
并且 (4) $\Rightarrow$ (3)。
\end{lemma}

\begin{proof}
证明 (3) $\Rightarrow$ (4)。若 $\mathcal{F}$ 和 $\mathcal{G}$
在 $X$ 上有限局部自由，则考虑
$\mathcal{H} = \SheafHom_{\mathcal{O}_X}(\mathcal{G}, \mathcal{F})$，
并使用概形上同调，引理
\ref{coherent-lemma-completion-internal-hom}，
可知 (3) 推出 (4)。

\medskip\noindent
证明 (2) $\rightarrow$ (3)。设 $\mathcal{L}$ 在 $X$ 上充足，
并设 $\mathcal{E}$ 为有限局部自由 $\mathcal{O}_X$-模。
我们断言可以找到泛正合序列
$$
0 \to \mathcal{E} \to
(\mathcal{L}^{\otimes p})^{\oplus r} \to
(\mathcal{L}^{\otimes q})^{\oplus s}
$$
，其中某些 $r, s \geq 0$ 且 $0 \ll p \ll q$。若此断言成立，则使用正合序列
$$
0 \to \lim \Gamma(\mathcal{E}|_{Y_n}) \to
\lim \Gamma((\mathcal{L}^{\otimes p})^{\oplus r}|_{Y_n}) \to
\lim \Gamma((\mathcal{L}^{\otimes q})^{\oplus s}|_{Y_n})
$$
以及 (2) 中的同构，就得到 (3) 中的同构。
为证明该断言，考虑对偶局部自由模
$\SheafHom_{\mathcal{O}_X}(\mathcal{E}, \mathcal{O}_X)$
，并应用性质，命题 \ref{properties-proposition-characterize-ample}，
得到满射
$$
(\mathcal{L}^{\otimes -p})^{\oplus r}
\longrightarrow
\SheafHom_{\mathcal{O}_X}(\mathcal{E}, \mathcal{O}_X)
$$
取对偶便得到正合序列中的第一个映射
（它是泛单射，因为满射性是泛的）。对余核重复这一过程便得到第二个映射。
略去若干细节。

\medskip\noindent
证明 (1) $\Rightarrow$ (2)。这是因为若 $X$ 是拟仿射的，
则 $\mathcal{O}_X$ 是充足可逆模；参见性质，引理
\ref{properties-lemma-quasi-affine-O-ample}。

\medskip\noindent
略去 (4) $\Rightarrow$ (3) 的证明。
\end{proof}

\noindent
给定诺特概形和拟凝聚理想层
$\mathcal{I} \subset \mathcal{O}_X$，若对象 $(\mathcal{F}_n)$ 属于
$\textit{Coh}(X, \mathcal{I})$，且每个 $\mathcal{F}_n$ 都是有限局部自由
$\mathcal{O}_X/\mathcal{I}^n$-模，则称该对象为{\it 有限局部自由的}。

\begin{lemma}
\label{algebraization-lemma-completion-fully-faithful-general}
设 $X$ 为诺特概形，$Y \subset X$ 为闭子概形，
其理想层为 $\mathcal{I} \subset \mathcal{O}_X$。
设 $Y_n \subset X$ 为第 $n$ 个无穷小邻域；这是 $Y$ 在 $X$ 中的邻域。
设 $\mathcal{V}$ 为开子概形 $V \subset X$ 所成的集合，这些开子概形包含
$Y$；该集合按反包含排序。
\begin{enumerate}
\item $X$ 是拟仿射的，并且
$$
\colim_\mathcal{V} \Gamma(V, \mathcal{O}_V)
\longrightarrow
\lim \Gamma(Y_n, \mathcal{O}_{Y_n})
$$
为同构；
\item $X$ 有充足可逆模 $\mathcal{L}$，并且
$$
\colim_\mathcal{V} \Gamma(V, \mathcal{L}^{\otimes m})
\longrightarrow
\lim \Gamma(Y_n, \mathcal{L}^{\otimes m}|_{Y_n})
$$
对所有 $m \gg 0$ 都是同构；
\item 对每个 $V \in \mathcal{V}$ 和每个有限局部自由
$\mathcal{O}_V$-模 $\mathcal{E}$，映射
$$
\colim_{V' \geq V} \Gamma(V', \mathcal{E}|_{V'})
\longrightarrow
\lim \Gamma(Y_n, \mathcal{E}|_{Y_n})
$$
为同构；
\item 完备化函子
$$
\colim_\mathcal{V} \textit{Coh}(\mathcal{O}_V)
\longrightarrow
\textit{Coh}(X, \mathcal{I}),
\quad
\mathcal{F} \longmapsto \mathcal{F}^\wedge
$$
在有限局部自由对象所成的全子范畴上是全忠实的
（参见上面的说明）。
\end{enumerate}
则 (1) $\Rightarrow$ (2) $\Rightarrow$ (3) $\Rightarrow$ (4)，
并且 (4) $\Rightarrow$ (3)。
\end{lemma}

\begin{proof}
注意 $\mathcal{V}$ 是有向集，所以这些余极限就是范畴，第
\ref{categories-section-directed-colimits} 节中的余极限。
其余论证几乎与引理 \ref{algebraization-lemma-completion-fully-faithful} 的证明
中的论证完全相同；我们建议读者跳过它。

\medskip\noindent
证明 (3) $\Rightarrow$ (4)。若 $\mathcal{F}$ 和 $\mathcal{G}$
在 $V \in \mathcal{V}$ 上有限局部自由，则考虑
$\mathcal{H} = \SheafHom_{\mathcal{O}_V}(\mathcal{G}, \mathcal{F})$，
并使用概形上同调，引理
\ref{coherent-lemma-completion-internal-hom}，
可知 (3) 推出 (4)。

\medskip\noindent
证明 (2) $\Rightarrow$ (3)。设 $\mathcal{L}$ 在 $X$ 上充足，
并设 $\mathcal{E}$ 是有限局部自由 $\mathcal{O}_V$-模，
其中某个 $V \in \mathcal{V}$。
我们断言可以找到泛正合序列
$$
0 \to \mathcal{E} \to
(\mathcal{L}^{\otimes p})^{\oplus r}|_{V} \to
(\mathcal{L}^{\otimes q})^{\oplus s}|_{V}
$$
，其中某些 $r, s \geq 0$ 且 $0 \ll p \ll q$。若此断言成立，
则 (2) 中的同构将推出 (3) 中的同构。
为证明该断言，回忆 $\mathcal{L}|_V$ 是充足的；参见性质，引理
\ref{properties-lemma-ample-on-locally-closed}。
考虑对偶局部自由模
$\SheafHom_{\mathcal{O}_V}(\mathcal{E}, \mathcal{O}_V)$
，并应用性质，命题 \ref{properties-proposition-characterize-ample}，
得到满射
$$
(\mathcal{L}^{\otimes -p})^{\oplus r}|_V \longrightarrow
\SheafHom_{\mathcal{O}_V}(\mathcal{E}, \mathcal{O}_V)
$$
（它是泛单射，因为满射性是泛的）。
取对偶便得到正合序列中的第一个映射。
对余核重复这一过程便得到第二个映射。略去若干细节。

\medskip\noindent
证明 (1) $\Rightarrow$ (2)。这是因为若 $X$ 是拟仿射的，
则 $\mathcal{O}_X$ 是充足可逆模；参见性质，引理
\ref{properties-lemma-quasi-affine-O-ample}。

\medskip\noindent
略去 (4) $\Rightarrow$ (3) 的证明。
\end{proof}

\begin{lemma}
\label{algebraization-lemma-recognize-formal-coherent-modules}
设 $X$ 为诺特概形。设 $\mathcal{I} \subset \mathcal{O}_X$
为拟凝聚理想层。函子
$$
\textit{Coh}(X, \mathcal{I}) \longrightarrow \text{Pro-}\QCoh(\mathcal{O}_X)
$$
是全忠实的；参见范畴，注记 \ref{categories-remark-pro-category}。
\end{lemma}

\begin{proof}
设 $(\mathcal{F}_n)$ 和 $(\mathcal{G}_n)$ 为
$\textit{Coh}(X, \mathcal{I})$ 的对象。pro-对象态射 $\alpha$ 从
$(\mathcal{F}_n)$ 到 $(\mathcal{G}_n)$，它由一族与转移映射相容的映射
$\alpha_n : \mathcal{F}_{m(n)} \to \mathcal{G}_n$ 给出，其中
$\mathbf{N} \to \mathbf{N}$，$n \mapsto m(n)$ 是递增函数
（特别地，$m(n) \geq n$）。由于
$\mathcal{F}_n = \mathcal{F}_{m(n)}/\mathcal{I}^n\mathcal{F}_{m(n)}$，
并且 $\mathcal{G}_n$ 被 $\mathcal{I}^n$ 零化，
可知 $\alpha_n$ 诱导映射 $\mathcal{F}_n \to \mathcal{G}_n$。
\end{proof}

\noindent
下面给出若干例子，说明利用本章前面完成的工作能够证明何种全忠实性结果。

\begin{lemma}
\label{algebraization-lemma-fully-faithful}
设 $I \subset \mathfrak a$ 为诺特环 $A$ 的理想。
令 $U = \Spec(A) \setminus V(\mathfrak a)$。假设
\begin{enumerate}
\item $A$ 是 $I$-进完备的，并且有对偶化复形；
\item 对任意满足下列条件的相伴素理想 $\mathfrak p \subset A$：
$\mathfrak p \not \in V(I)$、
$V(\mathfrak p) \cap V(I) \not \subset V(\mathfrak a)$，且
$\mathfrak q \in V(\mathfrak p) \cap V(\mathfrak a)$，都有
$\dim((A/\mathfrak p)_\mathfrak q) > \text{cd}(A, I) + 1$；
\item 对 $\mathfrak p \subset A$，若 $\mathfrak p \not \in V(I)$ 且
$V(\mathfrak p) \cap V(I) \subset V(\mathfrak a)$，
有 $\text{depth}(A_\mathfrak p) \geq 2$。
\end{enumerate}
则完备化函子
$$
\textit{Coh}(\mathcal{O}_U)
\longrightarrow
\textit{Coh}(U, I\mathcal{O}_U),
\quad
\mathcal{F} \longmapsto \mathcal{F}^\wedge
$$
在有限局部自由对象所成的全子范畴上是全忠实的。
\end{lemma}

\begin{proof}
由引理 \ref{algebraization-lemma-completion-fully-faithful}，
只需证明
$$
\Gamma(U, \mathcal{O}_U) =
\lim \Gamma(U, \mathcal{O}_U/I^n\mathcal{O}_U)
$$
。这立即由引理 \ref{algebraization-lemma-application-H0} 得到。
\end{proof}

\begin{lemma}
\label{algebraization-lemma-fully-faithful-alternative}
设 $A$ 为诺特环。设 $f \in \mathfrak a \subset A$
为 $A$ 的一个理想中的元素。令 $U = \Spec(A) \setminus V(\mathfrak a)$。
假设
\begin{enumerate}
\item $A$ 是 $f$-进完备的；
\item $H^1_\mathfrak a(A)$ 和 $H^2_\mathfrak a(A)$ 被 $f$ 的某个幂零化。
\end{enumerate}
则完备化函子
$$
\textit{Coh}(\mathcal{O}_U)
\longrightarrow
\textit{Coh}(U, I\mathcal{O}_U),
\quad
\mathcal{F} \longmapsto \mathcal{F}^\wedge
$$
在有限局部自由对象所成的全子范畴上是全忠实的。
\end{lemma}

\begin{proof}
由引理 \ref{algebraization-lemma-completion-fully-faithful}，
只需证明
$$
\Gamma(U, \mathcal{O}_U) =
\lim \Gamma(U, \mathcal{O}_U/I^n\mathcal{O}_U)
$$
。这立即由引理 \ref{algebraization-lemma-alternative-H0} 得到。
\end{proof}

\begin{lemma}
\label{algebraization-lemma-fully-faithful-simple-two}
设 $A$ 为诺特环。设 $f \in \mathfrak a$ 为 $A$ 的一个理想中的元素。
令 $U = \Spec(A) \setminus V(\mathfrak a)$。假设
\begin{enumerate}
\item $A$ 有对偶化复形，并且关于 $f$ 完备；
\item 对每个素理想 $\mathfrak p \subset A$，若 $f \not \in \mathfrak p$，
且 $\mathfrak q \in V(\mathfrak p) \cap V(\mathfrak a)$，则
$\text{depth}(A_\mathfrak p) + \dim((A/\mathfrak p)_\mathfrak q) > 2$。
\end{enumerate}
则完备化函子
$$
\textit{Coh}(\mathcal{O}_U)
\longrightarrow
\textit{Coh}(U, I\mathcal{O}_U),
\quad
\mathcal{F} \longmapsto \mathcal{F}^\wedge
$$
在有限局部自由对象所成的全子范畴上是全忠实的。
\end{lemma}

\begin{proof}
结论由引理 \ref{algebraization-lemma-fully-faithful-alternative} 和
局部上同调，命题 \ref{local-cohomology-proposition-annihilator} 得到。
\end{proof}

\begin{lemma}
\label{algebraization-lemma-fully-faithful-general}
设 $I \subset \mathfrak a \subset A$ 为诺特环 $A$ 的理想。
令 $U = \Spec(A) \setminus V(\mathfrak a)$。设 $\mathcal{V}$ 为
$U$ 的开子概形所成的集合，这些开子概形包含 $U \cap V(I)$，
并按反包含排序。假设
\begin{enumerate}
\item $A$ 是 $I$-进完备的，并且有对偶化复形；
\item 对任意满足下列条件的相伴素理想
$\mathfrak p \subset A$：
$I \not \subset \mathfrak p$，
$V(\mathfrak p) \cap V(I) \not \subset V(\mathfrak a)$，
且 $\mathfrak q \in V(\mathfrak p) \cap V(\mathfrak a)$，都有
$\dim((A/\mathfrak p)_\mathfrak q) > \text{cd}(A, I) + 1$。
\end{enumerate}
则完备化函子
$$
\colim_\mathcal{V} \textit{Coh}(\mathcal{O}_V)
\longrightarrow
\textit{Coh}(U, I\mathcal{O}_U),
\quad
\mathcal{F} \longmapsto \mathcal{F}^\wedge
$$
在有限局部自由对象所成的全子范畴上是全忠实的。
\end{lemma}

\begin{proof}
由引理 \ref{algebraization-lemma-completion-fully-faithful-general}，
只需证明
$$
\colim_\mathcal{V} \Gamma(V, \mathcal{O}_V) =
\lim \Gamma(U, \mathcal{O}_U/I^n\mathcal{O}_U)
$$
。这立即由命题 \ref{algebraization-proposition-application-H0} 得到。
\end{proof}

\begin{lemma}
\label{algebraization-lemma-fully-faithful-general-alternative}
设 $A$ 为诺特环。设 $f \in \mathfrak a \subset A$
为 $A$ 的一个理想中的元素。令 $U = \Spec(A) \setminus V(\mathfrak a)$。
设 $\mathcal{V}$ 为 $U$ 的开子概形所成的集合，这些开子概形包含
$U \cap V(f)$，并按反包含排序。假设
\begin{enumerate}
\item $A$ 是 $f$-进完备的；
\item $f$ 是非零因子；
\item $H^1_\mathfrak a(A/fA)$ 是有限 $A$-模。
\end{enumerate}
则完备化函子
$$
\colim_\mathcal{V} \textit{Coh}(\mathcal{O}_V)
\longrightarrow
\textit{Coh}(U, f\mathcal{O}_U),
\quad
\mathcal{F} \longmapsto \mathcal{F}^\wedge
$$
在有限局部自由对象所成的全子范畴上是全忠实的。
\end{lemma}

\begin{proof}
由引理 \ref{algebraization-lemma-completion-fully-faithful-general}，
只需证明
$$
\colim_\mathcal{V} \Gamma(V, \mathcal{O}_V) =
\lim \Gamma(U, \mathcal{O}_U/I^n\mathcal{O}_U)
$$
。这立即由引理 \ref{algebraization-lemma-alternative-colim-H0} 得到。
\end{proof}

\begin{lemma}
\label{algebraization-lemma-fully-faithful-very-general}
设 $I \subset \mathfrak a \subset A$ 为诺特环 $A$ 的理想。
令 $U = \Spec(A) \setminus V(\mathfrak a)$。设 $\mathcal{V}$ 为
$U$ 的开子概形所成的集合，这些开子概形包含 $U \cap V(I)$，
并按反包含排序。设 $\mathcal{F}$ 和
$\mathcal{G}$ 是凝聚 $\mathcal{O}_V$-模，其中某个 $V \in \mathcal{V}$。
映射
$$
\colim_{V' \geq V} \Hom_V(\mathcal{G}|_{V'}, \mathcal{F}|_{V'})
\longrightarrow
\Hom_{\textit{Coh}(U, I\mathcal{O}_U)}(\mathcal{G}^\wedge, \mathcal{F}^\wedge)
$$
在下列假设成立时是双射：
\begin{enumerate}
\item $A$ 是 $I$-进完备的，并且有对偶化复形；
\item 若 $x \in \text{Ass}(\mathcal{F})$，$x \not \in V(I)$，
$\overline{\{x\}} \cap V(I) \not \subset V(\mathfrak a)$，
且 $z \in \overline{\{x\}} \cap V(\mathfrak a)$，则
$\dim(\mathcal{O}_{\overline{\{x\}}, z}) > \text{cd}(A, I) + 1$。
\end{enumerate}
\end{lemma}

\begin{proof}
可以选取凝聚 $\mathcal{O}_U$-模
$\mathcal{F}'$ 和 $\mathcal{G}'$，它们在 $V$ 上的限制
分别为 $\mathcal{F}$ 和 $\mathcal{G}$；参见性质，引理
\ref{properties-lemma-lift-finite-presentation}。
可以修改对 $\mathcal{F}'$ 的选择，以保证
$\text{Ass}(\mathcal{F}') \subset V$；例如参见局部上同调，引理
\ref{local-cohomology-lemma-get-depth-1-along-Z}。
因此可以而且确实把 $V$ 换成 $U$，并把 $\mathcal{F}$ 和 $\mathcal{G}$
换成 $\mathcal{F}'$ 和 $\mathcal{G}'$。
令 $\mathcal{H} = \SheafHom_{\mathcal{O}_U}(\mathcal{G}, \mathcal{F})$。
这是凝聚 $\mathcal{O}_U$-模。我们有
$$
\Hom_V(\mathcal{G}|_V, \mathcal{F}|_V) =
H^0(V, \mathcal{H})
\quad\text{且}\quad
\lim H^0(U, \mathcal{H}/\mathcal{I}^n\mathcal{H}) =
\Mor_{\textit{Coh}(U, I\mathcal{O}_U)}
(\mathcal{G}^\wedge, \mathcal{F}^\wedge)
$$
；参见概形上同调，引理 \ref{coherent-lemma-completion-internal-hom}。
因此，如果能够证明命题 \ref{algebraization-proposition-application-H0}
的假设对 $\mathcal{H}$ 成立，那么证明就完成了。
这是因为
$\text{Ass}(\mathcal{H}) \subset \text{Ass}(\mathcal{F})$；
参见概形上同调，引理
\ref{coherent-lemma-hom-into-depth}。
\end{proof}








\section{凝聚形式模的代数化，I}
\label{algebraization-section-algebraization-modules}

\noindent
完备化函子的本质满性（见下文）在
\cite{SGA2}、\cite{MRaynaud-book} 和 \cite{MRaynaud-paper}
中得到了系统研究。我们在下列仿射情形中工作。

\begin{situation}
\label{algebraization-situation-algebraize}
这里 $A$ 是诺特环，$I \subset \mathfrak a \subset A$ 是理想。
令 $X = \Spec(A)$，$Y = V(I) = \Spec(A/I)$，并令
$Z = V(\mathfrak a) = \Spec(A/\mathfrak a)$。此外，$U = X \setminus Z$。
\end{situation}

\noindent
本节试图寻找条件，以保证 $\textit{Coh}(U, I\mathcal{O}_U)$ 中的对象
位于完备化函子
$\textit{Coh}(\mathcal{O}_U) \to \textit{Coh}(U, I\mathcal{O}_U)$
的像中。参见概形上同调，第 \ref{coherent-section-coherent-formal} 节和
第 \ref{algebraization-section-completion} 节。

\begin{lemma}
\label{algebraization-lemma-system-of-modules}
在情形 \ref{algebraization-situation-algebraize} 中，考虑满足下列条件的
逆系统 $(M_n)$，它由 $A$-模组成：
\begin{enumerate}
\item $M_n$ 是有限 $A$-模；
\item $M_n$ 被 $I^n$ 零化；
\item $M_{n + 1}/I^nM_{n + 1} \to M_n$ 的核与余核
都是 $\mathfrak a$-幂挠的。
\end{enumerate}
则 $(\widetilde{M}_n|_U)$ 属于 $\textit{Coh}(U, I\mathcal{O}_U)$。
反之，$\textit{Coh}(U, I\mathcal{O}_U)$ 的每个对象都以这种方式得到。
\end{lemma}

\begin{proof}
略去对 $(\widetilde{M}_n|_U)$ 属于
$\textit{Coh}(U, I\mathcal{O}_U)$ 的验证。设 $(\mathcal{F}_n)$
为 $\textit{Coh}(U, I\mathcal{O}_U)$ 的对象。
由局部上同调，引理
\ref{local-cohomology-lemma-finiteness-pushforwards-and-H1-local}
可知 $\mathcal{F}_n = \widetilde{M_n}$，其中 $A/I^n$-模 $M_n$ 有限。
将 $M_n$ 除以 $H^0_\mathfrak a(M_n)$ 后，
可以假设 $M_n \subset H^0(U, \mathcal{F}_n)$；参见对偶化复形，引理
\ref{dualizing-lemma-divide-by-torsion} 以及已经引用的引理。
归纳地把 $M_{n + 1}$ 替换为 $M_n$ 在映射
$M_{n + 1} \to H^0(U, \mathcal{F}_{n + 1})
\to H^0(U, \mathcal{F}_n)$ 下的逆像之后，可以假设 $M_{n + 1}$ 映入
$M_n$。这给出满足 (1) 和 (2) 的逆系统 $(M_n)$，使得
$\mathcal{F}_n = \widetilde{M_n}$。为说明 (3) 成立，使用如下事实：
$M_{n + 1}/I^nM_{n + 1} \to M_n$ 是有限 $A$-模之间的映射，
在应用 $\widetilde{\ }$ 并限制到 $U$ 后诱导同构
（这里再次使用第一个所引引理）。
\end{proof}

\noindent
在情形 \ref{algebraization-situation-algebraize} 中，可以研究概形上同调，
等式 (\ref{coherent-equation-completion-functor}) 的完备化函子
\begin{equation}
\label{algebraization-equation-completion}
\textit{Coh}(\mathcal{O}_U)
\longrightarrow
\textit{Coh}(U, I\mathcal{O}_U),\quad
\mathcal{F} \longmapsto \mathcal{F}^\wedge
\end{equation}
若 $A$ 是 $I$-进完备的，则由前面对形式截面代数化的研究，
这个函子在适当的子范畴上全忠实；一些示例性结论见第
\ref{algebraization-section-completion} 节和引理 \ref{algebraization-lemma-fully-faithful-inequalities}。
接着，设 $(\mathcal{F}_n)$ 为 $\textit{Coh}(U, I\mathcal{O}_U)$ 的对象。
仍假设 $A$ 是 $I$-进完备的，可以问：$(\mathcal{F}_n)$ 何时位于
上面所示完备化函子的本质像中？

\begin{lemma}
\label{algebraization-lemma-essential-image-completion}
在情形 \ref{algebraization-situation-algebraize} 中，设 $(\mathcal{F}_n)$
为 $\textit{Coh}(U, I\mathcal{O}_U)$ 的对象。考虑下列条件：
\begin{enumerate}
\item $(\mathcal{F}_n)$ 位于函子 (\ref{algebraization-equation-completion}) 的本质像中；
\item $(\mathcal{F}_n)$ 是某个凝聚 $\mathcal{O}_U$-模的完备化；
\item $(\mathcal{F}_n)$ 是某个凝聚 $\mathcal{O}_V$-模的完备化，
其中 $U \cap Y \subset V \subset U$ 为开子概形；
\item $(\mathcal{F}_n)$ 是某个模在 $U$ 上的限制的完备化，而该模是凝聚
$\mathcal{O}_X$-模；
\item $(\mathcal{F}_n)$ 是某个完备化在 $U$ 上的限制，而被完备化者是凝聚
$\mathcal{O}_X$-模；
\item 存在对象 $(\mathcal{G}_n)$，它属于 $\textit{Coh}(X, I\mathcal{O}_X)$，
其在 $U$ 上的限制为 $(\mathcal{F}_n)$。
\end{enumerate}
则条件 (1)、(2)、(3)、(4) 和 (5) 等价，并且蕴含 (6)。
若 $A$ 是 $I$-进完备的，则条件 (6) 蕴含其余条件。
\end{lemma}

\begin{proof}
条件 (1) 和 (2) 等价，因为凝聚 $\mathcal{O}_U$-模 $\mathcal{F}$ 的完备化，
按定义就是 $\mathcal{F}$ 在函子 (\ref{algebraization-equation-completion}) 下的像。
若 $V \subset U$ 是包含 $U \cap Y$ 的开子概形，则有
$$
\textit{Coh}(V, I\mathcal{O}_V) =
\textit{Coh}(U, I\mathcal{O}_U)
$$
，因为凝聚 $\mathcal{O}_V$-模的范畴（这些模支撑在 $V \cap Y$ 上），
与凝聚 $\mathcal{O}_U$-模的范畴（这些模支撑在 $U \cap Y$ 上）相同。
因此，凝聚 $\mathcal{O}_V$-模的完备化是
$\textit{Coh}(U, I\mathcal{O}_U)$ 的对象。由此，(2)、(3)、(4) 和 (5)
等价，因为函子
$\textit{Coh}(\mathcal{O}_X) \to \textit{Coh}(\mathcal{O}_U) \to
\textit{Coh}(\mathcal{O}_V)$ 是本质满的。
参见性质，引理 \ref{properties-lemma-lift-finite-presentation}。

\medskip\noindent
(5) 总是蕴含 (6)。假设 $A$ 是 $I$-进完备的。
于是由概形上同调，引理 \ref{coherent-lemma-inverse-systems-affine}，
$\textit{Coh}(X, I\mathcal{O}_X)$ 的任意对象都对应一个有限 $A$-模。
所以在此情形下，(6) 蕴含 (5)。
\end{proof}

\begin{example}
\label{algebraization-example-not-algebraizable}
设 $k$ 为域。令 $A = k[x, y][[t]]$，其中 $I = (t)$ 且
$\mathfrak a = (x, y, t)$。使用情形 \ref{algebraization-situation-algebraize} 中的记号。
注意
$U \cap Y = (D(x) \cap Y) \cup (D(y) \cap Y)$ 是仿射开覆盖。
对 $n \geq 1$，考虑可逆模 $\mathcal{L}_n$；它是
$\mathcal{O}_U/t^n\mathcal{O}_U$-模，由 $A_x/t^nA_x$ 与 $A_y/t^nA_y$ 粘合而成，
粘合使用 $A_{xy}/t^nA_{xy}$ 中的可逆元，该元是任意如下形式幂级数的像：
$$
u = 1 + \frac{t}{xy} + \sum_{n \geq 2} a_n \frac{t^n}{(xy)^{\varphi(n)}}
$$
其中 $a_n \in k[x, y]$ 且 $\varphi(n) \in \mathbf{N}$。
于是 $(\mathcal{L}_n)$ 是 $\textit{Coh}(U, I\mathcal{O}_U)$ 的可逆对象，
但它不是某个凝聚 $\mathcal{O}_U$-模 $\mathcal{L}$ 的完备化。
这里只概述论证，并略去大部分细节。
设 $y \in U \cap Y$。那么茎 $\mathcal{L}_y$ 的完备化将是可逆模，
因而 $\mathcal{L}_y$ 可逆。于是将存在开子概形
$V \subset U$，它包含 $U \cap Y$，并且 $\mathcal{L}|_V$ 可逆。
由除子，引理 \ref{divisors-lemma-extend-invertible-module}，
可得一个可逆 $A$-模 $M$，满足
$\widetilde{M}|_V \cong \mathcal{L}|_V$。然而环 $A$ 是唯一分解整环，
所以 $M \cong A$，这将推出
$\mathcal{L}_n \cong \mathcal{O}_U/I^n\mathcal{O}_U$。
但由构造有 $\mathcal{L}_2 \not \cong \mathcal{O}_U/I^2\mathcal{O}_U$，
于是得到所需矛盾。

\medskip\noindent
注意，若取 $a_n = 0$，其中 $n \geq 2$，则可见
$\lim H^0(U, \mathcal{L}_n)$ 非零：在这种情形下，函数 $x$
（在 $D(x)$ 上）与函数 $x + t/y$（在 $D(y)$ 上）可以粘合。
另一方面，若取 $a_n = 1$ 且
$\varphi(n) = 2^n$，甚至取 $\varphi(n) = n^2$，则读者可以证明
$\lim H^0(U, \mathcal{L}_n)$ 为零；这给出了在此情形下
$(\mathcal{L}_n)$ 不可代数化的另一证明。
\end{example}

\noindent
若在情形 \ref{algebraization-situation-algebraize} 中环 $A$ 不是
$I$-进完备的，则引理 \ref{algebraization-lemma-essential-image-completion} 表明，
正确的问题是 $(\mathcal{F}_n)$ 是否位于限制函子的本质像中：
$$
\textit{Coh}(X, I\mathcal{O}_X)
\longrightarrow
\textit{Coh}(U, I\mathcal{O}_U)
$$
然而，这时不能再说这意味着 $(\mathcal{F}_n)$ 可代数化。
因此引入下列术语。

\begin{definition}
\label{algebraization-definition-algebraizable}
在情形 \ref{algebraization-situation-algebraize} 中，设 $(\mathcal{F}_n)$ 为
$\textit{Coh}(U, I\mathcal{O}_U)$ 的对象。称
{\it $(\mathcal{F}_n)$ 可延拓到 $X$}，如果存在对象
$(\mathcal{G}_n)$ 属于 $\textit{Coh}(X, I\mathcal{O}_X)$，且其在
$U$ 上的限制同构于 $(\mathcal{F}_n)$。
\end{definition}

\noindent
这个概念等价于在完备化上可代数化。

\begin{lemma}
\label{algebraization-lemma-algebraizable}
在情形 \ref{algebraization-situation-algebraize} 中，设 $(\mathcal{F}_n)$ 为
$\textit{Coh}(U, I\mathcal{O}_U)$ 的对象。设 $A', I', \mathfrak a'$
为 $I$-进完备化，所完备化的对象分别是 $A, I, \mathfrak a$。令 $X' = \Spec(A')$，
且令 $U' = X' \setminus V(\mathfrak a')$。下列条件等价：
\begin{enumerate}
\item $(\mathcal{F}_n)$ 可延拓到 $X$；
\item $(\mathcal{F}_n)$ 到 $U'$ 的拉回是某个凝聚
$\mathcal{O}_{U'}$-模的完备化。
\end{enumerate}
\end{lemma}

\begin{proof}
回忆 $A \to A'$ 是平坦环映射，并诱导同构
$A/I \to A'/I'$。参见代数，引理 \ref{algebra-lemma-completion-flat} 和
\ref{algebra-lemma-completion-complete}。
因此 $X' \to X$ 是平坦态射，并诱导同构 $Y' \to Y$。
从而 $U' \to U$ 是平坦态射，并诱导同构
$U' \cap Y' \to U \cap Y$。这说明在交换图
$$
\xymatrix{
\textit{Coh}(X', I\mathcal{O}_{X'}) \ar[r] &
\textit{Coh}(U', I\mathcal{O}_{U'}) \\
\textit{Coh}(X, I\mathcal{O}_X) \ar[u] \ar[r] &
\textit{Coh}(U, I\mathcal{O}_U) \ar[u]
}
$$
中，竖直函子是等价。参见概形上同调，引理
\ref{coherent-lemma-inverse-systems-pullback-equivalence}。
该引理由此以及引理 \ref{algebraization-lemma-essential-image-completion} 的结论形式地推出。
\end{proof}

\noindent
在情形 \ref{algebraization-situation-algebraize} 中，设 $(\mathcal{F}_n)$ 为
$\textit{Coh}(U, I\mathcal{O}_U)$ 的对象。为了判断
$(\mathcal{F}_n)$ 是否可延拓到 $X$，考察下列 $A$-模是自然的：
\begin{equation}
\label{algebraization-equation-guess}
M = \lim H^0(U, \mathcal{F}_n)
\end{equation}
注意 $M$ 有极限拓扑；它（先验地）比 $I$-进拓扑更粗，因为
$M \to H^0(U, \mathcal{F}_n)$ 零化 $I^nM$。有典范映射
$$
\widetilde{M}|_U \to \widetilde{M/I^nM}|_U \to
\widetilde{H^0(U, \mathcal{F}_n)}|_U \to
\mathcal{F}_n
$$
可以希望 $\widetilde{M}$ 在 $U$ 上限制为凝聚模，且
$(\mathcal{F}_n)$ 是该模的完备化。如果 $A$ 不完备，
这几乎不可能成立，因此这种想法过于天真。但即使 $A$ 是 $I$-进完备的，
这个概念也很难使用。一种不那么天真的做法是考虑下列要求。

\begin{definition}
\label{algebraization-definition-canonically-algebraizable}
在情形 \ref{algebraization-situation-algebraize} 中，设 $(\mathcal{F}_n)$ 为
$\textit{Coh}(U, I\mathcal{O}_U)$ 的对象。称
{\it $(\mathcal{F}_n)$ 可典范地延拓到 $X$}，如果逆系统
$$
\{\widetilde{H^0(U, \mathcal{F}_n)}\}_{n \geq 1}
$$
在 $\QCoh(\mathcal{O}_X)$ 中 pro-同构于某个对象
$(\mathcal{G}_n)$，而该对象属于 $\textit{Coh}(X, I\mathcal{O}_X)$。
\end{definition}

\noindent
引理 \ref{algebraization-lemma-canonically-algebraizable} 将说明，
定义 \ref{algebraization-definition-canonically-algebraizable} 中的条件强于
定义 \ref{algebraization-definition-algebraizable} 中的条件。

\begin{lemma}
\label{algebraization-lemma-canonically-algebraizable}
在情形 \ref{algebraization-situation-algebraize} 中，设 $(\mathcal{F}_n)$ 为
$\textit{Coh}(U, I\mathcal{O}_U)$ 的对象。若 $(\mathcal{F}_n)$
可典范地延拓到 $X$，则
\begin{enumerate}
\item $(\widetilde{H^0(U, \mathcal{F}_n)})$ pro-同构于
一个对象 $(\mathcal{G}_n)$，该对象属于 $\textit{Coh}(X, I \mathcal{O}_X)$，
该对象在唯一同构意义下唯一；
\item $(\mathcal{G}_n)$ 在 $U$ 上的限制同构于
$(\mathcal{F}_n)$；也就是说，$(\mathcal{F}_n)$ 可延拓到 $X$；
\item 逆系统 $\{H^0(U, \mathcal{F}_n)\}$ 满足 Mittag--Leffler 条件；
\item (\ref{algebraization-equation-guess}) 中的模 $M$ 在 $I$-进完备化上有限，
这里所完备化的环是 $A$，
且 $M$ 上的极限拓扑是 $I$-进拓扑。
\end{enumerate}
\end{lemma}

\begin{proof}
由定义 \ref{algebraization-definition-canonically-algebraizable}，(1) 中的
$(\mathcal{G}_n)$ 存在。由引理
\ref{algebraization-lemma-recognize-formal-coherent-modules}，(1) 中的
$(\mathcal{G}_n)$ 唯一。写成 $\mathcal{G}_n = \widetilde{M_n}$。
则 $\{M_n\}$ 是有限 $A$-模的逆系统，且
$M_n = M_{n + 1}/I^n M_{n + 1}$。
由定义 \ref{algebraization-definition-canonically-algebraizable}，逆系统
$\{H^0(U, \mathcal{F}_n)\}$ pro-同构于 $\{M_n\}$。
因此，逆系统 $\{H^0(U, \mathcal{F}_n)\}$ 满足 Mittag--Leffler 条件，且
$M = \lim M_n$（作为拓扑模）。所以，(4) 中 $M$ 的性质由代数，引理
\ref{algebra-lemma-limit-complete}、
\ref{algebra-lemma-finite-over-complete-ring}、
\ref{algebra-lemma-hathat-finitely-generated} 得到。
由于 $U$ 是拟仿射的，典范映射
$$
\widetilde{H^0(U, \mathcal{F}_n)}|_U \to \mathcal{F}_n
$$
是同构（性质，引理
\ref{properties-lemma-quasi-coherent-quasi-affine}）。
由此可知 $(\mathcal{G}_n|_U)$ 和 $(\mathcal{F}_n)$ pro-同构，
因而由引理 \ref{algebraization-lemma-recognize-formal-coherent-modules} 可知它们同构。
\end{proof}

\begin{lemma}
\label{algebraization-lemma-canonically-extend-base-change}
在情形 \ref{algebraization-situation-algebraize} 中，设 $(\mathcal{F}_n)$ 为
$\textit{Coh}(U, I\mathcal{O}_U)$ 的对象。设 $A \to A'$ 为平坦环映射。
令 $X' = \Spec(A')$，设 $U' \subset X'$ 为 $U$ 的逆像，
并记诱导态射为 $g : U' \to U$。令
$(\mathcal{F}'_n) = (g^*\mathcal{F}_n)$；参见概形上同调，引理
\ref{coherent-lemma-inverse-systems-pullback}。
若 $(\mathcal{F}_n)$ 可典范地延拓到 $X$，则
$(\mathcal{F}'_n)$ 可典范地延拓到 $X'$。
此外，引理 \ref{algebraization-lemma-canonically-algebraizable} 为
$(\mathcal{F}_n)$ 给出的延拓拉回为 $(\mathcal{F}'_n)$ 的延拓。
\end{lemma}

\begin{proof}
设 $f : X' \to X$ 为诱导态射。由平坦基变换，有
$H^0(U', \mathcal{F}'_n) = H^0(U, \mathcal{F}_n) \otimes_A A'$；
参见概形上同调，引理
\ref{coherent-lemma-flat-base-change-cohomology}。
因此，若 $(\mathcal{G}_n)$ 属于 $\textit{Coh}(X, I\mathcal{O}_X)$，并且
pro-同构于 $(\widetilde{H^0(U, \mathcal{F}_n)})$，则
$(f^*\mathcal{G}_n)$ pro-同构于
$$
(f^*\widetilde{H^0(U, \mathcal{F}_n)}) =
(\widetilde{H^0(U, \mathcal{F}_n) \otimes_A A'}) =
(\widetilde{H^0(U', \mathcal{F}'_n)})
$$
。证明完成。
\end{proof}

\begin{lemma}
\label{algebraization-lemma-when-done}
在情形 \ref{algebraization-situation-algebraize} 中，设 $(\mathcal{F}_n)$ 为
$\textit{Coh}(U, I\mathcal{O}_U)$ 的对象。设 $M$ 如 (\ref{algebraization-equation-guess})。
假设
\begin{enumerate}
\item[(a)] 逆系统 $H^0(U, \mathcal{F}_n)$ 具有 Mittag--Leffler 性质；
\item[(b)] $M$ 上的极限拓扑与 $I$-进拓扑一致；
\item[(c)] $M \to H^0(U, \mathcal{F}_n)$ 的像是有限 $A$-模，
对所有 $n$ 均如此。
\end{enumerate}
则 $(\mathcal{F}_n)$ 可典范地延拓到 $X$。
特别地，若 $A$ 是 $I$-进完备的，则
$(\mathcal{F}_n)$ 是某个凝聚 $\mathcal{O}_U$-模的完备化。
\end{lemma}

\begin{proof}
由于 $H^0(U, \mathcal{F}_n)$ 具有 Mittag--Leffler 性质，
且 $M$ 上的极限拓扑是 $I$-进拓扑，可知
$\{M/I^nM\}$ 与 $\{H^0(U, \mathcal{F}_n)\}$
是 pro-同构的 $A$-模逆系统。因此，若令
$$
\mathcal{G}_n = \widetilde{M/I^n M}
$$
则为了验证定义 \ref{algebraization-definition-canonically-algebraizable} 中的条件，
只需证明 $M$ 在 $I$-进完备化上为有限模，其中被完备化的环是 $A$。
这由上述结论、条件 (c) 给出的 $M/I^n M$ 的有限性，以及
代数，引理 \ref{algebra-lemma-finite-over-complete-ring} 得到。
\end{proof}

\noindent
下述结论在某种意义上是上面引理 \ref{algebraization-lemma-when-done}
最直接的可能应用。

\begin{lemma}
\label{algebraization-lemma-algebraization-principal-variant}
在情形 \ref{algebraization-situation-algebraize} 中，设
$(\mathcal{F}_n)$ 为 $\textit{Coh}(U, I\mathcal{O}_U)$ 的对象。假设
\begin{enumerate}
\item $I = (f)$ 是由非零因子 $f \in \mathfrak a$ 生成的主理想；
\item $\mathcal{F}_n$ 是有限局部自由
$\mathcal{O}_U/f^n\mathcal{O}_U$-模；
\item $H^1_\mathfrak a(A/fA)$ 和 $H^2_\mathfrak a(A/fA)$
是有限 $A$-模。
\end{enumerate}
则 $(\mathcal{F}_n)$ 可典范地延拓到 $X$。特别地，若 $A$
完备，则 $(\mathcal{F}_n)$ 是某个凝聚 $\mathcal{O}_U$-模的完备化。
\end{lemma}

\begin{proof}
我们将通过验证引理 \ref{algebraization-lemma-when-done} 的假设 (a)、(b) 和 (c)
来证明此结论。

\medskip\noindent
由于 $\mathcal{F}_n$ 在 $\mathcal{O}_U/f^n\mathcal{O}_U$ 上局部自由，
所以有短正合序列
$0 \to \mathcal{F}_n \to \mathcal{F}_{n + 1} \to \mathcal{F}_1 \to 0$
，对所有 $n$ 均如此。因此，由上同调，引理
\ref{cohomology-lemma-topology-I-adic-f}，条件 (b) 成立。

\medskip\noindent
由于 $f$ 是非零因子，得到短正合序列
$$
0 \to A/f^nA \xrightarrow{f} A/f^{n + 1}A \to A/fA \to 0
$$
以及相应的短正合序列
$0 \to \mathcal{F}_n \to \mathcal{F}_{n + 1} \to \mathcal{F}_1 \to 0$.
下文将不再另行说明地使用局部上同调，引理
\ref{local-cohomology-lemma-finiteness-pushforwards-and-H1-local}
。我们的假设蕴含
$H^0(U, \mathcal{O}_U/f\mathcal{O}_U)$ 和
$H^1(U, \mathcal{O}_U/f\mathcal{O}_U)$
都是有限 $A$-模。因此，对 $\mathcal{F}_1$ 也有相同结论；参见局部上同调，引理
\ref{local-cohomology-lemma-finiteness-for-finite-locally-free}。
利用归纳法和这些短正合序列，可知
$H^0(U, \mathcal{F}_n)$ 是有限 $A$-模，对所有 $n$ 均如此。
由此可见假设 (c) 成立。

\medskip\noindent
最后，由于 $H^1(U, \mathcal{F}_1)$ 是有限 $A$-模，
可以应用上同调，引理 \ref{cohomology-lemma-ML}，
从而看出假设 (a) 成立。
\end{proof}

\begin{remark}
\label{algebraization-remark-interesting-case-variant}
在引理 \ref{algebraization-lemma-algebraization-principal-variant} 中，
若 $A$ 是具有 Cohen--Macaulay 形式纤维的泛链状环
（例如，若 $A$ 有对偶化复形），则条件
$H^1_\mathfrak a(A/fA)$ 和 $H^2_\mathfrak a(A/fA)$
是有限 $A$-模，等价于
$$
\text{depth}((A/f)_\mathfrak p) + \dim((A/\mathfrak p)_\mathfrak q) > 2
$$
对所有 $\mathfrak p \in V(f) \setminus V(\mathfrak a)$
和 $\mathfrak q \in V(\mathfrak p) \cap V(\mathfrak a)$ 成立；
这是局部上同调，定理 \ref{local-cohomology-theorem-finiteness} 的结论。

\medskip\noindent
例如，若 $A/fA$ 满足 $(S_2)$，且 $Z = V(\mathfrak a)$ 的每个不可约分支
的余维均 $\geq 3$（余维在 $Y = \Spec(A/fA)$ 中取），则可得
$H^1_\mathfrak a(A/fA)$ 和 $H^2_\mathfrak a(A/fA)$ 的有限性。
应将此与引理 \ref{algebraization-lemma-algebraization-principal} 中稍弱的条件作比较
（另见注记 \ref{algebraization-remark-interesting-case}）。
\end{remark}






\section{凝聚形式模的代数化，II}
\label{algebraization-section-algebraization-modules-general}

\noindent
我们继续第 \ref{algebraization-section-algebraization-modules} 节中开始的讨论。
初次阅读时可以跳过本节。

\begin{lemma}
\label{algebraization-lemma-map-kernel-cokernel-on-closed}
在情形 \ref{algebraization-situation-algebraize} 中，设
$(\mathcal{F}_n) \to (\mathcal{F}'_n)$ 是
$\textit{Coh}(U, I\mathcal{O}_U)$ 中的一个态射，
其核与余核均被 $I$ 的某个幂零化。则
\begin{enumerate}
\item $(\mathcal{F}_n)$ 可延拓到 $X$，当且仅当
$(\mathcal{F}'_n)$ 可延拓到 $X$；
\item $(\mathcal{F}_n)$ 是某个凝聚 $\mathcal{O}_U$-模的完备化，
当且仅当 $(\mathcal{F}'_n)$ 也是如此。
\end{enumerate}
\end{lemma}

\begin{proof}
第 (2) 项直接由概形上同调，引理
\ref{coherent-lemma-existence-easy} 推出。
为证明第 (1) 项，先用引理 \ref{algebraization-lemma-algebraizable}
约化到 $A$ 为 $I$-进完备的情形。
但在这种情形下，引理 \ref{algebraization-lemma-essential-image-completion}
把第 (1) 项约化为第 (2) 项。
\end{proof}

\noindent
下面两个引理原本用于引理 \ref{algebraization-lemma-when-done} 的证明。
我们将其保留在这里，供有兴趣了解可得到哪些中间结果的读者参考。

\begin{lemma}
\label{algebraization-lemma-when-ML}
在情形 \ref{algebraization-situation-algebraize} 中，设 $(\mathcal{F}_n)$ 是
$\textit{Coh}(U, I\mathcal{O}_U)$ 的一个对象。若逆系统
$H^0(U, \mathcal{F}_n)$ 满足 Mittag--Leffler 条件，则典范映射
$$
\widetilde{M/I^nM}|_U \to \mathcal{F}_n
$$
对所有 $n$ 均为满射，其中 $M$ 如 (\ref{algebraization-equation-guess}) 所示。
\end{lemma}

\begin{proof}
满性可在某点 $y \in Y \setminus Z$ 的茎上检验。
若 $y$ 对应于素理想 $\mathfrak q \subset A$，则可以选取
$f \in \mathfrak a$，使得 $f \not \in \mathfrak q$。于是只需证明
$$
M_f \longrightarrow H^0(U, \mathcal{F}_n)_f = H^0(D(f), \mathcal{F}_n)
$$
为满射，因为 $D(f)$ 是仿射的（等号由性质，引理
\ref{properties-lemma-invert-f-sections} 成立）。由于满足
Mittag--Leffler 条件，可知
$$
\Im(M \to H^0(U, \mathcal{F}_n)) =
\Im(H^0(U, \mathcal{F}_m) \to H^0(U, \mathcal{F}_n))
$$
对某个 $m \geq n$ 成立。利用上同调长正合序列可见
$$
\Coker(H^0(U, \mathcal{F}_m) \to H^0(U, \mathcal{F}_n))
\subset
H^1(U, \Ker(\mathcal{F}_m \to \mathcal{F}_n))
$$
由于 $U = X \setminus V(\mathfrak a)$，此 $H^1$ 是
$\mathfrak a$-幂挠的。因此将 $f$ 可逆化后，余核变为零。
\end{proof}

\begin{lemma}
\label{algebraization-lemma-when-topology}
在情形 \ref{algebraization-situation-algebraize} 中，设 $(\mathcal{F}_n)$ 是
$\textit{Coh}(U, I\mathcal{O}_U)$ 的一个对象，并设 $M$
如 (\ref{algebraization-equation-guess}) 所示。令
$$
\mathcal{G}_n = \widetilde{M/I^nM}.
$$
若 $M$ 上的极限拓扑与 $I$-进拓扑一致，则
$\mathcal{G}_n|_U$ 是凝聚
$\mathcal{O}_U$-模，并且逆系统的映射
$$
(\mathcal{G}_n|_U) \longrightarrow (\mathcal{F}_n)
$$
在阿贝尔范畴 $\textit{Coh}(U, I\mathcal{O}_U)$ 中是单射。
\end{lemma}

\begin{proof}
注意，$\mathcal{G}_n$ 是拟凝聚 $\mathcal{O}_X$-模，
并且被 $I^n$ 零化，且
$\mathcal{G}_{n + 1}/I^n\mathcal{G}_{n + 1} = \mathcal{G}_n$.
考虑
$$
M_n = \Im(M \longrightarrow H^0(U, \mathcal{F}_n))
$$
该假设说明逆系统 $(M_n)$ 与
$(M/I^nM)$ 作为 $\text{Mod}_A$ 的 pro-对象同构。
选取 $f \in \mathfrak a$，使得 $D(f) \subset U$ 是仿射开集。于是有
$$
(M_n)_f \subset H^0(U, \mathcal{F}_n)_f = H^0(D(f), \mathcal{F}_n)
$$
由性质，引理 \ref{properties-lemma-invert-f-sections}，这里等号成立。
因此 $\widetilde{M_n}|_U \to \mathcal{F}_n$ 是单射。
由此推出 $\widetilde{M_n}|_U$ 是凝聚模
（概形上同调，引理
\ref{coherent-lemma-coherent-Noetherian-quasi-coherent-sub-quotient}）。
由于 $M \to M/I^nM$ 是满射，并且它分解为
$M_{n'} \to M/I^nM$，其中某个 $n' \geq n$，所以 $\mathcal{G}_n|_U$
作为凝聚模的商是凝聚的。
结合证明开头的观察，可知
$(\mathcal{G}_n|_U)$ 确实构成
$\textit{Coh}(U, I\mathcal{O}_U)$ 的一个对象。
最后，为证明该映射的单性，只需证明
$$
\lim (M/I^nM)_f = \lim H^0(D(f), \mathcal{G}_n) \to
\lim H^0(D(f), \mathcal{F}_n)
$$
为单射；参见概形上同调，引理
\ref{coherent-lemma-inverse-systems-abelian} 和
\ref{coherent-lemma-inverse-systems-affine}。
$\lim (M_n)_f \to \lim H^0(D(f), \mathcal{F}_n)$ 的单性是清楚的
（见上文），而由我们关于 pro-系统的说明，有
$\lim (M_n)_f = \lim (M/I^nM)_f$。证明完成。
\end{proof}





\section{距离函数}
\label{algebraization-section-distance}

\noindent
设 $Y$ 是诺特概形，$Z \subset Y$ 是闭子集。我们定义函数
\begin{equation}
\label{algebraization-equation-delta-Z}
\delta^Y_Z = \delta_Z : Y \longrightarrow \mathbf{Z}_{\geq 0} \cup \{\infty\}
\end{equation}
来度量 $Y$ 的一点到 $Z$ 的“距离”。
非正式讨论见注记 \ref{algebraization-remark-discussion}。
设 $y \in Y$。规定 $\delta_Z(y) = \infty$，如果 $y$ 位于
$Y$ 的一个与 $Z$ 不相交的连通分支中。
若 $y$ 位于 $Y$ 的一个与 $Z$ 相交的连通分支中，
则可找到 $k \geq 0$ 以及系统
$$
V_0 \subset W_0 \supset V_1 \subset W_1 \supset \ldots \supset
V_k \subset W_k
$$
，其中各项都是 $Y$ 的整的闭子概形，$V_0 \subset Z$，
且 $y \in W_k$ 是泛点。令
$c_i = \text{codim}(V_i, W_i)$，其中 $i = 0, \ldots, k$；并令
$b_i = \text{codim}(V_{i + 1}, W_i)$，其中 $i = 0, \ldots, k - 1$。
对这样的系统，令
$$
\delta(V_0, W_0, V_1, \ldots, W_k) =
k +
\max_{i = 0, 1, \ldots, k}
(c_i + c_{i + 1} + \ldots + c_k - b_i - b_{i + 1} - \ldots - b_{k - 1})
$$
这个数 $\geq k$，因为可以取 $i = k$，且 $c_k \geq 0$。
最后令
$$
\delta_Z(y) = \min \delta(V_0, W_0, V_1, \ldots, W_k)
$$
，其中最小值遍历上述所有由 $Y$ 的整的闭子概形组成的系统。

\begin{lemma}
\label{algebraization-lemma-discussion}
设 $Y$ 是诺特概形，$Z \subset Y$ 是闭子集。
\begin{enumerate}
\item 对 $y \in Y$，有 $\delta_Z(y) = 0 \Leftrightarrow y \in Z$。
\item 子集 $\{y \in Y \mid \delta_Z(y) \leq k\}$
对特化稳定。
\item 对 $y \in Y$ 和 $z \in \overline{\{y\}} \cap Z$，有
$\dim(\mathcal{O}_{\overline{\{y\}}, z}) \geq \delta_Z(y)$。
\item 若 $\delta$ 是 $Y$ 上的维数函数，则
$\delta_Z(y) \geq \delta(y) - \delta_{max}$，其中 $\delta_{max}$
是 $\delta$ 在 $Z$ 上的最大值。
\item 若 $Y = \Spec(A)$ 是链状诺特局部环的谱，
其极大理想为 $\mathfrak m$，且 $Z = \{\mathfrak m\}$，则
$\delta_Z(y) = \dim(\overline{\{y\}})$。
\item 若 $Y' \subset Y$ 是开子概形，则
$\delta^{Y'}_{Y' \cap Z}(y') \geq \delta^Y_Z(y')$ 对 $y' \in Y'$ 成立。
\end{enumerate}
假设 $Y$ 是链状的。则
\begin{enumerate}
\setcounter{enumi}{6}
\item 设 $y' \leadsto y$ 是 $Y$ 中点的直接特化。
若 $Y$ 是链状的，则 $\delta_Z(y') \leq \delta_Z(y) + 1$。
\item 给定如下特化图式
$$
\xymatrix{
& y'_0 \ar@{~>}[ld] \ar@{~>}[rd] &
& y'_1 \ar@{~>}[ld] & \ldots
& y'_{k - 1} \ar@{~>}[rd] &
\\
y_0 & &
y_1 & &
\ldots & &
y_k = y
}
$$
，其中各点都在 $Y$ 中，$y_0 \in Z$，并且
$y_i' \leadsto y_i$ 是直接特化，则 $\delta_Z(y_k) \leq k$。
\end{enumerate}
\end{lemma}

\begin{proof}
第 (1) 项的证明。若 $y \in Z$，则可取 $k = 0$ 以及
$V_0 = W_0 = \overline{\{y\}}$，从而得到 $\delta(V_0, W_0) = 0$，故
$\delta_Z(y) = 0$。若 $y \not \in Z$，则对每个系统
$V_0 \subset W_0 \supset V_1 \subset W_1 \supset \ldots \subset W_k$
，只要它是 $y$ 的系统，就或者有 $k = 0$ 且
$V_0 \not = W_0$，或者有 $k > 0$。
在两种情形下都有 $\delta(V_0, W_0, \ldots, W_k) > 0$。因此
$\delta_Z(y) > 0$。

\medskip\noindent
第 (2) 项的证明。设 $y \leadsto y'$ 是非平凡特化，并设
$V_0 \subset W_0 \supset V_1 \subset W_1 \supset \ldots \subset W_k$
是 $y$ 的一个系统。分两种情形。
情形 I：$V_k = W_k$，即 $c_k = 0$。在此情形下可令
$V'_k = W'_k = \overline{\{y'\}}$。
直接计算表明
$\delta(V_0, W_0, \ldots, V'_k, W'_k) \leq
\delta(V_0, W_0, \ldots, V_k, W_k)$
，因为只有 $b_{k - 1}$ 被换成了一个更大的整数。
情形 II：$V_k \not = W_k$，即 $c_k > 0$。
在此情形下，令 $V_{k + 1} = W_{k + 1} = \overline{\{y'\}}$，
可见 $V_0 \subset W_0 \supset \ldots \subset W_k \supset V_{k + 1}
\subset W_{k + 1}$ 是 $y'$ 的一个系统。此时 $c_{k + 1} = 0$，且
$b_k > 0$，所以得到
\begin{align*}
& \delta(V_0, \ldots, W_{k + 1}) \\
& =
k + 1 +
\max_{i = 0, 1, \ldots, k + 1}
(c_i + c_{i + 1} + \ldots + c_k + c_{k + 1}
- b_i - b_{i + 1} - \ldots - b_{k - 1} - b_k) \\
& =
k + 1 +
\max_{i = 0, 1, \ldots, k + 1}
(c_i + c_{i + 1} + \ldots + c_k
- b_i - b_{i + 1} - \ldots - b_{k - 1} - b_k) \\
& \leq
k + \max_{i = 0, 1, \ldots, k}
(c_i + c_{i + 1} + \ldots + c_k - b_i - b_{i + 1} - \ldots - b_{k - 1}) \\
& = \delta(V_0, \ldots, W_k)
\end{align*}
该不等式成立，是因为 $c_k > 0$ 且 $b_k > 0$；这尤其蕴含
$\delta(V_0, \ldots, W_k) \geq k + c_k \geq k + 1$。

\medskip\noindent
第 (3) 项的证明。给定 $y \in Y$ 和 $z \in \overline{\{y\}} \cap Z$，
得到系统
$$
V_0 = \overline{\{z\}} \subset W_0 = \overline{\{y\}}
$$
，并且由性质，引理 \ref{properties-lemma-codimension-local-ring}，有
$c_0 = \text{codim}(V_0, W_0) = \dim(\mathcal{O}_{\overline{\{y\}}, z})$。
因此 $\delta(V_0, W_0) = 0 + c_0 = c_0$，这就证明了所需结论。

\medskip\noindent
第 (4) 项的证明。设 $\delta$ 是 $Y$ 上的维数函数。
设 $V_0 \subset W_0 \supset V_1 \subset W_1 \supset \ldots \subset W_k$
是 $y$ 的一个系统。令 $y'_i \in W_i$ 和 $y_i \in V_i$ 为相应泛点，
于是 $y_0 \in Z$ 且 $y_k = y$。这时
$$
\delta(y_i) - \delta(y_{i - 1}) =
\delta(y'_{i - 1}) - \delta(y_{i - 1}) - \delta(y'_{i - 1}) + \delta(y_i) =
c_{i - 1} - b_{i - 1}
$$
最后有 $\delta(y'_k) - \delta(y_{k - 1}) = c_k$。因此
$$
\delta(y) - \delta(y_0) =
c_0 + \ldots + c_k - b_0 - \ldots - b_{k - 1}
$$
由此得出
$\delta(V_0, W_0, \ldots, W_k) \geq k + \delta(y) - \delta(y_0)$
，这就证明了所需结论。

\medskip\noindent
第 (5) 项的证明。函数 $\delta(y) = \dim(\overline{\{y\}})$
是维数函数。因此由第 (4) 项有 $\delta(y) \leq \delta_Z(y)$。
由第 (3) 项有 $\delta_Z(y) \leq \delta(y)$，证明完成。

\medskip\noindent
第 (6) 项的证明。这是清楚的，因为计算
$\delta^{Y'}_{Y' \cap Z}$ 时需要考虑的系统更少。

\medskip\noindent
第 (7) 项的证明。设
$V_0 \subset W_0 \supset V_1 \subset W_1 \supset \ldots \subset W_k$
是 $y$ 的一个系统。令 $W'_k = \overline{\{y'\}}$。
由于 $Y$ 是链状的，可知
$\text{codim}(V_k, W'_k) = \text{codim}(V_k, W_k) + 1$.
容易推出 $\delta(V_0, \ldots, V_k, W'_k) \leq
\delta(V_0, \ldots, V_k, W_k) + 1$，这就证明了所需结论。

\medskip\noindent
第 (8) 项的证明：结合第 (7) 项与第 (2) 项即可。
\end{proof}

\begin{lemma}
\label{algebraization-lemma-change-distance-function}
设 $Y$ 是泛链状诺特概形，$Z \subset Y$ 是闭子概形。
设 $f : Y' \to Y$ 是有限型态射，其所有纤维的维数均
$\leq e$。令 $Z' = f^{-1}(Z)$。则
$$
\delta_Z(y) \leq \delta_{Z'}(y') + e - \text{trdeg}_{\kappa(y)}(\kappa(y'))
$$
对 $y' \in Y'$ 成立，其中其像为 $y \in Y$。
\end{lemma}

\begin{proof}
若 $\delta_{Z'}(y') = \infty$，则无须证明。
若 $\delta_{Z'}(y') < \infty$，则选取整的闭子概形系统
$$
V'_0 \subset W'_0 \supset V'_1 \subset W'_1 \supset \ldots \subset W'_k
$$
，其中各项在 $Y'$ 中，$V'_0 \subset Z'$，且 $y'$ 是 $W'_k$ 的泛点，
并使得 $\delta_{Z'}(y') = \delta(V'_0, W'_0, \ldots, W'_k)$。
记
$$
V_0 \subset W_0 \supset V_1 \subset W_1 \supset \ldots \subset W_k
$$
为上述概形在 $Y$ 中的概形论像。注意，$y$ 是 $W_k$ 的泛点，
并且 $V_0 \subset Z$。对每个 $i$，考察图表
$$
\xymatrix{
V'_i \ar[r] \ar[d] & W'_i \ar[d] & V'_{i + 1} \ar[l] \ar[d] \\
V_i \ar[r] & W_i & V_{i + 1} \ar[l]
}
$$
记 $n_i$ 为 $V'_i/V_i$ 的相对维数，记
$m_i$ 为 $W'_i/W_i$ 的相对维数；更确切地说，
它们是相应函数域扩张的超越次数。令
$c_i = \text{codim}(V_i, W_i)$,
$c'_i = \text{codim}(V'_i, W'_i)$,
$b_i = \text{codim}(V_{i + 1}, W_i)$，以及
$b'_i = \text{codim}(V'_{i + 1}, W'_i)$。
由维数公式有
$$
c_i = c'_i + n_i - m_i
\quad\text{且}\quad
b_i = b'_i + n_{i + 1} - m_i
$$
参见态射，引理 \ref{morphisms-lemma-dimension-formula}。
因此 $c_i - b_i = c'_i - b'_i + n_i - n_{i + 1}$。于是
\begin{align*}
& c_i + c_{i + 1} + \ldots + c_k - b_i - b_{i + 1} - \ldots - b_{k - 1} \\
& =
c'_i + c'_{i + 1} + \ldots + c'_k - b'_i - b'_{i + 1} - \ldots - b'_{k - 1}
+ n_i - n_k + c_k - c'_k \\
& =
c'_i + c'_{i + 1} + \ldots + c'_k - b'_i - b'_{i + 1} - \ldots - b'_{k - 1}
+ n_i - m_k
\end{align*}
由此可见
\begin{align*}
\max_{i = 0, \ldots, k}
& (c_i + c_{i + 1} + \ldots + c_k - b_i - b_{i + 1} - \ldots - b_{k - 1}) \\
& =
\max_{i = 0, \ldots, k}
(c'_i + c'_{i + 1} + \ldots + c'_k - b'_i - b'_{i + 1} - \ldots - b'_{k - 1}
+ n_i - m_k) \\
& =
\max_{i = 0, \ldots, k}
(c'_i + c'_{i + 1} + \ldots + c'_k - b'_i - b'_{i + 1} - \ldots - b'_{k - 1}
+ n_i) - m_k \\
& \leq
\max_{i = 0, \ldots, k}
(c'_i + c'_{i + 1} + \ldots + c'_k - b'_i - b'_{i + 1} - \ldots - b'_{k - 1})
+ e - m_k
\end{align*}
由于 $m_k = \text{trdeg}_{\kappa(y)}(\kappa(y'))$，可得
$$
\delta(V_0, W_0, \ldots, W_k) \leq
\delta(V'_0, W'_0, \ldots, W'_k) + e - \text{trdeg}_{\kappa(y)}(\kappa(y'))
$$
，即为所需结论。
\end{proof}

\begin{remark}
\label{algebraization-remark-discussion}
设 $Y$ 是链状诺特概形，$Z \subset Y$ 是闭子集。
由引理 \ref{algebraization-lemma-discussion} 有
$$
\delta_Z(y) \leq \min
\left\{ k \middle|
\begin{matrix}
\text{ there exist specializations in }Y \\
y_0 \leftarrow y'_0 \rightarrow y_1 \leftarrow y'_1 \rightarrow \ldots
\leftarrow y'_{k - 1} \rightarrow y_k = y \\
\text{ 其中 }y_0 \in Z\text{ 且 }y_i' \leadsto y_i
\text{ immediate}
\end{matrix}
\right\}
$$
我们断言，若 $Y$ 在某个域上是有限型的，则等号成立。
如果以后需要这一结果，我们会在这里给出精确表述并加以证明。
然而，一般而言，若用该不等式的右端定义 $\delta_Z$，
则我们不知道引理 \ref{algebraization-lemma-change-distance-function} 是否仍然成立。
\end{remark}

\begin{example}
\label{algebraization-example-distance}
设 $k$ 是域，且 $Y = \mathbf{A}^n_k$。记
$\delta : Y \to \mathbf{Z}_{\geq 0}$ 为通常的维数函数。
\begin{enumerate}
\item 若 $Z = \{z\}$，其中 $z$ 是某个闭点，则
\begin{enumerate}
\item $\delta_Z(y) = \delta(y)$ 在 $y \leadsto z$ 时成立；
\item $\delta_Z(y) = \delta(y) + 1$ 在 $y \not \leadsto z$ 时成立。
\end{enumerate}
\item 若 $Z$ 是闭子簇，且 $W = \overline{\{y\}}$，则
\begin{enumerate}
\item $\delta_Z(y) = 0$ 在 $W \subset Z$ 时成立；
\item $\delta_Z(y) = \dim(W) - \dim(Z)$ 在 $Z$ 包含于 $W$ 时成立；
\item $\delta_Z(y) = 1$ 在 $\dim(W) \leq \dim(Z)$ 且
$W \not \subset Z$ 时成立；
\item $\delta_Z(y) = \dim(W) - \dim(Z) + 1$ 在
$\dim(W) > \dim(Z)$ 且 $Z \not \subset W$ 时成立。
\end{enumerate}
\end{enumerate}
情形 (1) 的推广如下：设 $Y$ 在某个域上是有限型的，
且 $Z = \{z\}$ 是闭点。则 $\delta_Z(y) = \delta(y) + t$，
其中 $t$ 是连接 $z$ 与 $\overline{\{y\}}$ 的某个闭点的曲线链的最小长度。
\end{example}







\section{凝聚形式模的代数化，III}
\label{algebraization-section-uniqueness}

\noindent
我们继续第 \ref{algebraization-section-algebraization-modules} 节和
第 \ref{algebraization-section-algebraization-modules-general} 节中开始的讨论。
我们将使用第 \ref{algebraization-section-distance} 节的距离函数，
表述情形 \ref{algebraization-situation-algebraize} 中凝聚形式模所满足的一些自然条件。

\medskip\noindent
在情形 \ref{algebraization-situation-algebraize} 中，给定一点 $y \in U \cap Y$，
可以考虑 $I$-进完备化
$$
\mathcal{O}_{X, y}^\wedge = \lim \mathcal{O}_{X, y}/I^n\mathcal{O}_{X, y}
$$
。这是关于 $I\mathcal{O}_{X, y}^\wedge$ 完备的诺特局部环，
其极大理想为 $\mathfrak m_y^\wedge$；参见代数，
第 \ref{algebra-section-completion-noetherian} 节。
设 $(\mathcal{F}_n)$ 是 $\textit{Coh}(U, I\mathcal{O}_U)$ 的一个对象。
将 $(\mathcal{F}_n)$ 在 $y$ 处的“茎”定义为
$$
\mathcal{F}_y^\wedge = \lim \mathcal{F}_{n, y}
$$
。这是 $\mathcal{O}_{X, y}^\wedge$ 上的有限模。参见代数，引理
\ref{algebra-lemma-limit-complete} 和
\ref{algebra-lemma-finite-over-complete-ring}。

\begin{definition}
\label{algebraization-definition-s-d-inequalities}
在情形 \ref{algebraization-situation-algebraize} 中，设 $(\mathcal{F}_n)$ 是
$\textit{Coh}(U, I\mathcal{O}_U)$ 的一个对象。设 $a, b$ 是整数，
并设 $\delta^Y_Z$ 如 (\ref{algebraization-equation-delta-Z}) 所示。
称 {\it $(\mathcal{F}_n)$ 满足 $(a, b)$-不等式组}，如果对
$y \in U \cap Y$ 以及素理想
$\mathfrak p \subset \mathcal{O}_{X, y}^\wedge$，只要
$\mathfrak p \not \in V(I\mathcal{O}_{X, y}^\wedge)$，就有
\begin{enumerate}
\item 若 $V(\mathfrak p) \cap V(I\mathcal{O}_{X, y}^\wedge) \not =
\{\mathfrak m_y^\wedge\}$，则
$$
\text{depth}((\mathcal{F}^\wedge_y)_\mathfrak p) + \delta^Y_Z(y) \geq a
\quad\text{或}\quad
\text{depth}((\mathcal{F}^\wedge_y)_\mathfrak p) +
\dim(\mathcal{O}_{X, y}^\wedge/\mathfrak p) + \delta^Y_Z(y) > b
$$
\item 若 $V(\mathfrak p) \cap V(I\mathcal{O}_{X, y}^\wedge) =
\{\mathfrak m_y^\wedge\}$，则
$$
\text{depth}((\mathcal{F}^\wedge_y)_\mathfrak p) + \delta^Y_Z(y) > a
$$
\end{enumerate}
称 {\it $(\mathcal{F}_n)$ 满足严格 $(a, b)$-不等式组}，
如果对 $y \in U \cap Y$ 以及素理想
$\mathfrak p \subset \mathcal{O}_{X, y}^\wedge$，只要
$\mathfrak p \not \in V(I\mathcal{O}_{X, y}^\wedge)$，就有
$$
\text{depth}((\mathcal{F}^\wedge_y)_\mathfrak p) + \delta^Y_Z(y) > a
\quad\text{或}\quad
\text{depth}((\mathcal{F}^\wedge_y)_\mathfrak p) +
\dim(\mathcal{O}_{X, y}^\wedge/\mathfrak p) + \delta^Y_Z(y) > b
$$
\end{definition}

\noindent
下面是一些初等观察。

\begin{lemma}
\label{algebraization-lemma-elementary}
在情形 \ref{algebraization-situation-algebraize} 中，设 $(\mathcal{F}_n)$ 是
$\textit{Coh}(U, I\mathcal{O}_U)$ 的一个对象，并设 $a, b$ 是整数。
\begin{enumerate}
\item 若 $(\mathcal{F}_n)$ 被 $I$ 的某个幂零化，则
$(\mathcal{F}_n)$ 满足 $(a, b)$-不等式组，对任意 $a, b$ 都如此。
\item 若 $(\mathcal{F}_n)$ 满足 $(a + 1, b)$-不等式组，则
$(\mathcal{F}_n)$ 满足严格 $(a, b)$-不等式组。
\end{enumerate}
若 $\text{cd}(A, I) \leq d$，且 $A$ 有对偶化复形，则
\begin{enumerate}
\item[(3)] $(\mathcal{F}_n)$ 满足 $(s, s + d)$-不等式组，
当且仅当对所有 $y \in U \cap Y$，元组
$\mathcal{O}_{X, y}^\wedge, I\mathcal{O}_{X, y}^\wedge,
\{\mathfrak m_y^\wedge\}, \mathcal{F}_y^\wedge, s - \delta^Y_Z(y), d$
如情形 \ref{algebraization-situation-bootstrap} 所示。
\item[(4)]
若 $(\mathcal{F}_n)$ 满足严格 $(s, s + d)$-不等式组，则
$(\mathcal{F}_n)$ 满足 $(s, s + d)$-不等式组。
\end{enumerate}
\end{lemma}

\begin{proof}
除第 (4) 项外均为立即的；第 (4) 项由引理
\ref{algebraization-lemma-bootstrap-bis-bis} 以及第 (3) 项中的对应关系推出。
\end{proof}

\begin{lemma}
\label{algebraization-lemma-explain-2-3-cd-1}
在情形 \ref{algebraization-situation-algebraize} 中，设 $(\mathcal{F}_n)$ 是
$\textit{Coh}(U, I\mathcal{O}_U)$ 的一个对象。若
$\text{cd}(A, I) = 1$，则 $\mathcal{F}$ 满足 $(2, 3)$-不等式组，
当且仅当
$$
\text{depth}((\mathcal{F}^\wedge_y)_\mathfrak p) +
\dim(\mathcal{O}_{X, y}^\wedge/\mathfrak p) + \delta^Y_Z(y) > 3
$$
对所有 $y \in U \cap Y$ 和
$\mathfrak p \subset \mathcal{O}_{X, y}^\wedge$，只要
$\mathfrak p \not \in V(I\mathcal{O}_{X, y}^\wedge)$ 就成立。
\end{lemma}

\begin{proof}
注意，对素理想 $\mathfrak p \subset \mathcal{O}_{X, y}^\wedge$，
若 $\mathfrak p \not \in V(I\mathcal{O}_{X, y}^\wedge)$，则有
$V(\mathfrak p) \cap V(I\mathcal{O}_{X, y}^\wedge) =
\{\mathfrak m_y^\wedge\}
\Leftrightarrow \dim(\mathcal{O}_{X, y}^\wedge/\mathfrak p) = 1$
，因为 $\text{cd}(A, I) = 1$。参见局部上同调，引理
\ref{local-cohomology-lemma-cd-change-rings} 和
\ref{local-cohomology-lemma-cd-bound-dim-local}。
现在考虑三个数
$\alpha = \text{depth}((\mathcal{F}^\wedge_y)_\mathfrak p) \geq 0$、
$\beta = \dim(\mathcal{O}_{X, y}^\wedge/\mathfrak p) \geq 1$ 和
$\gamma = \delta^Y_Z(y) \geq 1$。
由定义 \ref{algebraization-definition-s-d-inequalities}，所需条件为
\begin{enumerate}
\item 若 $\beta > 1$，则
$\alpha + \gamma \geq 2$ 或 $\alpha + \beta + \gamma > 3$；
\item 若 $\beta = 1$，则 $\alpha + \gamma > 2$。
\end{enumerate}
显然，这等价于 $\alpha + \beta + \gamma > 3$。
\end{proof}

\noindent
在本节余下部分中（建议读者初次阅读时跳过），我们将证明：
当 $A$ 是 $I$-进完备的时，由 $(\mathcal{F}_n)$ 构成如下范畴：
它们属于 $\textit{Coh}(U, I\mathcal{O}_U)$、可延拓到 $X$，并满足严格
$(1, 1 + \text{cd}(A, I))$-不等式组。该范畴等价于凝聚
$\mathcal{O}_U$-模范畴的一个满子范畴。

\begin{lemma}
\label{algebraization-lemma-sanity}
在情形 \ref{algebraization-situation-algebraize} 中，设 $\mathcal{F}$ 是凝聚
$\mathcal{O}_U$-模，且 $d \geq 1$。假设
\begin{enumerate}
\item $A$ 是 $I$-进完备的，有对偶化复形，并且
$\text{cd}(A, I) \leq d$；
\item 完备化 $\mathcal{F}^\wedge$（它来自 $\mathcal{F}$）
满足严格 $(1, 1 + d)$-不等式组。
\end{enumerate}
设 $x \in X$ 是一点，并令 $W = \overline{\{x\}}$。
若 $W \cap Y$ 有一个包含于 $Z$ 的不可约分支，
并且还有一个不包含于其中的不可约分支，则
$\text{depth}(\mathcal{F}_x) \geq 1$。
\end{lemma}

\begin{proof}
设 $W \cap Y = W_1 \cup \ldots \cup W_n$ 是不可约分支分解。
由假设，重新编号后可找到 $0 < m < n$，使得
$W_1, \ldots, W_m  \subset Z$ 且
$W_{m + 1}, \ldots, W_n \not \subset Z$。于是
$$
W \cap Y \setminus
\left((W_1 \cup \ldots \cup W_m) \cap (W_{m + 1} \cup \ldots \cup W_n)\right)
$$
不连通。由引理 \ref{algebraization-lemma-connected}，可找到
$1 \leq i \leq m < j \leq n$ 以及
$z \in W_i \cap W_j$，使得 $\dim(\mathcal{O}_{W, z}) \leq d + 1$。
选取直接特化 $y \leadsto z$，其中
$y \in W_j$ 且 $y \not \in Z$；$y$ 的存在性由性质，引理
\ref{properties-lemma-complement-closed-point-Jacobson} 得到。
注意，$\delta^Y_Z(y) = 1$ 且 $\dim(\mathcal{O}_{W, y}) \leq d$。
设 $\mathfrak p \subset \mathcal{O}_{X, y}$ 是对应于 $x$ 的素理想。
设 $\mathfrak p' \subset \mathcal{O}_{X, y}^\wedge$ 是
$\mathfrak p\mathcal{O}_{X, y}^\wedge$ 上方的一个极小素理想。则
$$
\text{depth}(\mathcal{F}_x) =
\text{depth}((\mathcal{F}^\wedge_y)_{\mathfrak p'})
\quad\text{且}\quad
\dim(\mathcal{O}_{W, y}) = \dim(\mathcal{O}_{X, y}^\wedge/\mathfrak p')
$$
参见代数，引理 \ref{algebra-lemma-apply-grothendieck-module}，以及
局部上同调，引理 \ref{local-cohomology-lemma-change-completion}。
现在直接从所给不等式组读出结论。
\end{proof}

\begin{lemma}
\label{algebraization-lemma-recover}
在情形 \ref{algebraization-situation-algebraize} 中，设 $\mathcal{F}$ 是凝聚
$\mathcal{O}_U$-模，且 $d \geq 1$。假设
\begin{enumerate}
\item $A$ 是 $I$-进完备的，有对偶化复形，并且
$\text{cd}(A, I) \leq d$；
\item 完备化 $\mathcal{F}^\wedge$（它来自 $\mathcal{F}$）
满足严格 $(1, 1+ d)$-不等式组；
\item 对 $x \in U$，若 $\overline{\{x\}} \cap Y \subset Z$，
则有 $\text{depth}(\mathcal{F}_x) \geq 2$。
\end{enumerate}
则 $H^0(U, \mathcal{F}) \to \lim H^0(U, \mathcal{F}/I^n\mathcal{F})$
是同构。
\end{lemma}

\begin{proof}
我们将通过证明引理 \ref{algebraization-lemma-application-H0} 可以应用来得到结论。
为此，取 $x \in \text{Ass}(\mathcal{F})$，其中 $x \not \in Y$，并令
$W = \overline{\{x\}}$。
由条件 (3)，有 $W \cap Y \not \subset Z$。
由引理 \ref{algebraization-lemma-sanity}，$W \cap Y$ 的任何不可约分支都不包含于 $Z$。
因此，若 $z \in W \cap Z$，则存在直接特化
$y \leadsto z$，其中 $y \in W \cap Y$ 且 $y \not \in Z$。
$y$ 的存在性使用性质，引理
\ref{properties-lemma-complement-closed-point-Jacobson}。
这时 $\delta^Y_Z(y) = 1$，且假设蕴含
$\dim(\mathcal{O}_{W, y}) > d$。
因此 $\dim(\mathcal{O}_{W, z}) > 1 + d$，结论得证。
\end{proof}

\begin{lemma}
\label{algebraization-lemma-fully-faithful-inequalities}
在情形 \ref{algebraization-situation-algebraize} 中，设 $\mathcal{F}$ 是凝聚
$\mathcal{O}_U$-模，且 $d \geq 1$。假设
\begin{enumerate}
\item $A$ 是 $I$-进完备的，有对偶化复形，并且
$\text{cd}(A, I) \leq d$；
\item 完备化 $\mathcal{F}^\wedge$（它来自 $\mathcal{F}$）
满足严格 $(1, 1 + d)$-不等式组；
\item 对 $x \in U$，若 $\overline{\{x\}} \cap Y \subset Z$，
则有 $\text{depth}(\mathcal{F}_x) \geq 2$。
\end{enumerate}
则映射
$$
\Hom_U(\mathcal{G}, \mathcal{F})
\longrightarrow
\Hom_{\textit{Coh}(U, I\mathcal{O}_U)}(\mathcal{G}^\wedge, \mathcal{F}^\wedge)
$$
对每个凝聚 $\mathcal{O}_U$-模 $\mathcal{G}$ 都是双射。
\end{lemma}

\begin{proof}
令 $\mathcal{H} = \SheafHom_{\mathcal{O}_U}(\mathcal{G}, \mathcal{F})$。
利用概形上同调，引理
\ref{coherent-lemma-hom-into-depth}，或更多代数，引理
\ref{more-algebra-lemma-hom-into-depth}，可知 $\mathcal{H}$ 的完备化
满足严格 $(1, 1 + d)$-不等式组，并且对
$x \in U$，若 $\overline{\{x\}} \cap Y \subset Z$，则有
$\text{depth}(\mathcal{H}_x) \geq 2$。细节从略。
因此由引理 \ref{algebraization-lemma-recover}，有
$$
\Hom_U(\mathcal{G}, \mathcal{F}) =
H^0(U, \mathcal{H}) =
\lim H^0(U, \mathcal{H}/\mathcal{I}^n\mathcal{H}) =
\Mor_{\textit{Coh}(U, I\mathcal{O}_U)}
(\mathcal{G}^\wedge, \mathcal{F}^\wedge)
$$
最后一个等号见概形上同调，引理
\ref{coherent-lemma-completion-internal-hom}。
\end{proof}

\begin{lemma}
\label{algebraization-lemma-construct-unique}
在情形 \ref{algebraization-situation-algebraize} 中，设 $(\mathcal{F}_n)$ 是
$\textit{Coh}(U, I\mathcal{O}_U)$ 的一个对象，且 $d \geq 1$。假设
\begin{enumerate}
\item $A$ 是 $I$-进完备的，有对偶化复形，并且
$\text{cd}(A, I) \leq d$；
\item $(\mathcal{F}_n)$ 是某个凝聚 $\mathcal{O}_U$-模的完备化；
\item $(\mathcal{F}_n)$ 满足严格 $(1, 1 + d)$-不等式组。
\end{enumerate}
则存在唯一的凝聚 $\mathcal{O}_U$-模 $\mathcal{F}$，
其完备化是 $(\mathcal{F}_n)$，并且对
$x \in U$，若 $\overline{\{x\}} \cap Y \subset Z$，则有
$\text{depth}(\mathcal{F}_x) \geq 2$。
\end{lemma}

\begin{proof}
选取一个凝聚 $\mathcal{O}_U$-模 $\mathcal{F}$，其完备化为
$(\mathcal{F}_n)$。令
$T = \{x \in U \mid \overline{\{x\}} \cap Y \subset Z\}$.
我们将构造 $\mathcal{F}$：把局部上同调，引理
\ref{local-cohomology-lemma-make-S2-along-T-simple}
应用于 $\mathcal{F}$ 和 $T$。
唯一性则由引理 \ref{algebraization-lemma-fully-faithful-inequalities} 的映射性质推出。

\medskip\noindent
由于 $T$ 在 $U$ 中对特化稳定，只需检验下述条件。
若 $x' \leadsto x$ 是 $U$ 中点的直接特化，且
$x \in T$、$x' \not \in T$，则
$\text{depth}(\mathcal{F}_{x'}) \geq 1$。
令 $W = \overline{\{x\}}$，并令 $W' = \overline{\{x'\}}$。
由于 $x' \not \in T$，可知 $W' \cap Y$ 不包含于 $Z$。
若 $W' \cap Y$ 含有一个包含于 $Z$ 的不可约分支，
则由引理 \ref{algebraization-lemma-sanity} 得证。
否则，选取一个不可约分支 $W_1$（它属于 $W \cap Y$），以及
一个不可约分支 $W'_1$（它属于 $W' \cap Y$），使得
$W_1 \subset W'_1$。
设 $z \in W_1$ 是泛点。设 $y \leadsto z$、$y \in W'_1$
构成直接特化，且 $y \not \in Z$；$y$ 的存在性由
$W'_1 \not \subset Z$（见上文）以及性质，引理
\ref{properties-lemma-complement-closed-point-Jacobson} 推出。
于是有 $z \in Z$、$x \leadsto z$、
$x' \leadsto y \leadsto z$、$y \in Y \setminus Z$，以及
$\delta^Y_Z(y) = 1$。
由局部上同调，引理 \ref{local-cohomology-lemma-cd-bound-dim-local}
以及 $z$ 是 $W \cap Y$ 的泛点这一事实，有
$\dim(\mathcal{O}_{W, z}) \leq d$。
由于 $x' \leadsto x$ 是直接特化，有
$\dim(\mathcal{O}_{W', z}) \leq d + 1$。
由于 $y \not = z$，可得
$\dim(\mathcal{O}_{W', y}) \leq d$。
若 $\text{depth}(\mathcal{F}_{x'}) = 0$，则与假设 (3) 矛盾；
从 $\mathcal{O}_{X, y}$ 过渡到其完备化的细节从略。
证明完成。
\end{proof}








\section{凝聚形式模的代数化，IV}
\label{algebraization-section-algebraization-modules-local}

\noindent
在本节中，我们证明引理 \ref{algebraization-lemma-algebraization-principal-variant}
在局部情形下的两个更强版本，即引理
\ref{algebraization-lemma-algebraization-principal} 和
\ref{algebraization-lemma-algebraization-principal-bis}。
虽然更一般的命题 \ref{algebraization-proposition-cd-1} 将使这些引理不再需要，
但它们的证明要容易得多。

\begin{lemma}
\label{algebraization-lemma-algebraization-principal}
在情形 \ref{algebraization-situation-algebraize} 中，设 $(\mathcal{F}_n)$ 是
$\textit{Coh}(U, I\mathcal{O}_U)$ 的一个对象。假设
\begin{enumerate}
\item $A$ 是局部环，且 $\mathfrak a = \mathfrak m$ 是极大理想；
\item $A$ 有对偶化复形；
\item $I = (f)$ 是由非零因子 $f \in \mathfrak m$ 生成的主理想；
\item $\mathcal{F}_n$ 是有限局部自由
$\mathcal{O}_U/f^n\mathcal{O}_U$-模；
\item 若 $\mathfrak p \in V(f) \setminus \{\mathfrak m\}$，则
$\text{depth}((A/f)_\mathfrak p) + \dim(A/\mathfrak p) > 1$；
\item 若 $\mathfrak p \not \in V(f)$ 且
$V(\mathfrak p) \cap V(f) \not = \{\mathfrak m\}$，则
$\text{depth}(A_\mathfrak p) + \dim(A/\mathfrak p) > 3$。
\end{enumerate}
则 $(\mathcal{F}_n)$ 可典范地延拓到 $X$。特别地，若 $A$ 完备，
则 $(\mathcal{F}_n)$ 是某个凝聚 $\mathcal{O}_U$-模的完备化。
\end{lemma}

\begin{proof}
我们将通过验证引理 \ref{algebraization-lemma-when-done} 的假设 (a)、(b) 和 (c)
来证明该结论。

\medskip\noindent
由于 $\mathcal{F}_n$ 在 $\mathcal{O}_U/f^n\mathcal{O}_U$ 上局部自由，
所以有短正合序列
$0 \to \mathcal{F}_n \to \mathcal{F}_{n + 1} \to \mathcal{F}_1 \to 0$
，对所有 $n$ 均成立。因此由上同调，引理
\ref{cohomology-lemma-topology-I-adic-f}，条件 (b) 成立。

\medskip\noindent
对 $n$ 归纳，并使用短正合序列
$0 \to A/f^n \to A/f^{n + 1} \to A/f \to 0$，可知
$A/f^nA$ 的相伴素理想与 $A/fA$ 的相伴素理想相同。
由于假设 (3)，$\mathcal{F}_n$ 的相伴点对应于
$A/f^nA$ 的不等于 $\mathfrak m$ 的相伴素理想。因此由 (5) 以及
局部上同调，命题 \ref{local-cohomology-proposition-kollar}，可知
$M_n = H^0(U, \mathcal{F}_n)$ 是有限 $A$-模。
故假设 (c) 成立。

\medskip\noindent
为完成证明，只需证明存在 $n > 1$，使映射
$$
H^1(U, \mathcal{F}_n) \longrightarrow H^1(U, \mathcal{F}_1)
$$
的像作为 $A$-模具有有限长度。事实上，由上同调，引理
\ref{cohomology-lemma-ML-better}，这将蕴含假设 (a)。由引理
\ref{algebraization-lemma-ML-local}，当 $n$ 充分大时，该像与 $n$ 无关。
设 $\omega_A^\bullet$ 是 $A$ 的归一化对偶化复形。
由局部对偶定理和 Matlis 对偶
（对偶化复形，引理 \ref{dualizing-lemma-special-case-local-duality}
以及命题 \ref{dualizing-proposition-matlis}），我们的断言等价于：
映射
$$
\text{Ext}^{-2}_A(M_1, \omega_A^\bullet) \to
\text{Ext}^{-2}_A(M_n, \omega_A^\bullet)
$$
的像对 $n \gg 1$ 具有有限长度。所涉及的模是有限 $A$-模，
其支集位于 $V(f)$。因此，只需证明在素理想 $\mathfrak q$
处局部化后该映射为零，其中它包含 $f$ 且不等于 $\mathfrak m$。
设 $\omega_{A_\mathfrak q}^\bullet$ 是 $A_\mathfrak q$ 上的
归一化对偶化复形，并回忆
$\omega_{A_\mathfrak q}^\bullet =
(\omega_A^\bullet)_\mathfrak q[\dim(A/\mathfrak q)]$；这是由
对偶化复形，引理 \ref{dualizing-lemma-dimension-function} 得到的。
利用 (4) 中给出的 $\mathcal{F}_n$ 的局部结构，只需证明
$$
\text{Ext}^{-2 + \dim(A/\mathfrak q)}_{A_\mathfrak q}(
A_\mathfrak q/f, \omega_{A_\mathfrak q}^\bullet)
\to
\text{Ext}^{-2 + \dim(A/\mathfrak q)}_{A_\mathfrak q}(
A_\mathfrak q/f^n, \omega_{A_\mathfrak q}^\bullet)
$$
在 $n$ 充分大时为零。若 $\dim(A/\mathfrak q) > 3$，则由局部上同调，
引理 \ref{local-cohomology-lemma-sitting-in-degrees} 立即得到。
在其他情形下，使用长正合序列
$$
\ldots
\xrightarrow{f^n}
H^{-1}(\omega_{A_\mathfrak q}^\bullet)
\to
\text{Ext}^{-1}_{A_\mathfrak q}(
A_\mathfrak q/f^n, \omega_{A_\mathfrak q}^\bullet) \to
H^0(\omega_{A_\mathfrak q}^\bullet)
\xrightarrow{f^n}
H^0(\omega_{A_\mathfrak q}^\bullet)
\to
\text{Ext}^0_{A_\mathfrak q}(
A_\mathfrak q/f^n, \omega_{A_\mathfrak q}^\bullet) \to 0
$$
若 $\dim(A/\mathfrak q) = 2$，则
$H^0(\omega_{A_\mathfrak q}^\bullet) = 0$
，因为 $\text{depth}(A_\mathfrak q) \geq 1$，这是由于
$f$ 是非零因子。
因此长正合序列表明，所需条件是映射
$$
f^{n - 1} :
H^{-1}(\omega_{A_\mathfrak q}^\bullet)/f \to
H^{-1}(\omega_{A_\mathfrak q}^\bullet)/f^n
$$
为零。现在 $H^{-1}(\omega^\bullet_\mathfrak q)$ 是有限模，
其支集包含在满足下式的素理想 $\mathfrak p \subset A_\mathfrak q$ 中：
$\text{depth}(A_\mathfrak p) + \dim((A/\mathfrak p)_\mathfrak q) \leq 1$.
由于 $\dim((A/\mathfrak p)_\mathfrak q) = \dim(A/\mathfrak p) - 2$，
条件 (6) 表明这些素理想包含于 $V(f)$。因此，当 $n$ 充分大时，
得到所需的零化。
最后，若 $\dim(A/\mathfrak q) = 1$，则条件 (5) 与 $f$ 为非零因子
共同保证 $A_\mathfrak q$ 的深度至少为 $2$。因此
$H^0(\omega_{A_\mathfrak q}^\bullet) =
H^{-1}(\omega_{A_\mathfrak q}^\bullet) = 0$
，且长正合序列表明，该断言等价于下述映射为零：
$$
f^{n - 1} :
H^{-2}(\omega_{A_\mathfrak q}^\bullet)/f \to
H^{-2}(\omega_{A_\mathfrak q}^\bullet)/f^n
$$
现在 $H^{-2}(\omega^\bullet_\mathfrak q)$ 是有限模，
其支集包含在满足 $\mathfrak p \subset A_\mathfrak q$ 以及
$\text{depth}(A_\mathfrak p) + \dim((A/\mathfrak p)_\mathfrak q)
\leq 2$ 的素理想中。由条件 (6)，所有这些素理想都包含于 $V(f)$。
因此，当 $n$ 充分大时，得到所需的零化。
\end{proof}

\begin{remark}
\label{algebraization-remark-interesting-case}
设 $(A, \mathfrak m)$ 是维数 $\geq 4$ 的完备诺特正规局部整环，
并设 $f \in \mathfrak m$ 非零。此时引理
\ref{algebraization-lemma-algebraization-principal} 的假设 (1)、(2)、(3)、(5) 和 (6)
都成立。因此，$U$ 沿 $U \cap V(f)$ 的形式完备化上的向量丛
可以代数化。在引理 \ref{algebraization-lemma-algebraization-principal-bis} 中，
我们会将此推广到更一般的凝聚形式模；另请比较注记
\ref{algebraization-remark-interesting-case-bis}。
\end{remark}

\begin{lemma}
\label{algebraization-lemma-helper-algebraize}
在情形 \ref{algebraization-situation-algebraize} 中，设 $(M_n)$ 是如引理
\ref{algebraization-lemma-system-of-modules} 所示的 $A$-模逆系统，并设
$(\mathcal{F}_n)$ 是 $\textit{Coh}(U, I\mathcal{O}_U)$ 中相应的对象。
设 $d \geq \text{cd}(A, I)$ 和 $s \geq 0$ 为整数。
沿用上述记号，假设
\begin{enumerate}
\item $A$ 是局部环，其极大理想为 $\mathfrak m = \mathfrak a$；
\item $A$ 有对偶化复形；
\item $(\mathcal{F}_n)$ 满足 $(s, s + d)$-不等式组
（定义 \ref{algebraization-definition-s-d-inequalities}）。
\end{enumerate}
设 $E$ 是 $A$ 的剩余域的内射包。则对 $i \leq s$，
存在有限 $A$-模 $N$，它被 $I$ 的某个幂零化，并且对 $n \gg 0$
存在相容映射
$$
H^i_\mathfrak m(M_n) \to \Hom_A(N, E)
$$
，其余核是有限长度 $A$-模，其核 $K_n$ 构成逆系统，并且
$\Im(K_{n''} \to K_{n'})$ 对 $n'' \gg n' \gg 0$ 具有有限长度。
\end{lemma}

\begin{proof}
设 $\omega_A^\bullet$ 是归一化对偶化复形。则
$\delta^Y_Z = \delta$ 是与该对偶化复形相伴的维数函数。
注意，$\Ext^{-i}_A(M_n, \omega_A^\bullet)$ 是有限 $A$-模，
并被 $I^n$ 零化。固定 $0 \leq i \leq s$。
下面将找到 $n_1 > n_0 > 0$，使得若令
$$
N = \Im(\Ext^{-i}_A(M_{n_0}, \omega_A^\bullet) \to
\Ext^{-i}_A(M_{n_1}, \omega_A^\bullet))
$$
，则映射
$$
N \to \Ext^{-i}_A(M_n, \omega_A^\bullet),\quad n \geq n_1
$$
的核是有限长度 $A$-模，而余核 $Q_n$ 构成一个系统，使得
$\Im(Q_{n'} \to Q_{n''})$ 对 $n'' \gg n' \gg n_1$ 具有有限长度。
这等价于系统
$\{\Ext^{-i}_A(M_n, \omega_A^\bullet)\}_{n \geq 1}$
在有限 $A$-模范畴模去有限长度 $A$-模的 Serre 子范畴所得的商中
本质上为常值。由局部对偶定理
（对偶化复形，引理 \ref{dualizing-lemma-special-case-local-duality}）
以及 Matlis 对偶
（对偶化复形，命题 \ref{dualizing-proposition-matlis}），可知存在映射
$$
H^i_\mathfrak m(M_n) \to \Hom_A(N, E),\quad n \geq n_1
$$
，具有引理陈述中的性质。

\medskip\noindent
选取 $f \in \mathfrak m$。令 $B = A_f^\wedge$ 为 $I$-进完备化，
它取自局部化 $A_f$。回忆
$\omega_{A_f}^\bullet = \omega_A^\bullet \otimes_A A_f$
和 $\omega_B^\bullet = \omega_A^\bullet \otimes_A B$ 都是对偶化复形
（对偶化复形，引理 \ref{dualizing-lemma-dualizing-localize} 和
\ref{dualizing-lemma-completion-henselization-dualizing}）。
令 $M$ 为有限 $B$-模 $\lim M_{n, f}$（比较概形上同调，引理
\ref{coherent-lemma-inverse-systems-affine} 中的讨论）。则
$$
\Ext^{-i}_A(M_n, \omega_A^\bullet)_f =
\Ext^{-i}_{A_f}(M_{n, f}, \omega_{A_f}^\bullet) =
\Ext^{-i}_B(M/I^n M, \omega_B^\bullet)
$$
由于 $\mathfrak m$ 可以由有限多个 $f \in \mathfrak m$ 生成，
只需证明对每个 $f$，系统
$$
\{\Ext^{-i}_B(M/I^n M, \omega_B^\bullet)\}_{n \geq 1}
$$
本质上为常值。若干细节从略。

\medskip\noindent
设 $\mathfrak q \subset IB$ 是素理想。则 $\mathfrak q$ 对应于一点
$y \in U \cap Y$。注意，
$\delta(\mathfrak q) = \dim(\overline{\{y\}})$
也是与 $\omega_B^\bullet$ 相伴的维数函数之值
（细节从略；使用 $\omega_B^\bullet$ 是由
$\omega_A^\bullet$ 与 $B$ 作张量积得到的这一事实）。
假设 (3) 通过引理 \ref{algebraization-lemma-elementary} 保证引理
\ref{algebraization-lemma-algebraize-local-cohomology-bis} 可应用于
$B_\mathfrak q, IB_\mathfrak q, \mathfrak qB_\mathfrak q, M_\mathfrak q$
，其中将 $s$ 换为 $s - \delta(y)$。于是得到
$$
H^{i - \delta(\mathfrak q)}_{\mathfrak qB_\mathfrak q}(M_\mathfrak q) =
\lim H^{i - \delta(\mathfrak q)}_{\mathfrak qB_\mathfrak q}(
(M/I^nM)_\mathfrak q)
$$
，且该模被 $I$ 的某个幂零化。
由引理 \ref{algebraization-lemma-terrific}，逆系统
$H^{i - \delta(\mathfrak q)}_{\mathfrak qB_\mathfrak q}((M/I^nM)_\mathfrak q)$
本质上为常值，其值为
$H^{i - \delta(\mathfrak q)}_{\mathfrak qB_\mathfrak q}(M_\mathfrak q)$.
由于 $(\omega_B^\bullet)_\mathfrak q[-\delta(\mathfrak q)]$ 是
$B_\mathfrak q$ 上的归一化对偶化复形，局部对偶定理表明系统
$$
\Ext^{-i}_B(M/I^n M, \omega_B^\bullet)_\mathfrak q
$$
本质上为常值，其值为
$\Ext^{-i}_B(M, \omega_B^\bullet)_\mathfrak q$.

\medskip\noindent
为完成证明，我们像引理 \ref{algebraization-lemma-bootstrap} 的证明那样作整体化；
这里的论证更容易，因为已经知道该系统的值。具体而言，考虑映射
$$
\alpha_n :
\Ext^{-i}_B(M/I^n M, \omega_B^\bullet)
\longrightarrow
\Ext^{-i}_B(M, \omega_B^\bullet)
$$
，其中 $n$ 变化。由上文，对每个 $\mathfrak q$，可找到一个
$n$，使得 $\alpha_n$ 在 $\mathfrak q$ 处局部化后为满射。
由于 $B$ 是诺特的，且 $\Ext^{-i}_B(M, \omega_B^\bullet)$
是有限模，可找到一个 $n$，使得 $\alpha_n$ 为满射。
对任意 $n$，若 $\alpha_n$ 为满射，则给定素理想
$\mathfrak q \in V(IB)$，可找到 $n' > n$，使得
$\Ker(\alpha_n)$ 到 $\Ext^{-i}(M/I^{n'}M, \omega_B^\bullet)$ 的映射为零，
至少在 $\mathfrak q$ 处局部化后如此。
由于 $\Ker(\alpha_n)$ 是有限 $A$-模，且截面的支集是拟紧的，
可找到一个 $n'$，使得 $\Ker(\alpha_n)$ 到
$\Ext^{-i}(M/I^{n'}M, \omega_B^\bullet)$ 的映射为零。
由此可见，$\Ext^{-i}(M/I^n M, \omega_B^\bullet)$ 本质上为常值，
其值为 $\Ext^{-i}(M, \omega_B^\bullet)$。证明完成。
\end{proof}

\noindent
下面是引理 \ref{algebraization-lemma-algebraization-principal} 的一个更一般版本。

\begin{lemma}
\label{algebraization-lemma-algebraization-principal-bis}
在情形 \ref{algebraization-situation-algebraize} 中，设 $(\mathcal{F}_n)$ 是
$\textit{Coh}(U, I\mathcal{O}_U)$ 的一个对象。假设
\begin{enumerate}
\item $A$ 是局部环，且 $\mathfrak a = \mathfrak m$ 是极大理想；
\item $A$ 有对偶化复形；
\item $I = (f)$ 是主理想；
\item $(\mathcal{F}_n)$ 满足 $(2, 3)$-不等式组。
\end{enumerate}
则 $(\mathcal{F}_n)$ 可延拓到 $X$。特别地，若 $A$ 是
$I$-进完备的，则 $(\mathcal{F}_n)$ 是某个凝聚
$\mathcal{O}_U$-模的完备化。
\end{lemma}

\begin{proof}
回忆 $\textit{Coh}(U, I\mathcal{O}_U)$ 是阿贝尔范畴；参见概形上同调，
引理 \ref{coherent-lemma-inverse-systems-abelian}。
在 $U$ 的仿射开集上，对象 $(\mathcal{F}_n)$ 对应于诺特环上的有限模
（概形上同调，引理 \ref{coherent-lemma-inverse-systems-affine}）。
因此，映射 $f^N : (\mathcal{F}_n) \to (\mathcal{F}_n)$ 的核在
$N$ 充分大时稳定。由引理
\ref{algebraization-lemma-map-kernel-cokernel-on-closed} 和
\ref{algebraization-lemma-essential-image-completion}，为证明本引理，
可以用这样一个映射的像替换 $(\mathcal{F}_n)$。
因此可假设 $f$ 在 $(\mathcal{F}_n)$ 上是单射。
替换后，引理 \ref{algebraization-lemma-equivalent-f-good} 的等价条件对逆系统
$(\mathcal{F}_n)$ 在 $U$ 上成立。证明余下部分将不再另行说明地使用这一点。

\medskip\noindent
我们将检验引理 \ref{algebraization-lemma-when-done} 的假设 (a)、(b) 和 (c)。
由上同调，引理 \ref{cohomology-lemma-topology-I-adic-f}，假设 (b) 成立。

\medskip\noindent
选取如引理 \ref{algebraization-lemma-system-of-modules} 所示的模逆系统 $\{M_n\}$。
可以假设 $H^0_\mathfrak m(M_n) = 0$：必要时把 $M_n$ 替换为
$M_n/H^0_\mathfrak m(M_n)$ 即可。于是得到短正合序列
$$
0 \to M_n \to H^0(U, \mathcal{F}_n) \to H^1_\mathfrak m(M_n) \to 0
$$
，对所有 $n$ 均成立。设 $E$ 是 $A$ 的剩余域的内射包。
由引理 \ref{algebraization-lemma-helper-algebraize} 和当前假设 (4)，可以选取
一个整数 $m \geq 0$、有限 $A$-模 $N_1$ 和 $N_2$；后二者被
$f^c$ 零化，其中某个 $c \geq 0$，并且存在相容映射系统
$$
H^i_\mathfrak m(M_n) \to \Hom_A(N_i, E), \quad i = 1, 2
$$
，对 $n \geq m$ 成立，并具有引理中陈述的性质。

\medskip\noindent
已知 $M = \lim H^0(U, \mathcal{F}_n)$ 是 $A$-模，
其极限拓扑为 $f$-进拓扑。因此，给定 $n$，模
$M/f^nM$ 是 $H^0(U, \mathcal{F}_N)$ 的子商，其中 $N \gg n$。
由上面得到的信息可知，$f^cM/f^nM$ 是有限 $A$-模。
由于 $f$ 在 $M$ 上是非零因子，可知 $M/f^{n - c}M$ 是有限 $A$-模。
由此可见，引理 \ref{algebraization-lemma-when-done} 的假设 (c) 成立。

\medskip\noindent
接下来研究模
$$
Ob = \lim H^1(U, \mathcal{F}_n) = \lim H^2_\mathfrak m(M_n)
$$
对 $n \geq m$，令 $K_n$ 为映射
$H^2_\mathfrak m(M_n) \to \Hom_A(N_2, E)$.
的核。令 $K = \lim K_n$。得到正合序列
$$
0 \to K \to Ob \to \Hom_A(N_2, E)
$$
由上文，$Ob = \lim H^2_\mathfrak m(M_n)$ 上的极限拓扑是
$f$-进拓扑。由于 $N_2$ 被 $f^c$ 零化，可知
$K = \lim K_n$ 上的极限拓扑也同样如此。
因此 $K/fK$ 是 $K_n$ 的子商，其中 $n \gg 1$。
然而，由于 $\{K_n\}$ 与有限长度 $A$-模的某个逆系统 pro-同构
（由引理 \ref{algebraization-lemma-helper-algebraize} 的结论），可知
$K/fK$ 是某个有限长度 $A$-模的子商。因此 $K$ 是有限 $A$-模；
参见代数，引理 \ref{algebra-lemma-finite-over-complete-ring}。
（事实上，甚至可见 $\dim(\text{Supp}(K)) = 1$，但这里不需要这一点。）

\medskip\noindent
给定 $n \geq 1$，考虑边缘映射
$$
\delta_n :
H^0(U, \mathcal{F}_n)
\longrightarrow
\lim_N H^1(U, f^n\mathcal{F}_N) \xrightarrow{f^{-n}} Ob
$$
（第二个映射是同构），它来自短正合序列
$$
0 \to f^n\mathcal{F}_N \to \mathcal{F}_N \to \mathcal{F}_n \to 0
$$
对每个 $n$，令
$$
P_n = \Im(H^0(U, \mathcal{F}_{n + m}) \to H^0(U, \mathcal{F}_n))
$$
，其中 $m$ 如上。注意，$\{P_n\}$ 是逆系统，并且映射
$f : \mathcal{F}_n \to \mathcal{F}_{n + 1}$ 在整体截面上把
$P_n$ 映入 $P_{n + 1}$。
若 $p \in P_n$，则 $\delta_n(p) \in K \subset Ob$，
因为 $\delta_n(p)$ 在
$H^1(U, f^n\mathcal{F}_{n + m}) = H^2_\mathfrak m(M_m)$
中的像为零，并且 $\delta_n$ 与 $Ob \to \Hom_A(N_2, E)$ 的复合
通过 $H^2_\mathfrak m(M_m)$ 分解；这是由 $m$ 的选取保证的。
因此
$$
\bigoplus\nolimits_{n \geq 0} \Im(P_n \to Ob)
$$
是有限生成分次 $A[T]$-模，其中 $T$ 通过乘以 $f$ 起作用。
事实上，它是 $K[T]$ 的分次子模，而 $K$ 在 $A$ 上有限。
仿照上同调，引理 \ref{cohomology-lemma-ML-general} 的证明进行论证
\footnote{为所显示的模选取形如 $\delta_{n_j}(p_j)$ 的齐次生成元。
若 $k = \max(n_j)$，则对 $n \geq k$ 和任意 $p \in P_n$，
可找到 $a_j \in A$，使得
$p - \sum a_j f^{n - n_j} p_j$ 位于 $\delta_n$ 的核中，
从而位于 $P_{n'}$ 的像中；这对所有 $n' \geq n$ 成立。
因此有
$\Im(P_n \to P_{n - k}) = \Im(P_{n'} \to P_{n - k})$，
对所有 $n' \geq n$ 均成立。}
可知逆系统 $\{P_n\}$ 满足 ML。
由于 $\{P_n\}$ 与 $\{H^0(U, \mathcal{F}_n)\}$ pro-同构，
可知 $\{H^0(U, \mathcal{F}_n)\}$ 满足 ML。
因此，引理 \ref{algebraization-lemma-when-done} 的假设 (a) 成立，证明完成。
\end{proof}

\noindent
引理 \ref{algebraization-lemma-algebraization-principal-bis} 的条件可以具体展开如下。

\begin{lemma}
\label{algebraization-lemma-unwinding-conditions}
在情形 \ref{algebraization-situation-algebraize} 中，设 $(\mathcal{F}_n)$ 为
$\textit{Coh}(U, I\mathcal{O}_U)$ 的一个对象。假设
\begin{enumerate}
\item $A$ 是局部环，其极大理想为 $\mathfrak a = \mathfrak m$；
\item $\text{cd}(A, I) = 1$.
\end{enumerate}
则 $(\mathcal{F}_n)$ 满足 $(2, 3)$-不等式，当且仅当对每个
$y \in U \cap Y$，若 $\dim(\{y\}) = 1$，以及对每个素理想
$\mathfrak p \subset \mathcal{O}_{X, y}^\wedge$，若
$\mathfrak p \not \in V(I\mathcal{O}_{X, y}^\wedge)$，便有
$$
\text{depth}((\mathcal{F}_y^\wedge)_\mathfrak p) +
\dim(\mathcal{O}_{X, y}^\wedge/\mathfrak p) > 2
$$
\end{lemma}

\begin{proof}
以下将不再另行说明地使用引理 \ref{algebraization-lemma-explain-2-3-cd-1}。
它立即表明该条件是必要的。反过来，假设该条件成立。
由引理 \ref{algebraization-lemma-discussion}，有
$\delta^Y_Z(y) = \dim(\overline{\{y\}})$；把这个函数简记为 $\delta$。
设 $y \in U \cap Y$。若 $\delta(y) > 2$，则引理
\ref{algebraization-lemma-explain-2-3-cd-1} 的不等式成立。
最后假设 $\delta(y) = 2$。我们必须证明
$$
\text{depth}((\mathcal{F}_y^\wedge)_\mathfrak p) +
\dim(\mathcal{O}_{X, y}^\wedge/\mathfrak p) > 1
$$
选取一个特化 $y \leadsto y'$，使得 $\delta(y') = 1$。此时存在环映射
$\mathcal{O}_{X, y'}^\wedge \to \mathcal{O}_{X, y}^\wedge$，
它把靶环识别为 $\mathcal{O}_{X, y'}^\wedge$ 在某个素理想
$\mathfrak q$ 处局部化之后的完备化，并且
$\dim(\mathcal{O}_{X, y'}^\wedge/\mathfrak q) = 1$。此外有
$$
\mathcal{F}_y^\wedge =
\mathcal{F}_{y'}^\wedge
\otimes_{\mathcal{O}_{X, y'}^\wedge}
\mathcal{O}_{X, y}^\wedge
$$
令 $\mathfrak p' \subset \mathcal{O}_{X, y'}^\wedge$ 为
$\mathfrak p$ 的像。由《局部上同调》，引理
\ref{local-cohomology-lemma-change-completion}，有
\begin{align*}
\text{depth}((\mathcal{F}_y^\wedge)_\mathfrak p) +
\dim(\mathcal{O}_{X, y}^\wedge/\mathfrak p)
& =
\text{depth}((\mathcal{F}_{y'}^\wedge)_{\mathfrak p'}) +
\dim((\mathcal{O}_{X, y}^\wedge/\mathfrak p)_{\mathfrak p'}) \\
& =
\text{depth}((\mathcal{F}_{y'}^\wedge)_{\mathfrak p'}) +
\dim(\mathcal{O}_{X, y}^\wedge/\mathfrak p') - 1
\end{align*}
最后一个等号成立，是因为这个特化是直接特化。
现在对 $y', \mathfrak p'$ 使用所假设的不等式，便证明了本引理。
\end{proof}

\begin{lemma}
\label{algebraization-lemma-unwinding-conditions-bis}
在情形 \ref{algebraization-situation-algebraize} 中，设 $(\mathcal{F}_n)$ 为
$\textit{Coh}(U, I\mathcal{O}_U)$ 的一个对象。假设
\begin{enumerate}
\item $A$ 是局部环，其极大理想为 $\mathfrak a = \mathfrak m$；
\item $A$ 有对偶化复形；
\item $\text{cd}(A, I) = 1$,
\item 对 $y \in U \cap Y$，模 $\mathcal{F}_y^\wedge$ 在
$V(I\mathcal{O}_{X, y}^\wedge)$ 外有限局部自由；例如，
若 $\mathcal{F}_n$ 是有限局部自由
$\mathcal{O}_U/I^n\mathcal{O}_U$-模，则此条件成立；并且
\item 下列条件之一成立：
\begin{enumerate}
\item $A_f$ 满足 $(S_2)$，且 $X$ 的每个不包含在 $Y$ 中的
不可约分支之维数都 $\geq 4$；或者
\item 若 $\mathfrak p \not \in V(f)$ 且
$V(\mathfrak p) \cap V(f) \not = \{\mathfrak m\}$，则
$\text{depth}(A_\mathfrak p) + \dim(A/\mathfrak p) > 3$。
\end{enumerate}
\end{enumerate}
则 $(\mathcal{F}_n)$ 满足 $(2, 3)$-不等式。
\end{lemma}

\begin{proof}
我们使用引理 \ref{algebraization-lemma-unwinding-conditions} 的判据。
取 $y \in U \cap Y$，它满足 $\dim(\overline{\{y\}} = 1$；并取素理想
$\mathfrak p$，它满足 $\mathfrak p \subset \mathcal{O}_{X, y}^\wedge$ 且
$\mathfrak p \not \in V(I\mathcal{O}_{X, y}^\wedge)$。
条件 (4) 表明
$\text{depth}((\mathcal{F}_y^\wedge)_\mathfrak p) =
\text{depth}((\mathcal{O}_{X, y}^\wedge)_\mathfrak p)$.
因此只需证明
$$
\text{depth}((\mathcal{O}_{X, y}^\wedge)_\mathfrak p) +
\dim(\mathcal{O}_{X, y}^\wedge/\mathfrak p) > 2
$$
令 $\mathfrak p_0 \subset A$ 为 $\mathfrak p$ 的像。
令 $\mathfrak q \subset A$ 为与 $y$ 对应的素理想。
由《局部上同调》，引理
\ref{local-cohomology-lemma-change-completion}，有
\begin{align*}
\text{depth}((\mathcal{O}_{X, y}^\wedge)_\mathfrak p) +
\dim(\mathcal{O}_{X, y}^\wedge/\mathfrak p)
& =
\text{depth}(A_{\mathfrak p_0}) + \dim((A/\mathfrak p_0)_\mathfrak q) \\
& =
\text{depth}(A_{\mathfrak p_0}) + \dim(A/\mathfrak p_0) - 1
\end{align*}
若 (5)(a) 成立，则这个量
$$
\geq \min(2, \dim(A_{\mathfrak p_0})) + \dim(A/\mathfrak p_0) - 1
$$
注意，无论如何都有 $\dim(A/\mathfrak p_0) \geq 2$。因此，若该最小值为
$2$，结论已经成立。否则，我们得到
$$
\dim(A_{\mathfrak p_0}) + \dim(A/\mathfrak p_0) - 1 \geq 4 - 1
$$
，因为 $\Spec(A)$ 的每个经过 $\mathfrak p_0$ 的分支之维数都
$\geq 4$。若 (5)(b) 成立，则立即得到结论。
\end{proof}

\begin{remark}
\label{algebraization-remark-interesting-case-bis}
设 $(A, \mathfrak m)$ 为诺特局部环，它有对偶化复形，并且关于
$f \in \mathfrak m$ 完备。设 $(\mathcal{F}_n)$ 为
$\textit{Coh}(U, f\mathcal{O}_U)$ 的一个对象，其中
$U$ 是 $A$ 的穿孔谱。置 $Y = V(f) \subset X = \Spec(A)$。
对 $y \in U \cap V(f)$，若它在 $U$ 中闭，也就是满足
$\dim(\overline{\{y\}}) = 1$ 的点，假设
$\mathcal{O}_{X, y}^\wedge$-模 $\mathcal{F}_y^\wedge$
满足下列两个条件：
\begin{enumerate}
\item $\mathcal{F}_y^\wedge[1/f]$ 满足 $(S_2)$；这里把它视为
$\mathcal{O}_{X, y}^\wedge[1/f]$-模；并且
\item 对 $\mathfrak p \in \text{Ass}(\mathcal{F}_y^\wedge[1/f])$，有
$\dim(\mathcal{O}_{X, y}^\wedge/\mathfrak p) \geq 3$。
\end{enumerate}
则 $(\mathcal{F}_n)$ 是 $U$ 上某个凝聚模的完备化。
这由引理 \ref{algebraization-lemma-algebraization-principal-bis}
和 \ref{algebraization-lemma-unwinding-conditions} 得出。
\end{remark}
































\section{改进凝聚形式模}
\label{algebraization-section-improving-formal-modules}

\noindent
设 $X$ 为诺特概形。设 $Y \subset X$ 为闭子概形，
其拟凝聚理想层为 $\mathcal{I} \subset \mathcal{O}_X$。
设 $(\mathcal{F}_n)$ 为 $\textit{Coh}(X, \mathcal{I})$ 的一个对象。
本节构造映射
$(\mathcal{F}_n) \to (\mathcal{F}'_n)$
，它们类似于《局部上同调》，第
\ref{local-cohomology-section-improve} 节中为凝聚模构造的映射。
对一点 $y \in Y$，置
$$
\mathcal{O}_{X, y}^\wedge = \lim \mathcal{O}_{X, y}/\mathcal{I}^n_y, \quad
\mathcal{I}_y^\wedge = \lim \mathcal{I}_y/\mathcal{I}^n_y
\quad\text{且}\quad
\mathfrak m_y^\wedge = \lim \mathfrak m_y/\mathcal{I}_y^n
$$
则 $\mathcal{O}_{X, y}^\wedge$ 是诺特局部环，其极大理想为
$\mathfrak m_y^\wedge$，并且关于
$\mathcal{I}_y^\wedge = \mathcal{I}_y\mathcal{O}_{X, y}^\wedge$ 完备。
还置
$$
\mathcal{F}_y^\wedge = \lim \mathcal{F}_{n, y}
$$
则 $\mathcal{F}_y^\wedge$ 是
$\mathcal{O}_{X, y}^\wedge$ 上的有限模，并且
$\mathcal{F}_y^\wedge/(\mathcal{I}_y^\wedge)^n\mathcal{F}_y^\wedge =
\mathcal{F}_{n, y}$ 对所有 $n$ 成立；参见《代数》，引理
\ref{algebra-lemma-limit-complete} 和
\ref{algebra-lemma-finite-over-complete-ring}.

\begin{lemma}
\label{algebraization-lemma-divide-torsion-formal-coherent-module}
在上述情形中，假设 $X$ 局部有对偶化复形。
设 $T \subset Y$ 为对特化稳定的子集。
假设对 $y \in T$ 以及满足
$\mathfrak p \subset \mathcal{O}_{X, y}^\wedge$ 的非极大素理想，
若 $V(\mathfrak p) \cap V(\mathcal{I}^\wedge_y) = \{\mathfrak m_y^\wedge\}$，
则
$$
\text{depth}_{(\mathcal{O}_{X, y})_\mathfrak p}
((\mathcal{F}^\wedge_y)_\mathfrak p) > 0
$$
于是存在典范映射
$(\mathcal{F}_n) \to (\mathcal{F}_n')$
，它是凝聚 $\mathcal{O}_X$-模逆系统之间的映射，并满足下列性质：
\begin{enumerate}
\item 对 $y \in T$，有 $\text{depth}(\mathcal{F}'_{n, y}) \geq 1$；
\item $(\mathcal{F}'_n)$ 作为 pro-系统同构于一个对象
$(\mathcal{G}_n)$，后者属于 $\textit{Coh}(X, \mathcal{I})$；
\item 所诱导的态射
$(\mathcal{F}_n) \to (\mathcal{G}_n)$ 在
$\textit{Coh}(X, \mathcal{I})$ 中是满射，其核被
$\mathcal{I}$ 的某个幂零化。
\end{enumerate}
\end{lemma}

\begin{proof}
对每个 $n$，令 $\mathcal{F}_n \to \mathcal{F}'_n$ 为
《局部上同调》，引理
\ref{local-cohomology-lemma-get-depth-1-along-Z} 中构造的满射。
这是 $\mathcal{F}_n$ 对支撑于 $T$ 的截面子层取商所得，故得到典范映射
$\mathcal{F}'_{n + 1} \to \mathcal{F}'_n$，从而得到映射
$(\mathcal{F}_n) \to (\mathcal{F}_n')$；它是凝聚
$\mathcal{O}_X$-模逆系统之间的映射。
性质 (1) 由构造成立。

\medskip\noindent
为证明性质 (2) 和 (3)，可以假设 $X = \Spec(A_0)$ 是仿射的，
且 $A_0$ 有对偶化复形。令 $I_0 \subset A_0$ 为与 $Y$ 对应的理想。
令 $A, I$ 为关于 $I_0$ 的进完备化，分别对应 $A_0, I_0$。
供后文使用，注意 $A$ 有对偶化复形
（《对偶化复形》，引理 \ref{dualizing-lemma-ubiquity-dualizing}）。
令 $M$ 为有限 $A$-模，它对应于 $(\mathcal{F}_n)$；参见
《概形上同调》，引理 \ref{coherent-lemma-inverse-systems-affine}。
于是 $\mathcal{F}_n$ 对应于 $M_n = M/I^nM$。回忆
$\mathcal{F}'_n$ 对应于商模 $M'_n = M_n / H^0_T(M_n)$；参见
《局部上同调》，引理 \ref{local-cohomology-lemma-get-depth-1-along-Z}
及其证明。

\medskip\noindent
置 $s = 0$ 和 $d = \text{cd}(A, I)$。
我们断言 $A, I, T, M, s, d$ 满足情形
\ref{algebraization-situation-bootstrap} 的假设 (1)、(3)、(4)、(6)。
事实上，(1) 和 (3) 由上文立即成立；由于 $s = 0$，(4) 是空条件；
而 (6) 正是本引理陈述中的假设。

\medskip\noindent
由定理 \ref{algebraization-theorem-final-bootstrap} 可知，$\{H^0_T(M_n)\}$
本质常值，其值为 $H^0_T(M)$。因此，短正合序列
$0 \to H^0_T(M_n) \to M_n \to M'_n \to 0$ 的极限是短正合序列
$$
0 \to H^0_T(M) \to M \to \lim M'_n \to 0
$$
置 $M' = \lim M'_n = M/H^0_T(M)$。这是有限 $A$-模。
令 $(\mathcal{G}_n)$ 为 $\textit{Coh}(X, \mathcal{I})$ 中与
$M'$ 对应的对象。为完成证明，必须证明典范映射
$\{M'/I^nM'\} \to \{M'_n\}$ 是 pro-同构。
这等价于断言
$\{H^0_T(M) + I^nM\} \to \{\Ker(M \to M'_n)\}$ 是 pro-同构。
除以 $I^nM$ 后，只需证明
$\{H^0_T(M)/H^0_T(M) \cap I^nM\} \to \{H^0_T(M_n)\}$
是 pro-同构。由于 $H^0_T(M)$ 被 $I$ 的某个幂零化，
由 Artin--Rees 引理可知，$H^0_T(M) \cap I^nM = 0$，
对所有 $n \gg 0$ 均成立。而上面已经看到
$\{H^0_T(M_n)\}$ 本质常值，其值为 $H^0_T(M)$；
故得到所需的 pro-同构。
\end{proof}

\begin{lemma}
\label{algebraization-lemma-improvement-formal-coherent-module-better}
在上述情形中，假设 $X$ 局部有对偶化复形。
设 $T' \subset T \subset Y$ 为对特化稳定的子集。
设 $d \geq 0$ 为整数。假设
\begin{enumerate}
\item[(a)] 在仿射局部有 $X = \Spec(A_0)$ 和 $Y = V(I_0)$，
并且 $\text{cd}(A_0, I_0) \leq d$；
\item[(b)] 对 $y \in T$ 以及满足
$\mathfrak p \subset \mathcal{O}_{X, y}^\wedge$ 的非极大素理想，若
$V(\mathfrak p) \cap V(\mathcal{I}_y^\wedge) = \{\mathfrak m_y^\wedge\}$，
则
$$
\text{depth}_{(\mathcal{O}_{X, y})_\mathfrak p}
((\mathcal{F}^\wedge_y)_\mathfrak p) > 0
$$
\item[(c)] 对 $y \in T'$ 以及满足
$\mathfrak p \subset \mathcal{O}_{X, y}^\wedge$ 的素理想，若
$\mathfrak p \not \in V(\mathcal{I}_y^\wedge)$ 且
$V(\mathfrak p) \cap V(\mathcal{I}_y^\wedge) \not =
\{\mathfrak m_y^\wedge\}$，则
$$
\text{depth}_{(\mathcal{O}_{X, y})_\mathfrak p}
((\mathcal{F}^\wedge_y)_\mathfrak p) \geq 1
\quad\text{或}\quad
\text{depth}_{(\mathcal{O}_{X, y})_\mathfrak p}
((\mathcal{F}^\wedge_y)_\mathfrak p) +
\dim(\mathcal{O}_{X, y}^\wedge/\mathfrak p) > 1 + d
$$
\item[(d)] 对 $y \in T'$ 以及满足
$\mathfrak p \subset \mathcal{O}_{X, y}^\wedge$ 的非极大素理想，若
$V(\mathfrak p) \cap V(\mathcal{I}_y^\wedge) = \{\mathfrak m_y^\wedge\}$，
则
$$
\text{depth}_{(\mathcal{O}_{X, y})_\mathfrak p}
((\mathcal{F}^\wedge_y)_\mathfrak p) > 1
$$
\item[(e)] 若 $y \leadsto y'$ 是直接特化且
$y' \in T'$，则 $y \in T$。
\end{enumerate}
于是存在典范映射 $(\mathcal{F}_n) \to (\mathcal{F}_n'')$，
它是凝聚 $\mathcal{O}_X$-模逆系统之间的映射，并满足下列性质：
\begin{enumerate}
\item 对 $y \in T$，有 $\text{depth}(\mathcal{F}''_{n, y}) \geq 1$；
\item 对 $y' \in T'$，有 $\text{depth}(\mathcal{F}''_{n, y'}) \geq 2$；
\item $(\mathcal{F}''_n)$ 作为 pro-系统同构于一个对象
$(\mathcal{H}_n)$，它属于 $\textit{Coh}(X, \mathcal{I})$；
\item 所诱导的态射 $(\mathcal{F}_n) \to (\mathcal{H}_n)$ 在
$\textit{Coh}(X, \mathcal{I})$ 中的核与余核被
$\mathcal{I}$ 的某个幂零化。
\end{enumerate}
\end{lemma}

\begin{proof}
仿照引理 \ref{algebraization-lemma-divide-torsion-formal-coherent-module} 及其证明，
对每个 $n$，令 $\mathcal{F}_n \to \mathcal{F}'_n$ 为
《局部上同调》，引理
\ref{local-cohomology-lemma-get-depth-1-along-Z} 中利用 $T$ 构造的满射。
接着，令 $\mathcal{F}'_n \to \mathcal{F}''_n$ 为
《局部上同调》，引理
\ref{local-cohomology-lemma-make-S2-along-T} 及其证明中构造的单射。
这些构造表明，得到典范映射
$\mathcal{F}''_{n + 1} \to \mathcal{F}''_n$，从而得到映射
$$
(\mathcal{F}_n) \longrightarrow (\mathcal{F}_n') \longrightarrow
(\mathcal{F}''_n)
$$
，它们是凝聚 $\mathcal{O}_X$-模逆系统之间的映射。
性质 (1) 和 (2) 由构造成立。

\medskip\noindent
为证明性质 (3) 和 (4)，可以假设 $X = \Spec(A_0)$ 是仿射的，
且 $A_0$ 有对偶化复形。令 $I_0 \subset A_0$ 为与 $Y$ 对应的理想。
令 $A, I$ 为关于 $I_0$ 的进完备化，分别对应 $A_0, I_0$。
供后文使用，注意 $A$ 有对偶化复形
（《对偶化复形》，引理 \ref{dualizing-lemma-ubiquity-dualizing}）。
令 $M$ 为有限 $A$-模，它对应于 $(\mathcal{F}_n)$；参见
《概形上同调》，引理 \ref{coherent-lemma-inverse-systems-affine}。
于是 $\mathcal{F}_n$ 对应于 $M_n = M/I^nM$。回忆
$\mathcal{F}'_n$ 对应于商模 $M'_n = M_n / H^0_T(M_n)$。
还回忆，$M' = \lim M'_n$ 是 $M$ 对 $H^0_T(M)$ 取的商，
$H^0_T(M)$ 是 $I$-幂挠模，并且
$\{M'_n\}$ 与 $\{M'/I^nM'\}$ 作为 pro-系统同构。
最后，$\mathcal{F}''_n$ 对应于扩张
$$
0 \to M'_n \to M''_n \to H^1_{T'}(M'_n) \to 0
$$
；参见《局部上同调》，引理
\ref{local-cohomology-lemma-make-S2-along-T} 的证明。

\medskip\noindent
置 $s = 1$。我们断言 $A, I, T', M', s, d$ 满足情形
\ref{algebraization-situation-bootstrap} 的假设 (1)、(3)、(4)、(6)。事实上，
(1) 和 (3) 立即成立，(c) 蕴含 (4)，而 (d) 蕴含 (6)。
我们省略 (c) $\Rightarrow$ (4) 和 (d) $\Rightarrow$ (6) 的证明。

\medskip\noindent
由定理 \ref{algebraization-theorem-final-bootstrap} 可知，$\{H^1_{T'}(M'/I^nM')\}$
本质常值，其值为 $H^1_{T'}(M')$。由此推出
$\{H^1_{T'}(M'_n)\}$ 本质常值，其值为 $H^1_{T'}(M')$。
因此，上面所显示的短正合序列之极限是短正合序列
$$
0 \to M' \to \lim M''_n \to H^1_{T'}(M') \to 0
$$
置 $M'' = \lim M''_n$。由《局部上同调》，命题
\ref{local-cohomology-proposition-finiteness} 可知，
$H^1_{T'}(M')$ 以及因而 $M''$ 都是有限 $A$-模
（验证条件时，使用 $M'$ 的深度至少为 $1$，这是对
$T - T'$ 中各素理想而言）。
令 $(\mathcal{H}_n)$ 为 $\textit{Coh}(X, \mathcal{I})$ 中与有限
$A$-模 $M''$ 对应的对象。为完成证明，必须证明典范映射
$\{M''/I^nM''\} \to \{M''_n\}$ 是 pro-同构。
已知 $\{M'/I^nM'\}$ 与 $\{M'_n\}$ pro-同构；读者可验证
（此处省略），这等价于要求
$\{H^1_{T'}(M')/I^nH^1_{T'}(M')\} \to \{H^1_{T'}(M'_n)\}$
是 pro-同构。由于
$H^1_{T'}(M')/I^nH^1_{T'}(M') = H^1_{T'}(M')$
在 $n$ 充分大时成立，
而上面已经看到 $\{H^1_{T'}(M'_n)\}$
本质常值，其值为 $H^1_{T'}(M')$，故结论成立。
\end{proof}

\begin{lemma}
\label{algebraization-lemma-improvement-application}
在情形 \ref{algebraization-situation-algebraize} 中，假设 $A$ 有对偶化复形。
设 $d \geq \text{cd}(A, I)$。设 $(\mathcal{F}_n)$ 为
$\textit{Coh}(U, I\mathcal{O}_U)$ 的一个对象。假设
$(\mathcal{F}_n)$ 满足 $(2, 2 + d)$-不等式；参见定义
\ref{algebraization-definition-s-d-inequalities}。
于是存在典范映射 $(\mathcal{F}_n) \to (\mathcal{F}_n'')$，
它是凝聚 $\mathcal{O}_U$-模逆系统之间的映射，并满足下列性质：
\begin{enumerate}
\item $\text{depth}(\mathcal{F}''_{n, y}) + \delta^Y_Z(y) \geq 3$
对所有 $y \in U \cap Y$ 成立；
\item $(\mathcal{F}''_n)$ 作为 pro-系统同构于一个对象
$(\mathcal{H}_n)$，后者属于 $\textit{Coh}(U, I\mathcal{O}_U)$；
\item 所诱导的态射 $(\mathcal{F}_n) \to (\mathcal{H}_n)$ 在
$\textit{Coh}(U, I\mathcal{O}_U)$ 中的核与余核被
$I$ 的某个幂零化；
\item 模 $H^0(U, \mathcal{F}''_n)$ 和 $H^1(U, \mathcal{F}''_n)$
都是有限 $A$-模，对所有 $n$ 均成立。
\end{enumerate}
\end{lemma}

\begin{proof}
存在性以及性质 (1)、(2)、(3) 由引理
\ref{algebraization-lemma-improvement-formal-coherent-module-better} 得出：把该引理应用于
$U$、$U \cap Y$、
$T = \{y \in U \cap Y : \delta^Y_Z(y) \leq 2\}$、
$T' = \{y \in U \cap Y : \delta^Y_Z(y) \leq 1\}$ 和
$(\mathcal{F}_n)$ 即可。模 $H^0(U, \mathcal{F}''_n)$ 和
$H^1(U, \mathcal{F}''_n)$ 的有限性由《局部上同调》，引理
\ref{local-cohomology-lemma-finiteness-Rjstar}，以及函数
$\delta^Y_Z(-)$ 在引理 \ref{algebraization-lemma-discussion} 中证明的初等性质得出。
\end{proof}











\section{凝聚形式模的代数化，V}
\label{algebraization-section-algebraization-modules-conclusion}

\noindent
本节证明关于凝聚形式模代数化的最一般结果。先在理想的上同调维数为
$1$ 时证明，再把所得结果应用于一次吹起，以证明更一般的结论。

\begin{lemma}
\label{algebraization-lemma-cd-1-canonical}
在情形 \ref{algebraization-situation-algebraize} 中，设 $(\mathcal{F}_n)$ 为
$\textit{Coh}(U, I\mathcal{O}_U)$ 的一个对象。假设
\begin{enumerate}
\item $A$ 有对偶化复形，且 $\text{cd}(A, I) = 1$；
\item $(\mathcal{F}_n)$ pro-同构于逆系统
$(\mathcal{F}_n'')$，后者由凝聚 $\mathcal{O}_U$-模组成，并且
$\text{depth}(\mathcal{F}''_{n, y}) + \delta^Y_Z(y) \geq 3$
对所有 $y \in U \cap Y$ 成立。
\end{enumerate}
则 $(\mathcal{F}_n)$ 典范地延拓到 $X$；参见定义
\ref{algebraization-definition-canonically-algebraizable}。
\end{lemma}

\begin{proof}
我们检验引理 \ref{algebraization-lemma-when-done} 的假设 (a)、(b) 和 (c)。
开始之前，先指出模
$H^0(U, \mathcal{F}''_n)$ 和 $H^1(U, \mathcal{F}''_n)$
都是有限 $A$-模，对所有 $n$ 均如此；这是由《局部上同调》，引理
\ref{local-cohomology-lemma-finiteness-Rjstar} 得出的。

\medskip\noindent
注意，对每个 $p \geq 0$，由引理 \ref{algebraization-lemma-topology-I-adic}，
$\lim H^p(U, \mathcal{F}_n)$ 上的极限拓扑是 $I$-进拓扑。
特别地，假设 (b) 成立。

\medskip\noindent
已知 $M = \lim H^0(U, \mathcal{F}_n)$ 是 $A$-模，
其极限拓扑为 $I$-进拓扑。因此，给定 $n$，模
$M/I^nM$ 是 $H^0(U, \mathcal{F}_N)$ 的子商，其中 $N \gg n$。
逆系统 $\{H^0(U, \mathcal{F}_N)\}$ pro-同构于有限 $A$-模的逆系统，
即 $\{H^0(U, \mathcal{F}''_N)\}$，故 $M/I^nM$ 有限。
由此 $M$ 有限；参见《代数》，引理
\ref{algebra-lemma-finite-over-complete-ring}。
特别地，假设 (c) 成立。

\medskip\noindent
对每个 $n \geq 0$，记 $Ob_n = \lim_N H^1(U, I^n\mathcal{F}_N)$。
一个特例是 $Ob = Ob_0 = \lim_N H^1(U, \mathcal{F}_N)$。
与上一段完全相同的论证表明，$Ob$ 是有限 $A$-模。
（事实上，还知道 $Ob/I Ob$ 被 $\mathfrak a$ 的某个幂零化，
但这似乎不太容易使用。）

\medskip\noindent
置 $\mathcal{F} = \lim \mathcal{F}_n$，选取生成元
$f_1, \ldots, f_r$，它们生成 $I$；再选取 $c \geq 1$，并像引理
\ref{algebraization-lemma-cd-is-one-for-system} 中那样选取 $\Phi_\mathcal{F}$。
以下将不再另行说明地使用引理 \ref{algebraization-lemma-properties-system} 的结论。
特别地，对每个 $n \geq 1$ 都有映射
$$
\delta_n :
H^0(U, \mathcal{F}_n)
\longrightarrow
H^1(U, I^n\mathcal{F})
\longrightarrow
Ob_n
$$
第一个映射来自短正合序列
$0 \to I^n\mathcal{F} \to \mathcal{F} \to \mathcal{F}_n \to 0$
，第二个映射来自 $I^n\mathcal{F} = \lim I^n\mathcal{F}_N$。
稍后将使用如下事实：若 $\delta_n(s) = 0$，其中
$s \in H^0(U, \mathcal{F}_n)$，则对每个 $n' \geq n$，
都可找到 $s' \in H^0(U, \mathcal{F}_{n'})$ 映到 $s$。
注意有交换图
$$
\xymatrix{
H^0(U, \mathcal{F}_{nc}) \ar[r] \ar[dd] &
H^1(U, I^{nc}\mathcal{F}) \ar[dd] \ar[rd]^{\Phi_\mathcal{F}} \\
& &
\bigoplus_{e_1 + \ldots + e_r = n}
H^1(U, \mathcal{F}) \cdot T_1^{e_1} \ldots T_r^{e_r} \ar[ld] \\
H^0(U, \mathcal{F}_n) \ar[r] &
H^1(U, I^n\mathcal{F})
}
$$
由此可知，障碍映射
$H^0(U, \mathcal{F}_n) \to Ob_n$
把映射
$H^0(U, \mathcal{F}_{nc}) \to H^0(U, \mathcal{F}_n)$
的像送入子模
$$
Ob'_n =
\Im\left(
\bigoplus\nolimits_{e_1 + \ldots + e_r = n}
Ob \cdot T_1^{e_1} \ldots T_r^{e_r} \to Ob_n
\right)
$$
；在直和项 $Ob \cdot T_1^{e_1} \ldots T_r^{e_r}$ 上，
这里使用的上同调映射来自对 $I$ 的各次幂取模后的约化，
所约化的乘法映射为
$f_1^{e_1} \ldots f_r^{e_r} : \mathcal{F} \to I^n\mathcal{F}$。
由构造，
$$
\bigoplus\nolimits_{n \geq 0} Ob'_n
$$
是 Rees 代数 $\bigoplus_{n \geq 0} I^n$ 上的有限分次模。
对每个 $n$，置
$$
M_n = \{s \in H^0(U, \mathcal{F}_n) \mid \delta_n(s) \in Ob'_n\}
$$
注意，$\{M_n\}$ 是逆系统，并且
$f_j : \mathcal{F}_n \to \mathcal{F}_{n + 1}$ 在整体截面上把
$M_n$ 映入 $M_{n + 1}$。
仿照《上同调》，引理 \ref{cohomology-lemma-ML-general} 的证明作完全相同的论证，
可知 $\{M_n\}$ 满足 ML。具体而言，由于 Rees 代数是诺特的，
可以为分次子模选取有限多个形如 $\delta_{n_j}(z_j)$ 的齐次生成元，
其中 $z_j \in M_{n_j}$；该分次子模为
$\bigoplus_{n \geq 0} \Im(M_n \to Ob'_n)$.
若 $k = \max(n_j)$，则对 $n \geq k$ 和任意 $z \in M_n$，
可找到 $a_j \in I^{n - n_j}$，使得
$z - \sum a_j z_j$ 位于 $\delta_n$ 的核中，因而位于
$M_{n'}$ 的像中，对所有 $n' \geq n$ 均成立
（因为 $\delta_n$ 的消失意味着，可以把
$z - \sum a_j z_j$ 提升为元素
$z' \in H^0(U, \mathcal{F}_{n'c})$，对所有 $n' \ge n$ 均可如此；
而后 $z'$ 在 $H^0(U, \mathcal{F}_{n'})$ 中的像由上文可知属于
$M_{n'}$）。因此
$\Im(M_n \to M_{n - k}) = \Im(M_{n'} \to M_{n - k})$
对所有 $n' \geq n$ 均成立。

\medskip\noindent
选取 $n$。由刚刚证明的 $\{M_n\}$ 的 Mittag--Leffler 性质，
可以找到 $n' \geq n$，使得 $M_{n'} \to M_n$ 的像
等于 $M' \to M_n$ 的像。由上文可知，
$M' \to M_n$ 的像包含
$H^0(U, \mathcal{F}_{n'c}) \to H^0(U, \mathcal{F}_n)$.
的像。因此，$\{M_n\}$ 与 $\{H^0(U, \mathcal{F}_n)\}$
pro-同构。于是 $\{H^0(U, \mathcal{F}_n)\}$ 满足 ML，
最终得到假设 (a) 成立。证明完成。
\end{proof}

\begin{proposition}[上同调维数为 1 时的代数化]
\label{algebraization-proposition-cd-1}
\begin{reference}
本结果的局部情形见 \cite[IV Corollaire 2.9]{MRaynaud-book}。
\end{reference}
在情形 \ref{algebraization-situation-algebraize} 中，设 $(\mathcal{F}_n)$ 为
$\textit{Coh}(U, I\mathcal{O}_U)$ 的一个对象。假设
\begin{enumerate}
\item $A$ 有对偶化复形，且 $\text{cd}(A, I) = 1$；
\item $(\mathcal{F}_n)$ 满足 $(2, 3)$-不等式；参见定义
\ref{algebraization-definition-s-d-inequalities}。
\end{enumerate}
则 $(\mathcal{F}_n)$ 延拓到 $X$。特别地，若 $A$ 是
$I$-进完备的，则 $(\mathcal{F}_n)$ 是某个凝聚
$\mathcal{O}_U$-模的完备化。
\end{proposition}

\begin{proof}
由引理 \ref{algebraization-lemma-map-kernel-cokernel-on-closed}，可以把
$(\mathcal{F}_n)$ 替换为对象 $(\mathcal{H}_n)$，它属于
$\textit{Coh}(U, I\mathcal{O}_U)$，并且是在引理
\ref{algebraization-lemma-improvement-application} 中找到的。因此，可以假设 $(\mathcal{F}_n)$ pro-同构于
逆系统 $(\mathcal{F}_n'')$，后者具有引理
\ref{algebraization-lemma-improvement-application} 中所述的性质。
引理 \ref{algebraization-lemma-cd-1-canonical} 已经证明
$(\mathcal{F}_n)$ 典范地延拓到 $X$。
最后一句由引理 \ref{algebraization-lemma-canonically-algebraizable} 得出。
\end{proof}

\begin{lemma}
\label{algebraization-lemma-blowup}
在情形 \ref{algebraization-situation-algebraize} 中，设 $(\mathcal{F}_n)$ 为
$\textit{Coh}(U, I\mathcal{O}_U)$ 的一个对象。假设
\begin{enumerate}
\item $A$ 有对偶化复形；
\item 吹起 $b : X' \to X$（沿着 $I$）的所有纤维之维数都
$\leq d - 1$；
\item 下列条件之一成立：
\begin{enumerate}
\item $(\mathcal{F}_n)$ 满足 $(d + 1, d + 2)$-不等式
（定义 \ref{algebraization-definition-s-d-inequalities}）；或者
\item 对 $y \in U \cap Y$ 以及素理想
$\mathfrak p \subset \mathcal{O}_{X, y}^\wedge$，若
$\mathfrak p \not \in V(I\mathcal{O}_{X, y}^\wedge)$
，则
$$
\text{depth}((\mathcal{F}^\wedge_y)_\mathfrak p) +
\dim(\mathcal{O}_{X, y}^\wedge/\mathfrak p) + \delta^Y_Z(y) > d + 2
$$
\end{enumerate}
\end{enumerate}
则 $(\mathcal{F}_n)$ 延拓到 $X$。
\end{lemma}

\begin{proof}
令 $Y' \subset X'$ 为例外除子。
令 $Z' \subset Y'$ 为 $Z \subset Y$ 的逆像。
则 $U' = X' \setminus Z'$ 是 $U$ 的逆像。
使用 (\ref{algebraization-equation-delta-Z}) 中的 $\delta^{Y'}_{Z'}$，置
$$
T' = \{y' \in Y' \mid \delta^{Y'}_{Z'}(y') = 1\}
\subset
T = \{y' \in Y' \mid \delta^{Y'}_{Z'}(y') = 1\text{ 或 }2\}
$$
由引理 \ref{algebraization-lemma-discussion} 的 (1) 和 (2)，它们是
$U' \cap Y' = Y' \setminus Z'$ 中对特化稳定的子集。
考虑对象 $(b|_{U'}^*\mathcal{F}_n)$；它属于
$\textit{Coh}(U', I\mathcal{O}_{U'})$。参见《概形上同调》，引理
\ref{coherent-lemma-inverse-systems-pullback}。
对 $y' \in U' \cap Y'$，记
$$
\mathcal{F}_{y'}^\wedge = \lim (b|_{U'}^*\mathcal{F}_n)_{y'}
$$
为该拉回在 $y'$ 处的“茎”。我们断言，引理
\ref{algebraization-lemma-improvement-formal-coherent-module-better} 的条件
(a)、(b)、(c)、(d)、(e) 对对象
$(b|_{U'}^*\mathcal{F}_n)$ 成立；该对象定义在 $U'$ 上，其中把 $d$ 换成 $1$，
并取子集 $T' \subset T \subset U' \cap Y'$。
条件 (a) 成立，因为 $Y'$ 是有效 Cartier 除子，故局部由
$1$ 个方程截出。条件 (e) 由引理 \ref{algebraization-lemma-discussion} 的 (7) 成立。
为证明 (b)、(c)、(d)，需要作一些准备。

\medskip\noindent
设 $y' \in U' \cap Y'$，并设
$\mathfrak p' \subset \mathcal{O}_{X', y'}^\wedge$
为不包含在 $V(I\mathcal{O}_{X', y'}^\wedge)$ 中的素理想。
记 $y = b(y') \in U \cap Y$。选取 $f \in I$，使得
$y'$ 包含在仿射吹起代数 $A[\frac{I}{f}]$ 的谱中；参见《除子》，引理
\ref{divisors-lemma-blowing-up-affine}。
对任意 $A$-代数 $B$，记 $B' = B[\frac{IB}{f}]$ 为相应的仿射吹起代数。
把 $I$-进完备化记作 ${\ }^\wedge$。由 $f$ 的选取，得到环映射
$(\mathcal{O}_{X, y}^\wedge)' \to \mathcal{O}_{X', y'}^\wedge$.
若令 $\mathfrak q' \subset (\mathcal{O}_{X, y}^\wedge)'$
为 $\mathfrak m_{y'}^\wedge$ 的逆像，则
$((\mathcal{O}_{X, y}^\wedge)'_{\mathfrak q'})^\wedge =
\mathcal{O}_{X', y'}^\wedge$.
令 $\mathfrak p \subset \mathcal{O}_{X, y}^\wedge$ 为相应的素理想。
此时有交换图
$$
\xymatrix{
\mathcal{O}_{X, y}^\wedge \ar[d] \ar[r] &
(\mathcal{O}_{X, y}^\wedge)' \ar[d]_\alpha \ar[r] &
(\mathcal{O}_{X, y}^\wedge)'_{\mathfrak q'} \ar[d] \ar[r]_\beta &
\mathcal{O}_{X', y'}^\wedge \ar[d] \\
\mathcal{O}_{X, y}^\wedge/\mathfrak p \ar[r] &
(\mathcal{O}_{X, y}^\wedge/\mathfrak p)' \ar[r] &
(\mathcal{O}_{X, y}^\wedge/\mathfrak p)'_{\mathfrak q'} \ar[r]^\gamma &
((\mathcal{O}_{X, y}^\wedge/\mathfrak p)'_{\mathfrak q'})^\wedge \ar[d] \\
 & & &
\mathcal{O}_{X', y'}^\wedge/\mathfrak p'
}
$$
，其中竖直箭头都是满射。由《代数续篇》，引理
\ref{more-algebra-lemma-completion-dimension} 和维数公式
（《代数》，引理 \ref{algebra-lemma-dimension-formula}），有
$$
\dim(((\mathcal{O}_{X, y}^\wedge/\mathfrak p)'_{\mathfrak q'})^\wedge) =
\dim((\mathcal{O}_{X, y}^\wedge/\mathfrak p)'_{\mathfrak q'}) =
\dim(\mathcal{O}_{X, y}^\wedge/\mathfrak p)
- \text{trdeg}(\kappa(y')/\kappa(y))
$$
沿着拉回、茎、局部化和完备化的定义逐一追踪，可得
$$
(\mathcal{F}_y^\wedge)_{\mathfrak p}
\otimes_{(\mathcal{O}_{X, y}^\wedge)_\mathfrak p}
(\mathcal{O}_{X', y'}^\wedge)_{\mathfrak p'}
=
(\mathcal{F}_{y'}^\wedge)_{\mathfrak p'}
$$
细节从略。图中的环映射 $\beta$ 和 $\gamma$ 都是平坦的，
其纤维为 Gorenstein（从而 Cohen--Macaulay）环，因为这些映射是
对具有对偶化复形的环取完备化所得。参见《对偶化复形》，引理
\ref{dualizing-lemma-formal-fibres-gorenstein} 和
\ref{dualizing-lemma-dualizing-gorenstein-formal-fibres}，以及
《代数续篇》，第 \ref{more-algebra-section-properties-formal-fibres} 节中的讨论。
注意
$(\mathcal{O}_{X, y}^\wedge)_\mathfrak p =
(\mathcal{O}_{X, y}^\wedge)'_{\tilde{\mathfrak p}}$
，其中 $\tilde{\mathfrak p}$ 是图中 $\alpha$ 的核。另一方面，
$(\mathcal{O}_{X, y}^\wedge)'_{\tilde{\mathfrak p}}
\to (\mathcal{O}_{X', y'}^\wedge)_{\mathfrak p'}$
由上文可知是纤维为 CM 环的平坦映射。因此
$(\mathcal{O}_{X, y}^\wedge)_\mathfrak p  \to
(\mathcal{O}_{X', y'}^\wedge)_{\mathfrak p'}$ 也是纤维为 CM 环的平坦映射。
使用《代数》，引理 \ref{algebra-lemma-apply-grothendieck-module}，可得
$$
\text{depth}((\mathcal{F}_{y'}^\wedge)_{\mathfrak p'}) =
\text{depth}((\mathcal{F}_y^\wedge)_{\mathfrak p}) +
\dim(F_\mathfrak r)
$$
，其中 $F$ 是
$(\mathcal{O}_{X, y}^\wedge/\mathfrak p)'_{\mathfrak q'}$
的泛形式纤维，而 $\mathfrak r$ 是与 $\mathfrak p'$ 对应的素理想。
由于 $(\mathcal{O}_{X, y}^\wedge/\mathfrak p)'_{\mathfrak q'}$
是泛链状局部整环，由 Ratliff 定理，其 $I$-进完备化
是等维且（泛）链状的
（《代数续篇》，命题 \ref{more-algebra-proposition-ratliff}）。
从而
$$
\dim(((\mathcal{O}_{X, y}^\wedge/\mathfrak p)'_{\mathfrak q'})^\wedge) =
\dim(F_\mathfrak r) + \dim(\mathcal{O}_{X', y'}^\wedge/\mathfrak p')
$$
结合引理 \ref{algebraization-lemma-change-distance-function}，得到
\begin{equation}
\label{algebraization-equation-one}
\begin{aligned}
&
\text{depth}((\mathcal{F}_{y'}^\wedge)_{\mathfrak p'}) +
\delta^{Y'}_{Z'}(y') \\
& =
\text{depth}((\mathcal{F}_y^\wedge)_{\mathfrak p}) +
\dim(F_\mathfrak r) + \delta^{Y'}_{Z'}(y') \\
& \geq
\text{depth}((\mathcal{F}_y^\wedge)_{\mathfrak p}) + \delta^Y_Z(y) +
\dim(F_\mathfrak r) + \text{trdeg}(\kappa(y')/\kappa(y)) - (d - 1) \\
& =
\text{depth}((\mathcal{F}_y^\wedge)_{\mathfrak p}) + \delta^Y_Z(y) - (d - 1)
+ \dim(\mathcal{O}_{X, y}^\wedge/\mathfrak p) -
\dim(\mathcal{O}_{X', y'}^\wedge/\mathfrak p')
\end{aligned}
\end{equation}
请注意
$\dim(\mathcal{O}_{X, y}^\wedge/\mathfrak p) \geq
\dim(\mathcal{O}_{X', y'}^\wedge/\mathfrak p')$。将其改写，得到
\begin{equation}
\label{algebraization-equation-two}
\begin{aligned}
&
\text{depth}((\mathcal{F}_{y'}^\wedge)_{\mathfrak p'}) +
\dim(\mathcal{O}_{X', y'}^\wedge/\mathfrak p') +
\delta^{Y'}_{Z'}(y') \\
& \geq
\text{depth}((\mathcal{F}_y^\wedge)_{\mathfrak p}) +
\dim(\mathcal{O}_{X, y}^\wedge/\mathfrak p) +
\delta^Y_Z(y) - (d - 1)
\end{aligned}
\end{equation}
这个不等式将使我们能够检验余下条件。

\medskip\noindent
检验引理 \ref{algebraization-lemma-improvement-formal-coherent-module-better} 的条件 (b) 和 (d)。
假设
$V(\mathfrak p') \cap V(I\mathcal{O}_{X', y'}^\wedge) =
\{\mathfrak m_{y'}^\wedge\}$.
这蕴含 $\dim(\mathcal{O}_{X', y'}^\wedge/\mathfrak p') = 1$，
因为 $Z'$ 是有效 Cartier 除子。
(b) 与 (d) 合起来等价于
$$
\text{depth}((\mathcal{F}_{y'}^\wedge)_{\mathfrak p'}) + \delta^{Y'}_{Z'}(y')
> 2
$$
若 $(\mathcal{F}_n)$ 满足 (3)(b) 中的不等式，
应用 (\ref{algebraization-equation-two}) 便立即得到这个结论。
若 $(\mathcal{F}_n)$ 满足 (3)(a)，即
$(d + 1, d + 2)$-不等式，则无论如何都有
$$
\text{depth}((\mathcal{F}_y^\wedge)_{\mathfrak p}) + \delta^Y_Z(y)
\geq d + 1
\quad\text{或}\quad
\text{depth}((\mathcal{F}_y^\wedge)_{\mathfrak p}) +
\dim(\mathcal{O}_{X, y}^\wedge/\mathfrak p) +
\delta^Y_Z(y) > d + 2
$$
考察上面的 (\ref{algebraization-equation-one}) 和 (\ref{algebraization-equation-two})，便得到所需结论，
唯一可能的例外是
$\dim(\mathcal{O}_{X, y}^\wedge/\mathfrak p) = 1$.
然而，若 $\dim(\mathcal{O}_{X, y}^\wedge/\mathfrak p) = 1$，则
$V(\mathfrak p) \cap V(I\mathcal{O}_{X, y}^\wedge) = \{\mathfrak m_y^\wedge\}$
，而且事实上
$$
\text{depth}((\mathcal{F}_y^\wedge)_{\mathfrak p}) + \delta^Y_Z(y) > d + 1
$$
，因为 $(\mathcal{F}_n)$ 满足 $(d + 1, d + 2)$-不等式；故仍得到结论。

\medskip\noindent
检验引理 \ref{algebraization-lemma-improvement-formal-coherent-module-better} 的条件 (c)。
假设
$V(\mathfrak p') \cap V(I\mathcal{O}_{X', y'}^\wedge) \not =
\{\mathfrak m_{y'}^\wedge\}$。则条件 (c) 等价于
$$
\text{depth}((\mathcal{F}_{y'}^\wedge)_{\mathfrak p'}) + \delta^{Y'}_{Z'}(y')
\geq 2
\quad\text{或}\quad
\text{depth}((\mathcal{F}_{y'}^\wedge)_{\mathfrak p'}) +
\dim(\mathcal{O}_{X', y'}^\wedge/\mathfrak p') +
\delta^{Y'}_{Z'}(y') > 3
$$
若 $(\mathcal{F}_n)$ 满足 (3)(b) 中的不等式，
应用 (\ref{algebraization-equation-two}) 可知所显示的两个不等式中的第二个成立。
若 $(\mathcal{F}_n)$ 满足 (3)(a)，即 $(d + 1, d + 2)$-不等式，
则由 (\ref{algebraization-equation-one}) 和 (\ref{algebraization-equation-two}) 立即得到结论。
这就证明了我们的断言。

\medskip\noindent
$(b|_{U'}^*\mathcal{F}_n) \to (\mathcal{F}_n'')$
以及 $(\mathcal{H}_n)$，后者在 $\textit{Coh}(U', I\mathcal{O}_{U'})$ 中；
像引理 \ref{algebraization-lemma-improvement-formal-coherent-module-better} 中那样选取它们。
对任意仿射开集 $W \subset X'$，注意到
$\delta^{W \cap Y'}_{W \cap Z'}(y') \geq \delta^{Y'}_{Z'}(y')$；
这是由引理 \ref{algebraization-lemma-discussion} 的 (6) 得出的。因此，
$(\mathcal{H}_n|_W)$ 满足引理 \ref{algebraization-lemma-cd-1-canonical} 的假设。
故 $(\mathcal{H}_n|_W)$ 典范地延拓到 $W$。
令 $(\mathcal{G}_{W, n})$ 为 $\textit{Coh}(W, I\mathcal{O}_W)$ 中
如引理 \ref{algebraization-lemma-canonically-algebraizable} 所述的典范延拓。
由引理 \ref{algebraization-lemma-canonically-extend-base-change}，可知对
$W' \subset W$ 存在唯一同构
$$
(\mathcal{G}_{W, n}|_{W'}) \longrightarrow
(\mathcal{G}_{W', n})
$$
，它与给定同构
$(\mathcal{G}_{W, n}|_{W \cap U}) \cong (\mathcal{H}_n|_{W \cap U})$.
相容。由此得到一个对象
$(\mathcal{G}_n)$，它属于 $\textit{Coh}(X', I\mathcal{O}_{X'})$，
其在 $U$ 上的限制同构于 $(\mathcal{H}_n)$。

\medskip\noindent
若 $A$ 是 $I$-进完备的，则可如下完成证明。
由 Grothendieck 存在性定理
（《概形上同调》，引理 \ref{coherent-lemma-existence-projective}），
可知 $(\mathcal{G}_n)$ 是某个凝聚 $\mathcal{O}_{X'}$-模的完备化。
再由《概形上同调》，引理 \ref{coherent-lemma-existence-easy}，
可知 $(b|_{U'}^*\mathcal{F}_n)$ 是某个凝聚
$\mathcal{O}_{U'}$-模 $\mathcal{F}'$ 的完备化。
由《概形上同调》，引理 \ref{coherent-lemma-inverse-systems-push-pull}，
存在映射
$$
(\mathcal{F}_n) \longrightarrow ((b|_{U'})_*\mathcal{F}')^\wedge
$$
，其核与余核被 $I$ 的某个幂零化。
最后应用引理 \ref{algebraization-lemma-map-kernel-cokernel-on-closed} 即得结论。

\medskip\noindent
若 $A$ 不完备，则在开始证明之前可以把 $A$ 换成其完备化；
参见引理 \ref{algebraization-lemma-algebraizable}。
完备化后各假设仍成立：条件 (3) 是立即的；条件 (1) 由
《对偶化复形》，引理 \ref{dualizing-lemma-ubiquity-dualizing} 得出；
条件 (2) 由《除子》，引理
\ref{divisors-lemma-flat-base-change-blowing-up} 得出。
因此，完备情形蕴含一般情形。
\end{proof}

\begin{proposition}[生成元较少的理想之代数化]
\label{algebraization-proposition-d-generators}
在情形 \ref{algebraization-situation-algebraize} 中，设 $(\mathcal{F}_n)$ 为
$\textit{Coh}(U, I\mathcal{O}_U)$ 的一个对象。假设
\begin{enumerate}
\item $A$ 有对偶化复形；
\item $V(I) = V(f_1, \ldots, f_d)$，其中某个 $d \geq 1$，并且
$f_1, \ldots, f_d \in A$；
\item 下列条件之一成立：
\begin{enumerate}
\item $(\mathcal{F}_n)$ 满足 $(d + 1, d + 2)$-不等式
（定义 \ref{algebraization-definition-s-d-inequalities}）；或者
\item 对 $y \in U \cap Y$ 以及素理想
$\mathfrak p \subset \mathcal{O}_{X, y}^\wedge$，若
$\mathfrak p \not \in V(I\mathcal{O}_{X, y}^\wedge)$
，则
$$
\text{depth}((\mathcal{F}^\wedge_y)_\mathfrak p) +
\dim(\mathcal{O}_{X, y}^\wedge/\mathfrak p) + \delta^Y_Z(y) > d + 2
$$
\end{enumerate}
\end{enumerate}
则 $(\mathcal{F}_n)$ 延拓到 $X$。特别地，若 $A$ 是
$I$-进完备的，则 $(\mathcal{F}_n)$ 是某个凝聚
$\mathcal{O}_U$-模的完备化。
\end{proposition}

\begin{proof}
可以假设 $I = (f_1, \ldots, f_d)$；参见《概形上同调》，引理
\ref{coherent-lemma-inverse-systems-ideals-equivalence}。
于是，$X$ 沿 $I$ 的吹起之所有纤维的维数至多为 $d - 1$。
故由引理 \ref{algebraization-lemma-blowup} 得到延拓。
最后一句由引理 \ref{algebraization-lemma-essential-image-completion} 得出。
\end{proof}

\noindent
请把下一个引理与注
\ref{algebraization-remark-interesting-case-variant}、
\ref{algebraization-remark-interesting-case},
\ref{algebraization-remark-interesting-case-bis} 和
\ref{algebraization-remark-interesting-case-ter} 比较。

\begin{lemma}
\label{algebraization-lemma-interesting-case-final}
在情形 \ref{algebraization-situation-algebraize} 中，设 $(\mathcal{F}_n)$ 为
$\textit{Coh}(U, I\mathcal{O}_U)$ 的一个对象。假设
\begin{enumerate}
\item $A$ 是有对偶化复形的局部环；
\item $X$ 的所有不可约分支具有相同维数；
\item 概形 $X \setminus Y$ 是 Cohen--Macaulay 的；
\item $I$ 由 $d$ 个元素生成；
\item $\dim(X) - \dim(Z) > d + 2$；并且
\item 对 $y \in U \cap Y$，模 $\mathcal{F}_y^\wedge$ 在
$V(I\mathcal{O}_{X, y}^\wedge)$ 外有限局部自由；例如，若
$\mathcal{F}_n$ 是有限局部自由
$\mathcal{O}_U/I^n\mathcal{O}_U$-模，则此条件成立。
\end{enumerate}
则 $(\mathcal{F}_n)$ 延拓到 $X$。特别地，若 $A$ 是 $I$-进
完备的，则 $(\mathcal{F}_n)$ 是某个凝聚 $\mathcal{O}_U$-模的完备化。
\end{lemma}

\begin{proof}
我们证明命题 \ref{algebraization-proposition-d-generators} 的假设 (1)、(2)、(3)(b)
均成立。其中 (1) 和 (2) 是显然的。

\medskip\noindent
设 $y \in U \cap Y$，并设 $\mathfrak p$ 为素理想
$\mathfrak p \subset \mathcal{O}_{X, y}^\wedge$，且
$\mathfrak p \not \in V(I\mathcal{O}_{X, y}^\wedge)$.
最后一个条件表明
$\text{depth}((\mathcal{F}_y^\wedge)_\mathfrak p) =
\text{depth}((\mathcal{O}_{X, y}^\wedge)_\mathfrak p)$.
由于 $X \setminus Y$ 是 Cohen--Macaulay 的，可知
$(\mathcal{O}_{X, y}^\wedge)_\mathfrak p$ 是 Cohen--Macaulay 的。
因此
\begin{align*}
& \text{depth}((\mathcal{F}^\wedge_y)_\mathfrak p) +
\dim(\mathcal{O}_{X, y}^\wedge/\mathfrak p) + \delta^Y_Z(y) \\
& =
\dim((\mathcal{O}_{X, y}^\wedge)_\mathfrak p) +
\dim(\mathcal{O}_{X, y}^\wedge/\mathfrak p) + \delta^Y_Z(y) \\
& =
\dim(\mathcal{O}_{X, y}^\wedge) + \delta^Y_Z(y)
\end{align*}
最后一个等号成立，是因为第二个条件表明
$\mathcal{O}_{X, y}$ 等维。
令 $\delta(y) = \dim(\overline{\{y\}})$。由于 $A$ 是链状局部环，
这是维数函数。由引理 \ref{algebraization-lemma-discussion}，有
$\delta^Y_Z(y) \geq \delta(y) - \dim(Z)$。由于 $X$ 等维，得到
$$
\dim(\mathcal{O}_{X, y}^\wedge) + \delta^Y_Z(y)
\geq \dim(\mathcal{O}_{X, y}^\wedge) + \delta(y) - \dim(Z)
= \dim(X) - \dim(Z)
$$
这就得到所需不等式，从而结论成立。
\end{proof}

\begin{remark}
\label{algebraization-remark-question}
我们无法证明或否证命题 \ref{algebraization-proposition-d-generators} 的如下类似命题：
把 $I$ 有 $d$ 个生成元这一假设替换为
$\text{cd}(A, I) \leq d$。
如果您知道证明或有反例，请发送电子邮件至
\href{mailto:stacks.project@gmail.com}{stacks.project@gmail.com}.
另一个显然的问题是，命题 \ref{algebraization-proposition-d-generators} 中的条件
在何种程度上是必要的。
\end{remark}




\section{凝聚形式模的代数化，VI}
\label{algebraization-section-algebraization-modules-yet-more}

\noindent
本节再给出几个较易证明的情形。

\begin{proposition}
\label{algebraization-proposition-algebraization-regular-sequence}
在情形 \ref{algebraization-situation-algebraize} 中，令
$(\mathcal{F}_n)$ 为 $\textit{Coh}(U, I\mathcal{O}_U)$ 的一个对象。
假设
\begin{enumerate}
\item 存在 $f_1, \ldots, f_d \in I$，使得对
$y \in U \cap Y$，理想 $I\mathcal{O}_{X, y}$
由 $f_1, \ldots, f_d$ 生成，并且
$f_1, \ldots, f_d$ 构成一个 $\mathcal{F}_y^\wedge$-正则序列；
\item $H^0(U, \mathcal{F}_1)$ 和 $H^1(U, \mathcal{F}_1)$
是有限 $A$-模。
\end{enumerate}
那么 $(\mathcal{F}_n)$ 可典范地延拓到 $X$。特别地，若 $A$
完备，则 $(\mathcal{F}_n)$ 是某个凝聚
$\mathcal{O}_U$-模的完备化。
\end{proposition}

\begin{proof}
我们通过验证引理 \ref{algebraization-lemma-when-done} 的假设 (a)、(b) 和 (c)
来证明这一点。对每个 $n$，都有短正合列
$$
0 \to I^n\mathcal{F}_{n + 1} \to \mathcal{F}_{n + 1} \to \mathcal{F}_n \to 0
$$
由于 $f_1, \ldots, f_d$ 在每个“茎”
$\mathcal{F}_y^\wedge$ 上构成正则序列（因而是拟正则序列，见
代数，引理 \ref{algebra-lemma-regular-quasi-regular}），并且我们有
$I\mathcal{F}_n = (f_1, \ldots, f_d)\mathcal{F}_n$ 对所有 $n$ 成立，
逐茎检验便得到
$$
I^n\mathcal{F}_{n + 1} =
\bigoplus\nolimits_{e_1 + \ldots + e_d = n} \mathcal{F}_1 \cdot
f_1^{e_1} \ldots f_d^{e_d}
$$
利用 $H^0(U, \mathcal{F}_1)$ 的有限性假设并作归纳，我们先得出
$M_n = H^0(U, \mathcal{F}_n)$ 是有限 $A$-模，且这对所有 $n$ 成立。
由此可见引理 \ref{algebraization-lemma-when-done} 的条件 (c) 成立。
我们还看到
$$
\bigoplus\nolimits_{n \geq 0} H^1(U, I^n\mathcal{F}_{n + 1})
$$
是有限分次 $R = \bigoplus I^n/I^{n +1}$-模。
由上同调，引理 \ref{cohomology-lemma-ML-general}，
可知引理 \ref{algebraization-lemma-when-done} 的条件 (a) 成立。最后，引理
\ref{algebraization-lemma-when-done} 的条件 (b) 也成立，因为
$\bigoplus H^0(U, I^n\mathcal{F}_{n + 1})$ 是有限分次 $R$-模，
而且可以应用上同调，引理
\ref{cohomology-lemma-topology-I-adic-general}。
\end{proof}

\begin{remark}
\label{algebraization-remark-interesting-case-ter}
在命题 \ref{algebraization-proposition-algebraization-regular-sequence} 的情形中，
若假设 $A$ 有对偶化复形，则
$H^0(U, \mathcal{F}_1)$ 和
$H^1(U, \mathcal{F}_1)$ 有限这一条件等价于
$$
\text{depth}(\mathcal{F}_{1, y}) +
\dim(\mathcal{O}_{\overline{\{y\}}, z}) > 2
$$
对所有 $y \in U \cap Y$ 以及 $z \in Z \cap \overline{\{y\}}$ 均成立。
见局部上同调，引理 \ref{local-cohomology-lemma-finiteness-Rjstar}。
例如，若 $\mathcal{F}_1$ 是有限局部自由
$\mathcal{O}_{U \cap Y}$-模，$Y$ 满足 $(S_2)$，并且
$\text{codim}(Z', Y') \geq 3$ 对如下每一对不可约分支都成立：
$Y'$ 是 $Y$ 的不可约分支，$Z'$ 是 $Z$ 的不可约分支，且
$Z' \subset Y'$，则上述条件成立。
\end{remark}

\begin{proposition}
\label{algebraization-proposition-algebraization-flat}
在情形 \ref{algebraization-situation-algebraize} 中，令
$(\mathcal{F}_n)$ 为 $\textit{Coh}(U, I\mathcal{O}_U)$ 的一个对象。
假设存在一个 Noether 局部环 $(R, \mathfrak m)$ 以及环同态
$R \to A$，使得
\begin{enumerate}
\item $I = \mathfrak m A$,
\item 对 $y \in U \cap Y$，茎 $\mathcal{F}_y^\wedge$ 是 $R$-平坦的；
\item $H^0(U, \mathcal{F}_1)$ 和 $H^1(U, \mathcal{F}_1)$ 是有限
$A$-模。
\end{enumerate}
那么 $(\mathcal{F}_n)$ 可典范地延拓到 $X$。特别地，若 $A$
完备，则 $(\mathcal{F}_n)$ 是某个凝聚
$\mathcal{O}_U$-模的完备化。
\end{proposition}

\begin{proof}
证明与命题 \ref{algebraization-proposition-algebraization-regular-sequence} 的证明
完全相同。具体而言，若 $\kappa = R/\mathfrak m$，则对 $n \geq 0$
有同构
$$
I^n \mathcal{F}_{n + 1} \cong
\mathcal{F}_1 \otimes_\kappa \mathfrak m^n/\mathfrak m^{n + 1}
$$
而右端是 $\mathcal{F}_1$ 的有限个副本的直和。这一点可以逐茎检验。
其余部分完全相同。
\end{proof}

\begin{remark}
\label{algebraization-remark-interesting-case-quater}
命题 \ref{algebraization-proposition-algebraization-flat} 是
\cite[Theorem 2.10 (i)]{Baranovsky} 的一个局部版本。从局部结果推出
整体结果很直接；下面简述论证。设 $(R, \mathfrak m)$ 是完备
Noether 局部环，且 $X \to \Spec(R)$ 是固有态射。对 $n \geq 1$，令
$X_n = X \times_{\Spec(R)} \Spec(R/\mathfrak m^n)$。
令 $Z \subset X_1$ 为特殊纤维的一个闭子集。置
$U = X \setminus Z$，并以 $j : U \to X$ 表示嵌入态射。给定一个对象
$$
(\mathcal{F}_n) \text{ 属于 } \textit{Coh}(U, \mathfrak m\mathcal{O}_U)
$$
它在 $R$ 上平坦，其含义是 $\mathcal{F}_n$ 在
$R/\mathfrak m^n$ 上平坦，且这对所有 $n$ 成立。假设
$j_*\mathcal{F}_1$ 和
$R^1j_*\mathcal{F}_1$ 是凝聚模。那么，由命题
\ref{algebraization-proposition-algebraization-flat}，仿射局部地在 $X$ 上可得到
$(\mathcal{F}_n)$ 的典范延拓；而该延拓的构造与局部化可交换
（由引理 \ref{algebraization-lemma-algebraization-principal-variant}）。
于是得到一个典范的整体对象 $(\mathcal{G}_n)$，它属于
$\textit{Coh}(X, \mathfrak m\mathcal{O}_X)$
且其在 $U$ 上的限制是 $(\mathcal{F}_n)$。由 Grothendieck
存在性定理（概形上同调，命题
\ref{coherent-proposition-existence-proper}），可知存在凝聚
$\mathcal{O}_X$-模 $\mathcal{G}$，其完备化为 $(\mathcal{G}_n)$。
由此可见 $(\mathcal{F}_n)$ 可代数化，也就是说，它是某个凝聚
$\mathcal{O}_U$-模的完备化。

\medskip\noindent
再补充一点：$j_*\mathcal{F}_1$ 和 $R^1j_*\mathcal{F}_1$
的凝聚性是特殊纤维上的条件。具体而言，若以
$j_1 : U_1 \to X_1$ 表示 $j : U \to X$ 的特殊纤维，则可以把
$\mathcal{F}_1$ 视为 $U_1$ 上的凝聚层，并且有
$j_*\mathcal{F}_1 = j_{1, *}\mathcal{F}_1$，且
$R^1j_*\mathcal{F}_1 = R^1j_{1, *}\mathcal{F}_1$。
因此，例如，若 $X_1$ 满足 $(S_2)$ 且不可约，并且
$\dim(X_1) - \dim(Z) \geq 3$，而 $\mathcal{F}_1$ 是局部自由
$\mathcal{O}_{U_1}$-模，则 $j_{1, *}\mathcal{F}_1$ 和
$R^1j_{1, *}\mathcal{F}_1$ 都是凝聚模。
\end{remark}








\section{完备化函子的应用}
\label{algebraization-section-completion-application}

\noindent
本节只是把若干已得结果组合起来，以便引用。这里还可以陈述许多
（更强的）结果。

\begin{lemma}
\label{algebraization-lemma-equivalence-better}
在情形 \ref{algebraization-situation-algebraize} 中，假设
\begin{enumerate}
\item $A$ 有对偶化复形，并且是 $I$-进完备的；
\item $I = (f)$ 由一个元素生成；
\item $A$ 是局部环，其极大理想为 $\mathfrak a = \mathfrak m$；
\item 下列条件之一成立：
\begin{enumerate}
\item $A_f$ 满足 $(S_2)$，并且对满足 $\mathfrak p \subset A$、
$f \not \in \mathfrak p$ 的极小素理想，有
$\dim(A/\mathfrak p) \geq 4$；或者
\item 若 $\mathfrak p \not \in V(f)$ 且
$V(\mathfrak p) \cap V(f) \not = \{\mathfrak m\}$，则
$\text{depth}(A_\mathfrak p) + \dim(A/\mathfrak p) > 3$。
\end{enumerate}
\end{enumerate}
那么，置 $U_0 = U \cap V(f)$，完备化函子
$$
\colim_{U_0 \subset U' \subset U\text{ 开}}
\textit{Coh}(\mathcal{O}_{U'})
\longrightarrow
\textit{Coh}(U, f\mathcal{O}_U)
$$
在有限局部自由对象所成的满子范畴上是等价。
\end{lemma}

\begin{proof}
由引理 \ref{algebraization-lemma-fully-faithful-general} 可知该函子充分忠实
（略去细节）。下面证明本质满。令 $(\mathcal{F}_n)$ 为
$\textit{Coh}(U, f\mathcal{O}_U)$ 中的有限局部自由对象。由引理
\ref{algebraization-lemma-algebraization-principal-bis} 或命题
\ref{algebraization-proposition-cd-1}，存在凝聚 $\mathcal{O}_U$-模 $\mathcal{F}$，
使得 $(\mathcal{F}_n)$ 是 $\mathcal{F}$ 的完备化。事实上，要应用
任一结果，唯一需要检验的是 $(\mathcal{F}_n)$ 满足 $(2, 3)$-不等式。
这由引理 \ref{algebraization-lemma-unwinding-conditions-bis} 给出。若 $y \in U_0$，
则 $f$-进完备化后的茎 $\mathcal{F}_y$ 同构于某个有限自由模，该模
定义在 $f$-进完备化后的 $\mathcal{O}_{U, y}$ 上。因此
$\mathcal{F}$ 在某个开邻域 $U'$ 上有限局部自由，而该邻域包含
$U_0$。证毕。
\end{proof}

\begin{lemma}
\label{algebraization-lemma-equivalence}
在情形 \ref{algebraization-situation-algebraize} 中，假设
\begin{enumerate}
\item $I = (f)$ 是主理想；
\item $A$ 是 $f$-进完备的；
\item $f$ 是非零因子；
\item $H^1_\mathfrak a(A/fA)$ 和 $H^2_\mathfrak a(A/fA)$
是有限 $A$-模。
\end{enumerate}
那么，置 $U_0 = U \cap V(f)$，完备化函子
$$
\colim_{U_0 \subset U' \subset U\text{ 开}}
\textit{Coh}(\mathcal{O}_{U'})
\longrightarrow
\textit{Coh}(U, f\mathcal{O}_U)
$$
在有限局部自由对象所成的满子范畴上是等价。
\end{lemma}

\begin{proof}
由引理 \ref{algebraization-lemma-fully-faithful-general-alternative}，该函子充分忠实。
本质满性由引理 \ref{algebraization-lemma-algebraization-principal-variant} 得到。
\end{proof}








\section{凝聚三元组}
\label{algebraization-section-coherent-triples}

\noindent
令 $(A, \mathfrak m)$ 为 Noether 局部环，令
$f \in \mathfrak m$ 为非零因子。置
$X = \Spec(A)$、$X_0 = \Spec(A/fA)$、$U = X \setminus V(\mathfrak m)$，以及
$U_0 = U \cap X_0$。若满足下列条件，就称
$(\mathcal{F}, \mathcal{F}_0, \alpha)$ 为一个{\it 凝聚三元组}：
\begin{enumerate}
\item $\mathcal{F}$ 是凝聚 $\mathcal{O}_U$-模，且
$f : \mathcal{F} \to \mathcal{F}$ 是单射；
\item $\mathcal{F}_0$ 是凝聚 $\mathcal{O}_{X_0}$-模；
\item $\alpha : \mathcal{F}/f\mathcal{F} \to \mathcal{F}_0|_{U_0}$
是同构。
\end{enumerate}
凝聚三元组之间有显然的{\it 态射}概念，它使所有凝聚三元组组成一个范畴。

\medskip\noindent
凝聚三元组范畴是加性范畴，但不是 Abel 范畴。不过，凝聚三元组的
短正合列应当是什么意思是清楚的。

\medskip\noindent
给定两个凝聚三元组 $(\mathcal{F}, \mathcal{F}_0, \alpha)$ 和
$(\mathcal{G}, \mathcal{G}_0, \beta)$，则
$(\mathcal{F} \otimes_{\mathcal{O}_U} \mathcal{G},
\mathcal{F}_0  \otimes_{\mathcal{O}_{X_0}} \mathcal{G}_0,
\alpha \otimes \beta)$ 未必是凝聚三元组\footnote{也就是说，$f$
在 $\mathcal{F} \otimes_{\mathcal{O}_U} \mathcal{G}$ 上未必是单射。}。
不过，若茎 $\mathcal{G}_x$ 对所有 $x \in U_0$ 都自由，则它确实是。

\medskip\noindent
称凝聚三元组 $(\mathcal{G}, \mathcal{G}_0, \beta)$ 为
{\it 局部自由的}（相应地，{\it 可逆的}），如果
$\mathcal{G}$ 和 $\mathcal{G}_0$ 分别是局部自由模（相应地，可逆模）。
在这种情形下，与
$(\mathcal{G}, \mathcal{G}_0, \beta)$ 作张量积是有意义的（见上文），
并且它把凝聚三元组的短正合列变为凝聚三元组的短正合列。

\begin{lemma}
\label{algebraization-lemma-prepare-chi-triple}
对任意凝聚三元组 $(\mathcal{F}, \mathcal{F}_0, \alpha)$，存在凝聚
$\mathcal{O}_X$-模 $\mathcal{F}'$，使得
$f : \mathcal{F}' \to \mathcal{F}'$ 是单射；还存在同构
$\alpha' : \mathcal{F}'|_U \to \mathcal{F}$ 以及映射
$\alpha'_0 : \mathcal{F}'/f\mathcal{F}' \to \mathcal{F}_0$
使得 $\alpha \circ (\alpha' \bmod f) = \alpha'_0|_{U_0}$。
\end{lemma}

\begin{proof}
选取有限 $A$-模 $M$。要求 $\mathcal{F}$ 是某个层在 $U$ 上的限制；
这个层是凝聚 $\mathcal{O}_X$-模，即由 $M$ 给出的相伴层，见
局部上同调，引理
\ref{local-cohomology-lemma-finiteness-pushforwards-and-H1-local}。
由于 $\mathcal{F}$ 无 $f$-挠，可以用 $M$ 对其
$f$-幂挠子模的商来代替它。另一方面，令
$M_0 = \Gamma(X_0, \mathcal{F}_0)$，于是 $\mathcal{F}_0$ 是凝聚
$\mathcal{O}_{X_0}$-模，并且它与有限 $A/fA$-模 $M_0$ 相伴。
由概形上同调，引理 \ref{coherent-lemma-homs-over-open}，存在 $n$，
使同构 $\alpha_0$ 对应于一个 $A/fA$-模同态
$\mathfrak m^n M/fM \to M_0$（其核与余核均被 $\mathfrak m$ 的某个幂
零化，但这里不需要这一点）。因此，若取 $M' = \mathfrak m^n M$，
并令 $\mathcal{F}'$ 为凝聚 $\mathcal{O}_X$-模，即与 $M'$ 相伴的模，
则引理显然成立。
\end{proof}

\noindent
令 $(\mathcal{F}, \mathcal{F}_0, \alpha)$ 为凝聚三元组。按照
引理 \ref{algebraization-lemma-prepare-chi-triple} 选取
$\mathcal{F}', \alpha', \alpha'_0$。置
\begin{equation}
\label{algebraization-equation-chi-triple}
\chi(\mathcal{F}, \mathcal{F}_0, \alpha) =
\text{length}_A(\Coker(\alpha'_0)) - 
\text{length}_A(\Ker(\alpha'_0))
\end{equation}
右端表达式有意义，因为 $\alpha'_0$ 在 $U_0$ 上是同构，故其核与余核
都是支撑在 $\{\mathfrak m\}$ 上的凝聚模，从而具有有限长度
（代数，引理 \ref{algebra-lemma-support-point}）。

\begin{lemma}
\label{algebraization-lemma-well-defined-chi-triple}
式 (\ref{algebraization-equation-chi-triple}) 中的量
$\chi(\mathcal{F}, \mathcal{F}_0, \alpha)$ 与按照引理
\ref{algebraization-lemma-prepare-chi-triple} 对
$\mathcal{F}', \alpha', \alpha'_0$ 的选择无关。
\end{lemma}

\begin{proof}
令 $\mathcal{F}', \alpha', \alpha'_0$ 和
$\mathcal{F}'', \alpha'', \alpha''_0$ 为两组这样的选择。对 $n > 0$，置
$\mathcal{F}'_n = \mathfrak m^n \mathcal{F}'$。由概形上同调，引理
\ref{coherent-lemma-homs-over-open}，存在某个 $n$ 以及
$\mathcal{O}_X$-模映射 $\mathcal{F}'_n \to \mathcal{F}''$，它与等同
$\mathcal{F}''|_U = \mathcal{F}'|_U$ 相符；这个等同由 $\alpha'$ 和
$\alpha''$ 确定。于是图
$$
\xymatrix{
\mathcal{F}'_n/f\mathcal{F}'_n \ar[r] \ar[d] &
\mathcal{F}'/f\mathcal{F}' \ar[d]^{\alpha_0'} \\
\mathcal{F}''/f\mathcal{F}'' \ar[r]^{\alpha_0''} &
\mathcal{F}_0
}
$$
限制到 $U_0$ 后交换。因此，由概形上同调，引理
\ref{coherent-lemma-homs-over-open}，该图限制到
$\mathfrak m^l(\mathcal{F}'_n/f\mathcal{F}'_n)$ 后交换，其中
$l > 0$ 是某个整数。由于
$\mathcal{F}'_{n + l}/f\mathcal{F}'_{n + l} \to \mathcal{F}'_n/f\mathcal{F}'_n$
通过 $\mathfrak m^l(\mathcal{F}'_n/f\mathcal{F}'_n)$ 分解，可知把
$n$ 替换为 $n + l$ 后该图交换。换言之，我们找到了第三组选择
$\mathcal{F}''', \alpha''', \alpha'''_0$
以及映射 $\mathcal{F}''' \to \mathcal{F}''$ 和
$\mathcal{F}''' \to \mathcal{F}'$；它们定义在 $X$ 上，并与 $U$ 和
$X_0$ 上的映射相容。
于是问题化归为下一段讨论的情形。

\medskip\noindent
假设有映射 $\mathcal{F}'' \to \mathcal{F}'$，它定义在 $X$ 上，并分别
与 $\alpha', \alpha''$ 在 $U$ 上以及与 $\alpha'_0, \alpha''_0$
在 $X_0$ 上相容。注意 $\mathcal{F}'' \to \mathcal{F}'$ 是单射，
因为它在 $U$ 上是同构，而且
$f : \mathcal{F}'' \to \mathcal{F}''$ 是单射。显然
$\mathcal{F}'/\mathcal{F}''$ 支撑在 $\{\mathfrak m\}$ 上，因而长度有限。
我们有如下凝聚 $\mathcal{O}_{X_0}$-模的映射：
$$
\mathcal{F}''/f\mathcal{F}'' \to
\mathcal{F}'/f\mathcal{F}' \xrightarrow{\alpha'_0}
\mathcal{F}_0
$$
其复合为 $\alpha''_0$，并且这些映射在 $U_0$ 上都是同构。初等同调代数
给出 $6$-项正合列
$$
\begin{matrix}
0 \to
\Ker(\mathcal{F}''/f\mathcal{F}'' \to \mathcal{F}'/f\mathcal{F}') \to
\Ker(\alpha''_0) \to
\Ker(\alpha'_0) \to \\
\Coker(\mathcal{F}''/f\mathcal{F}'' \to \mathcal{F}'/f\mathcal{F}') \to
\Coker(\alpha''_0) \to
\Coker(\alpha'_0) \to 0
\end{matrix}
$$
由长度的可加性（代数，引理 \ref{algebra-lemma-length-additive}），
只需证明
$$
\text{length}_A(
\Coker(\mathcal{F}''/f\mathcal{F}'' \to \mathcal{F}'/f\mathcal{F}')) -
\text{length}_A(
\Ker(\mathcal{F}''/f\mathcal{F}'' \to \mathcal{F}'/f\mathcal{F}')) = 0
$$
对下图应用蛇形引理即可得到这一点：
$$
\xymatrix{
0 \ar[r] &
\mathcal{F}'' \ar[r]_f \ar[d] &
\mathcal{F}'' \ar[r] \ar[d] &
\mathcal{F}''/f\mathcal{F}'' \ar[r] \ar[d] &
0 \\
0 \ar[r] &
\mathcal{F}' \ar[r]^f &
\mathcal{F}' \ar[r] &
\mathcal{F}'/f\mathcal{F}' \ar[r] &
0
}
$$
这里还使用了 $\mathcal{F}'/\mathcal{F}''$ 长度有限这一事实。
\end{proof}

\begin{lemma}
\label{algebraization-lemma-ses-chi-triple}
下式成立：
$\chi(\mathcal{G}, \mathcal{G}_0, \beta) =
\chi(\mathcal{F}, \mathcal{F}_0, \alpha) +
\chi(\mathcal{H}, \mathcal{H}_0, \gamma)$，其条件是
$$
0 \to
(\mathcal{F}, \mathcal{F}_0, \alpha) \to
(\mathcal{G}, \mathcal{G}_0, \beta) \to
(\mathcal{H}, \mathcal{H}_0, \gamma)
\to 0
$$
是凝聚三元组的短正合列。
\end{lemma}

\begin{proof}
选取 $\mathcal{G}', \beta', \beta'_0$，如引理
\ref{algebraization-lemma-prepare-chi-triple} 所述；这是针对三元组
$(\mathcal{G}, \mathcal{G}_0, \beta)$ 的选择。以
$j : U \to X$ 表示嵌入态射。令 $\mathcal{F}' \subset \mathcal{G}'$
为如下复合的核：
$$
\mathcal{G}' \xrightarrow{\beta'} j_*\mathcal{G} \to j_*\mathcal{H}
$$
注意 $\mathcal{H}' = \mathcal{G}'/\mathcal{F}'$ 是
$j_*\mathcal{H}$ 的凝聚子层，因而
$f : \mathcal{H}' \to \mathcal{H}'$ 是单射。因此，由蛇形引理得到短正合列
$$
0 \to \mathcal{F}'/f\mathcal{F}' \to
\mathcal{G}'/f\mathcal{G}' \to
\mathcal{H}'/f\mathcal{H}' \to 0
$$
由构造可得同构
$\alpha' : \mathcal{F}'|_U \to \mathcal{F}$,
$\beta' : \mathcal{G}'|_U \to \mathcal{G}$，以及
$\gamma' : \mathcal{H}'|_U \to \mathcal{H}$。为完成证明，需要构造映射
$\alpha'_0 : \mathcal{F}'/f\mathcal{F}' \to \mathcal{F}_0$ 和
$\gamma'_0 : \mathcal{H}'/f\mathcal{H}' \to \mathcal{H}_0$，使其如
引理 \ref{algebraization-lemma-prepare-chi-triple} 所述，并嵌入下列交换图：
$$
\xymatrix{
0 \ar[r] &
\mathcal{F}'/f\mathcal{F}' \ar[r] \ar@{..>}[d]^{\alpha'_0} &
\mathcal{G}'/f\mathcal{G}' \ar[r] \ar[d]^{\beta'_0} &
\mathcal{H}'/f\mathcal{H}' \ar[r] \ar@{..>}[d]^{\gamma'_0} &
0 \\
0 \ar[r] &
\mathcal{F}_0 \ar[r] &
\mathcal{G}_0 \ar[r] &
\mathcal{H}_0 \ar[r] &
0
}
$$
不过，对 $\mathcal{G}'$ 的初始选择，未必能做到这一点。从上图可见，
障碍恰好是复合
$$
\delta :
\mathcal{F}'/f\mathcal{F}' \to
\mathcal{G}'/f\mathcal{G}' \xrightarrow{\beta'_0}
\mathcal{G}_0 \to
\mathcal{H}_0
$$
注意，$\delta$ 限制到 $U_0$ 后为零，这是由我们对 $\mathcal{F}'$ 和
$\mathcal{H}'$ 的选择所致。因此，由概形上同调，引理
\ref{coherent-lemma-homs-over-open}，存在 $k > 0$，使得
$\delta$ 在 $\mathfrak m^k \cdot (\mathcal{F}'/f\mathcal{F}')$ 上为零。
对 $n > k$，置 $\mathcal{G}'_n = \mathfrak m^n \mathcal{G}'$、
$\mathcal{F}'_n = \mathcal{G}'_n \cap \mathcal{F}'$，以及
$\mathcal{H}'_n = \mathcal{G}'_n/\mathcal{F}'_n$。注意，$\beta'_0$
可与
$\mathcal{G}'_n/f\mathcal{G}'_n \to \mathcal{G}'/f\mathcal{G}'$
复合，从而给出映射
$\beta'_{n, 0} : \mathcal{G}'_n/f\mathcal{G}'_n \to \mathcal{G}_0$
，如引理 \ref{algebraization-lemma-prepare-chi-triple} 所述。由 Artin--Rees
（代数，引理 \ref{algebra-lemma-Artin-Rees}），可以选取 $n$，使得
$\mathcal{F}'_n \subset \mathfrak m^k \mathcal{F}'$。如上，映射
$f : \mathcal{F}'_n \to \mathcal{F}'_n$,
$f : \mathcal{G}'_n \to \mathcal{G}'_n$，以及
$f : \mathcal{H}'_n \to \mathcal{H}'_n$ 都是单射；再次使用蛇形引理，
得到短正合列
$$
0 \to \mathcal{F}'_n/f\mathcal{F}'_n \to
\mathcal{G}'_n/f\mathcal{G}'_n \to
\mathcal{H}'_n/f\mathcal{H}'_n \to 0
$$
如上，有同构
$\alpha'_n : \mathcal{F}'_n|_U \to \mathcal{F}$,
$\beta'_n : \mathcal{G}'_n|_U \to \mathcal{G}$，以及
$\gamma'_n : \mathcal{H}'_n|_U \to \mathcal{H}$。与前面一样，考虑障碍
$$
\delta_n :
\mathcal{F}'_n/f\mathcal{F}'_n \to
\mathcal{G}'_n/f\mathcal{G}'_n
\xrightarrow{\beta'_{n, 0}}
\mathcal{G}_0 \to
\mathcal{H}_0
$$
不过，交换图
$$
\xymatrix{
\mathcal{F}'_n/f\mathcal{F}'_n \ar[r] \ar[d] &
\mathcal{G}'_n/f\mathcal{G}'_n \ar[r]_{\beta'_{n, 0}} \ar[d] &
\mathcal{G}_0 \ar[r] \ar[d] &
\mathcal{H}_0 \ar[d] \\
\mathcal{F}'/f\mathcal{F}' \ar[r] &
\mathcal{G}'/f\mathcal{G}' \ar[r]^{\beta'_0} &
\mathcal{G}_0 \ar[r] &
\mathcal{H}_0
}
$$
结合对 $n$ 的选择和关于 $\delta$ 的观察，表明 $\delta_n = 0$。
这就产生了所需映射
$\alpha'_{n, 0} : \mathcal{F}'_n/f\mathcal{F}'_n \to \mathcal{F}_0$ 和
$\gamma'_{n, 0} : \mathcal{H}'_n/f\mathcal{H}'_n \to \mathcal{H}_0$.
因此，可以使用
$\mathcal{F}'_n, \alpha'_n, \alpha'_{n, 0}$,
$\mathcal{G}'_n, \beta'_n, \beta'_{n, 0}$，以及
$\mathcal{H}'_n, \gamma'_n, \gamma'_{n, 0}$
来计算
$\chi(\mathcal{F}, \mathcal{F}_0, \alpha)$,
$\chi(\mathcal{G}, \mathcal{G}_0, \beta)$，以及
$\chi(\mathcal{H}, \mathcal{H}_0, \gamma)$。最后，对下图应用蛇形引理
$$
\xymatrix{
0 \ar[r] &
\mathcal{F}'_n/f\mathcal{F}'_n \ar[r] \ar[d] &
\mathcal{G}'_n/f\mathcal{G}'_n \ar[r] \ar[d] &
\mathcal{H}'_n/f\mathcal{H}'_n \ar[r] \ar[d] &
0 \\
0 \ar[r] &
\mathcal{F}_0 \ar[r] &
\mathcal{G}_0 \ar[r] &
\mathcal{H}_0 \ar[r] &
0
}
$$
并使用长度的可加性（代数，引理
\ref{algebra-lemma-length-additive}），即可得到本引理。
\end{proof}

\begin{proposition}
\label{algebraization-proposition-hilbert-triple}
令 $(\mathcal{F}, \mathcal{F}_0, \alpha)$ 为凝聚三元组，并令
$(\mathcal{L}, \mathcal{L}_0, \lambda)$ 为可逆凝聚三元组。那么函数
$$
\mathbf{Z} \longrightarrow \mathbf{Z},\quad
n \longmapsto
\chi((\mathcal{F}, \mathcal{F}_0, \alpha) \otimes
(\mathcal{L}, \mathcal{L}_0, \lambda)^{\otimes n})
$$
是次数 $\leq \dim(\text{Supp}(\mathcal{F}))$ 的多项式。
\end{proposition}

\noindent
更确切地说，若 $\mathcal{F} = 0$，则该函数为常数。若
$\mathcal{F}$ 在 $U$ 中的支撑有限，则该函数为常数。若
$\mathcal{F}$ 在 $U$ 中的支撑维数为 $1$，即 $\mathcal{F}$ 的支撑
在 $X$ 中的闭包维数为 $2$，则该函数为线性函数，依此类推。

\begin{proof}
我们对 $\mathcal{F}$ 的支撑维数作归纳来证明。

\medskip\noindent
基本情形是 $\mathcal{F} = 0$。此时 $\mathcal{F}_0$ 或者为零，
或者其支撑为 $\{\mathfrak m\}$。在这种情形下，有
$$
(\mathcal{F}, \mathcal{F}_0, \alpha) \otimes
(\mathcal{L}, \mathcal{L}_0, \lambda)^{\otimes n} =
(0, \mathcal{F}_0 \otimes \mathcal{L}_0^{\otimes n}, 0) \cong
(0, \mathcal{F}_0, 0)
$$
因此命题中的函数为常数，其值等于 $\mathcal{F}_0$ 的长度。

\medskip\noindent
归纳步骤。假设 $\mathcal{F}$ 的支撑非空。以
$\mathcal{G}_0 \subset \mathcal{F}_0$ 表示由支撑在
$\{\mathfrak m\}$ 上的截面组成的子模。于是得到短正合列
$$
0 \to (0, \mathcal{G}_0, 0) \to
(\mathcal{F}, \mathcal{F}_0, \alpha) \to
(\mathcal{F}, \mathcal{F}_0/\mathcal{G}_0, \alpha) \to 0
$$
与可逆凝聚三元组 $(\mathcal{L}, \mathcal{L}_0, \lambda)$ 作张量积后，
此列仍然正合，见上文讨论。因此，由 $\chi$ 的可加性
（引理 \ref{algebraization-lemma-ses-chi-triple}）以及上述基本情形，只需对
$(\mathcal{F}, \mathcal{F}_0/\mathcal{G}_0, \alpha)$ 证明归纳步骤。
由此可假设 $\mathfrak m$ 不是 $\mathcal{F}_0$ 的相伴点。

\medskip\noindent
令 $T = \text{Ass}(\mathcal{F}) \cup \text{Ass}(\mathcal{F}/f\mathcal{F})$。
由于 $U$ 拟仿射，可以找到 $s \in \Gamma(U, \mathcal{L})$，它在每个
$u \in T$ 处均不消失，见性质，引理
\ref{properties-lemma-quasi-affine-invertible-nonvanishing-section}。
将 $s$ 乘以 $\mathfrak m$ 的适当元素后，可以假设
$\lambda(s \bmod f) = s_0|_{U_0}$，其中
$s_0 \in \Gamma(X_0, \mathcal{L}_0)$；略去细节。于是得到态射
$$
(s, s_0) :
(\mathcal{O}_U, \mathcal{O}_{X_0}, 1)
\longrightarrow
(\mathcal{L}, \mathcal{L}_0, \lambda)
$$
它属于凝聚三元组范畴。令
$\mathcal{G} = \Coker(s : \mathcal{F} \to \mathcal{F} \otimes \mathcal{L})$
并令
$\mathcal{G}_0 = \Coker(s_0 : \mathcal{F}_0 \to
\mathcal{F}_0 \otimes \mathcal{L}_0)$。注意，
$s_0 : \mathcal{F}_0 \to
\mathcal{F}_0 \otimes \mathcal{L}_0$ 是单射：它在 $U_0$ 上单射，
这是由 $s$ 的选择得到的；而且 $\mathfrak m$ 不是
$\mathcal{F}_0$ 的相伴点。因此存在同构
$\beta : \mathcal{G}/f\mathcal{G} \to \mathcal{G}_0|_{U_0}$，使得得到
短正合列
$$
0 \to
(\mathcal{F}, \mathcal{F}_0, \alpha) \to
(\mathcal{F}, \mathcal{F}_0, \alpha) \otimes
(\mathcal{L}, \mathcal{L}_0, \lambda) \to
(\mathcal{G}, \mathcal{G}_0, \beta) \to 0
$$
由支撑维数的归纳假设，命题对凝聚三元组
$(\mathcal{G}, \mathcal{G}_0, \beta)$ 成立。利用引理
\ref{algebraization-lemma-ses-chi-triple} 的可加性，可知
$$
n \longmapsto 
\chi((\mathcal{F}, \mathcal{F}_0, \alpha) \otimes
(\mathcal{L}, \mathcal{L}_0, \lambda)^{\otimes n + 1})
-
\chi((\mathcal{F}, \mathcal{F}_0, \alpha) \otimes
(\mathcal{L}, \mathcal{L}_0, \lambda)^{\otimes n})
$$
是多项式。对定义在全体整数上的函数应用代数，引理
\ref{algebra-lemma-numerical-polynomial} 的一个变体，即得结论
（略去细节）。
\end{proof}

\begin{lemma}
\label{algebraization-lemma-nonnegative-chi-triple}
假设 $\text{depth}(A) \geq 3$，等价地，
$\text{depth}(A/fA) \geq 2$。令 $(\mathcal{L}, \mathcal{L}_0, \lambda)$
为可逆凝聚三元组。那么
$$
\chi(\mathcal{L}, \mathcal{L}_0, \lambda) =
\text{length}_A \Coker(\Gamma(U, \mathcal{L}) \to \Gamma(U_0, \mathcal{L}_0))
$$
特别地，该数 $\geq 0$。此外，
$\chi(\mathcal{L}, \mathcal{L}_0, \lambda) = 0$ 当且仅当
$\mathcal{L} \cong \mathcal{O}_U$。
\end{lemma}

\begin{proof}
深度条件的等价性由代数，引理
\ref{algebra-lemma-depth-drops-by-one} 得到。由深度条件可知
$\Gamma(U, \mathcal{O}_U) = A$ 且
$\Gamma(U_0, \mathcal{O}_{U_0}) = A/fA$，见对偶化复形，引理
\ref{dualizing-lemma-depth} 以及局部上同调，引理
\ref{local-cohomology-lemma-finiteness-pushforwards-and-H1-local}。
利用局部上同调，引理
\ref{local-cohomology-lemma-finiteness-for-finite-locally-free}，可知
$M = \Gamma(U, \mathcal{L})$ 是有限 $A$-模。进而有
$\text{depth}(M) \geq 2$；例如，可用局部上同调，引理
\ref{local-cohomology-lemma-finiteness-pushforwards-and-H1-local} 的第 (4) 部分，
或除子，引理 \ref{divisors-lemma-depth-pushforward}。此外，
$\mathcal{L}_0 \cong \mathcal{O}_{X_0}$，因为 $X_0$ 是局部概形。因此还有
$M_0 = \Gamma(X_0, \mathcal{L}_0) = \Gamma(U_0, \mathcal{L}_0|_{U_0})$
，且该模同构于 $A/fA$。

\medskip\noindent
由上可知，$\mathcal{F}' = \widetilde{M}$ 是凝聚
$\mathcal{O}_X$-模，其在 $U$ 上的限制同构于 $\mathcal{L}$。同构
$\lambda : \mathcal{L}/f\mathcal{L} \to \mathcal{L}_0|_{U_0}$
在整体截面上确定映射 $M/fM \to M_0$，该映射在 $U_0$ 上是同构。
由于 $\text{depth}(M) \geq 2$，可知
$H^0_\mathfrak m(M/fM) = 0$，从而
$M/fM \to M_0$ 是单射。因此，由定义
$$
\chi(\mathcal{L}, \mathcal{L}_0, \lambda) =
\text{length}_A \Coker(M/fM \to M_0)
$$
这就给出引理的第一个断言。

\medskip\noindent
最后，若此长度为 $0$，则 $M \to M_0$ 是满射。因此可以找到
$s \in M = \Gamma(U, \mathcal{L})$，它映到 $\mathcal{L}_0$
的一个平凡化截面。考虑由下列正合列定义的有限 $A$-模 $K$、$Q$：
$$
0 \to K \to A \xrightarrow{s} M \to Q \to 0
$$
$K$ 和 $Q$ 的支撑均不与 $U_0$ 相交，因为 $s$ 在 $U_0$ 的各点
都非零。利用代数，引理 \ref{algebra-lemma-depth-in-ses}，可知
$\text{depth}(K) \geq 2$（注意，$As \subset M$ 有
$\text{depth} \geq 1$，因为它是 $M$ 的子模）。因此，由代数，引理
\ref{algebra-lemma-bound-depth}，$K$ 的支撑若非空，则其维数
$\geq 2$。这与
$\text{Supp}(M) \cap V(f) \subset \{\mathfrak m\}$ 矛盾，除非
$K = 0$。当 $K = 0$ 时，可知
$\text{depth}(Q) \geq 2$，并且如前得出 $Q = 0$。因此
$A \cong M$，且 $\mathcal{L}$ 平凡。
\end{proof}







\section{穿孔谱上的可逆模，I}
\label{algebraization-section-local-lefschetz-for-pic}

\noindent
本节证明关于 Picard 群的若干局部 Lefschetz 定理。其中一些思路取自
\cite{Kollar-pic}、\cite{Bhatt-local} 和 \cite{Kollar-map-pic}。

\begin{lemma}
\label{algebraization-lemma-injective-torsion-in-pic}
令 $(A, \mathfrak m)$ 为 Noether 局部环。令 $f \in \mathfrak m$
为非零因子，并假设 $\text{depth}(A/fA) \geq 2$，等价地，
$\text{depth}(A) \geq 3$。令 $U$、相应地 $U_0$ 分别为
$A$、相应地 $A/fA$ 的穿孔谱。映射
$$
\Pic(U) \to \Pic(U_0)
$$
在挠部分上是单射。
\end{lemma}

\begin{proof}
令 $\mathcal{L}$ 为可逆 $\mathcal{O}_U$-模。注意，$\mathcal{L}$
映到 $0$（在 $\Pic(U_0)$ 中），当且仅当能把 $\mathcal{L}$ 延拓为
可逆凝聚三元组 $(\mathcal{L}, \mathcal{L}_0, \lambda)$，如
第 \ref{algebraization-section-coherent-triples} 节所述。由命题
\ref{algebraization-proposition-hilbert-triple}，函数
$$
n \longmapsto \chi((\mathcal{L}, \mathcal{L}_0, \lambda)^{\otimes n})
$$
是多项式。由引理 \ref{algebraization-lemma-nonnegative-chi-triple}，该多项式的值为零
当且仅当 $\mathcal{L}^{\otimes n}$ 平凡。因此，若 $\mathcal{L}$ 为挠元，
则该多项式有无穷多个零点，从而恒为零，于是 $\mathcal{L}$ 平凡。
\end{proof}

\begin{proposition}[Koll\'ar]
\label{algebraization-proposition-injective-pic}
\begin{reference}
\cite[Theorem 1.9]{Kollar-map-pic}
\end{reference}
令 $(A, \mathfrak m)$ 为 Noether 局部环，并令 $f \in \mathfrak m$。假设
\begin{enumerate}
\item $A$ 有对偶化复形；
\item $f$ 是非零因子；
\item $\text{depth}(A/fA) \geq 2$，等价地，
$\text{depth}(A) \geq 3$；
\item 若 $f \in \mathfrak p \subset A$ 是素理想且
$\dim(A/\mathfrak p) = 2$，则 $\text{depth}(A_\mathfrak p) \geq 2$。
\end{enumerate}
令 $U$、相应地 $U_0$ 分别为 $A$、相应地 $A/fA$ 的穿孔谱。映射
$$
\Pic(U) \to \Pic(U_0)
$$
是单射。最后，若 (1)、(2)、(3) 成立，$A$ 满足 $(S_2)$，且
$\dim(A) \geq 4$，则 (4) 成立。
\end{proposition}

\begin{proof}
令 $\mathcal{L}$ 为可逆 $\mathcal{O}_U$-模。注意，$\mathcal{L}$
映到 $0$（在 $\Pic(U_0)$ 中），当且仅当能把 $\mathcal{L}$ 延拓为
可逆凝聚三元组 $(\mathcal{L}, \mathcal{L}_0, \lambda)$，如
第 \ref{algebraization-section-coherent-triples} 节所述。由命题
\ref{algebraization-proposition-hilbert-triple}，函数
$$
n \longmapsto \chi((\mathcal{L}, \mathcal{L}_0, \lambda)^{\otimes n})
$$
是多项式 $P$。由引理 \ref{algebraization-lemma-nonnegative-chi-triple}，不等式
$P(n) \geq 0$ 对所有 $n \in \mathbf{Z}$ 都成立，且等号成立当且仅当
$\mathcal{L}^{\otimes n}$ 平凡。特别地，$P(0) = 0$，而 $P$
或者恒为零，此时结论成立；或者 $P$ 的次数为偶数且 $\geq 2$。

\medskip\noindent
置 $M = \Gamma(U, \mathcal{L})$ 以及
$M_0 = \Gamma(X_0, \mathcal{L}_0) = \Gamma(U_0, \mathcal{L}_0)$。
于是 $M$ 是有限 $A$-模，其深度 $\geq 2$，并且
$M_0 \cong A/fA$，见引理 \ref{algebraization-lemma-nonnegative-chi-triple} 的证明。
注意，$H^2_\mathfrak m(M)$ 是有限 $A$-模；这是由局部上同调，引理
\ref{local-cohomology-lemma-local-finiteness-for-finite-locally-free}
以及下述事实得到的：$H^i_\mathfrak m(A) = 0$ 对 $i = 0, 1, 2$ 成立，
因为 $\text{depth}(A) \geq 3$。考虑短正合列
$$
0 \to M/fM \to M_0 \to Q \to 0
$$
引理 \ref{algebraization-lemma-nonnegative-chi-triple} 告诉我们，$Q$ 具有有限长度，
且该长度等于 $\chi(\mathcal{L}, \mathcal{L}_0, \lambda)$。由上式的
局部上同调长正合列，得到 $Q = H^1_\mathfrak m(M/fM)$ 以及
$H^i_\mathfrak m(M/fM) = H^i_\mathfrak m(M_0) \cong H^i_\mathfrak m(A/fA)$
对 $i > 1$ 成立。再考虑与序列
$0 \to M \to M \to M/fM \to 0$ 相伴的局部上同调长正合列，其开头为
$$
0 \to Q \to H^2_\mathfrak m(M) \to H^2_\mathfrak m(M) \to
H^2_\mathfrak m(A/fA)
$$
利用长度的可加性，可知
$\chi(\mathcal{L}, \mathcal{L}_0, \lambda)$
等于映射 $H^2_\mathfrak m(M) \to H^2_\mathfrak m(A/fA)$ 的像的长度。

\medskip\noindent
先在一个特殊情形中证明命题，以阐明余下证明。暂且假设
$H^2_\mathfrak m(A/fA)$ 是有限长度模。此时有
$P(1) \leq \text{length}_A H^2_\mathfrak m(A/fA)$。把完全相同的论证
应用于 $\mathcal{L}^{\otimes n}$，可知
$P(n) \leq \text{length}_A H^2_\mathfrak m(A/fA)$ 对所有 $n$ 成立。
因此 $P$ 不可能有正次数，结论成立。在余下证明中，我们将修改该论证，
给出 $P(n)$ 的线性上界，这已足够。

\medskip\noindent
研究映射
$H^2_\mathfrak m(M) \to H^2_\mathfrak m(M_0) \cong H^2_\mathfrak m(A/fA)$.
以 $\omega_A^\bullet$ 表示为 $A$ 选取的正规化对偶化复形。由局部对偶
（对偶化复形，引理
\ref{dualizing-lemma-special-case-local-duality}），此映射是下列映射的
Matlis 对偶：
$$
\text{Ext}^{-2}_A(M, \omega_A^\bullet)
\longleftarrow
\text{Ext}^{-2}_A(M_0, \omega_A^\bullet)
$$
因而其像具有相同的有限长度。有限 $A$-模
$\text{Ext}^{-2}_A(M_0, \omega_A^\bullet)$
的支撑若非空，则由 $\mathfrak m$ 以及有限多个素理想
$\mathfrak p_1, \ldots, \mathfrak p_r$ 组成；这些素理想包含 $f$，并且
$\dim(A/\mathfrak p_i) = 1$。事实上，由局部上同调，引理
\ref{local-cohomology-lemma-sitting-in-degrees}，该支撑包含于满足
$\mathfrak p \subset A$ 以及
$\text{depth}_{A_\mathfrak p}(M_{0, \mathfrak p}) + \dim(A/\mathfrak p) \leq 2$.
的素理想之集。因此，只需证明不存在素理想 $\mathfrak p$，它包含 $f$，
使得 $\dim(A/\mathfrak p) = 2$ 且
$\text{depth}_{A_\mathfrak p}(M_{0, \mathfrak p}) = 0$。然而，由
$M_{0, \mathfrak p} \cong (A/fA)_\mathfrak p$，这会推出
$\text{depth}(A_\mathfrak p) = 1$，与假设 (4) 矛盾。
选择截面 $t \in \Gamma(U, \mathcal{L}^{\otimes -1})$，使其在各点
$\mathfrak p_1, \ldots, \mathfrak p_r$ 均不消失，见性质，引理
\ref{properties-lemma-quasi-affine-invertible-nonvanishing-section}。
整体截面上乘以 $t$ 确定映射 $t : M \to A$；它给出同构
$M_{\mathfrak p_i} \to A_{\mathfrak p_i}$，这对
$i = 1, \ldots, r$ 成立。以
$t_0 = t|_{U_0}$ 表示对应截面
$\Gamma(U_0, \mathcal{L}_0^{\otimes -1})$；它同样确定映射
$t_0 : M_0 \to A/fA$，且与 $t$ 相容。于是有交换图
$$
\xymatrix{
\text{Ext}^{-2}_A(M, \omega_A^\bullet) &
\text{Ext}^{-2}_A(M_0, \omega_A^\bullet) \ar[l] \\
\text{Ext}^{-2}_A(A, \omega_A^\bullet) \ar[u]^t &
\text{Ext}^{-2}_A(A/fA, \omega_A^\bullet) \ar[l] \ar[u]_{t_0}
}
$$
由此可知，顶部水平映射之像的长度至多等于
$\text{Ext}^{-2}_A(A/fA, \omega_A^\bullet)$ 的长度加上 $t_0$ 的余核长度。

\medskip\noindent
不过，若把 $\mathcal{L}$ 替换为 $\mathcal{L}^n$，其中 $n > 1$，
则可以使用
$$
t^n :
M_n = \Gamma(U, \mathcal{L}^{\otimes n})
\longrightarrow
\Gamma(U, \mathcal{O}_U) = A
$$
来代替 $t$。这会把 $t_0$ 替换为它的 $n$ 次幂。因此，映射
$\text{Ext}^{-2}_A(M_n, \omega_A^\bullet) \leftarrow
\text{Ext}^{-2}_A(M_{n, 0}, \omega_A^\bullet)$
的像的长度至多等于 $\text{Ext}^{-2}_A(A/fA, \omega_A^\bullet)$ 的长度
加上下列映射的余核长度：
$$
t_0^n : 
\text{Ext}^{-2}_A(A/fA, \omega_A^\bullet)
\longrightarrow
\text{Ext}^{-2}_A(M_{n, 0}, \omega_A^\bullet)
$$
经由同构 $M_0 \cong A/fA$，映射 $t_0$ 变为
$g : A/fA \to A/fA$，其中 $g \in A/fA$；经由相应同构
$M_{n, 0} \cong A/fA$，映射 $t_0^n$ 变为
$g^n : A/fA \to A/fA$。因此，上述余核的长度就是
$\text{Ext}^{-2}_A(A/fA, \omega_A^\bullet)$ 对 $g^n$ 的商的长度。
由于 $\text{Ext}^{-2}_A(A/fA, \omega_A^\bullet)$ 是有限 $A$-模，
其支撑 $T$ 的维数为 $1$，并且 $V(g) \cap T$ 只含闭点，
这是由对 $t$ 的选择得到的。所以由代数，引理
\ref{algebra-lemma-support-dimension-d}，该长度关于 $n$ 线性增长。

\medskip\noindent
最后证明末尾断言。假设 $f \in \mathfrak m \subset A$ 满足
(1)、(2)、(3)，$A$ 满足 $(S_2)$，并且 $\dim(A) \geq 4$。
条件 (1) 蕴含 $A$ 是链状环，见对偶化复形，引理
\ref{dualizing-lemma-universally-catenary}。于是
$\Spec(A)$ 等维，见局部上同调，引理
\ref{local-cohomology-lemma-catenary-S2-equidimensional}。因此
$\dim(A_\mathfrak p) + \dim(A/\mathfrak p) \geq 4$ 对每个素理想
$\mathfrak p$ 都成立，而它是 $A$ 的素理想。于是
$\text{depth}(A_\mathfrak p) \geq \min(2, \dim(A_\mathfrak p))
\geq \min(2, 4 - \dim(A/\mathfrak p))$，从而 (4) 成立。
\end{proof}

\begin{remark}
\label{algebraization-remark-compare-SGA2}
在 SGA2 中可找到下列结果。令 $(A, \mathfrak m)$ 为 Noether 局部环，
并令 $f \in \mathfrak m$。假设 $A$ 是正则环的商，元素
$f$ 是非零因子，并且
\begin{enumerate}
\item[(a)] 若 $\mathfrak p \subset A$ 是素理想且
$\dim(A/\mathfrak p) = 1$，则 $\text{depth}(A_\mathfrak p) \geq 2$；
\item[(b)] $\text{depth}(A/fA) \geq 3$，等价地，
$\text{depth}(A) \geq 4$。
\end{enumerate}
令 $U$、相应地 $U_0$ 分别为 $A$、相应地 $A/fA$ 的穿孔谱。那么映射
$$
\Pic(U) \to \Pic(U_0)
$$
是单射。这就是 \cite[Exposee XI, Lemma 3.16]{SGA2}\footnote{条件
(a) 由条件 (b) 推出，见代数，引理
\ref{algebra-lemma-depth-localization}。}。SGA2 的这一结果由命题
\ref{algebraization-proposition-injective-pic} 推出，因为
\begin{enumerate}
\item 正则环的商有对偶化复形（见对偶化复形，引理
\ref{dualizing-lemma-regular-gorenstein} 和命题
\ref{dualizing-proposition-dualizing-essentially-finite-type}）；
\item 若 $\text{depth}(A) \geq 4$，则
$\text{depth}(A_\mathfrak p) \geq 2$ 对所有素理想 $\mathfrak p$ 都成立，
只要 $\dim(A/\mathfrak p) = 2$，见代数，引理
\ref{algebra-lemma-depth-localization}。
\end{enumerate}
\end{remark}






\section{穿孔谱上的可逆模，II}
\label{algebraization-section-local-lefschetz-for-pic-surjective}

\noindent
下面转向 Picard 群局部 Lefschetz 定理中的满性。首先，为了把 $U_0$
上的可逆模延拓到一个开邻域，有如下简单判据。

\begin{lemma}
\label{algebraization-lemma-surjective-Pic-first}
令 $(A, \mathfrak m)$ 为 Noether 局部环，且 $f \in \mathfrak m$。假设
\begin{enumerate}
\item $A$ 是 $f$-进完备的；
\item $f$ 是非零因子；
\item $H^1_\mathfrak m(A/fA)$ 和 $H^2_\mathfrak m(A/fA)$
是有限 $A$-模；
\item $H^3_\mathfrak m(A/fA) = 0$\footnote{注意，若
$\text{depth}(A/fA) \geq 4$，等价地
$\text{depth}(A) \geq 5$，则 (3) 和 (4) 成立。}。
\end{enumerate}
令 $U$、相应地 $U_0$ 分别为 $A$、相应地 $A/fA$ 的穿孔谱。那么
$$
\colim_{U_0 \subset U' \subset U\text{ 开}} \Pic(U')
\longrightarrow
\Pic(U_0)
$$
是满射。
\end{lemma}

\begin{proof}
令 $U_0 \subset U_n \subset U$ 为第 $n$ 个无穷小邻域，即 $U_0$
的无穷小邻域。注意，
$U_n$ 在 $U_{n + 1}$ 中的理想层同构于 $\mathcal{O}_{U_0}$，因为
$U_0 \subset U$ 是由非零因子 $f$ 截出的主闭子概形。因此有 Abel 群正合列
$$
\Pic(U_{n + 1}) \to \Pic(U_n) \to
H^2(U_0, \mathcal{O}_{U_0}) = H^3_\mathfrak m(A/fA) = 0
$$
见态射进阶，引理
\ref{more-morphisms-lemma-picard-group-first-order-thickening}。
因此，每个可逆 $\mathcal{O}_{U_0}$-模都是某个可逆凝聚形式模的限制，
也就是说，它是 $\textit{Coh}(U, f\mathcal{O}_U)$ 的可逆对象。
应用引理 \ref{algebraization-lemma-equivalence} 即得结论。
\end{proof}

\begin{remark}
\label{algebraization-remark-surjective-Pic-second}
令 $(A, \mathfrak m)$ 为 Noether 局部环，且 $f \in \mathfrak m$。
若假设下列条件，则引理 \ref{algebraization-lemma-surjective-Pic-first} 的结论成立：
\begin{enumerate}
\item $A$ 有对偶化复形；
\item $A$ 是 $f$-进完备的；
\item $f$ 是非零因子；
\item 下列条件之一成立：
\begin{enumerate}
\item $A_f$ 满足 $(S_2)$，并且对满足 $\mathfrak p \subset A$、
$f \not \in \mathfrak p$ 的极小素理想，有
$\dim(A/\mathfrak p) \geq 4$；或者
\item 若 $\mathfrak p \not \in V(f)$ 且
$V(\mathfrak p) \cap V(f) \not = \{\mathfrak m\}$，则
$\text{depth}(A_\mathfrak p) + \dim(A/\mathfrak p) > 3$。
\end{enumerate}
\item $H^3_{\mathfrak m}(A/fA) = 0$.
\end{enumerate}
证明与引理 \ref{algebraization-lemma-surjective-Pic-first} 的证明完全相同，只是用引理
\ref{algebraization-lemma-equivalence-better} 代替引理 \ref{algebraization-lemma-equivalence}。
这里需要说明两点：(a) 似乎很难找到能够知道
$H^3_{\mathfrak m}(A/fA) = 0$ 的例子，而又不假设
$\text{depth}(A/fA) \geq 4$；(b) 引理
\ref{algebraization-lemma-equivalence-better} 的证明比引理
\ref{algebraization-lemma-equivalence} 的证明难得多。
\end{remark}

\begin{lemma}
\label{algebraization-lemma-surjective-Pic-first-better}
令 $(A, \mathfrak m)$ 为 Noether 局部环，且 $f \in \mathfrak m$。假设
\begin{enumerate}
\item 引理 \ref{algebraization-lemma-surjective-Pic-first} 的条件成立；
\item 对每个极大理想 $\mathfrak p \subset A_f$，局部环
$(A_f)_\mathfrak p$ 的穿孔谱具有平凡 Picard 群。
\end{enumerate}
令 $U$、相应地 $U_0$ 分别为 $A$、相应地 $A/fA$ 的穿孔谱。那么
$$
\Pic(U) \longrightarrow \Pic(U_0)
$$
是满射。
\end{lemma}

\begin{proof}
令 $\mathcal{L}_0 \in \Pic(U_0)$。由引理
\ref{algebraization-lemma-surjective-Pic-first}，存在开集 $U_0 \subset U' \subset U$
以及 $\mathcal{L}' \in \Pic(U')$，其在 $U_0$ 上的限制为
$\mathcal{L}_0$。由于 $U' \supset U_0$，可知 $U \setminus U'$ 由与
(2) 中素理想 $\mathfrak p_1, \ldots, \mathfrak p_n$ 对应的点组成。
由假设，可以找到可逆模 $\mathcal{L}'_i$，它定义在
$\Spec(A_{\mathfrak p_i})$ 上，并与 $\mathcal{L}'$ 在穿孔谱
$U' \times_U \Spec(A_{\mathfrak p_i})$ 上相符，
因为平凡可逆模总能延拓。把极限，引理
\ref{limits-lemma-glueing-near-closed-point-modules} 应用 $n$ 次，
可知 $\mathcal{L}'$ 延拓为 $U$ 上的可逆模。
\end{proof}

\begin{lemma}
\label{algebraization-lemma-local-pic-to-completion}
令 $(A, \mathfrak m)$ 为深度 $\geq 2$ 的 Noether 局部环，并令
$A^\wedge$ 为其完备化。令 $U$、相应地 $U^\wedge$ 分别为
$A$、相应地 $A^\wedge$ 的穿孔谱。那么
$\Pic(U) \to \Pic(U^\wedge)$ 是单射。
\end{lemma}

\begin{proof}
令 $\mathcal{L}$ 为可逆 $\mathcal{O}_U$-模，其拉回为
$\mathcal{L}^\wedge$，定义在 $U^\wedge$ 上。由深度假设、对偶化复形，引理
\ref{dualizing-lemma-depth} 以及局部上同调，引理
\ref{local-cohomology-lemma-finiteness-pushforwards-and-H1-local}，有
$H^0(U, \mathcal{O}_U) = A$。因此，$\mathcal{L}$ 平凡当且仅当
$M = H^0(U, \mathcal{L})$ 作为 $A$-模同构于 $A$（略去细节）。
由于 $A \to A^\wedge$ 平坦，由平坦基变换有
$M \otimes_A A^\wedge = \Gamma(U^\wedge, \mathcal{L}^\wedge)$，见
概形上同调，引理 \ref{coherent-lemma-flat-base-change-cohomology}。
最后，容易看出 $M \cong A$ 当且仅当
$M \otimes_A A^\wedge \cong A^\wedge$。
\end{proof}

\begin{lemma}
\label{algebraization-lemma-trivial-local-pic-regular}
令 $(A, \mathfrak m)$ 为正则局部环。那么 $A$ 的穿孔谱的 Picard 群平凡。
\end{lemma}

\begin{proof}
结合除子，引理 \ref{divisors-lemma-extend-invertible-module} 与
代数进阶，引理 \ref{more-algebra-lemma-regular-local-UFD}。
\end{proof}

\noindent
现在可以自举前述结果，证明维数 $\geq 4$ 的完全交的穿孔谱具有平凡
Picard 群。回忆，若一个 Noether 局部环的完备化是某个正则局部环
关于一个由正则序列生成的理想的商，则称该环为完全交。见除幂代数，
第 \ref{dpa-section-lci} 节的讨论。

\begin{proposition}[Grothendieck]
\label{algebraization-proposition-trivial-local-pic-complete-intersection}
令 $(A, \mathfrak m)$ 为 Noether 局部环。若 $A$ 是维数
$\geq 4$ 的完全交，则 $A$ 的穿孔谱的 Picard 群平凡。
\end{proposition}

\begin{proof}
由引理 \ref{algebraization-lemma-local-pic-to-completion}，可以假设 $A$ 是完备局部环。
由假设可写成 $A = B/(f_1, \ldots, f_r)$，其中 $B$ 是完备正则局部环，
而 $f_1, \ldots, f_r$ 是正则序列。我们对 $r$ 作归纳以完成证明。
基本情形 $r = 0$ 由引理 \ref{algebraization-lemma-trivial-local-pic-regular} 得到。

\medskip\noindent
假设 $A = B/(f_1, \ldots, f_r)$，且命题对 $r - 1$ 成立。置
$A' = B/(f_1, \ldots, f_{r - 1})$，并把引理
\ref{algebraization-lemma-surjective-Pic-first-better} 应用于 $f_r \in A'$。
这是允许的：
\begin{enumerate}
\item 引理 \ref{algebraization-lemma-surjective-Pic-first} 的条件 (1) 成立，因为所用局部环
均完备；
\item 引理 \ref{algebraization-lemma-surjective-Pic-first} 的条件 (2) 成立，因为
$f_1, \ldots, f_r$ 是正则序列；
\item 引理 \ref{algebraization-lemma-surjective-Pic-first} 的条件 (3) 和 (4) 成立，因为
$A = A'/f_r A'$ 是维数 $\dim(A) \geq 4$ 的 Cohen--Macaulay 环；
\item 引理 \ref{algebraization-lemma-surjective-Pic-first-better} 的条件 (2) 由归纳假设
成立，因为 $\dim((A'_{f_r})_\mathfrak p) \geq 4$ 对极大素理想
$\mathfrak p$ 成立，而该素理想属于 $A'_{f_r}$；并且
$(A'_{f_r})_\mathfrak p = B_\mathfrak q/(f_1, \ldots, f_{r - 1})$
对某个素理想 $\mathfrak q \subset B$ 成立，且 $B_\mathfrak q$ 正则。
\end{enumerate}
证毕。
\end{proof}

\begin{example}
\label{algebraization-example-grothendieck-sharp}
命题 \ref{algebraization-proposition-trivial-local-pic-complete-intersection} 中的维数界是
尖锐的。例如，$A = k[[x, y, z, w]]/(xy - zw)$ 的穿孔谱的 Picard 群
非平凡。事实上，理想 $I = (x, z)$ 在穿孔谱上截出一个有效 Cartier
除子 $D$；该穿孔谱记为 $U$，它是 $A$ 的穿孔谱。容易看出
$I_x, I_y, I_z, I_w$ 分别是 $A_x, A_y, A_z, A_w$ 中的可逆理想。
另一方面，$A/I$ 的深度 $\geq 1$（事实上为 $2$），故 $I$ 的深度
$\geq 2$（事实上为 $3$），从而
$I = \Gamma(U, \mathcal{O}_U(-D))$。因此，若
$\mathcal{O}_U(-D)$ 平凡，则会有
$I \cong \Gamma(U, \mathcal{O}_U) = A$，但这是不成立的，因为 $I$
不是由 $1$ 个元素生成的。
\end{example}

\begin{example}
\label{algebraization-example-grothendieck-sharp-bis}
命题 \ref{algebraization-proposition-trivial-local-pic-complete-intersection} 不能推广到如下商：
$$
A = B/(f_1, \ldots, f_r)
$$
其中 $B$ 正则，且 $\dim(B) - r \geq 4$。换言之，
$f_1, \ldots, f_r$ 构成正则序列这一条件，一般是 $A$ 的穿孔谱的
Picard 群消失所必需的。
具体而言，令 $k$ 为域，并置
$$
A = k[[a, b, x, y, z, u, v, w]]/(a^3, b^3, xa^2 + yab + zb^2, w^2)
$$
注意 $A = A_0[w]/(w^2)$，其中
$A_0 = k[[a, b, x, y, z, u, v]]/(a^3, b^3, xa^2 + yab + zb^2)$。
下面将证明 $A_0$ 的深度为 $2$。以 $U$ 表示 $A$ 的穿孔谱，以
$U_0$ 表示 $A_0$ 的穿孔谱。注意有短正合列
$0 \to A_0 \to A \to A_0 \to 0$，其中第一个箭头由乘以 $w$ 给出。
由态射进阶，引理
\ref{more-morphisms-lemma-picard-group-first-order-thickening}，得到正合列
$$
H^0(U, \mathcal{O}_U^*) \to
H^0(U_0, \mathcal{O}_{U_0}^*) \to
H^1(U_0, \mathcal{O}_{U_0}) \to
\Pic(U)
$$
由于 $A_0$ 的深度、从而 $A$ 的深度为 $2$，可知
$H^0(U_0, \mathcal{O}_{U_0}) = A_0$ 且 $H^0(U, \mathcal{O}_U) = A$，
并且 $H^1(U_0, \mathcal{O}_{U_0})$ 非零，见对偶化复形，引理
\ref{dualizing-lemma-depth} 以及局部上同调，引理
\ref{local-cohomology-lemma-local-cohomology}。因此，上述最后一个箭头非零，
从而 $\Pic(U)$ 非零。

\medskip\noindent
要证明 $A_0$ 的深度为 $2$，只需证明
$A_1 = k[[a, b, x, y, z]]/(a^3, b^3, xa^2 + yab + zb^2)$
的深度为 $0$。这是因为 $a^2b^2$ 映到 $A_1$ 中的一个非零元素，
且该元素被每个变量 $a, b, x, y, z$ 零化。例如
$ya^2b^2 = (yab)(ab) = - (xa^2 + zb^2)(ab) = -xa^3b - yab^3 = 0$ 在 $A_1$ 中。
其他情形类似。
\end{example}










\section{Lefschetz 定理的应用}
\label{algebraization-section-lefschetz}

\noindent
本节讨论投射概形 $P$ 上的凝聚层与沿丰沛除子 $Q$
的形式完备化上的凝聚模之间的关系。

\medskip\noindent
设 $k$ 为域，$P$ 为 $k$ 上的固有概形。
设 $\mathcal{L}$ 为丰沛的可逆 $\mathcal{O}_P$-模。
取截面 $s \in \Gamma(P, \mathcal{L})$\footnote{我们不要求
$s$ 是正则截面。因此，$Q$ 仅为 $P$ 的局部主闭子概形，
而未必为有效 Cartier 除子。}，并令
$Q = Z(s)$ 为其零点概形，见
《除子》，定义 \ref{divisors-definition-zero-scheme-s}。
对所有 $n \geq 1$，记 $Q_n = Z(s^n)$
为第 $n$ 个无穷小邻域，其中心为 $Q$。
若 $\mathcal{F}$ 是凝聚 $\mathcal{O}_P$-模，则记
$\mathcal{F}_n = \mathcal{F}|_{Q_n}$ 为其限制；换言之，
它是 $\mathcal{F}$ 沿闭浸入 $Q_n \to P$ 的拉回。

\begin{proposition}
\label{algebraization-proposition-lefschetz}
在上述情形中，假设存在整数 $\sigma$，使得对每一点
$p \in P \setminus Q$ 都有
$$
\text{depth}(\mathcal{F}_p) + \dim(\overline{\{p\}}) > \sigma
$$
则映射
$$
H^i(P, \mathcal{F}) \longrightarrow \lim H^i(Q_n, \mathcal{F}_n)
$$
对 $0 \leq i < \sigma$ 为同构。
\end{proposition}

\begin{proof}
我们将使用《态射进阶》，引理 \ref{more-morphisms-lemma-apply-proj-spec}，
并在不再另行说明的情况下使用
《态射进阶》，第 \ref{more-morphisms-section-proj-spec} 节
中的记号和结果；若至少不了解该引理的陈述，本证明便无法理解。
注意，在我们的情形中
$A = \bigoplus_{m \geq 0} \Gamma(P, \mathcal{L}^{\otimes m})$
是有限型 $k$-代数，且它的每个分次部分都是有限维
$k$-向量空间，见《概形上同调》，
引理 \ref{coherent-lemma-coherent-proper-ample}。

\medskip\noindent
我们可以并且确实把 $s$ 看成元素 $f \in A_1 \subset A$，
即次数为 $1$ 的一个 $A$-齐次元素。记
$Y = V(f) \subset X$ 为由 $f$ 定义的闭子概形。
于是从概形论意义看，$U \cap Y = (\pi|_U)^{-1}(Q)$。
回顾记号
$\mathcal{F}_U = \pi^*\mathcal{F}|_U = (\pi|_U)^*\mathcal{F}$.
这是凝聚 $\mathcal{O}_U$-模。选取有限 $A$-模 $M$，使得
$\mathcal{F}_U = \widetilde{M}|_U$
（其存在性见《局部上同调》，引理
\ref{local-cohomology-lemma-finiteness-pushforwards-and-H1-local}）。
我们断言：$H^i_Z(M)$ 被 $f$ 的某个幂零化，只要
$i \leq \sigma + 1$。

\medskip\noindent
为证明该断言，我们应用《局部上同调》，
命题 \ref{local-cohomology-proposition-annihilator}。
将条件译成几何语言可知，只须证明：对
$u \in U$、$u \not \in Y$ 以及
$z \in \overline{\{u\}} \cap Z$，有
$$
\text{depth}(\mathcal{F}_{U, u}) +
\dim(\mathcal{O}_{\overline{\{u\}}, z}) > \sigma + 1
$$
这只需稍加思考。

\medskip\noindent
注意，$Z = \Spec(A_0)$ 是 $X$ 的有限个闭点所成的集合，
因为 $A_0$ 是有限维 $k$-代数。
（若读者希望 $Z$ 只有一个点，可以把有限 $k$-代数
$A_0$ 换成 $k$；这不会影响证明的其余部分。）

\medskip\noindent
态射 $\pi : L \to P$ 及其限制 $\pi|_U : U \to P$
都是相对维数为 $1$ 的光滑态射。
取 $u \in U$、$u \not \in Y$ 以及
$z \in \overline{\{u\}} \cap Z$，并令
$p = \pi(u) \in P \setminus Q$ 为其像。
于是 $u$ 或为 $\pi$ 在 $p$ 上纤维的泛点，或为该纤维的闭点。
若 $u$ 是纤维的泛点，则
$\text{depth}(\mathcal{F}_{U, u}) = \text{depth}(\mathcal{F}_p)$
且
$\dim(\overline{\{u\}}) = \dim(\overline{\{p\}}) + 1$.
若 $u$ 是纤维的闭点，则
$\text{depth}(\mathcal{F}_{U, u}) = \text{depth}(\mathcal{F}_p) + 1$
且
$\dim(\overline{\{u\}}) = \dim(\overline{\{p\}})$.
在这两种情形下都有
$\dim(\overline{\{u\}}) = \dim(\mathcal{O}_{\overline{\{u\}}, z})$
，因为 $Z$ 的每一点都是闭点。因此所需不等式由命题陈述中的
假设推出。

\medskip\noindent
令 $A'$ 为由 $f$ 所定义的 $A$ 的进完备化。由《代数》，
引理 \ref{algebra-lemma-completion-flat}，$A \to A'$ 是平坦的。
记 $U' \subset X' = \Spec(A')$ 为 $U$ 的逆像，并类似地定义
$Y'$ 与 $Z'$。令 $\mathcal{F}'$ 为 $U'$ 上由
$\mathcal{F}_U$ 拉回所得的模，并令
$M' = M \otimes_A A'$.
由局部上同调的平坦基变换
（《局部上同调》，引理
\ref{local-cohomology-lemma-torsion-change-rings}），有
$$
H^i_{Z'}(M') = H^i_Z(M) \otimes_A A'
$$
，故对 $i \leq \sigma + 1$，这些模均被 $f$ 的某个幂零化。
考虑交换图
$$
\xymatrix{
& H^i(U, \mathcal{F}_U) \ar[ld] \ar[d] \ar[r] &
\lim H^i(U, \mathcal{F}_U/f^n\mathcal{F}_U) \ar@{=}[d] \\
H^i(U, \mathcal{F}_U) \otimes_A A' \ar@{=}[r] &
H^i(U', \mathcal{F}') \ar[r] &
\lim H^i(U', \mathcal{F}'/f^n\mathcal{F}')
}
$$
由引理 \ref{algebraization-lemma-alternative-higher} 以及刚证明的挠性，
当 $i < \sigma$ 时，下方水平箭头为同构。
水平等号来自平坦基变换
（《概形上同调》，引理
\ref{coherent-lemma-flat-base-change-cohomology}）；
竖直等号则是因为 $U \cap Y$ 与 $U' \cap Y'$，
连同它们的第 $n$ 个无穷小邻域，都彼此同构地映射
（这是因为我们关于 $f$ 作完备化）。

\medskip\noindent
应用《态射进阶》，等式
(\ref{more-morphisms-equation-cohomology-torsor})
，得到相容的直和分解
$$
\lim H^i(U, \mathcal{F}_U/f^n\mathcal{F}_U) =
\lim
\left(
\bigoplus\nolimits_{m \in \mathbf{Z}}
H^i(Q_n, \mathcal{F}_n \otimes \mathcal{L}^{\otimes m})
\right)
$$
以及
$$
H^i(U, \mathcal{F}_U) =
\bigoplus\nolimits_{m \in \mathbf{Z}}
H^i(P, \mathcal{F} \otimes \mathcal{L}^{\otimes m})
$$
因此结论由《代数》，引理 \ref{algebra-lemma-daniel-litt} 得出。
\end{proof}

\begin{lemma}
\label{algebraization-lemma-lefschetz-addendum}
设 $k$ 为域，$X$ 为 $k$ 上的固有概形。
设 $\mathcal{L}$ 为丰沛的可逆 $\mathcal{O}_X$-模，
$s \in \Gamma(X, \mathcal{L})$。令 $Y = Z(s)$ 为 $s$ 的零点概形，
其第 $n$ 个无穷小邻域为 $Y_n = Z(s^n)$。
设 $\mathcal{F}$ 为凝聚 $\mathcal{O}_X$-模。
假设对所有 $x \in X \setminus Y$ 都有
$$
\text{depth}(\mathcal{F}_x) + \dim(\overline{\{x\}}) > 1
$$
则映射
$\Gamma(V, \mathcal{F}) \to \lim \Gamma(Y_n, \mathcal{F}|_{Y_n})$
对任意开子概形 $V \subset X$，只要它包含 $Y$，都是同构。
\end{lemma}

\begin{proof}
由命题 \ref{algebraization-proposition-lefschetz}，当 $V = X$ 时结论成立。
因此只须证明映射
$\Gamma(V, \mathcal{F}) \to \lim \Gamma(Y_n, \mathcal{F}|_{Y_n})$
是单射。若 $\sigma \in \Gamma(V, \mathcal{F})$
映为零，则其支撑与 $Y$ 不交（略去细节；提示：
使用 Krull 交定理）。于是，令 $T \subset X$ 为
$\text{Supp}(\sigma)$ 的闭包；它与 $Y$ 不交。
因此 $T$ 在 $k$ 上固有（因为它在 $X$ 中闭），
又是仿射的（因为它在仿射概形 $X \setminus Y$ 中闭，见
《态射》，引理 \ref{morphisms-lemma-proper-ample-delete-affine}），
从而在 $k$ 上有限
（《态射》，引理 \ref{morphisms-lemma-finite-proper}）。
故 $T$ 是 $X$ 的有限个闭点所成的集合。
由假设，$\text{depth}(\mathcal{F}_x) \geq 2$，因而深度至少为
$1$；这里 $x \in T$。
于是
$\Gamma(V, \mathcal{F}) \to \Gamma(V \setminus T, \mathcal{F})$
为单射，故如所需 $\sigma = 0$。
\end{proof}


\begin{example}
\label{algebraization-example-lefschetz}
设 $k$ 为域，$X$ 为 $k$ 上的固有簇。
设 $Y \subset X$ 为有效 Cartier 除子，且
$\mathcal{O}_X(Y)$ 丰沛；记 $Y_n$ 为它的第 $n$ 个无穷小邻域。
设 $\mathcal{E}$ 为有限局部自由 $\mathcal{O}_X$-模。
下面列出命题 \ref{algebraization-proposition-lefschetz} 的若干特殊情形。
\begin{enumerate}
\item 若 $X$ 是曲线，则得不到任何新结论。
\item 若 $X$ 是 Cohen-Macaulay 曲面（例如正规曲面），则
$$
H^0(X, \mathcal{E}) \to \lim H^0(Y_n, \mathcal{E}|_{Y_n})
$$
为同构。
\item 若 $X$ 是 Cohen-Macaulay 三维簇，则
$$
H^0(X, \mathcal{E}) \to \lim H^0(Y_n, \mathcal{E}|_{Y_n})
\quad\text{且}\quad
H^1(X, \mathcal{E}) \to \lim H^1(Y_n, \mathcal{E}|_{Y_n})
$$
均为同构。
\end{enumerate}
其规律应当已经清楚。若 $X$ 是正规三维簇，则可以对 $H^0$
得出结论，但不能对 $H^1$ 得出结论。
\end{example}

\noindent
在证明下一个主要结果之前，我们需要一个引理。

\begin{lemma}
\label{algebraization-lemma-Gm-equivariant-extend-canonically}
在情形 \ref{algebraization-situation-algebraize} 中，设 $(\mathcal{F}_n)$ 是
$\textit{Coh}(U, I\mathcal{O}_U)$ 的一个对象。假设
\begin{enumerate}
\item $A$ 是分次环，$\mathfrak a = A_+$，且 $I$ 是齐次理想；
\item $(\mathcal{F}_n) = (\widetilde{M_n}|_U)$，其中 $(M_n)$
是分次 $A$-模的逆系统；以及
\item $(\mathcal{F}_n)$ 可典范延拓到 $X$。
\end{enumerate}
则存在有限分次 $A$-模 $N$，使得
\begin{enumerate}
\item[(a)] 逆系统 $(N/I^nN)$ 与 $(M_n)$ 在分次 $A$-模
模去 $A_+$-幂挠模所得的范畴中为 pro-同构；并且
\item[(b)] $(\mathcal{F}_n)$ 是与 $N$ 相伴的凝聚模的完备化。
\end{enumerate}
\end{lemma}

\begin{proof}
令 $(\mathcal{G}_n)$ 为引理
\ref{algebraization-lemma-canonically-algebraizable} 中的典范延拓。
$A$ 与 $M_n$ 上的分次决定如下作用
$$
a : \mathbf{G}_m \times X \longrightarrow X
$$
；这是群概形 $\mathbf{G}_m$ 在 $X$ 上的作用，使
$(\widetilde{M_n})$ 成为
$\mathbf{G}_m$-等变拟凝聚 $\mathcal{O}_X$-模的逆系统，见
《群胚》，例 \ref{groupoids-example-Gm-on-affine}。
由于 $\mathfrak a$ 与 $I$ 是齐次理想，闭子概形 $Z$、$Y$
以及开子概形 $U$ 都是 $\mathbf{G}_m$-不变的。
$(\mathcal{F}_n)$ 是 $(\widetilde{M_n})$ 的限制，
并且是 $\mathbf{G}_m$-等变凝聚 $\mathcal{O}_U$-模的逆系统。
换言之，$(\mathcal{F}_n)$ 是 $\mathbf{G}_m$-等变凝聚形式模；
这意味着存在同构
$$
\alpha : (a^*\mathcal{F}_n) \longrightarrow (p^*\mathcal{F}_n)
$$
，它定义在 $\mathbf{G}_m \times U$ 上，并满足适当的余循环条件。
由于 $a$ 与 $p$ 是仿射概形之间的平坦态射，由引理
\ref{algebraization-lemma-canonically-extend-base-change} 可知，
存在唯一同构
$$
\beta : (a^*\mathcal{G}_n) \longrightarrow (p^*\mathcal{G}_n)
$$
，它定义在 $\mathbf{G}_m \times X$ 上，并将 $\alpha$
作为在 $\mathbf{G}_m \times U$ 上的限制。
唯一性保证 $\beta$ 满足相应的余循环条件。
由此，每个 $\mathcal{G}_n$ 都以与转移映射相容的方式
成为 $\mathbf{G}_m$-等变凝聚 $\mathcal{O}_X$-模。

\medskip\noindent
由《群胚》，引理 \ref{groupoids-lemma-Gm-equivariant-module}，
$\mathcal{G}_n$ 连同其 $\mathbf{G}_m$-等变结构
对应于分次 $A$-模 $N_n$。转移映射
$N_{n + 1} \to N_n$ 是分次模同态。注意，$N_n$ 是有限
$A$-模，并且 $N_n = N_{n + 1}/I^n N_{n + 1}$，因为
$(\mathcal{G}_n)$ 是 $\textit{Coh}(X, I\mathcal{O}_X)$ 的对象。
令 $N$ 为《代数》，引理
\ref{algebra-lemma-finiteness-graded} 中得到的有限分次 $A$-模。
于是 $N_n = N/I^nN$，故 $(\mathcal{G}_n)$
是与 $N$ 相伴的凝聚模的完备化，尤其 (b) 成立。

\medskip\noindent
为证明 (a)，还需更细致地展开上述情形。
首先，注意映射 $M_n \to H^0(U, \mathcal{F}_n)$ 的核与余核
都是 $A_+$-幂挠的（《局部上同调》，引理
\ref{local-cohomology-lemma-finiteness-pushforwards-and-H1-local}）。
注意，$H^0(U, \mathcal{F}_n)$ 带有自然分次，使这些映射及系统的
转移映射都是分次 $A$-模同态；例如，可以利用
$(U \to X)_*\mathcal{F}_n$ 是 $\mathbf{G}_m$-等变模，
并且定义在 $X$ 上，再应用《群胚》，
引理 \ref{groupoids-lemma-Gm-equivariant-module}。
其次，回顾定义 \ref{algebraization-definition-canonically-algebraizable}
和引理 \ref{algebraization-lemma-canonically-algebraizable}：
$(N_n)$ 与 $(H^0(U, \mathcal{F}_n))$ 为 pro-同构。
我们略去验证定义此 pro-同构的映射均为分次模同态。
因此，$(N_n)$ 与 $(M_n)$ 在分次 $A$-模模去
$A_+$-幂挠模所得的范畴中为 pro-同构。
\end{proof}

\noindent
设 $k$ 为域，$P$ 为 $k$ 上的固有概形。
设 $\mathcal{L}$ 为丰沛的可逆 $\mathcal{O}_P$-模。
取截面 $s \in \Gamma(P, \mathcal{L})$，并令
$Q = Z(s)$ 为其零点概形，见《除子》，
定义 \ref{divisors-definition-zero-scheme-s}。
令 $\mathcal{I} \subset \mathcal{O}_P$ 为 $Q$ 的理想层。
我们用 $\textit{Coh}(P, \mathcal{I})$ 表示
《概形上同调》，第 \ref{coherent-section-coherent-formal} 节
中引入的凝聚形式模范畴。

\begin{proposition}
\label{algebraization-proposition-lefschetz-existence}
在上述情形中，设 $(\mathcal{F}_n)$ 是
$\textit{Coh}(P, \mathcal{I})$ 的一个对象。假设对所有
$q \in Q$ 以及所有素理想
$\mathfrak p \in \mathcal{O}_{P, q}^\wedge$，若
$\mathfrak p \not \in V(\mathcal{I}_q^\wedge)$，都有
$$
\text{depth}((\mathcal{F}_q^\wedge)_\mathfrak p) +
\dim(\mathcal{O}_{P, q}^\wedge/\mathfrak p) +
\dim(\overline{\{q\}}) > 2
$$
则 $(\mathcal{F}_n)$ 是某个凝聚 $\mathcal{O}_P$-模的完备化。
\end{proposition}

\begin{proof}
由《概形上同调》，引理 \ref{coherent-lemma-existence-easy}，
为证明本命题，可以把 $(\mathcal{F}_n)$ 换成一个与它仅相差
$\mathcal{I}$-挠的对象（更精确的说明见下文）。
令 $T' = \{q \in Q \mid \dim(\overline{\{q\}}) = 0\}$，
并令 $T = \{q \in Q \mid \dim(\overline{\{q\}}) \leq 1\}$。
命题中的假设恰好说明
$Q \subset P$、$(\mathcal{F}_n)$ 以及
$T' \subset T \subset Q$ 在 $d = 1$ 时满足
引理 \ref{algebraization-lemma-improvement-formal-coherent-module-better}
的条件；除了直接整理不等式外，还要用到
$V(\mathfrak p) \cap V(\mathcal{I}^\wedge_y) = \{\mathfrak m^\wedge_y\}
\Leftrightarrow \dim(\mathcal{O}_{P, q}^\wedge/\mathfrak p) = 1$
，因为 $\mathcal{I}_y^\wedge$ 由 $1$ 个元素生成。
结合这两点，可以把 $(\mathcal{F}_n)$ 换成对象
$(\mathcal{H}_n)$；它属于 $\textit{Coh}(P, \mathcal{I})$，
并由引理 \ref{algebraization-lemma-improvement-formal-coherent-module-better}
给出。因此可以并且确实假设
$(\mathcal{F}_n)$ 与逆系统 $(\mathcal{F}_n'')$ 为
pro-同构；后者由凝聚 $\mathcal{O}_P$-模组成，并且满足
$\text{depth}(\mathcal{F}''_{n, q}) + \dim(\overline{\{q\}}) \geq 2$，
对所有 $q \in Q$ 都如此。

\medskip\noindent
我们将使用《态射进阶》，引理
\ref{more-morphisms-lemma-apply-proj-spec}，
并在不再另行说明的情况下使用
《态射进阶》，第 \ref{more-morphisms-section-proj-spec} 节
中的记号和结果；若至少不了解该引理的陈述，本证明便无法理解。
注意，在我们的情形中
$A = \bigoplus_{m \geq 0} \Gamma(P, \mathcal{L}^{\otimes m})$
是有限型 $k$-代数，且它的每个分次部分都是有限维
$k$-向量空间，见《概形上同调》，
引理 \ref{coherent-lemma-coherent-proper-ample}。

\medskip\noindent
由《概形上同调》，引理
\ref{coherent-lemma-inverse-systems-pullback}，
沿 $\pi|_U : U \to P$ 的拉回
$(\pi|_U^*\mathcal{F}_n)$ 是
$\textit{Coh}(U, f\mathcal{O}_U)$ 的对象，并且它与逆系统
$(\pi|_U^*\mathcal{F}_n'')$ 为 pro-同构；后者由凝聚
$\mathcal{O}_U$-模组成。
我们断言
$$
\text{depth}(\pi|_U^*\mathcal{F}''_{n, y}) + \delta_Z^Y(y) \geq 3
$$
，对所有 $y \in U \cap Y$ 都成立。由于 $Z$ 的所有点都是闭点，由引理
\ref{algebraization-lemma-discussion} 可知，
$\delta_Z^Y(y) \geq \dim(\overline{\{y\}})$ 对所有
$y \in U \cap Y$ 都成立。
令 $q \in Q$ 为 $y$ 的像。因为态射 $\pi : U \to P$
是相对维数为 $1$ 的光滑态射，所以 $y$ 或是 $\pi$
某条纤维的闭点，或是其泛点。因此
$$
\text{depth}(\pi^*\mathcal{F}''_{n, y}) + \delta_Z^Y(y)
\geq
\text{depth}(\pi^*\mathcal{F}''_{n, y}) + \dim(\overline{\{y\}}) =
\text{depth}(\mathcal{F}''_{n, q}) + \dim(\overline{\{q\}}) + 1
$$
，因为深度或维数二者之一增加 $1$。这就证明了断言。

\medskip\noindent
由引理 \ref{algebraization-lemma-cd-1-canonical} 可知，
$(\pi|_U^*\mathcal{F}_n)$ 可典范延拓到 $X$。注意
$$
M_n = \Gamma(U, \pi|_U^*\mathcal{F}_n) =
\bigoplus\nolimits_{m \in \mathbf{Z}}
\Gamma(P, \mathcal{F}_n \otimes_{\mathcal{O}_P} \mathcal{L}^{\otimes m})
$$
典范地成为分次 $A$-模，见《态射进阶》，等式
(\ref{more-morphisms-equation-cohomology-torsor})。
由《性质》，引理 \ref{properties-lemma-quasi-coherent-quasi-affine}，
有 $\pi|_U^*\mathcal{F}_n = \widetilde{M_n}|_U$。
因此可以应用引理
\ref{algebraization-lemma-Gm-equivariant-extend-canonically}，
得到有限分次 $A$-模 $N$，使 $(M_n)$ 与
$(N/I^nN)$ 在分次 $A$-模模去 $A_+$-挠模所得的范畴中
为 pro-同构。令 $\mathcal{F}$ 为凝聚 $\mathcal{O}_P$-模，
并要求它与 $N$ 相伴，见《概形上同调》，命题
\ref{coherent-proposition-coherent-modules-on-proj-general}。
同一命题还说明 $(\mathcal{F}/\mathcal{I}^n\mathcal{F})$
与 $(\mathcal{F}_n)$ 为 pro-同构。由于二者都是
$\textit{Coh}(P, \mathcal{I})$ 的对象，结论由引理
\ref{algebraization-lemma-recognize-formal-coherent-modules} 得出。
\end{proof}

\begin{example}
\label{algebraization-example-lefschetz-existence}
设 $k$ 为域，$X$ 为 $k$ 上的固有簇。
设 $Y \subset X$ 为有效 Cartier 除子，且
$\mathcal{O}_X(Y)$ 丰沛；记
$\mathcal{I} \subset \mathcal{O}_X$ 为相应的理想层。
设 $(\mathcal{E}_n)$ 是 $\textit{Coh}(X, \mathcal{I})$
的一个对象，且 $\mathcal{E}_n$ 有限局部自由。
下面列出命题 \ref{algebraization-proposition-lefschetz-existence} 的若干特殊情形。
\begin{enumerate}
\item 若 $X$ 是曲线或曲面，则得不到任何新结论。
\item 若 $X$ 是 Cohen-Macaulay 三维簇，则
$(\mathcal{E}_n)$ 是某个凝聚 $\mathcal{O}_X$-模
$\mathcal{E}$ 的完备化。
\item 更一般地，若 $\dim(X) \geq 3$ 且 $X$ 满足 $(S_3)$，
则 $(\mathcal{E}_n)$ 是某个凝聚 $\mathcal{O}_X$-模
$\mathcal{E}$ 的完备化。
\end{enumerate}
当然，若 $\mathcal{E}$ 存在，则 $\mathcal{E}$ 在 $Y$ 的某个开邻域上
有限局部自由。
\end{example}

\begin{proposition}
\label{algebraization-proposition-lefschetz-equivalence}
设 $k$ 为域，$X$ 为 $k$ 上的固有概形。
设 $\mathcal{L}$ 为丰沛的可逆 $\mathcal{O}_X$-模，并取
$s \in \Gamma(X, \mathcal{L})$。令 $Y = Z(s)$
为 $s$ 的零点概形，记 $\mathcal{I} \subset \mathcal{O}_X$
为相应的理想层。令 $\mathcal{V}$ 为 $X$ 中所有包含 $Y$
的开子概形所成的集合，并按反包含关系排序。
假设对所有 $x \in X \setminus Y$ 都有
$$
\text{depth}(\mathcal{O}_{X, x}) + \dim(\overline{\{x\}}) > 2
$$
则完备化函子
$$
\colim_\mathcal{V}
\textit{Coh}(\mathcal{O}_V)
\longrightarrow
\textit{Coh}(X, \mathcal{I})
$$
在有限局部自由对象所成的满子范畴之间给出等价。
\end{proposition}

\begin{proof}
为证明充分忠实，只须证明映射
$$
\colim_\mathcal{V} \Gamma(V, \mathcal{L}^{\otimes m})
\longrightarrow
\lim \Gamma(Y_n, \mathcal{L}^{\otimes m}|_{Y_n})
$$
对所有 $m$ 为同构，见引理
\ref{algebraization-lemma-completion-fully-faithful-general}。
这由引理 \ref{algebraization-lemma-lefschetz-addendum} 推出。

\medskip\noindent
本质满性。设 $(\mathcal{F}_n)$ 是
$\textit{Coh}(X, \mathcal{I})$ 的有限局部自由对象。
则对 $y \in Y$，有
$\mathcal{F}_y^\wedge = \lim \mathcal{F}_{n, y}$，且它是有限自由
$\mathcal{O}_{X, y}^\wedge$-模。
设 $\mathfrak p \subset \mathcal{O}_{X, y}^\wedge$
为满足 $\mathfrak p \not \in V(\mathcal{I}_y^\wedge)$ 的素理想。
那么 $\mathfrak p$ 位于某个素理想
$\mathfrak p_0 \subset \mathcal{O}_{X, y}$ 之上，而后者对应于
特化 $x \leadsto y$，其中 $x \not \in Y$。
由《局部上同调》，引理
\ref{local-cohomology-lemma-change-completion} 以及一些维数理论
（见《簇》，第 \ref{varieties-section-algebraic-schemes} 节），有
$$
\text{depth}((\mathcal{O}_{X, y}^\wedge)_\mathfrak p) +
\dim(\mathcal{O}_{X, y}^\wedge/\mathfrak p) =
\text{depth}(\mathcal{O}_{X, x}) +
\dim(\overline{\{x\}}) - \dim(\overline{\{y\}})
$$
因此，我们的假设蕴含命题
\ref{algebraization-proposition-lefschetz-existence} 的假设均成立，
从而 $(\mathcal{F}_n)$ 是某个凝聚
$\mathcal{O}_X$-模 $\mathcal{F}$ 的完备化。
于是 $\mathcal{F}_y$ 对所有 $y \in Y$ 都是有限自由的；
因此 $\mathcal{F}$ 在某个开邻域 $V$ 上有限局部自由，
且该邻域包含 $Y$。
证明完毕。
\end{proof}
